% \iffalse meta-comment
% arara: luatex: { files: [ spintent.ins] }
% arara: lualatex: { branch: developer, draft: true }
% arara: lualatex: { branch: developer, interaction: batchmode, shell: false }
% arara: clean: { extensions: [ aux, log, out, glo, gls, hd, ind, ilg, idx, toc, atfi ] }
%! arara: clean: { files: [ userdoc.idx, userdoc.ilg, userdoc.ind, mylhmc.sty, mylhmc.lua] }
% arara: halt
%
% File: spintent.dtx
%
% Copyright (C) 2026 by Pablo González L <pablgonz@educarchile.cl>
%
% This work may be distributed and/or modified under the conditions of the
% LaTeX Project Public License, either version 1.3c of this license or (at
% your option) any later version. The latest version of this license is in
%
%    https://www.latex-project.org/lppl.txt
%
% and version 1.3c or later is part of all distributions of LaTeX version
% 2005/12/01 or later.
%
% This work is "maintained" (as per the LPPL maintenance status)
% by Pablo González L.
%
% This work consists of the files spintent.dtx, spintent.lua and spintent.sty.
%
% \fi
%
% \iffalse
%<*driver>
\documentclass[full]{l3doc}
\usepackage[svgnames]{xcolor}
\usepackage[sf,bf,compact,medium,pagestyles]{titlesec}
\usepackage{lastpage,microtype,attachfile2,multicol,listings,titletoc,booktabs, longtable, fontawesome5}
\usepackage{scalerel} % for \spmQ sample
\usepackage[english]{babel} % \selectlanguage{english}
\usepackage[top=0.5in,bottom=0.3in,left=1.70in,right=0.7in,footskip=0.2in,headheight=1cm,headsep=0.27cm]{geometry}
% Set \parindent
\setlength{\parindent}{0pt}
% Only for arara... I like arara :)
\usepackage[scale=0.85]{comfortaa}
\definecolor{araracolor}{rgb}{0, 0.72, 0.28}
\newcommand{\araratext}[1]{{\small\normalfont\comfortaa\color{araracolor}\bfseries#1}}
\newcommand*\arara{\araratext{ar\kern-.03emar\kern-.03ema}}
% Fonts
\usepackage[osf,nomath,mono=false,ScaleSF=0.95,ScaleRM=0.95]{libertinus-otf}
\usepackage{sourcecodepro}
\defaultfontfeatures[\ttfamily]
  {
    Numbers   = OldStyle,
    Scale     = 0.84249,
    Extension = .otf,
  }
\setmonofont[
    UprightFont    = *-Regular,
    ItalicFont     = *-RegularIt,
    BoldFont       = *-Medium,
    BoldItalicFont = *-MediumIt,
    RawFeature     = {+zero,+ss01}]{SourceCodePro}
\newfontfamily\mysourcecodeprolight{SourceCodePro-Regular.otf}[
    Scale = 0.80,FakeStretch =0.75,
    FontFace = {m}{it}{Font =SourceCodePro-RegularIt.otf,Color=FF0000}]
\RenewDocumentCommand{\MacroLongFont}{}
  {
    \mysourcecodeprolight\small
  }
% The character of visible space is now taken from Latin Modern Mono
% to prevent fonts in T1. The original definition for xetex/luatex is
% \def\verbvisiblespace{\usefont{OT1}{cmtt}{m}{n}\asciispace}
\def\verbvisiblespace{{\fontfamily{lmtt}\selectfont\char"2423}}
% Math font
\typeout{ \fileversion }
\global\let\danger\relax
\usepackage[no-text,Scale = 0.95]{fourier-otf}
\typeout{ \fileversion }
\setmathfont{Erewhon-Math}[range={cal,bfcal},StylisticSet={1,4}]
\newfontfamily\fetamontotf{ffmw10.otf}[
   Scale             = 0.95,%
   RawFeature        = {+latn,+rand,+kern,+size},%
   FontFace          = {b}{n}{ffmw10.otf},% fix replace font
   ]
% Remove `"` and `|` for short verbatim defined in l3doc, then set short |...|
\AtBeginDocument
  {
    \DeleteShortVerb \"
    \DeleteShortVerb \|
    \lstMakeShortInline[language=spintent-doc,basicstyle=\ttfamily]\|
  }
% Index
\usepackage{imakeidx}
\indexsetup{level=\section,firstpagestyle=myheader,othercode=\pagestyle{myheader}}
% imakeidx also opens a write stream for the .idx file, and that conflicts
% with the one opened by l3doc. Here we close that write stream:
\ExplSyntaxOn
\iow_close:c { @indexfile }
% and later we will copy the write stream opened by imakeidx into \@indexfile
% so that entries written to both streams end up in the same file.
\makeindex[name=userdoc,options={-q ~ -s ~ gind.ist},columnsep=15pt,title={Índice ~ de ~ la ~ Documentación}]
\makeindex[options={-q ~ -s ~ gind.ist},columnsep=15pt,title={Índice ~ de ~ la ~ Implementación}]
\EnableCrossrefs
% \PageIndex, \CodelineIndex undoes what \PageIndex does \CodelineIndex
% tries to open another write stream for the index file. We don't want
% that, so we temporarily make \makeindex a no-op:
\cs_set_eq:NN \l__spintent_tmp_makeindex_cs \makeindex
\cs_set_nopar:Npn \makeindex { }
\CodelineIndex
\cs_set_eq:NN \makeindex \l__spintent_tmp_makeindex_cs
\DoNotIndex{\ , \1,\^}
\exp_args:Ne \DoNotIndex { \tl_to_str:n { \{ } }
\exp_args:Ne \DoNotIndex { \tl_to_str:n { \} } }
\exp_args:Ne \DoNotIndex { \tl_to_str:n { \begin } }
% \StartImplementation
\cs_new_protected:Npn \StartImplementation
  {
    \bool_set_true:N \l__codedoc_in_implementation_bool
  }
\bool_set_false:N \l__codedoc_in_implementation_bool
% Links and index
\newcommand{\HP}[1]{\emph{\hyperpage{#1}}\normalsize}
\cs_new_protected:Nn \__mydoc_sort_index:nn
  {
    \bool_if:NTF \l__codedoc_in_implementation_bool
      {
        \index{#1\actualchar#2|HP}
      }
      { \index[userdoc]{#1\actualchar#2|HP} }
  }
% Now, after imakeidx opens the write stream for the index file, we copy
% the reference to \@indexfile:
\cs_set_eq:cc { @indexfile } { \c_sys_jobname_str @idxfile }
\ExplSyntaxOff
% LuaLaTeX syntax highlighter for l3doc
\usepackage{mylhmc}
\setmylhmc
  {
    provide =
      { setspintent, spnum, spunit, spnewunit, spunitalias, spmoney,
        spshort, spsiglo, sptime, spdate, spread, spdsym, spmsym,
        spusep, spdiv, spmul, spnM, spnD, spmcd, spmcm, spnZ, spang,
        spmQ, spPoint, sprecta, sprayo, spside, spmside, spangle, spmangle,
        sparc, spmarc, spPoly, spAfig, spPfig, %
      }
  }
% Colors
\definecolor{clspkg}{HTML}{8DA0F5} % Azul claro
\definecolor{rulecolor}{rgb}{0.96,0.96,0.96}
\definecolor{nvdapurple}{rgb}{0.353, 0.176, 0.506}   % Púrpura NVDA (#5A2D81)
\definecolor{mathcatgreen}{rgb}{0.165, 0.478, 0.157} % Verde MathCAT (#2A7A28)
\definecolor{catgray}{rgb}{0.4, 0.4, 0.4}            % Gris para "CAT" o "ML"
% File info
\def\fileversion{v0.98}
\def\filedate{2026-08-20}
\def\fileinfo{Spanish parse intents}

% Config hyperref
\definecolor{linkcolor}{rgb}{0.04,0.38,0.04}
\definecolor{newlinkcolor}{HTML}{0A602C} % (0A6048 bonito Deep Teal) 029352
\hypersetup
  {
    allcolors          = newlinkcolor, %%%araracolor,%%%;linkcolor,
    colorlinks         = true,
    linktoc            = all,
    pdftitle           = {.:: The spintent package --- \fileinfo{} ::.},
    pdfauthor          = {Pablo González L},
    pdfsubject         = {Documentation for \fileversion{} [\filedate] },
    pdfkeywords        = {MathCAT, NVDA, tagged PDF, spanish, accesibilidad},
    bookmarksopenlevel = 2,
  }
% User level commands
\ExplSyntaxOn
% Logo
\NewDocumentCommand{\pkglogo}{}
  {
    \group_begin:
      \LibertinusSansOsF{\textcolor{module_name}{sp}\textcolor{OrangeRed}{intent}}
    \group_end:
  }
% Custom \meta[...]{...}, \marg[...]{...} and \oarg[...]{...} with color
\NewDocumentCommand{\mymeta}{O{}m}
  {
    \userdoc_meta_generic:Nnn \userdoc_meta:n { #1 } { #2 }
  }
\NewDocumentCommand{\mymarg}{O{}m}
  {
    \userdoc_meta_generic:Nnn \userdoc_marg:n { #1 } { #2 }
  }
\NewDocumentCommand{\myoarg}{O{}m}
  {
    \userdoc_meta_generic:Nnn \userdoc_oarg:n { #1 } { #2 }
  }
\NewDocumentCommand{\myparg}{O{}m}
  {
    \userdoc_meta_generic:Nnn \userdoc_parg:n { #1 } { #2 }
  }
% Variables and keys
\tl_new:N \l_userdoc_meta_font_tl
\keys_define:nn { userdoc / mymeta }
  {
    type .choice:,
    type / tt .code:n = \tl_set:Nn \l_userdoc_meta_font_tl { \ttfamily },
    type / rm .code:n = \tl_set:Nn \l_userdoc_meta_font_tl { \rmfamily },
    type .initial:n   = tt,
    cf .tl_set:N      = \l_userdoc_meta_color_tl,
    cf .initial:n     = keyname,
    ac .tl_set:N      = \l_userdoc_meta_anglecolor_tl,
    ac .initial:n     = lightgray,
    sbc .tl_set:N     = \l_userdoc_meta_brackcolor_tl,
    sbc .initial:n    = gray,
    cbc .tl_set:N     = \l_userdoc_meta_bracecolor_tl,
    cbc .initial:n    = gray,
    pbc .tl_set:N     = \l_userdoc_meta_parencolor_tl,
    pbc .initial:n    = gray,
  }
% Internal commands
\cs_new_protected:Npn \userdoc_meta_generic:Nnn #1 #2 #3
  {
    \group_begin:
      \keys_set:nn { userdoc / mymeta } { #2 }
      \color{ \l_userdoc_meta_color_tl }
      \l_userdoc_meta_font_tl
      #1 { #3 } % #1 is \userdoc_meta:n, \userdoc_marg:n or \userdoc_oarg:n
    \group_end:
  }
\cs_new_protected:Npn \userdoc_meta:n #1
  {
    \userdoc_meta_angle:n { \textlangle }
    \userdoc_meta_meta:n { #1 }
    \userdoc_meta_angle:n { \textrangle }
  }
\cs_new_protected:Npn \userdoc_marg:n #1
  {
    \userdoc_meta_brace:n { \textbraceleft }
    \userdoc_meta:n { #1 }
    \userdoc_meta_brace:n { \textbraceright }
  }
\cs_new_protected:Npn \userdoc_oarg:n #1
  {
    \userdoc_meta_brack:n { [ }
    \userdoc_meta:n { #1 }
    \userdoc_meta_brack:n { ] }
  }
\cs_new_protected:Npn \userdoc_parg:n #1
  {
    \userdoc_meta_brace:n { ( }
    \userdoc_meta:n { #1 }
    \userdoc_meta_brace:n { ) }
  }
\cs_new_protected:Npn \userdoc_meta_meta:n #1
  {
    \textnormal{\textit{#1}}
  }
\cs_new_protected:Npn \userdoc_meta_angle:n #1
  {
    \group_begin:
      \fontfamily{lmr}\selectfont
      \textcolor{\l_userdoc_meta_anglecolor_tl}{#1}
    \group_end:
  }
\cs_new_protected:Npn \userdoc_meta_brace:n #1
  {
    \group_begin:
      \color{\l_userdoc_meta_bracecolor_tl}
        #1
    \group_end:
  }
\cs_new_protected:Npn \userdoc_meta_brack:n #1
  {
    \textcolor{\l_userdoc_meta_brackcolor_tl}{#1}
  }
% \cmdexamp{o m o m o}
\tl_new:N \l_tempd_tl
\NewDocumentCommand \cmdexamp { o m o m o }
  {
    \group_begin:
    \small\ttfamily
    \textcolor{module_name}{\textbackslash#2}
    \IfValueT{#1}{ \textcolor{MediumOrchid}{#1} }
    \IfValueT{#3}
      {
        % \textnormal{\ttfamily } or \c{texttt}\cB\{  \cE\}
        \tl_set:Nn \l_tempd_tl { #3 }
        % \regex_replace_once:nnN { (\*) } { \c{textcolor}\cB\{MediumOrchid\cE\}\cB\{\1\cE\} } \l_tempd_tl
        \regex_replace_once:nnN { (\*) } { \c{texttt}\cB\{\c{textcolor}\cB\{MediumOrchid\cE\}\cB\{\1\cE\} \cE\} } \l_tempd_tl
        \myoarg[type=tt,sbc=gray,ac=lightgray,cf=keyname]{\tl_use:N \l_tempd_tl}
      }
    \IfValueTF{#5}
      {
        \mymeta[ac=gray,type=tt,cf=MediumOrchid]{#5}%
        \mymeta[type=tt,ac=gray,cf=keyname]{#4}
        \mymeta[ac=gray,type=tt,cf=MediumOrchid]{#5}%
      }
      { \mymarg[type=tt,cbc=gray,ac=lightgray,cf=keyname]{#4} }
    \par
    \group_end:
  }
% \keyexampcmd{s m m m o}
\NewDocumentCommand \keyexampcmd { s m m m o }
  {
    \noindent
    \mode_leave_vertical:
    \hbox_to_wd:nn { -\marginparsep }
      {
        \tex_hss:D
        \textcolor { keyname } { \small \ttfamily {#2} }
      }
    \IfBooleanTF {#1}
      {
        \hphantom { \textcolor { white } { \, \bfseries \texttt{=} } \, {} }
        \mymeta [ type=tt, cbc=gray, ac=lightgray, cf=lightgray ] { \small {#3} }
      }
      {
        \textcolor { gray } { \, \bfseries \texttt{=} } \, {}
        \mymarg [ type=tt, cbc=gray, ac=lightgray, cf=keyname ] { \small {#3} }
      }
    \hfill
    \textcolor { gray } { \small \textsf{default}:~ \emph{#4} } \par
    \IfValueT {#5}
      {
        \__mydoc_sort_index:nn { Keys }
          {
            Keys~for~ \texttt { \textbackslash #5 } ~ provide~by~ \pkglogo :> \texttt {#2}
          }
      }
  }

% \mykey{o m}
\NewDocumentCommand \mykey { o m }
  {
    \group_begin:
      \textcolor{keyname}{\ttfamily #2}
      \IfValueT {#1}
        {
          \clist_map_inline:nn {#1}
            {
              \__mydoc_sort_index:nn {Keys} {Keys~for~\texttt{\textbackslash{}##1}~provide~by~\pkglogo :>\texttt{#2}}
            }
        }
    \group_end:
  }
% \mypkg{sm}
\NewDocumentCommand \mypkg { s O{} m }
  {
    \group_begin:
    \IfBooleanTF {#1}
      {
        \pkglogo
      }
      {
        \textcolor{pkgname} { \textbf{\texttt{#3} } } %\texttt{#3}
      }
    \str_if_eq:nnF {#2} { no-index }
      {
        \__mydoc_sort_index:nn { packages } { Packages~:>~\texttt{#3} }
      }
    \group_end:
  }
% \myclass{m}
\NewDocumentCommand \myclass { m }
  {
    \textcolor{clsname} { \textbf{\texttt{#1} } } % \texttt{#3}
    \__mydoc_sort_index:nn {class} {Document~class :>\texttt{#1}}
  }
% \mynvmath
\tl_new:N \l__mynvmath_input_tl
\NewDocumentCommand \mynvmath { m }
  {
    \group_begin:
      \tl_set:Nn \l__mynvmath_input_tl { #1 }
      \regex_replace_all:nnN { NVDA } { \c{textcolor}\cB\{percent\cE\}\cB\{NVDA\cE\} } \l__mynvmath_input_tl
      \regex_replace_all:nnN { Math } { \c{textcolor}\cB\{percent\cE\}\cB\{Math\cE\} } \l__mynvmath_input_tl
      \regex_replace_all:nnN { (CAT|ML) } { \c{textcolor}\cB\{comment\cE\}\cB\{\1\cE\} }  \l__mynvmath_input_tl
      \regex_replace_all:nnN { / } { \c{textcolor}\cB\{backtick\cE\}\cB\{/\cE\} } \l__mynvmath_input_tl
      \textbf{\texttt{ \tl_use:N \l__mynvmath_input_tl } }
    \group_end:
  }
% \myenv{sm}
\NewDocumentCommand \myenv { s m O{}}
  {
    \IfBooleanTF{#1}
      {
        \tl_set:Nn \l_tmpa_tl { #2 }
        \regex_replace_once:nnN { (\*) } { \c{textcolor}\cB\{MediumOrchid\cE\}\cB\{\1\cE\} } \l_tmpa_tl
        \textcolor{module_name}{\ttfamily{\tl_use:N \l_tmpa_tl}}%
        \__mydoc_sort_index:nn {environment} {Environments~provide~by ~\pkglogo :>\texttt{\tl_use:N \l_tmpa_tl}}
      }
      {
        \textcolor{SlateBlue}{\ttfamily{#2}}%
        \__mydoc_sort_index:nn {environment} {Environments :>\texttt{#2}}
      }
  }
% \ics{sm}
\NewDocumentCommand \ics {s m}
  {
    \IfBooleanTF{#1}
      {
        \tl_set:Nn \l_tmpa_tl { #2 }
        \regex_replace_once:nnN { (\*) } { \c{textcolor}\cB\{MediumOrchid\cE\}\cB\{\1\cE\} } \l_tmpa_tl
        \regex_replace_once:nnN { (!) } { \c{textcolor}\cB\{MediumOrchid\cE\}\cB\{\1\cE\} } \l_tmpa_tl
        \textcolor{module_name}{\ttfamily\textbackslash{\tl_use:N \l_tmpa_tl}}
        \__mydoc_sort_index:nn {Commands} {Commands~provide ~by~\pkglogo :>\texttt{\textbackslash#2}}
      }
      {
        \textcolor{latex_cmd}{\ttfamily\textbackslash{#2}}
        \__mydoc_sort_index:nn {#2} {\texttt{\textbackslash#2}}
      }
  }
% \mydim{m}
\NewDocumentCommand \mydim { m }
  {
    \group_begin:
      \tl_set:Nn \l_tmpa_tl { #1 }
      \regex_replace_all:nnN { (\d) } { \c{textcolor}\cB\{number\cE\}\cB\{\1\cE\} } \l_tmpa_tl
      %\regex_replace_all:nnN  { (\.) } { \c{textcolor}\cB\{unit\cE\}\cB\{\1\cE\} } \l_tmpa_tl
    \textcolor{unit}{ \ttfamily \tl_use:N \l_tmpa_tl }
    %%\texttt { \tl_use:N \l_tmpa_tl }
    \group_end:
  }
% \myluamod{}
\tl_new:N \l__mydoc_luamod_tl
\NewDocumentCommand \myluamod { }
  {
    \group_begin:
      \tl_set:Nn \l__mydoc_luamod_tl { spintent.lua }
      \regex_replace_all:nnN { spintent }
        {
          \c{textcolor}\cB\{module_name\cE\}\cB\{sp\cE\}
          \c{textcolor}\cB\{OrangeRed\cE\}\cB\{intent\cE\}
        }
        \l__mydoc_luamod_tl
      \regex_replace_once:nnN { \. ( .* ) \Z }
        {
          \c{textcolor}\cB\{MediumOrchid\cE\}\cB\{\.\cE\}
          \c{textcolor}\cB\{gray\cE\}\cB\{\1\cE\}
        }
        \l__mydoc_luamod_tl
      \texttt{ \l__mydoc_luamod_tl }
      \__mydoc_sort_index:nn {Module} {Lua~module~provide ~by~\pkglogo :>\texttt{spintent.lua}}
    \group_end:
  }
% \myfile{m}
\NewDocumentCommand \myfile { m }
  {
    \group_begin:
      \tl_clear:N \l_tmpa_tl
      \tl_set:Nn \l_tmpa_tl { #1 }
      \regex_replace_all:nnN { (\d) } { \c{textcolor}\cB\{MediumOrchid\cE\}\cB\{\1\cE\} } \l_tmpa_tl
      \regex_replace_all:nnN  { (\.) } { \c{textcolor}\cB\{MediumOrchid\cE\}\cB\{\1\cE\} } \l_tmpa_tl
      \textcolor{filename}{ \ttfamily \tl_use:N \l_tmpa_tl }
    \group_end:
  }

% Iconos al margen para entorno important (uso interno)
\cs_new_protected:Npn \__mydoc_margin_icon:n #1
  {
    \par
    \mode_leave_vertical:
    \skip_horizontal:n { -0.5 \marginparsep }
    \hbox_overlap_left:n
      {
        \hbox_to_wd:nn { \marginparsep }
          {
            \tex_hss:D
            \textcolor { module_at } { #1 } % privatefun
          }
      }
    \skip_horizontal:n { 0.5 \marginparsep }
  }
% Entorno de notas
\NewDocumentEnvironment { important } { s }
  {
    \small\sffamily
    \IfBooleanTF { #1 }
      {
        \__mydoc_margin_icon:n { \footnotesize \faIcon { bomb } }
      }
      {
        \__mydoc_margin_icon:n { \scriptsize \reflectbox { \faIcon { tag } } }
      }
    \ignorespaces
  }
  { }
% Análisis y formato de correo electrónico
\seq_new:N \l__spintent_email_seq
\cs_new_protected:Npn \__spintent_format_email:n #1
  {
    \seq_set_split:Nnn \l__spintent_email_seq { @ } { #1 }
    \href { mailto:#1 }
      {
        \ttfamily
        \seq_item:Nn \l__spintent_email_seq { 1 }
        \textcolor { arroba } { \footnotesize @ }
        \seq_item:Nn \l__spintent_email_seq { 2 }
      }
  }
\NewDocumentCommand \email { m }
  {
    \__spintent_format_email:n { #1 }
  }
% Sobrescritura de la Tabla de Contenidos (TOC)
\cs_new_eq:Nc \__spintent_mkboth:nn { @mkboth }
\cs_new_eq:Nc \__spintent_starttoc:n { @starttoc }
\RenewDocumentCommand \tableofcontents { }
  {
    \group_begin:
    \section*
      {
        \contentsname \quad
          {
            \color { gray }
            \tex_leaders:D \tex_hrule:D height 5pt depth -4.4pt \tex_hfill:D
          }
      }
    \__spintent_mkboth:nn { \text_uppercase:n { \contentsname } } { \text_uppercase:n { \contentsname } }
    \vspace* { -14pt }
    \dim_set:Nn \columnsep { 10pt }
    \begin{multicols}{2}
      \__spintent_starttoc:n { toc }
    \end{multicols}
    \vspace* { -3pt }
      {
        \color { gray }
        \tex_hrule:D height 0.6pt
      }
    \vspace* { 5pt }
    \group_end:
  }
% \thanks
\cs_new:Npn \__spintent_fnsymbol:n #1
  {
    \ensuremath
      {
        \int_case:nnF { #1 }
          {
            { 1 } { \textcolor { module_name } { * } }
            { 2 } { \textcolor { module_name } { \dagger } }
            { 3 } { \ddagger }
            { 4 } { \mathsection }
            { 5 } { \mathparagraph }
            { 6 } { \| }
            { 7 } { ** }
            { 8 } { \dagger \dagger }
            { 9 } { \ddagger \ddagger }
          }
          { \@ctrerr }
      }
  }
\hook_gput_code:nnn { begindocument } { spintent }
  {
    \cs_set_eq:cN { @fnsymbol } \__spintent_fnsymbol:n
  }

% Don't copy numbers in code example
\newcommand*{\noaccsupp}[1]
  {
    \__mylhmc_actual_text_begin: #1 \__mylhmc_actual_text_end:
  }
\ExplSyntaxOff

% Create a language for documentation
\lstdefinelanguage{spintent-doc}{
    texcsstyle    = *,%
    escapechar    =`,%
    extendedchars = true,%
    keepspaces    = true,%
    stringstyle   = {\color{red}},%
    basicstyle    = \small\ttfamily,% \small
    alsoletter    = {-,.},% @,
% comments
    morecomment   = [l]{\%},%
    commentstyle  = \ttfamily\small\color{comment},%
% Important \NeedsTeXFormat \ProvidesExplPackage \documentclass \DocumentMetadata
    keywordstyle  = [1]{\bfseries\color{structure}},%
    keywords      = [1]{documentclass, makeatletter, makeatother, %
                    },%
% Other words \usepackage,
    keywordstyle  = [2]{\color{pkgstruct}},%
    keywords      = [2]{usepackage,},%,
% Other words 3 MathML/tagpdf
    keywordstyle  = [3]{\color{mathtag}},%
    keywords      = [3]{MathMLintent,MathMLarg, invisibletimes, luamml_attribute, luamml_annotate, DocumentMetadata,},%,
% Other words 4
    keywordstyle  = [4]{\color{latex_cmd}},%
    keywords      = [4]{frac, dfrac, section, input, listparindent, linewidth, emph, textbf, footnote,%
                    triangleright, diamond, circ, parleft, star, centering, ast, sptext, , href, label, ref, %
                    mathord, parallel, mathstyle,.choice, .code, .meta, text, textsuperscript, cs_to_str, c_zero_skip, keys_define, %
                    mode_leave_vertical, Umathchardef,mathrm, mathbf, textstyle, addfontfeatures, %
                    left, right, Uleft, Uright,  },%
% Other words 5
    keywordstyle  = [5]{\bfseries\color{pkgname}},%
    keywords      = [5]{babel,spintent,libertinus-otf,sourcecodepro,fourier-otf,Erewhon-Math,hyperref,},%,
% Reserved words 6 (document class)
    keywordstyle  = [6]{\bfseries\color{clsname}},%
    keywords      = [6]{article, l3doc,},%
% Reserved words 7 (spintent cmd)
    keywordstyle  = [7]{\color{module_name}},%
    keywords      = [7]{spnum, spunit, spmoney, spshort, spdiv, spang, sptime, spdate, spsiglo, spmcm, spnD, spnM, %
                    spmcd, spmul, spread, spdsym, spmsym, spnewunit, spunitalias, spnZ, spmQ , setspintent, spreadsym,
                    spusep, },%
% Reserved words 8 (keyname) KEYS
    keywordstyle  = [8]{\color{keyname}},%
    keywords      = [8]{print-era, none, ac, dc, currency, sub-units, clear, terse, times, cdot, true, false, unknown, valid, %
                    multislash, coma, punto, result, join-word, entero, enteros, y, error, periódico, postfix,%
                    es-paralela, es-paralelo, two-dots, old-style, slash, dividido, dividido-por, dividido-entre, por,%
                    multiplicado-por, veces, cross, dot, asterisk, sample-arg, cifras-sep, read-cmd, %
                    osf, nomath, mono, ScaleSF, ScaleRM, spanish, shorthands, no-text, semibold, colorlinks, pdftitle, pdfdisplaydoctitle, %
                    pdfstartview, FitH, pdfinfo, Title, Subject, range, StylisticSet, small-caps, read-arg, píxeles, px, caudal,%
                    read-period-sep, read-sign, notation-sym, partido-por, Numbers, Lining, OldStyle, size,%
                    en-US, math, mathml-SE, es-CL, on, off, setup, tagging, tagging-setup, lang, pdfversion, pdfstandard,},%
% Reserved words 9 (intent values)
    keywordstyle  = [9]{\color{ltcmd_arg}},%
    keywords      = [9]{number,intent,test,sample,above,below, ActualText, Span, E, cal,bfcal, silent, number, %
                    intent, title, table, firstaid, on,XXI, setup,},%
% Reserved words 10 (signature)
    keywordstyle  = [10]{\color{sigargs}},%
    keywords      = [10]{een, },%
% Reserved words 11 (support)
    keywordstyle  = [11]{\color{support}},%
    keywords      = [11]{hypersetup, setmathfont,},%
% literate
    literate      =*{\{}{{{\color{brace}{\char`\{}}}}{1}
                    {\}}{{{\color{brace}{\char`\}}}}}{1}
                    {\|}{{{\color{gray}{\char`\|}}}}{1}
                    {\$}{{{\color{math}{\char`\$}}}}{1}
                    {\#}{{{\color{argument}{\char`\#}}}}{1}
                    {(}{{\textcolor{gray}{(}}}{1}
                    {)}{{\textcolor{gray}{)}}}{1}
                    {[}{{\textcolor{bracket}{[}}}{1}
                    {]}{{\textcolor{bracket}{]}}}{1}
                    {,}{{\textcolor{gray}{,}}}{1}
                    {;}{{\textcolor{gray}{;}}}{1}
                    {.}{{\textcolor{signature}{.}}}{1}
                    {:}{{\textcolor{signature}{:}}}{1}
                    {`}{{\textcolor{quote}{\textasciigrave}}}{1}
                    {\&}{{\textcolor{gray}{\&}}}{1}
                    {>}{{\textcolor{module_eq}{>}}}{1}
                    {<}{{\textcolor{module_eq}{<}}}{1}
                    {*}{{\textcolor{star}{*}}}{1}
                    {ua-2}{{\textcolor{keyname}{ua-2}}}{4} % UA-2
                    {!}{{\textcolor{MediumOrchid}{!}}}{1}
                    {@}{{\textcolor{arroba}{@}}}{1}%
                    {\^}{{\textcolor{keyname}{\textasciicircum}}}{1}
                    {0}{{\textcolor{argnum}{0}}}{1}
                    {1}{{\textcolor{argnum}{1}}}{1}
                    {2}{{\textcolor{argnum}{2}}}{1}
                    {3}{{\textcolor{argnum}{3}}}{1}
                    {4}{{\textcolor{argnum}{4}}}{1}
                    {5}{{\textcolor{argnum}{5}}}{1}
                    {6}{{\textcolor{argnum}{6}}}{1}
                    {7}{{\textcolor{argnum}{7}}}{1}
                    {8}{{\textcolor{argnum}{8}}}{1}
                    {9}{{\textcolor{argnum}{9}}}{1}
                    {.0}{{\textcolor{argnum}{.0}}}{2}% Following is to ensure that only periods
                    {.1}{{\textcolor{argnum}{.1}}}{2}% followed by a digit are changed.
                    {.2}{{\textcolor{argnum}{.2}}}{2}
                    {.3}{{\textcolor{argnum}{.3}}}{2}
                    {.4}{{\textcolor{argnum}{.4}}}{2}
                    {.5}{{\textcolor{argnum}{.5}}}{2}
                    {.6}{{\textcolor{argnum}{.6}}}{2}
                    {.7}{{\textcolor{argnum}{.7}}}{2}
                    {.8}{{\textcolor{argnum}{.8}}}{2}
                    {.9}{{\textcolor{argnum}{.9}}}{2}
                    {=}{{\textcolor{module_eq}{=}}}{1},%
}[keywords,tex,comments,strings]% end language

% \begin{examplecode}[key=val]...\end{examplecode}
\lstnewenvironment{examplecode}[1][]{%
\lstset{
    language    = spintent-doc,%
    stringstyle = {\color{red}},%
    basicstyle  = \ttfamily\small,%
    columns     = fullflexible,%
    fontadjust  = true,%
    aboveskip   = 8pt,%
    belowskip   = 5pt,%
    rulecolor   = \color{gray!50},%
    framesep    = \fboxsep,%
    framerule   = \fboxrule,%
    xleftmargin = \dimexpr\fboxsep+\fboxrule\relax,%
    xrightmargin= \dimexpr\fboxsep+\fboxrule\relax,%
           #1,%
    }% close lstset
}{\vspace{-\parskip}}% close examplecode
% \lstinline[style=inline]|...|
\lstdefinestyle{inline}
  {
    language    = spintent-doc,%
    upquote     = true,%
    numbersep   = 1em,%
    escapechar  =`,%
    basicstyle  =\ttfamily,%\small
    numberstyle = \tiny\color{lightgray}\noaccsupp,%
    literate    =*{\%}{{\textcolor{percent}{\%}}}{1}
                  {i}{{\textcolor{keyname}{i}}}{1}
                  {ii}{{\textcolor{keyname}{ii}}}{2}
                  {iii}{{\textcolor{keyname}{iii}}}{3}
                  {iv}{{\textcolor{keyname}{iv}}}{2}
                  {v}{{\textcolor{keyname}{v}}}{1}
                  {vi}{{\textcolor{keyname}{vi}}}{1}
  }
% Set default \lstinline style
\lstset{style=inline}

% Configuration titleps and titlesec
\settitlemarks{section}
\renewpagestyle{plain}[\color{keyname}\small\sffamily]{%
\setfoot{\rlap{\hskip\dimexpr-\oddsidemargin-1in\relax%
         \parbox{1.93\paperwidth}{\hfil\thepage\,/\,\pageref{LastPage}}}}%
        {}%
        {}%
}
\newpagestyle{myheader}[\color{keyname}\small\sffamily]{%
\renewcommand\makeheadrule{%
\rlap{\hskip\dimexpr-\oddsidemargin-1in\relax
      \color{rulecolor}\rule[0.3cm]{\paperwidth}{0.7cm}}\hss
}
\setfoot{\rlap{\hskip\dimexpr-\oddsidemargin-1in\relax%
         \parbox{1.93\paperwidth}{\hfil\thepage\,/\,\pageref{LastPage}}}}%
        {\parbox{\textwidth}{\raggedright \textcolor{gray}{\raisebox{-1pt}{\textcopyright}{}2026 by Pablo González L}}}%
        {}%
\sethead{\llap{\raisebox{0.55cm}{\parbox{\dimexpr\oddsidemargin+1in\relax}{\makebox[0pt][l]{\hspace{15pt}\pkglogo\space\fileversion}}}}}
        {\raisebox{0.55cm}{\parbox{\textwidth}{\hspace*{-\oddsidemargin}\centering\small\S.\thesection\space\sectiontitle}}}%
        {}%
}
\titlecontents{section}[0mm]{}%
    {\bfseries\contentspush{\makebox[5mm][l]{\thecontentslabel\hfill}}}%
    {\hspace*{-5mm}}% numberless
    {\hspace{0.25em}\titlerule*[6pt]{.}\contentspage}%
\titlecontents{subsection}[5mm]{}%
    {\contentspush{\makebox[6mm][l]{\thecontentslabel\hfill}}}
    {\hspace*{-11mm}}% numberless
    {\hspace{0.25em}\titlerule*[6pt]{.}\contentspage}%
\titlecontents{subsubsection}[11mm]{}%
    {\contentspush{\makebox[8mm][l]{\thecontentslabel\hfill}}}
    {\hspace*{-19mm}}% numberless
    {\hspace{0.25em}\titlerule*[6pt]{.}\contentspage}%
\begin{document}
  \DocInput{\jobname.dtx}
\end{document}
%</driver>
% \fi
%
% \selectlanguage{spanish}
%
% \title{
%    \scalebox{2.67}{\pkglogo}\\[0pt]
%    {\Large\textcolor{guard_angle}{\fetamontotf{spanish parse intents}}}\\[-5pt]
%    \fileversion{} \quad \filedate\thanks{
%     Este archivo describe la documentación para la versión \fileversion{}
%     revisada por última vez el
%    \filedate.}\\[30pt]
%    \author{
%    \large
%    \raisebox{-1pt}{\textcolor{guard_angle}{\textcopyright{}2026 by Pablo González L}}\thanks{
%    E-mail: \textcolor{guard_angle}{\textsf{\guillemotleft}}\email{pablgonz@educarchile.cl}\textcolor{guard_angle}{\textsf{\guillemotright}}.
%       }%
%    }
% \small
% \vspace*{-5pt}
% \textsc{ctan}: \url{https://www.ctan.org/pkg/spintent}\\
% \textcolor{gray}{\scriptsize\faIcon[regular]{github}}\,\,\,\url{https://github.com/pablgonz/spintent}
% \vspace*{-2cm}
% }
% \date{}
% \maketitle
%
% \begin{abstract}
%  El paquete \mypkg*[no-index]{spintent} proporciona una serie de utilidades para
%  docentes de educación primaria y secundaria que necesiten crean documentos
%  \texttt{PDF} \emph{accesibles} (\emph{tagged} \texttt{PDF}) en español
%  usando \hologo{LuaLaTeX}.
% \end{abstract}
%
% \tableofcontents
%
% \section*{Motivación y agradecimientos}
%
% Desde que desarrollé \mypkg{scontents}\cite{scontents} en 2019 y, más
% tarde, \mypkg{enumext}\cite{enumext} en 2024, me he ido interesando
% cada vez más por la accesibilidad en los documentos \texttt{PDF}. Al
% principio era un tema prácticamente desconocido para mí, pero a medida
% que fui aprendiendo sobre él comprendí la importancia que tiene,
% especialmente en el ámbito educativo.
%
% \smallskip
%
%
% El equipo del proyecto de \hologo{LaTeX}\cite{tagged-project} ha
% realizado un trabajo extraordinario en este campo. Hoy es posible
% generar documentos etiquetados, validarlos con herramientas como
% \href{https://dev.verapdf-rest.duallab.com}{\texttt{veraPDF}}\cite{verapdf}
% o inspeccionar su estructura con
% \href{https://ngpdf.com}{\texttt{ngpdf}}\cite{ngpdf}. Sin embargo,
% siempre me quedó una duda: cuando un estudiante utiliza \mynvmath{NVDA}\cite{nvda}
% junto con \mynvmath{MathCAT}\cite{mathcat}, ¿realmente escucha lo que
% el o la docente quiso expresar?
%
% \smallskip
%
% Basta con generar un documento accesible siguiendo las recomendaciones
% de \href{https://latex3.github.io/tagging-project/}{\texttt{The LaTeX
% Tagged PDF Project}} y escucharlo con \mynvmath{NVDA/MathCAT} para darse
% cuenta de que, en muchos casos, la \emph{verbalización} no coincide con lo
% que uno espera. Esto no resulta extraño: las reglas del español, y en
% particular el lenguaje matemático, dependen en gran medida del contexto
% en que se utilizan.
%
% \smallskip
%
% El objetivo de \mypkg*[no-index]{spintent} nace precisamente de esta necesidad:
% proporcionar herramientas para que el o la docente pueda indicar
% explícitamente el contexto de determinadas expresiones matemáticas
% desde el propio código fuente de \hologo{LaTeX}, de modo que el
% documento \texttt{PDF} resultante no solo sea \emph{accesible}, sino que
% también produzca una \emph{verbalización más natural} y adecuada para el
% estudiante al utilizar \mynvmath{NVDA/MathCAT}
%
% \smallskip
%
% El nuevo estándar \texttt{PDF} |2.0|\cite{pdf20,pdf20-impact,pdf20-errata},
% las especificaciones de \texttt{WTPDF}\cite{wtpdf} y la guía más reciente de la
% \texttt{PDF} Association, «Best Practice Guide: Math in
% \texttt{PDF}»\cite{bpg-math},
% coinciden en que el camino hacia una accesibilidad matemática de calidad
% pasa por el uso de \mynvmath{MathML}. En ese contexto,
% \hologo{LuaLaTeX} ofrece una plataforma especialmente adecuada para
% incorporar esta tecnología.
%
% \thispagestyle{plain}
%
% \section*{Licencia y requerimientos}
%
% El paquete \mypkg*[no-index]{spintent} carga y requiere el paquete \mypkg{unicode-math}\cite{unicode-math}
% y que se ejecute bajo \hologo{LuaLaTeX} por lo que es necesario contar con
% una distribución moderna de \hologo{TeX} como \hologo{TeX} Live 2026. Ha sido probado con las clases estándar
% proporcionadas por \hologo{LaTeX}: \myclass{book}, \myclass{report}, \myclass{article} y
% \myclass{letter} en tamaños \mydim{10pt}, \mydim{11pt} y \mydim{12pt}.
%
% \medskip
%
% Se otorga permiso para copiar, distribuir y/o modificar este paquete bajo
% los términos de \LaTeX{} Project Public License (\textsc{lppl}), versión
% 1.3 o posterior (\url{https://www.latex-project.org/lppl.txt}). El paquete
% tiene la condición de mantenido al estilo \emph{«A Work in Progress\footnote{%
% \emph{A Work in Progress} es un álbum en vivo de Neil Peart
% ($\ast$~1952--\textdied~2020), baterista y principal letrista de Rush.
% El nombre no pretende formar parte de la terminología de la
% \textsc{lppl}; simplemente describe el estado permanente de este paquete
% mientras sigo aprendiendo sobre \emph{accesibilidad} en \texttt{PDF}.}»}.
%
% \smallskip
%
% \begin{important}*
%  El requerimiento mínimo es \hologo{LaTeX} release 2026-11-01.
% \end{important}
%
% \smallskip
%
% Este paquete solo se puede utilizar solo con \hologo{LuaLaTeX} y
% \emph{tagged} \texttt{PDF} activo. Está presente en \hologo{TeX} Live y \hologo{MiKTeX},
% para instalarlo, utilice el gestor de paquetes de su distribución. Si
% desea una instalación manual, descargue \href{http://mirrors.ctan.org/macros/latex/contrib/spintent.zip}{spintent.zip}
% y descomprímalo, luego ejecute \lstinline[language=spintent-doc,basicstyle=\ttfamily]+luatex spintent.ins+,
% mueva todos los archivos a las ubicaciones correspondientes y luego
% ejecute |mktexlsr|. Para generar la documentación, ejecute
% \lstinline[language=spintent-doc,basicstyle=\ttfamily]+arara spintent.dtx+.
%
% \iffalse
%<*example>
% \fi
\begin{examplecode}[frame=single,language={}]
  spintent.sty  `$\Rightarrow$`  TDS:tex/lualatex/spintent/
  spintent.lua  `$\Rightarrow$`  TDS:tex/lualatex/spintent/
  spintent.pdf  `$\Rightarrow$`  TDS:doc/lualatex/spintent/
  README.md     `$\Rightarrow$`  TDS:doc/lualatex/spintent/
  spintent.dtx  `$\Rightarrow$`  TDS:source/lualatex/spintent/
  spintent.ins  `$\Rightarrow$`  TDS:source/lualatex/spintent/
\end{examplecode}
% \iffalse
%</example>
% \fi
%
% \pagestyle{myheader}
%
% \section{Descripción}
%
% El nombre \mypkg*[no-index]{spintent} puede entenderse como \emph{«Spanish Parse
% Intent»} o también como \emph{«School Parse Intent»}. Ambas
% interpretaciones reflejan el objetivo del paquete y el público para el
% que fue diseñado: docentes de educación primaria y secundaria.
%
% Está pensado para crear \emph{«apuntes»}, \emph{«hojas de ejercicios»},
% \emph{«clases escritas»} enfocadas en el estudiante. Solo un docente sabe cual es
% el nivel de aprendizaje del estudiante. El paquete intenta sobrescribir
% muchas de las reglas dictadas por la heurística de \mynvmath{NVDA/MathCAT}
% teniendo esto presente.
%
% \smallskip
%
% \begin{important}*
%  El paquete \mypkg*[no-index]{spintent} usa y abusa del atributo |:intent| de \mynvmath{MathML}
%  por medio del comando |\MathMLintent| definido en \mypkg{latex-lab-mathintent}\cite{latexlab},
%  un encapsulado de |\luamml_attribute:een| proveída por \mypkg{luamml}\cite{luamml}. !Gracias Marcel¡
% \end{important}
%
% \subsection{Carga del paquete}
%
% El paquete \mypkg*[no-index]{spintent} trabaja solo con \hologo{LuaLaTeX} y
% \emph{tagged} \texttt{PDF} activo usando |tagging=on| dentro del comando
% |\DocumentMetadata| al inicio del archivo.
% Las llaves (\mymeta{keys})  aceptadas por |\DocumentMetadata| y en
% general para el código de \emph{tagged} \texttt{PDF} están documentadas
% en los paquetes y \mypkg{tagpdf}\cite{tagpdf} \mypkg{documentmetadata-support}\cite{latexlab},
% es deber del usuario revisar la documentación asociada a ellos.
%
% \iffalse
%<*example>
% \fi
\begin{examplecode}[frame=single]
\DocumentMetadata
  {
    lang=es-CL,
    pdfversion=2.0,
    pdfstandard=ua-2,
    tagging=on,
    tagging-setup={math/setup={mathml-SE}},
  }
\documentclass{article}
\usepackage[spanish]{babel}
\usepackage{spintent}
\end{examplecode}
% \iffalse
%</example>
% \fi
%
% \subsection{El paquete \texttt{babel}}
%
% El paquete \mypkg{babel-spanish}\cite{babel-spanish} está marcado como
% «parcialmente compatible» para \emph{tagged} \texttt{PDF}, pero, es
% absolutamente necesario para la creación de documentos en español.
% En caso de encontrar problemas de compatibilidad con otros paquetes
% se puede desactivar usando:
%
% \iffalse
%<*example>
% \fi
\begin{examplecode}[frame=single]
\usepackage[spanish,shorthands=off]{babel}
\end{examplecode}
% \iffalse
%</example>
% \fi
%
% Muchas de las «shorthands» proveidas por \mypkg{babel-spanish} hoy no
% son necesarias, más allá de |\sptext| para las abreviaturas y ordinales
% el resto son solo atajos de teclas para escritura que hoy en día no se ocupan.
% El paquete \mypkg*[no-index]{spintent} provee el comando |\spshort| (\S\ref{sec:spshort})
% que cubre estos casos, de esta manera podemos cargar \mypkg{babel}\cite{babel}
% con todas sus funcionalidades sin problemas.
%
% \subsection{Fuentes tipográficas}
%
% Como \mypkg*[no-index]{spintent} es un paquete \hologo{LuaLaTeX} debemos
% distinguir entre las fuentes tipográficas \texttt{OpenType/TrueType}
% versus las fuentes de \texttt{8bit} estándar usadas por
% \hologo{LaTeX} para el \emph{modo texto} y el \emph{modo matemático}.
% Para evitar problemas con \emph{tagged} \texttt{PDF} preferiremos no
% combinar fuentes \texttt{OpenType/TrueType} con fuentes \texttt{8bit} y
% trabajaremos solo con fuentes \texttt{OpenType/TrueType} que tengan
% soporte para matemáticas.
%
% \smallskip
%
% La carga de las fuentes tipográficas para \hologo{LuaLaTeX} se
% maneja por medio del paquete \mypkg{fontspec}\cite{fontspec}, si no se espcifica
% ninguna se usarán las fuentes \texttt{Latin Modern} para el \emph{modo texto} y
% \texttt{Latin Modern Math} para el \emph{modo matemático}.
%
% \smallskip
%
% El tipo de fuente tipográfica escogida es preferencia de cada usuario, en lo
% personal utilizo las fuente \texttt{Libertinus} para el \emph{modo
% texto} por medio del paquete
% \mypkg{libertinus-otf}\cite{libertinus-otf} y la fuente
% \texttt{Erewhon-Math} para el modo matemático cargada por el paquete
% \mypkg{fourier-otf}\cite{erewhon-math}. Adicionalmente cargo la fuente \texttt{Source Code Pro}
% para el texto tipo \texttt{«máquina de escribir»} por medio del paquete
% \mypkg{sourcecodepro}\cite{sourcecodepro}.
%
% \iffalse
%<*example>
% \fi
\begin{examplecode}[frame=single]
% \DocumentMetadata{...}
\documentclass{article}
\usepackage[spanish]{babel}
\usepackage[osf,nomath,mono=false]{libertinus-otf}
\usepackage[osf,semibold]{sourcecodepro}
\usepackage[no-text]{fourier-otf}
\usepackage{spintent}
\end{examplecode}
% \iffalse
%</example>
% \fi
%
% \subsection{El paquete \texttt{hyperref}}
%
% El paquete \mypkg{hyperref}\cite{hyperref} es \emph{«casi obligatorio»}
% cuando se crean documentos \emph{tagged} \texttt{PDF}, se encarga de las
% referencias cruzadas |\label| y |\ref|, las notas al pie |\footnote|,
% los hipervínculos |\href| y muchos otros elementos como marcadores, índices
% y metadatos.
%
% \iffalse
%<*example>
% \fi
\begin{examplecode}[frame=single]
\DocumentMetadata
  {
    lang=es-CL,
    pdfversion=2.0,
    pdfstandard=ua-2,
    tagging=on,
    tagging-setup={math/setup={mathml-SE}},
  }
\documentclass{article}
\usepackage[spanish]{babel}
\usepackage[osf,nomath,mono=false]{libertinus-otf}
\usepackage[osf,semibold]{sourcecodepro}
\usepackage[no-text]{fourier-otf}
\setmathfont{Erewhon-Math}[range={cal,bfcal},StylisticSet=1]
\usepackage{spintent, hyperref}
\hypersetup
  {
    colorlinks         = true,
    pdfstartview       = FitH,
    pdfdisplaydoctitle = true,
    pdftitle           = {Test para el paquete spintent},
    pdfinfo            =
      {
        Title   = {Test},
        Subject = {spintent},
      },
  }
\end{examplecode}
% \iffalse
%</example>
% \fi
%
% \section{El comando \cs[no-index]{setspintent}}\label{sec:setspintent}
%
% \vspace*{-\baselineskip}
%
% \begin{function}{\setspintent}
%   \begin{syntax}
%     \ics*{setspintent}\mymarg{comando}\mymarg{key-uno \textnormal{\textcolor{gray}{=}} valor\textnormal{\textcolor{gray}{,}} key-dos \textnormal{\textcolor{gray}{=}} valor\textnormal{\textcolor{gray}{,}} key-tres \textnormal{\textcolor{gray}{=}} valor\textnormal{\textcolor{gray}{,}} \dots}
%   \end{syntax}
%   El comando \ics*{setspintent} establece las llaves (\mymeta{keys}) de
%   forma global para los comandos provistos por \mypkg*[no-index]{spintent}. Puede
%   utilizarse tanto en el preámbulo como en el cuerpo del documento
%   tantas veces como se desee.
%
% \smallskip
%
%   Las llaves (\mymeta{keys}) definidas en el \emph{argumento opcional}
%   \myoarg{key \textnormal{\textcolor{gray}{=}} val} de los comandos
%   tienen prioridad sobre las establecidas por este, anulando las
%   configuraciones globales de |\setspintent|.
%
% \smallskip
%
%   Muchas de las utilidades que provee el paquete \mypkg*[no-index]{spintent} implementan
%   el sistema \myoarg{key \textnormal{\textcolor{gray}{=}} val} y están pensadas
%   para ser utilizadas en la generación de \emph{«hojas de ejercicios»} o \emph{«clases
%   escritas»}. Se recomienda usar |\setspintent| y establecerlas de manera
%   global para cada documento, para mantener la coherencia en la verbalización
%   de \mynvmath{NVDA/MathCAT} a lo largo del mismo.
%
%  \smallskip
%
%   El sistema \myoarg{key \textnormal{\textcolor{gray}{=}} val} utilizado
%   por \mypkg*[no-index]{spintent} se implementa usando el módulo \mypkg{l3keys}
%   de \mypkg{expl3}\cite{expl3}. Por diseño, las llaves que reciben un
%   \emph{valor} después del signo `|=|' NO aceptan \emph{valores vacíos}
%   y es obligatorio pasar uno valido; las llaves documentadas como \mymeta{valor prohibido}
%   no reciben \emph{ningun valor}.
%   En caso de omitir un \emph{valor} valido o pasar un \emph{valor vacío}
%   a una llave que acepte \emph{valor} o pasar un \emph{valor} a una llave
%   del tipo \mymeta{valor prohibido} se devolverá un \enquote{error}
%   y se detendrá la compilación del documento.
% \end{function}
%
% \section{The big four}
%
% El paquete \mypkg*[no-index]{spintent} provee cuatro grandes comandos
% \emph{(«the big four\footnote{Los \emph{«cuatro grandes»} del
% Thrash Metal: Metallica, Megadeth, Slayer y Anthrax, que me acompañan
% mientras desarrollo \mypkg*[no-index]{spintent}.}»)}:
% |\spnum| para el manejo de números, |\spunit| para el manejo
% de unidades, |\spmoney| para el manejo de dinero y |\spshort|
% para el manejo de abreviaturas, acrónimos y números ordinales.
%
% \smallskip
%
% \begin{important}*
%  Tambien se incluyen las utilidades |\spnewunit| y |\spunitalias| asociadas
%  a al comando |\spunit|.
% \end{important}
%
% \subsection{El comando \cs[no-index]{spnum}}\label{sec:spnum}
%
% \vspace*{-\baselineskip}
%
% \begin{function}{\spnum}
%   \begin{minipage}[t]{0.35\linewidth}
%     \cmdexamp{spnum}{número}
%     \cmdexamp{spnum}[key \textnormal{\textcolor{gray}{=}} val]{\textnormal{\textcolor{red}{{+}}}número}
%     \cmdexamp{spnum}[key \textnormal{\textcolor{gray}{=}} val]{\textnormal{\textcolor{red}{{-}}}número}
%     \cmdexamp{spnum}[key \textnormal{\textcolor{gray}{=}} val]{número\textnormal{\textcolor{red}{{,}}}número}
%     \cmdexamp{spnum}[key \textnormal{\textcolor{gray}{=}} val]{número\textnormal{\textcolor{red}{{.}}}número}
%   \end{minipage}\hfill
%   \begin{minipage}[t]{0.55\linewidth}
%     \cmdexamp{spnum}[key \textnormal{\textcolor{gray}{=}} val]{número\textnormal{\textcolor{red}{{,}}}número\textnormal{\textcolor{red}{{;}}}número}
%     \cmdexamp{spnum}[key \textnormal{\textcolor{gray}{=}} val]{número\textnormal{\textcolor{red}{{.}}}número\textnormal{\textcolor{red}{{;}}}número}
%     \cmdexamp{spnum}[key \textnormal{\textcolor{gray}{=}} val]{número\textnormal{\textcolor{red}{{;}}}número}
%     \cmdexamp{spnum}[key \textnormal{\textcolor{gray}{=}} val]{número \textnormal{\textcolor{red}{{e}}} exponente}
%     \cmdexamp{spnum}[key \textnormal{\textcolor{gray}{=}} val]{número \textnormal{\textcolor{red}{{E}}} exponente}
%   \end{minipage}
%
%  \smallskip
%
%   El comando \ics*{spnum} trabaja en \emph{modo matemático} y recibe
%   el \emph{argumento obligatorio} \mymarg{número}, donde la \emph{«parte entera»}
%   puede iniciar con los signos `$\mathcolor{red}{+}$' o `$\mathcolor{red}{-}$',
%   la \emph{«parte decimal finita»} \textbf{debe} iniciar con una \emph{coma}
%   `\textcolor{red}{\texttt{,}}' o con un \emph{punto} `\textcolor{red}{\texttt{.}}',
%   y la \emph{«parte decimal infinita (periódica)»} \textbf{debe} iniciar
%   con un \emph{punto y coma} `\textcolor{red}{\texttt{;}}'.
%
%  \smallskip
%
%   Si no se proporciona el \emph{argumento opcional} \myoarg{key \textnormal{\textcolor{gray}{=}} val}
%   y se ejecuta |$\spnum{+23}$| se imprimirá `$+23$' en la salida \texttt{PDF}
%   y \mynvmath{NVDA/MathCAT} verbalizará \emph{«más veintitrés»}, si se
%   ejecuta |$\spnum{3,14}$| se imprimirá `$3{,}14$' en la salida \texttt{PDF}
%   y \mynvmath{NVDA/MathCAT} verbalizará \emph{«tres coma catorce»}, si
%   se ejecuta |$\spnum{3,1;4}$| se imprimirá `$3{,}1\overline{4}$' en la
%   salida \texttt{PDF} y \mynvmath{NVDA/MathCAT} verbalizará \emph{«tres coma
%   uno con cuatro periódico»}.
% \end{function}
%
% \smallskip
%
% Además, el \emph{argumento obligatorio} \mymarg{número} admite al final
% de este un sufijo de \emph{«notación científica»}. Este se indica mediante las
% letras `\textcolor{red}{\texttt{e}}' o `\textcolor{red}{\texttt{E}}',
% acompañadas de los signos opcionales `$\mathcolor{red}{+}$' o
% `$\mathcolor{red}{-}$' seguido por los dígitos del exponente.
% Por ejemplo |$\spnum{3,14e3}$| imprimirá `$3{,}14\cdot 10^{3}$' en la salida \texttt{PDF}
% y \mynvmath{NVDA/MathCAT} verbalizará \emph{«tres coma catorce por diez al cubo»}.
%
% \smallskip
%
% \begin{important}*
%  La \emph{coma} `\textcolor{red}{\texttt{,}}' o el \emph{punto} `\textcolor{red}{\texttt{.}}'
%  son \emph{separadores decimales} y NO de millares, por lo tanto solo
%  pueden aparecer una «única vez» dentro del \mymarg{argumento}.
%  El \mymarg{argumento} puede contener espacios matemáticos validos, pero,
%  NO puede contener comandos con o sin argumentos como |\textstyle|, |\mathrm{#1}| o
%  |\mathbf{#1}| por ejemplo. Cualquier entrada fuera de esta sintaxis
%  devolverá un \enquote{error} deteniendo la compilación del documento.
% \end{important}
%
% \subsubsection{Llaves para \cs[no-index]{spnum}}\label{subsec:key-spnum}
%
% \keyexampcmd{cifras-sep}{true \textnormal{\textcolor{lightgray}{\textbar}} false}{true}[spnum]
% \smallskip
% Establece el \emph{espaciado} utilizado para agrupar los dígitos en
% bloques de tres, tanto en la parte entera como en la parte decimal finita.
% Si se establece en |true| se siguen las normas de la \textsc{rae} y se
% utiliza un \emph{espacio fino} `\textcolor{blue!75}{\textbackslash}|,|'
% entre las cifras; si se establece en |false|, se respetan los espacios
% matemáticos introducidos en el argumento. Por ejemplo |$\spnum{100000}$|
% imprimirá `$100\,000$' en la salida \texttt{PDF} y |$\spnum[cifras-sep=false]{10\,00\,00}$|
% imprimirá `$10\,00\,00$' en la salida \texttt{PDF}.
%
% \smallskip
%
% \keyexampcmd{read-sign}{terse \textnormal{\textcolor{lightgray}{\textbar}} clear}{terse}[spnum]
% \smallskip
% Establece la \emph{forma} en la que \mynvmath{NVDA/MathCAT} verbalizará los signos
% `$\mathcolor{red}{+}$' o `$\mathcolor{red}{-}$'.
% Si se establece en |terse|, \mynvmath{NVDA/MathCAT} verbalizará \emph{«más»} o
% \emph{«menos»} \textbf{antes} de leer el número; si se establece en |clear|, \mynvmath{NVDA/MathCAT}
% verbalizará \emph{«positivo»} o \emph{«negativo»} \textbf{después} de leer el número.
% Por ejemplo |$\spnum{-23}$| imprimirá `$-23$' en la salida \texttt{PDF} y \mynvmath{NVDA/MathCAT}
% verbalizará \emph{«menos veintitrés»} y |$\spnum[read-sign=clear]{-23}$|
% imprimirá `$-23$' en la salida \texttt{PDF} y \mynvmath{NVDA/MathCAT}
% verbalizará \emph{«veintitrés negativo»}.
%
% \smallskip
%
% \keyexampcmd{read-period-sep}{coma \textnormal{\textcolor{lightgray}{\textbar}} punto}{coma}[spnum]
% Establece la \emph{forma} en la que \mynvmath{NVDA/MathCAT} verbalizará e imprimirá
% el signo de separación `\textcolor{red}{\texttt{;}}' de la \emph{«parte
% decimal infinita (periódica)»} cuando no se tiene la \emph{«parte
% decimal finita»}. Si se establece en |coma|, se imprimirá `\textcolor{red}{\texttt{,}}'
% en la salida \texttt{PDF} y \mynvmath{NVDA/MathCAT} verbalizará \emph{«coma»};
% si se establece en |punto|, se imprimirá `\textcolor{red}{\texttt{.}}'
% en la salida \texttt{PDF} y \mynvmath{NVDA/MathCAT} verbalizará \emph{«punto»}.
% Por ejemplo |$\spnum{3;4}$| imprimirá `$3{,}\overline{4}$' en la salida \texttt{PDF}
% y \mynvmath{NVDA/MathCAT} verbalizará \emph{«tres coma cuatro periódico»}
% y |$\spnum[read-period-sep=punto]{3;4}$| imprimirá `$3{.}\overline{4}$'
% en la salida \texttt{PDF}  y \mynvmath{NVDA/MathCAT} verbalizará
% \emph{«tres punto cuatro periódico»}
%
% \smallskip
%
% \begin{important}*
%  Esta llave es necesaria SOLO cuando se utilizan \emph{«decimales periódicos puros»},
%  es decir, cuando no se ha incluido en el \mymarg{argumento} una \emph{coma}
%  `\textcolor{red}{\texttt{,}}' o un \emph{punto} `\textcolor{red}{\texttt{.}}'
%  antes del signo de separación `\textcolor{red}{\texttt{;}}' periódica.
% \end{important}
%
% \smallskip
%
% \keyexampcmd{notation-sym}{times \textnormal{\textcolor{lightgray}{\textbar}} cdot}{cdot}[spnum]
% \smallskip
% Establece el \emph{símbolo} que se utilizará para imprimir
% la \emph{«notación científica»} al usar `\textcolor{red}{\texttt{e}}' o
% `\textcolor{red}{\texttt{E}}'. Si se establece en |times| imprimirá
% `$\mathcolor{red}{\times}$' en la salida \texttt{PDF} y \mynvmath{NVDA/MathCAT}
% verbalizará \emph{«por»}; si se establece en |cdot| imprimirá
% `$\mathcolor{red}{\cdot}$' en la salida \texttt{PDF} y \mynvmath{NVDA/MathCAT}
% verbalizará \emph{«por»}. Por ejemplo |$\spnum{3,14e3}$| imprimirá `$3{,}14\cdot 10^{3}$'
% en la salida \texttt{PDF} y \mynvmath{NVDA/MathCAT} verbalizará
% \emph{«tres coma catorce por diez al cubo»} y |$\spnum[notation-sym=times]{3,14e3}$|
% imprimirá `$3{,}14\times 10^{3}$' en la salida \texttt{PDF} y \mynvmath{NVDA/MathCAT}
% verbalizará \emph{«tres coma catorce por diez al cubo»}.
%
% \smallskip
%
% \begin{important}*
%  Decir \emph{«notación científica»} aquí es un poco exagerado, en realidad
%  es \emph{«multiplicación por potenicas de $10$»}, no se verfica si la
%  entrada cumple o no con la condición de estár escrita en \emph{«notación científica»}.
% \end{important}
%
% \subsection{El comando \cs[no-index]{spunit}}\label{sec:spunit}
%
% \vspace*{-\baselineskip}
%
% \begin{function}{\spunit}
%   \begin{minipage}[t]{0.35\linewidth}
%     \cmdexamp{spunit}{unidad}
%     \cmdexamp{spunit}[key \textnormal{\textcolor{gray}{=}} val]{unidad}
%     \cmdexamp{spunit}[key \textnormal{\textcolor{gray}{=}} val]{número unidad}
%   \end{minipage}\hfill
%   \begin{minipage}[t]{0.55\linewidth}
%     \cmdexamp{spunit}[key \textnormal{\textcolor{gray}{=}} val]{número unidad\textnormal{\textcolor{red}{{\textasteriskcentered}}}unidad}
%     \cmdexamp{spunit}[key \textnormal{\textcolor{gray}{=}} val]{número unidad\textnormal{\textcolor{red}{{\textasteriskcentered}}}unidad\textnormal{\textcolor{red}{{/}}}unidad}
%     \cmdexamp{spunit}[key \textnormal{\textcolor{gray}{=}} val]{número unidad\textnormal{\textcolor{red}{{\textasteriskcentered}}}unidad\textnormal{\textcolor{red}{{/}}}unidad\textnormal{\textcolor{red}{{\textasteriskcentered}}}unidad}
%   \end{minipage}
%
% \smallskip
%
%   El comando |\spunit| trabaja en \emph{modo matemático} y recibe
%   el \emph{argumento obligatorio} \mymarg{unidad}, el cual puede contener
%   una \emph{«parte numérica»} opcional la cual se procesa
%   internamente por el comando |\spnum| al inicio del \mymarg{argumento} seguido
%   de una \mymarg{unidad} o un \mymarg{alias} valido  descritos en~\S\ref{subsec:arg-spunit}
%   o una \mymarg{unidad} creada por el comando |\spnewunit| (\S\ref{subsec:spnewunit})
%   o un \mymarg{alias} creado por el comando |\spunitalias| (\S\ref{subsec:spunitalias}).
%
%   \smallskip
%
%   Si no se proporciona el \emph{argumento opcional} \myoarg{key \textnormal{\textcolor{gray}{=}} val}
%   y se ejecuta |$\spunit{23 m}$| se imprimirá `$23\,\mathrm{m}$' en la
%   salida \texttt{PDF} y \mynvmath{NVDA/MathCAT} verbalizará \emph{«veintitrés metros»}.
%
%   \smallskip
%
%   La multiplicación entre las \mymarg{unidades} se realiza mediante un
%   \emph{asterisco} `$\mathcolor{red}{\ast}$' y las potencias de estas
%   se indican con un \emph{circunflejo seguido del exponente entre llaves}
%   `\textcolor{red}{\texttt{\^{}}}\mymarg{exponente}'. Por ejemplo si
%   se ejecuta |$\spunit{23 m*s^{2}}$| se imprimirá `$23\,\mathrm{m}\,\mathrm{s}^{2}$'
%   en la salida \texttt{PDF} y \mynvmath{NVDA/MathCAT} verbalizará
%   \emph{«veintitrés metros segundos al cuadrado»}.
%
%   \smallskip
%
%   Se pueden usar \emph{fracciones lineales} con las \mymarg{unidades}
%   utilizando una \emph{«única barra»} `$\mathcolor{red}{/}$' para la
%   separación entre las \mymarg{unidades} del numerador y el denominador,
%   teniendo presente que el denominador NO puede contener números de ningún
%   tipo. Por ejemplo si se ejecuta |$\spunit{23 m/s^{2}}$| se imprimirá
%   `$23\,\mathrm{m}/\mathrm{s}^{2}$' en la salida \texttt{PDF} y
%   \mynvmath{NVDA/MathCAT} verbalizará \emph{«veintitrés metros por segundos al cuadrado»}.
% \end{function}
%
% \smallskip
%
% \begin{important}*
%  Solo imprimiremos \emph{fracciones lineales}, si se requiere otra forma
%  fraccionaria para explicaciones y otros usos se debe usar el comando
%  por separado. El \mymarg{argumento} puede contener espacios matemáticos
%  validos, pero, NO puede contener comandos con o sin argumentos como
%  |\textstyle|, |\mathrm{#1}| o |\mathbf{#1}| por ejemplo, NO puede contener
%  más de una \emph{«única barra»} `$\mathcolor{red}{/}$' y NO puede contener
%  números en el denominador. Cualquier entrada fuera de esta sintaxis
%  devolverá un \enquote{error} deteniendo la compilación del documento.
% \end{important}
%
% \subsubsection{Llaves para \cs[no-index]{spunit}}
%
% El comando |\spunit| posee las \mymeta{meta-keys} \mykey[spunit]{cifras-sep},
% \mykey[spunit]{read-period-sep} y \mykey[spunit]{notation-sym} pasadas
% al comando |\spnum| (\S\ref{subsec:key-spnum}).
%
% \smallskip
%
% \keyexampcmd{print-cdot}{true \textnormal{\textcolor{lightgray}{\textbar}} false}{false}[spunit]
% \smallskip
% Establece si se \emph{imprime} |\cdot| `$\mathcolor{red}{\cdot}$' o un \emph{espacio fino}
% `\textcolor{blue!75}{\textbackslash}|,|' para la multiplicación de \mymarg{unidades}.
% Si se establece en |true|, se imprimirá `$\mathcolor{red}{\cdot}$' entre las \mymarg{unidades}
% en la salida \texttt{PDF}; si se establece en |false| se imprimirá un \emph{espacio fino}
% `\textcolor{blue!75}{\textbackslash}|,|' entre las \mymarg{unidades} en
% la salida \texttt{PDF}.
%
% \keyexampcmd{read-cdot}{true \textnormal{\textcolor{lightgray}{\textbar}} false}{false}[spunit]
% \smallskip
% Establece la \emph{forma} en la que \mynvmath{NVDA/MathCAT} verbalizará la
% multiplicación de \mymarg{unidades}. Si se establece en |true|,
% \mynvmath{NVDA/MathCAT} verbalizará «por» entre las \mymarg{unidades}; si se
% establece en |false|, \mynvmath{NVDA/MathCAT} no verbalizará nada entre las
% \mymarg{unidades}.
%
% \keyexampcmd{read-slash}{por \textnormal{\textcolor{lightgray}{\textbar}} partido-por}{por}[spunit]
% \smallskip
% Establece la \emph{forma} en la que \mynvmath{NVDA/MathCAT} verbalizará la
% \emph{«barra»} `$\mathcolor{red}{/}$' de separación entre el
% numerador y el denominador. Si se establece en |por|, \mynvmath{NVDA/MathCAT}
% verbalizará \emph{«por»}; si se establece en |partido-por|, \mynvmath{NVDA/MathCAT}
% verbalizará \emph{«partido por»}.
%
% \subsubsection{Argumentos para \cs[no-index]{spunit}}\label{subsec:arg-spunit}
%
% Resumen de \mymarg{unidades} y \mymarg{alias} soportados por \ics*{spunit}.
%
% \smallskip
%
% \begin{longtable}{@{}cllcl@{}}\small
%   \label{spunit:valid-units} \\
%   \toprule
%   \textbf{Entrada} & \textbf{Alias soportados} & \textbf{Nombre unidad} & \textbf{Salida} & \textbf{NVDA/MathCAT} \\
%   \midrule
%   \endfirsthead
%   %
%   \multicolumn{5}{c}{{\bfseries Continuación de la tabla \ref{spunit:valid-units}}} \\
%   \toprule
%   \textbf{Entrada} & \textbf{Alias soportados} & \textbf{Nombre unidad} & \textbf{Salida} & \textbf{MathCAT} \\
%   \midrule
%   \endhead
%   %
%   \midrule
%   \multicolumn{5}{r}{\textit{Continúa en la siguiente página...}} \\
%   \endfoot
%   %
%   \bottomrule
%   \endlastfoot
%   %
%   m   & metro, metros             & Metro               & $\mathrm{m}$   & \emph{«metro»}               \\
%   g   & gramo, gramos             & Gramo               & $\mathrm{g}$   & \emph{«gramo»}               \\
%   s   & segundo, segundos         & Segundo (Tiempo)    & $\mathrm{s}$   & \emph{«segundo (Tiempo)»}    \\
%   l   & litro, l-litro            & Litro (minúscula)   & $\mathrm{l}$   & \emph{«litro (minúscula)»}   \\
%   h   & hora, horas               & Hora                & $\mathrm{h}$   & \emph{«hora»}                \\
%   d   & dia, dias                 & Día                 & $\mathrm{d}$   & \emph{«día»}                 \\
%   A   & amperio, ampere           & Amperio             & $\mathrm{A}$   & \emph{«amperio»}             \\
%   V   & voltio, volt              & Voltio              & $\mathrm{V}$   & \emph{«voltio»}              \\
%   W   & vatio, watt               & Vatio               & $\mathrm{W}$   & \emph{«vatio»}               \\
%   J   & julio, joule              & Julio               & $\mathrm{J}$   & \emph{«julio»}               \\
%   N   & newton                    & Newton              & $\mathrm{N}$   & \emph{«newton»}              \\
%   Pa  & pascal                    & Pascal              & $\mathrm{Pa}$  & \emph{«pascal»}              \\
%   Hz  & hercio, hertz             & Hercio              & $\mathrm{Hz}$  & \emph{«hercio»}              \\
%   Ω   & ohmio, ohm, Omega         & Ohmio               & $\mathrm{Ω}$   & \emph{«Ohmio»}               \\
%   F   & faradio                   & Faradio             & $\mathrm{F}$   & \emph{«Faradio»}             \\
%   C   & culombio                  & Culombio            & $\mathrm{C}$   & \emph{«Culombio»}            \\
%   K   & kelvin, Kelvin            & Kelvin              & $\mathrm{K}$   & \emph{«Kelvin»}              \\
%   °C  & celsius, degC, ℃          & Grados Celsius      & $\mathrm{°C}$  & \emph{«Grados Celsius»}      \\
%   °F  & fahrenheit, degF, ℉       & Grados Fahrenheit   & $\mathrm{°F}$  & \emph{«Grados Fahrenheit»}   \\
%   cal & caloria, calorias         & Caloría             & $\mathrm{cal}$ & \emph{«Caloría»}             \\
%   b   & bit                       & Bit                 & $\mathrm{b}$   & \emph{«Bit»}                 \\
%   B   & byte, bytes               & Byte                & $\mathrm{B}$   & \emph{«Byte»}                \\
%   in  & pulgada, pulgadas         & Pulgada             & $\mathrm{in}$  & \emph{«Pulgada»}             \\
%   ft  & pie, pies                 & Pie                 & $\mathrm{ft}$  & \emph{«Pie»}                 \\
%   yd  & yarda, yardas             & Yarda               & $\mathrm{yd}$  & \emph{«Yarda»}               \\
%   mi  & milla, millas             & Milla               & $\mathrm{mi}$  & \emph{«Milla»}               \\
%   lb  & libra, libras             & Libra               & $\mathrm{lb}$  & \emph{«Libra»}               \\
%   oz  & onza, onzas               & Onza                & $\mathrm{oz}$  & \emph{«Onza»}               \\
%   gal & galon, galones            & Galón               & $\mathrm{gal}$ & \emph{«Galón»}               \\
%   Å   & angstrom                  & Angstrom            & $\mathrm{Å}$   & \emph{«Angstrom»}            \\
%   µ   & micro                     & Micro (Prefijo)     & $\mathrm{µ}$   & \emph{«Micro (Prefijo)»}     \\
%   ℧   & mho                       & Mho / Siemens       & $\mathrm{℧}$   & \emph{«Mho / Siemens»}       \\
%   ′   & minmin                    & Minuto (Arco)       & $\mathrm{′}$   & \emph{«Minuto (Arco)»}       \\
%   ″   & segseg                    & Segundo (Arco)      & $\mathrm{″}$   & \emph{«Segundo (Arco)»}      \\
%   u   & u-masa                    & Unidad masa atómica & $\mathrm{u}$   & \emph{«Unidad masa atómica»} \\
%   \midrule
%   \multicolumn{5}{c}{\textit{Unidades directas (sin alias):}} \\
%   \midrule
%   L   & -                         & Litro (Mayúscula)   & $\mathrm{L}$   & \emph{«Litro (Mayúscula)»}   \\
%   ℓ   & -                         & Litro (Cursiva)     & $\mathrm{ℓ}$   & \emph{«Litro (Cursiva)»}     \\
%   w   & -                         & Vatio (Minúscula)   & $\mathrm{w}$   & \emph{«Vatio (Minúscula)»}   \\
%   °K  & -                         & Grado Kelvin        & $\mathrm{°K}$  & \emph{«Grado Kelvin»}        \\
%   Bq  & -                         & Becquerel           & $\mathrm{Bq}$  & \emph{«Becquerel»}           \\
%   Da  & -                         & Dalton              & $\mathrm{Da}$  & \emph{«Dalton»}              \\
%   Gy  & -                         & Gray                & $\mathrm{Gy}$  & \emph{«Gray»}                \\
%   S   & -                         & Siemens             & $\mathrm{S}$   & \emph{«Siemens»}             \\
%   Sv  & -                         & Sievert             & $\mathrm{Sv}$  & \emph{«Sievert»}             \\
%   T   & -                         & Tesla               & $\mathrm{T}$   & \emph{«Tesla»}               \\
%   Wb  & -                         & Weber               & $\mathrm{Wb}$  & \emph{«Weber»}               \\
%   atm & -                         & Atmósfera           & $\mathrm{atm}$ & \emph{«Atmósfera»}           \\
%   au  & -                         & Unidad astronómica  & $\mathrm{au}$  & \emph{«Unidad astronómica»}  \\
%   bar & -                         & Bar                 & $\mathrm{bar}$ & \emph{«Bar»}                 \\
%   cd  & -                         & Candela             & $\mathrm{cd}$  & \emph{«Candela»}             \\
%   dB  & -                         & Decibelio           & $\mathrm{dB}$  & \emph{«Decibelio»}           \\
%   dyn & -                         & Dina                & $\mathrm{dyn}$ & \emph{«Dina»}                \\
%   eV  & -                         & Electronvoltio      & $\mathrm{eV}$  & \emph{«Electronvoltio»}      \\
%   erg & -                         & Ergio               & $\mathrm{erg}$ & \emph{«Ergio»}               \\
%   ha  & -                         & Hectárea            & $\mathrm{ha}$  & \emph{«Hectárea»}            \\
%   kat & -                         & Katal               & $\mathrm{kat}$ & \emph{«Katal»}               \\
%   kg  & -                         & Kilogramo           & $\mathrm{kg}$  & \emph{«Kilogramo»}           \\
%   km  & -                         & Kilómetro           & $\mathrm{km}$  & \emph{«Kilómetro»}           \\
%   lm  & -                         & Lumen               & $\mathrm{lm}$  & \emph{«Lumen»}               \\
%   lx  & -                         & Lux                 & $\mathrm{lx}$  & \emph{«Lux»}                 \\
%   cm  & -                         & Centímetro          & $\mathrm{cm}$  & \emph{«Centímetro»}          \\
%   min & -                         & Minuto (Tiempo)     & $\mathrm{min}$ & \emph{«Minuto (Tiempo)»}     \\
%   mm  & -                         & Milímetro           & $\mathrm{mm}$  & \emph{«Milímetro»}           \\
%   mol & -                         & Mol                 & $\mathrm{mol}$ & \emph{«Mol»}                 \\
%   pc  & -                         & Pársec              & $\mathrm{pc}$  & \emph{«Pársec»}              \\
%   rad & -                         & Radián              & $\mathrm{rad}$ & \emph{«Radián»}              \\
%   rpm & -                         & Rev. por minuto     & $\mathrm{rpm}$ & \emph{«Rev. por minuto»}     \\
%   sr  & -                         & Estereorradián      & $\mathrm{sr}$  & \emph{«Estereorradián»}      \\
%   t   & -                         & Tonelada            & $\mathrm{t}$   & \emph{«Tonelada»}            \\
%   a   & -                         & Año / Área (are)    & $\mathrm{a}$   & \emph{«Año / Área (are)»}    \\
%   y   & -                         & Año (Year)          & $\mathrm{y}$   & \emph{«Año (Year)»}          \\
%   °   & -                         & Grado (Arco)        & $\mathrm{°}$   & \emph{«Grado (Arco)»}        \\
% \end{longtable}
%
% \subsubsection{El comando \cs[no-index]{spnewunit}}\label{subsec:spnewunit}
%
% \vspace*{-0.5\baselineskip}
%
% \begin{function}{\spnewunit}
%   \begin{syntax}
%     \ics*{spnewunit}\mymarg{nueva unidad}\mymarg{verbalización \textnormal{NVDA}}
%   \end{syntax}
%   El comando \ics*{spnewunit} trabaja en \emph{modo texto} y recibe dos
%   \emph{argumentos obligatorios}: \mymarg{nueva unidad} que establece
%   lo que se imprimirá en la salida \texttt{PDF} y \mymarg{verbalización \textnormal{NVDA}}
%   que establece como \mynvmath{NVDA/MathCAT} \emph{verbalizará} la \mymarg{nueva unidad}.
% \end{function}
%
% \smallskip
%
% La \mymarg{nueva unidad} no debe estar registrada como \mymarg{unidad}
% o \mymarg{alias} para |\spunit| descritos en \S\ref{subsec:arg-spunit}.
% Si la \mymarg{nueva unidad} ya está registrada se devolverá un \enquote{error}
% y se detendrá la compilación del documento.
%
% \smallskip
%
% La declaración de la \mymarg{nueva unidad} es «global» y se almacena
% en el módulo \myluamod{} como una \mymarg{unidad} valida para pasar como
% \mymarg{argumento} a |\spunit|. Por ejemplo: |\spnewunit{px}{píxeles}|
% creará la \mymarg{nueva unidad} `|px|' y cuando se utilice |\spunit{100 px}|
% se imprimirá `$100\,\mathrm{px}$' en la salida \texttt{PDF} y \mynvmath{NVDA/MathCAT}
% verbalizará \emph{«cien píxeles»}.
%
% \subsubsection{El comando \cs[no-index]{spunitalias}}\label{subsec:spunitalias}
%
% \vspace*{-0.5\baselineskip}
%
% \begin{function}{\spunitalias}
%   \begin{syntax}
%     \ics*{spunitalias}\mymarg{nuevo alias}\mymarg{expresión de unidades}
%   \end{syntax}
%   El comando \ics*{spunitalias} trabaja en \emph{modo texto} y recibe
%   dos \emph{argumentos obligatorios}: \mymarg{nuevo alias} que establece
%   es el nuevo \mymarg{alias} de \mymarg{argumento} para |\spunit| y
%   \mymarg{expresión de unidades} que puede contener multiplicación o
%   división de \mymarg{unidades} ya registradas que se imprimirán en la
%   salida \texttt{PDF}.
% \end{function}
%
% \smallskip
%
% El \mymarg{nuevo alias} no debe estar registrado como \mymarg{unidad}
% o \mymarg{alias} para |\spunit| descritos en \S\ref{subsec:arg-spunit}
% o estar definido como \mymarg{unidad} por medio del comando |\spnewunit|.
% Si el \mymarg{nuevo alias} ya está registrado se devolverá un \enquote{error}
% y se detendrá la compilación del documento.
%
% \smallskip
%
% La declaración de un \mymarg{nuevo alias} es «global» y se almacena
% en el módulo \myluamod{} como una \mymarg{alias} valido para pasar como
% \mymarg{argumento} a |\spunit|. Por ejemplo: |\spunitalias{caudal}{m^{3}/s}|
% creará el \mymarg{alias} `|caudal|' y cuando se utilice |$\spunit{50 caudal}$|
% se imprimirá `$50\,\mathrm{m}^{3} / \mathrm{s}$' en la salida \texttt{PDF}
% y \mynvmath{NVDA/MathCAT} verbalizará \emph{«cincuenta metros al cubo por segundos»}.
%
% \subsection{El comando \cs[no-index]{spmoney}}\label{subsec:spmoney}
%
% \vspace*{-0.5\baselineskip}
%
% \begin{function}{\spmoney}
%   \begin{syntax}
%     \cmdexamp{spmoney}{cifra}
%     \cmdexamp{spmoney}{cifra moneda}
%     \cmdexamp{spmoney}[key \textnormal{\textcolor{gray}{=}} val]{cifra moneda}
%   \end{syntax}
%   El comando \ics*{spmoney} trabaja en \emph{modo matemático} y recibe
%   el \emph{argumento obligatorio} \mymarg{cifra} que contiene la
%   \emph{«cifra numérica»} positiva o negativa la se cual procesa internamente
%   por el comando |\spnum| al inicio del \mymarg{argumento} acompañada
%   de un \emph{«símbolo de moneda»} o \emph{«alias de moneda»} opcional
%   descritos en \S\ref{subsec:arg-spmoney}.
%
%   \smallskip
%
%   Si no se proporciona el \emph{argumento opcional} \myoarg{key \textnormal{\textcolor{gray}{=}} val}
%   y se ejecuta |$\spmoney{500}$| se imprimirá `$\$500$' en la salida
%   \texttt{PDF} y \mynvmath{NVDA/MathCAT} verbalizará \emph{«quinientos pesos»}.
% \end{function}
%
% \smallskip
%
% \begin{important}*
%  El comportamiento de |\spmoney| depende de como se pase el \emph{argumento obligatorio}
%  que tiene prioridad por sobre los valores por defecto. Por ejemplo:
%  |$\spmoney{500 dolares}$| imprimirá `$\$500$' en la salida \texttt{PDF}
%  y \mynvmath{NVDA/MathCAT} verbalizará \emph{«quinientos dolares»}
%  aunque la moneda por defecto sean pesos.
% \end{important}
%
% \smallskip
%
% Si se pasa un \emph{«alias de moneda»} del tipo \texttt{ISO} (código de divisas)
% como |CLP| o |USD| (mayúsculas) se modifica la verbalización en \mynvmath{NVDA/MathCAT}
% agregando el país y se ajusta la salida \texttt{PDF} al estandar \texttt{ISO}. Por ejemplo:
% |$\spmoney{500 CLP}$| imprimirá `$\mathrm{CLP}500$' en la salida \texttt{PDF}
% y \mynvmath{NVDA/MathCAT} verbalizará \emph{«quinientos pesos chilenos»}.
%
% \smallskip
%
% Si se pasa un \emph{«alias de moneda»} del tipo \texttt{ISO} (código de divisas)
% como |clp| o |usd| (minúsculas) se modifica la verbalización en \mynvmath{NVDA/MathCAT}
% agregando el país, pero, NO se modifica la salida \texttt{PDF} imprimiendo
% el \emph{«símbolo de moneda»} correspondiente. Por ejemplo: |$\spmoney{500 clp}$|
% imprimirá `$\$500$' en la salida \texttt{PDF} y \mynvmath{NVDA/MathCAT} verbalizará
% \emph{«quinientos pesos chilenos»}
%
% \smallskip
%
% Pasar de manera directa en el \emph{argumento obligatorio} un \emph{«símbolo de moneda»}
% como `$\mathcolor{red}{\mathdollar}$', `$\mathcolor{red}{\euro}$', `$\mathcolor{red}{\mathsterling}$'
% o un \emph{«alias de moneda»} como |peso|, |euro| o |libra| produce
% la salida \texttt{PDF} y verbalización estándar esperada para \mynvmath{NVDA/MathCAT}.
%
% \smallskip
%
% \begin{important}*
%  La posición del \emph{«símbolo de moneda»} a la izquierda o a la derecha
%  de la \emph{«cifra numérica»} y la verbalización para \mynvmath{NVDA/MathCAT}
%  se manejan desde el módulo Lua \myluamod{} y no se proveen llaves para
%  modificar este comportamiento.
% \end{important}
%
% \smallskip
%
% \begin{important}*
%  El \emph{«símbolo de moneda»} es tomado de la fuente tipográfica definida
%  para el \emph{modo matemático} al momento de ejecutar el comando, se
%  debe tener presente en caso de querer usar un \emph{«símbolo de moneda»}
%  poco habitual que esté soportada por \mypkg*[no-index]{spintent}. El \mymarg{argumento}
%  puede contener espacios matemáticos validos, pero, NO puede contener
%  comandos con o sin argumentos como |\textstyle|, |\mathrm{#1}| o
%  |\mathbf{#1}| por ejemplo. La \emph{coma} `\textcolor{red}{\texttt{,}}'
%  o el \emph{punto} `\textcolor{red}{\texttt{.}}' como \emph{separador decimal} solo
%  pueden aparecer una «única vez». Cualquier entrada fuera de esta sintaxis
%  devolverá un \enquote{error} deteniendo la compilación del documento.
% \end{important}
%
% \subsubsection{Llaves para \cs[no-index]{spmoney}}
%
% El comando |\spmoney| posee la \mymeta{meta-key} \mykey[spmoney]{cifras-sep}
% pasada al comando |\spnum| (\S\ref{subsec:key-spnum}).
%
% \smallskip
%
% \keyexampcmd{currency}{moneda \textnormal{\textcolor{lightgray}{\textbar}} alias}{peso}[spmoney]
% \smallskip
% Establece el \emph{símbolo de moneda} por defecto que se imprimirá en la
% salida \texttt{PDF} y cómo lo verbalizará \mynvmath{NVDA/MathCAT} cuando
% SOLO se pasa una \emph{«cifra numérica»} como \mymarg{argumento} al comando |\spmoney|.
% El \emph{«símbolo de moneda»} o el \emph{«alias de moneda»} pasado a esta
% llave debe ser alguno de los descritos en \S\ref{subsec:arg-spmoney}.
%
% \smallskip
%
% \keyexampcmd{sub-units}{true \textnormal{\textcolor{lightgray}{\textbar}} false}{false}[spmoney]
% \smallskip
% Establece si se usarán \emph{subunidades} de la moneda en curso para la
% verbalización en \mynvmath{NVDA/MathCAT}. La moneda en curso debe aceptar
% \emph{subunidades}. Por ejemplo: |$\spmoney[sub-units=true]{100,5 dolares}$|
% imprimirá `$\$100{,}5$' en la salida \texttt{PDF} y \mynvmath{NVDA/MathCAT} verbalizará
% \emph{«cien dólares con cincuenta centavos»}.
%
% \smallskip
%
% \keyexampcmd{dolar-math}{true \textnormal{\textcolor{lightgray}{\textbar}} false}{true}[spmoney]
% \smallskip
% Establece si el \emph{símbolo} `|$|' se imprimirá en \emph{modo texto}
% o en \emph{modo matemático}. Si se establece en |true|, se usará la fuente tipográfica definida para
% el \emph{modo matemático} en curso para imprimir `|$|' en la salida \texttt{PDF};
% si se establece en |false|, se usará la la fuente tipográfica definida
% para el \emph{modo texto}  para imprimir `|$|' en la salida \texttt{PDF}.
%
% \subsubsection{Monedas y alias para \cs[no-index]{spmoney}}\label{subsec:arg-spmoney}
%
% Resumen de los \mymarg{símbolos de moneda} y \mymarg{alias de moneda} soportados
% como \mymarg{argumento} para |\spmoney|.
%
% \begin{table}[htpb]
%   \centering
%   \label{tab:spmoney-divisas}
%   \begin{tabular}{c p{2.75cm} l c l}
%     \toprule
%     \textbf{Entrada} & \textbf{Alias} & \textbf{Nombre} & \textbf{Salida} & \textbf{NVDA/MathCAT} \\
%     \midrule
%     \texttt{\$}   & \texttt{peso}, \texttt{pesos}                                       & peso          & \$ & \emph{«peso / pesos»} \\
%     \texttt{€}    & \texttt{euro}, \texttt{euros}                                       & euro          & €  & \emph{«euro / euros»} \\
%     \texttt{£}    & \texttt{libra}, \texttt{libras}                                     & libra         & £  & \emph{«libra / libras»} \\
%     \texttt{¥}    & \texttt{yen}, \texttt{yenes}                                        & yen           & ¥  & \emph{«yen / yenes»} \\
%     \texttt{¥}    & \texttt{yuan}, \texttt{yuanes}                                      & yuan          & ¥  & \emph{«yuan / yuanes»} \\
%     \texttt{₩}    & \texttt{won}, \texttt{wons}                                         & won           & \texttt{₩}  & \emph{«won / wons»} \\
%     \texttt{₪}    & \texttt{shekel}, \texttt{shekels}, \texttt{sequel}, \texttt{sequeles} & séquel      & ₪  & \emph{«séquel / séqueles»} \\
%     \texttt{₹}    & \texttt{rupia}, \texttt{rupias}                                     & rupia         & ₹  & \emph{«rupia / rupias»} \\
%     \texttt{₽}    & \texttt{rublo}, \texttt{rublos}                                     & rublo         & \texttt{₽}  & \emph{«rublo / rublos»} \\
%     \texttt{₺}    & \texttt{lira}, \texttt{liras}                                       & lira          & \texttt{₺}  & \emph{«lira / liras»} \\
%     \texttt{₴}    & \texttt{grivna}, \texttt{grivnas}                                   & grivna        & \texttt{₴}  & \emph{«grivna / grivnas»} \\
%     \texttt{¢}    & \texttt{centavo}, \texttt{centavos}                                 & centavo       & ¢  & \emph{«centavo / centavos»} \\
%     \midrule
%     \multicolumn{5}{c}{\textbf{Divisas en formato ISO (Mayúsculas)}} \\
%     \midrule
%     \texttt{CLP}  & ---                                                                 & peso chileno  & \texttt{CLP} & \emph{«peso chileno / pesos chilenos»} \\
%     \texttt{MXN}  & ---                                                                 & peso mexicano & \texttt{MXN} & \emph{«peso mexicano / pesos mexicanos»} \\
%     \texttt{USD}  & ---                                                                 & dólar         & \texttt{USD} & \emph{«dólar estadounidense / dólares estadounidenses»} \\
%     \texttt{EUR}  & ---                                                                 & euro          & \texttt{EUR} & \emph{«euro / euros»} \\
%     \texttt{GBP}  & ---                                                                 & libra         & \texttt{GBP} & \emph{«libra / libras»} \\
%     \texttt{JPY}  & ---                                                                 & yen           & \texttt{JPY} & \emph{«yen / yenes»} \\
%     \texttt{CNY}  & ---                                                                 & yuan          & \texttt{CNY} & \emph{«yuan / yuanes»} \\
%     \texttt{KRW}  & ---                                                                 & won           & \texttt{KRW} & \emph{«won / wons»} \\
%     \texttt{ILS}  & ---                                                                 & séquel        & \texttt{ILS} & \emph{«séquel / séqueles»} \\
%     \texttt{INR}  & ---                                                                 & rupia         & \texttt{INR} & \emph{«rupia / rupias»} \\
%     \texttt{RUB}  & ---                                                                 & rublo         & \texttt{RUB} & \emph{«rublo / rublos»} \\
%     \texttt{TRY}  & ---                                                                 & lira          & \texttt{TRY} & \emph{«lira / liras»} \\
%     \texttt{UAH}  & ---                                                                 & grivna        & \texttt{UAH} & \emph{«grivna / grivnas»} \\
%     \midrule
%     \multicolumn{5}{c}{\textbf{Códigos ISO en minúsculas}} \\
%     \midrule
%     \texttt{clp}  & ---                                                                 & peso chileno  & \$ & \emph{«peso chileno / pesos chilenos»} \\
%     \texttt{mxn}  & ---                                                                 & peso mexicano & \$ & \emph{«peso mexicano / pesos mexicanos»} \\
%     \texttt{usd}  & \texttt{dolar}, \texttt{dolares}                                    & dólar         & \$ & \emph{«dólar estadounidense / dólares estadounidenses»} \\
%     \texttt{eur}  & ---                                                                 & euro          & €  & \emph{«euro / euros»} \\
%     \texttt{gbp}  & ---                                                                 & libra         & £  & \emph{«libra / libras»} \\
%     \texttt{jpy}  & ---                                                                 & yen           & ¥  & \emph{«yen / yenes»} \\
%     \texttt{cny}  & ---                                                                 & yuan          & ¥  & \emph{«yuan / yuanes»} \\
%     \texttt{krw}  & ---                                                                 & won           & \texttt{₩} & \emph{«won / wons»} \\
%     \texttt{ils}  & ---                                                                 & séquel        & ₪  & \emph{«séquel / séqueles»} \\
%     \texttt{inr}  & ---                                                                 & rupia         & ₹  & \emph{«rupia / rupias»} \\
%     \texttt{rub}  & ---                                                                 & rublo         & \texttt{₽} & \emph{«rublo / rublos»} \\
%     \texttt{try}  & ---                                                                 & lira          & \texttt{₺} & \emph{«lira / liras»} \\
%     \texttt{uah}  & ---                                                                 & grivna        & \texttt{₴} & \emph{«grivna / grivnas»} \\
%     \bottomrule
%   \end{tabular}
% \end{table}
%
% \subsection{El comando \cs[no-index]{spshort}}\label{sec:spshort}
%
% \vspace*{-\baselineskip}
%
% \begin{function}{\spshort}
%   \cmdexamp{spshort}{abreviaturas \textnormal{\textcolor{lightgray}{\textbar}} acrónimos \textnormal{\textcolor{lightgray}{\textbar}} números ordinales}
%   \cmdexamp{spshort}[key \textnormal{\textcolor{gray}{=}} val]{abreviaturas \textnormal{\textcolor{lightgray}{\textbar}} acrónimos \textnormal{\textcolor{lightgray}{\textbar}} números ordinales}
%   \cmdexamp{spshort}[read-arg\textnormal{\textcolor{gray}{=}}\mymarg{verbalización del argumento para \textnormal{NVDA}}]{comandos \hologo{LaTeX}}
%
% \smallskip
%
%   El comando \ics*{spshort} trabaja \emph{solo en modo texto} y recibe
%   \emph{un argumento obligatorio} del tipo \mymarg{abreviaturas}, \mymarg{acrónimos} o \mymarg{números ordinales}
%   descritos en \S\ref{subsec:arg-spshort}.
%
%   \smallskip
%
%   Si no se proporciona el \emph{argumento opcional} \myoarg{key \textnormal{\textcolor{gray}{=}} val}
%   y se ejecuta |\spshort{dr.a}| se imprimirá `Dr.\textsuperscript{a}' en la salida
%   \texttt{PDF} y \mynvmath{NVDA/MathCAT} verbalizará \emph{«doctora»};
%   si se ejecuta |\spshort{onu}| se imprimirá `ONU' en la salida \texttt{PDF}
%   y \mynvmath{NVDA/MathCAT} verbalizará \emph{«organización de las naciones unidas»}
%   y si se ejecuta |\spshort{mineduc}| se imprimirá `Mineduc' en la salida \texttt{PDF}
%   y \mynvmath{NVDA/MathCAT} verbalizará \emph{«ministerio de educación»}.
% \end{function}
%
% \smallskip
%
% Para la escritura de \mymarg{números ordinales} se puede ejecutar
% |\spshort{1.a}| que imprimirá `{\addfontfeatures{Numbers=Lining}1.\textsuperscript{a}}'
% en la salida \texttt{PDF} y \mynvmath{NVDA/MathCAT} verbalizará \emph{«primera»};
% si se ejecuta |\spshort{1.o}| imprimirá `{\addfontfeatures{Numbers=Lining}1.\textsuperscript{o}}'
% en la salida \texttt{PDF} y \mynvmath{NVDA/MathCAT} verbalizará \emph{«primero»}.
%
% \smallskip
%
% Para crear un nueva \mymarg{abreviatura}, \mymarg{acrónimo} o \mymarg{número ordinal}
% que no se encuentre registrada por \mypkg*[no-index]{spintent} dentro del módulo Lua \myluamod{}
% descritas en \S\ref{subsec:arg-spshort} se debe utilizar la llave
% |read-arg| y definir esta dentro del \emph{argumento obligatorio}
% que puede contener \mymarg{comandos \hologo{LaTeX}} \emph{en modo texto}
% los cuales se imprimirán en salida \texttt{PDF}.
%
% \smallskip
%
% \begin{important}*
%   Para la escritura de \mymarg{números ordinales} se debe tener cuidado
%   y no confundir la \emph{«letra o voladita»} `\textsuperscript{o}' con
%   el símbolo de grado `°'. Si la fuente tipográfica del \emph{modo texto}
%   está utilizando |Numbers=OldStyle| los números pueden quedar muy separados
%   de la \emph{«letra voladita»}, la solución para esto es agregar |\addfontfeatures{Numbers=Lining}|
%   dentro de un grupo junto al comando: |{\addfontfeatures{Numbers=Lining}\spshort{1.o}}|.
% \end{important}
%
% \smallskip
%
% \begin{important}*
%   El comando |\spshort| está pensado como un reemplazo para \emph{shorthands}
%   y el comando |\sptext| provistos por el paquete \mypkg{babel} en español
%   compatible con \emph{tagged} PDF. Es de las pocas utilidades que provee
%   el paquete \mypkg*[no-index]{spintent} que trabaja \emph{solo en modo texto}. Internamente
%   utiliza la llave \mykey{E} \emph{(expansion text)} proveída por el paquete
%   \mypkg{tagpdf}\cite{tagpdf}.
% \end{important}
%
% \subsubsection{Llaves para \cs[no-index]{spshort}}
%
% \keyexampcmd{read-arg}{verbalización \textnormal{NVDA}}{no usado}[spshort]
% \smallskip
% Establece la \emph{verbalización del argumento} \mymarg{comandos \textnormal{\hologo{LaTeX}}}
% de la nueva \mymarg{abreviatura}, \mymarg{acrónimo} o \mymarg{número ordinal} para \mynvmath{NVDA}.
% El valor de esta llave \emph{siempre} debe pasarse entre `|{ }|’. Por ejemplo:
% |\spshort[read-arg={octavo}]{8.\textsuperscript{\emph{avo}}}|
% creará una nueva \mymarg{abreviatura} que imprimirá `8.\textsuperscript{\emph{avo}}'
% en la salida \texttt{PDF} y \mynvmath{NVDA} verbalizará \emph{«octavo»}.
%
% \smallskip
%
% \keyexampcmd{small-caps}{true \textnormal{\textcolor{lightgray}{\textbar}} false}{false}[spshort]
% \smallskip
% Establece si se usarán \emph{versalitas} para imprimir \mymarg{acrónimos}
% y para la \emph{«letra volada»} en \mymarg{abreviaturas} y \mymarg{números ordinales}
% respetando las reglas de la \textsc{rae} para el uso de estas. Si se establece en |true|,
% las abreviaturas con mayúscula como \texttt{EE.~UU.} y los acrónimos
% de cinco o más letras como \texttt{Unicef} no se verán afectados; si se
% establece en |false|, no se aplicarán \emph{versalitas}. Por ejemplo:
% |\spshort[small-caps=true]{dr.a}| imprimirá `Dr.\textsuperscript{\textsc{a}}'
% en la salida \texttt{PDF} y \mynvmath{NVDA} verbalizará \emph{«doctora»}.
%
% \subsubsection{Argumentos para \cs[no-index]{spshort}}\label{subsec:arg-spshort}
%
% El cuadro~\ref{tab:spshort-abreviaturas} resume las abreviaturas,
% expresiones latinas y títulos registrados en el módulo Lua \myluamod{}.
%
% \begin{longtable}{l p{3.5cm} c p{5cm}}
% \caption{Abreviaturas soportadas} \label{tab:spshort-abreviaturas} \\
% \toprule
% \textbf{Entrada} & \textbf{Alias / Variantes} & \textbf{Salida PDF} & \multicolumn{1}{c}{\textbf{Lectura NVDA}} \\
% \midrule
% \endfirsthead
% \caption[]{\emph{(Continuación - Abreviaturas soportadas)}} \\
% \toprule
% \textbf{Entrada} & \textbf{Alias / Variantes} & \textbf{Salida PDF} & \multicolumn{1}{c}{\textbf{Lectura NVDA}} \\
% \midrule
% \endhead
% \midrule
% \multicolumn{4}{r}{\emph{Continúa en la siguiente página...}} \\
% \endfoot
% \bottomrule
% \endlastfoot
% \multicolumn{4}{c}{\textbf{Abreviaturas con letras voladas}} \\
% \midrule
% \texttt{dr.a} & \texttt{dr.ª} & Dr.$^{\text{a}}$ & \emph{«doctora»} \\
% \texttt{dr.as} & --- & Dr.$^{\text{as}}$ & \emph{«doctoras»} \\
% \texttt{prof.a} & --- & Prof.$^{\text{a}}$ & \emph{«profesora»} \\
% \texttt{d.a} & \texttt{d.ª} & D.$^{\text{a}}$ & \emph{«doña»} \\
% \texttt{m.a} & --- & M.$^{\text{a}}$ & \emph{«maría»} \\
% \texttt{n.o} & \texttt{n.º} & N.$^{\text{o}}$ & \emph{«número»} \\
% \texttt{n.os} & --- & N.$^{\text{os}}$ & \emph{«números»} \\
% \texttt{c.ia} & --- & C.$^{\text{ia}}$ & \emph{«compañía»} \\
% \midrule
% \multicolumn{4}{c}{\textbf{Abreviaturas comunes y expresiones latinas}} \\
% \midrule
% \texttt{pág.} & \texttt{pag.} & pág. & \emph{«página»} \\
% \texttt{vol.} & --- & vol. & \emph{«volumen»} \\
% \texttt{etc.} & --- & etc. & \emph{«etcétera»} \\
% \texttt{et al.} & \texttt{et. al.} & et~al. & \emph{«y otros»} \\
% \texttt{ibíd.} & \texttt{ibid.} & ibíd. & \emph{«ibídem»} \\
% \texttt{op. cit.} & \texttt{op.cit.} & op.~cit. & \emph{«obra citada»} \\
% \texttt{loc. cit.} & \texttt{loc.cit.} & loc.~cit. & \emph{«lugar citado»} \\
% \texttt{v. gr.} & \texttt{v.gr.} & v.~gr. & \emph{«verbigracia»} \\
% \texttt{i. e.} & \texttt{i.e.} & i.~e. & \emph{«esto es»} \\
% \texttt{e. g.} & \texttt{e.g.} & e.~g. & \emph{«por ejemplo»} \\
% \texttt{p. ej.} & \texttt{p.ej.} & p.~ej. & \emph{«por ejemplo»} \\
% \midrule
% \multicolumn{4}{c}{\textbf{Títulos, profesiones y cargos}} \\
% \midrule
% \texttt{sr.} & --- & Sr. & \emph{«señor»} \\
% \texttt{sra.} & --- & Sra. & \emph{«señora»} \\
% \texttt{srta.} & --- & Srta. & \emph{«señorita»} \\
% \texttt{dr.} & --- & Dr. & \emph{«doctor»} \\
% \texttt{dra.} & --- & Dra. & \emph{«doctora»} \\
% \texttt{dres.} & --- & Dres. & \emph{«doctores»} \\
% \texttt{dras.} & --- & Dras. & \emph{«doctoras»} \\
% \texttt{prof.} & --- & Prof. & \emph{«profesor»} \\
% \texttt{profa.} & --- & Profa. & \emph{«profesora»} \\
% \texttt{profs.} & --- & Profs. & \emph{«profesores»} \\
% \texttt{ing.} & --- & Ing. & \emph{«ingeniero»} \\
% \texttt{ings.} & --- & Ings. & \emph{«ingenieros»} \\
% \texttt{lic.} & --- & Lic. & \emph{«licenciado»} \\
% \texttt{v. b.} & \texttt{v.b.} & V.~B. & \emph{«visto bueno»} \\
% \midrule
% \multicolumn{4}{c}{\textbf{Abreviaturas compuestas e institucionales}} \\
% \midrule
% \texttt{s. a.} & \texttt{s.a.} & S.~A. & \emph{«sociedad anónima»} \\
% \texttt{ee. uu.} & \texttt{ee.uu.} & EE.~UU. & \emph{«estados unidos»} \\
% \texttt{dd. hh.} & \texttt{dd.hh.} & DD.~HH. & \emph{«derechos humanos»} \\
% \texttt{jj. oo.} & \texttt{jj.oo.} & JJ.~OO. & \emph{«juegos olímpicos»} \\
% \texttt{ff. aa.} & \texttt{ff.aa.} & FF.~AA. & \emph{«fuerzas armadas»} \\
% \texttt{rr. ee.} & \texttt{rr.ee.} & RR.~EE. & \emph{«relaciones exteriores»} \\
% \texttt{cc. aa.} & \texttt{cc.aa.} & CC.~AA. & \emph{«comunidades autónomas»} \\
% \texttt{p. d.} & \texttt{p.d.} & P.~D. & \emph{«posdata»} \\
% \texttt{a. c.} & \texttt{a.c.} & a.~C. & \emph{«antes de Cristo»} \\
% \texttt{d. c.} & \texttt{d.c.} & d.~C. & \emph{«después de Cristo»} \\
% \texttt{d. n. i.} & \texttt{d.n.i.} & D.~N.~I. & \emph{«documento nacional de identidad»} \\
% \texttt{r. u. t.} & \texttt{r.u.t.} & R.~U.~T. & \emph{«rol único tributario»} \\
% \texttt{r. u. n.} & \texttt{r.u.n.} & R.~U.~N. & \emph{«rol único nacional»} \\
% \end{longtable}
%
% El cuadro~\ref{tab:spshort-siglas} detalla las siglas y acrónimos registrados en el módulo Lua \myluamod{}.
%
% \begin{longtable}{l p{3.5cm} c p{5cm}}
% \caption{Siglas y acrónimos soportados} \label{tab:spshort-siglas} \\
% \toprule
% \textbf{Entrada} & \textbf{Alias / Variantes} & \textbf{Salida PDF} & \multicolumn{1}{>{\centering\arraybackslash}p{5cm}}{\textbf{Lectura NVDA}} \\
% \midrule
% \endfirsthead
% \caption[]{\emph{(Continuación - Siglas y acrónimos)}} \\
% \toprule
% \textbf{Entrada} & \textbf{Alias / Variantes} & \textbf{Salida PDF} & \multicolumn{1}{>{\centering\arraybackslash}p{5cm}}{\textbf{Lectura NVDA}} \\
% \midrule
% \endhead
% \midrule
% \multicolumn{4}{r}{\emph{Continúa en la siguiente página...}} \\
% \endfoot
% \bottomrule
% \endlastfoot
% \multicolumn{4}{c}{\textbf{Siglas (Deletreo o lectura expandida)}} \\
% \midrule
% \texttt{dni} & --- & DNI & \emph{«documento nacional de identidad»} \\
% \texttt{rut} & --- & RUT & \emph{«rol único tributario»} \\
% \texttt{run} & --- & RUN & \emph{«rol único nacional»} \\
% \texttt{onu} & --- & ONU & \emph{«organización de las naciones unidas»} \\
% \texttt{rae} & --- & RAE & \emph{«real academia española»} \\
% \texttt{ong} & \texttt{ongs} & ONG & \emph{«organización no gubernamental»} \\
% \texttt{urss} & --- & URSS & \emph{«unión de repúblicas socialistas soviéticas»} \\
% \texttt{oea} & --- & OEA & \emph{«organización de los estados americanos»} \\
% \texttt{oms} & --- & OMS & \emph{«organización mundial de la salud»} \\
% \texttt{fmi} & --- & FMI & \emph{«fondo monetario internacional»} \\
% \texttt{bid} & --- & BID & \emph{«banco interamericano de desarrollo»} \\
% \texttt{otan} & --- & OTAN & \emph{«organización del tratado del atlántico norte»} \\
% \texttt{ue} & --- & UE & \emph{«unión europea»} \\
% \texttt{ru} & --- & RU & \emph{«reino unido»} \\
% \midrule
% \multicolumn{4}{c}{\textbf{Acrónimos silábicos}} \\
% \midrule
% \texttt{unicef} & --- & Unicef & \emph{«fondo de las naciones unidas para la infancia»} \\
% \texttt{unesco} & --- & Unesco & \emph{«organización de las naciones unidas para la educación, la ciencia y la cultura»} \\
% \texttt{mercosur} & --- & Mercosur & \emph{«mercado común del sur»} \\
% \texttt{cepal} & --- & Cepal & \emph{«comisión económica para américa latina»} \\
% \texttt{mineduc} & --- & Mineduc & \emph{«ministerio de educación»} \\
% \texttt{conicet} & --- & Conicet & \emph{«consejo nacional de investigaciones»} \\
% \end{longtable}
%
% El cuadro~\ref{tab:spshort-ordinales} muestra ejemplos de números ordinales
% registrados en el módulo Lua \myluamod{}..
%
% \begin{longtable}{l p{3.5cm} c p{5cm}}
% \caption{Números ordinales soportados} \label{tab:spshort-ordinales} \\
% \toprule
% \textbf{Entrada} & \textbf{Alias / Variantes} & \textbf{Salida PDF} & \multicolumn{1}{c}{\textbf{Lectura NVDA}} \\
% \midrule
% \endfirsthead
% \caption[]{\emph{(Continuación - Números ordinales)}} \\
% \toprule
% \textbf{Entrada} & \textbf{Alias / Variantes} & \textbf{Salida PDF} & \multicolumn{1}{c}{\textbf{Lectura NVDA}} \\
% \midrule
% \endhead
% \midrule
% \multicolumn{4}{r}{\emph{Continúa en la siguiente página...}} \\
% \endfoot
% \bottomrule
% \endlastfoot
% \multicolumn{4}{c}{\textbf{Sufijos para ordinales (Ejemplos)}} \\
% \midrule
% \texttt{1.a} & \texttt{1.ª} & 1.$^{\text{a}}$ & \emph{«primera»} \\
% \texttt{2.o} & \texttt{2.º} & 2.$^{\text{o}}$ & \emph{«segundo»} \\
% \texttt{3.er} & --- & 3.$^{\text{er}}$ & \emph{«tercer»} \\
% \texttt{4.os} & --- & 4.$^{\text{os}}$ & \emph{«cuartos»} \\
% \texttt{5.as} & --- & 5.$^{\text{as}}$ & \emph{«quintas»} \\
% \end{longtable}
%
% \section{Utilidades generales}\label{sec:utilidades}
%
% Además de los \emph{«cuatro grandes»} y sus asociados, el paquete \mypkg*[no-index]{spintent}
% provee tres utilidades de uso general: |\spsiglo| para la escritura y
% lectura de siglos, |\sptime| para la escritura y lectura de tiempo y
% |\spdate| para la escritura y lectura de fechas.
%
% \subsection*{El comando \cs[no-index]{spsiglo}}\label{subsec:spsiglo}
%
% \vspace*{-\baselineskip}
%
% \begin{function}{\spsiglo}
%   \begin{syntax}
%     \cmdexamp{spsiglo}{siglo}
%     \cmdexamp{spsiglo}[key \textnormal{\textcolor{gray}{=}} val]{siglo}
%   \end{syntax}
%   El comando \ics*{spsiglo} trabaja en \emph{modo texto} y en \emph{modo matemático},
%   recibe el \emph{argumento obligatorio} \mymarg{siglo}, el cual puede
%   ser un \emph{número arábigo} `|21|' o un \emph{número romano} `|XXI|'
%   mayor que `|0|' y menor que `|4000|'.
%
%   \smallskip
%
%   Si no se proporciona el \emph{argumento opcional} \myoarg{key \textnormal{\textcolor{gray}{=}} val}
%   y se ejecuta |\spsiglo{21}| o |\spsiglo{XXI}| se imprimirá `\textsc{xxi}' en la salida
%   \texttt{PDF} y \mynvmath{NVDA/MathCAT} verbalizará \emph{«veintiuno»}.
% \end{function}
%
% \smallskip
%
% \begin{important}*
%  El \emph{argumento obligatorio} \mymarg{siglo} \emph{«siempre»} es
%  convertido a \emph{número romano en versalitas} y es verbalizado por
%  \mynvmath{NVDA/MathCAT} como \emph{número arábigo}. Esta es la única utilidad
%  que provee el paquete \mypkg*[no-index]{spintent} que trabaja en \emph{modo texto}
%  y en \emph{modo matemático}. Internamente utiliza la llave \mykey{alt}
%  proveída por el paquete \mypkg{tagpdf} en \emph{modo texto}
%  y usa |\MathMLintent| para el \emph{modo matemático}.
% \end{important}
%
% \subsubsection*{Llaves para \cs[no-index]{spsiglo}}
%
% \keyexampcmd{print-era}{ac \textnormal{\textcolor{lightgray}{\textbar}} dc \textnormal{\textcolor{lightgray}{\textbar}} none}{none}[spsiglo]
% \smallskip
% Establece si se \emph{imprime} o no la abreviatura después de imprimir el argumento \mymarg{siglo}.
% Si se establece en |ac|, |\spsiglo[print-era=ac]{XXI}| imprimirá `\textsc{xi} a.~C.'
% en la salida \texttt{PDF} y \mynvmath{NVDA/MathCAT} verbalizará \emph{«veintiuno antes de cristo»};
% si se establece en |dc|, |\spsiglo[print-era=dc]{XXI}| imprimirá `\textsc{xi} d.~C.'
% en la salida \texttt{PDF} y \mynvmath{NVDA/MathCAT} verbalizará \emph{«veintiuno después de cristo»};
% si se establece en |none|, |\spsiglo[print-era=none]{XXI}| imprimirá `\textsc{xi}'
% en la salida \texttt{PDF} y \mynvmath{NVDA/MathCAT} verbalizará \emph{«veintiuno»}.
%
% \subsection*{El comando \cs[no-index]{sptime}}\label{subsec:sptime}
%
% \vspace*{-\baselineskip}
%
% \begin{function}{\sptime}
%   \begin{syntax}
%     \cmdexamp{sptime}{horas \textnormal{\textcolor{red}{:}} minutos}
%     \cmdexamp{sptime}{horas \textnormal{\textcolor{red}{:}} minutos \textnormal{\textcolor{red}{:}} segundos}
%   \end{syntax}
%   El comando \ics*{sptime} trabaja en \emph{modo matemático}
%   y recibe el \emph{argumento obligatorio} \mymarg{horas \textnormal{\textcolor{red}{:}} minutos}
%   o \mymarg{horas \textnormal{\textcolor{red}{:}} minutos \textnormal{\textcolor{red}{:}} segundos}.
%   Por ejemplo: |$\sptime{23:45}$| imprimirá `23:45' en \emph{modo texto} la salida \texttt{PDF}
%   y \mynvmath{NVDA/MathCAT} verbalizará \emph{«las veintitrés y cuarenta y cinco minutos»}.
% \end{function}
%
% \subsection*{El comando \cs[no-index]{spdate}}\label{subsec:spdate}
%
% \vspace*{-\baselineskip}
%
% \begin{function}{\spdate}
%   \begin{minipage}[t]{0.35\linewidth}
%     \cmdexamp{spdate}{\textnormal{DD\textcolor{red}{/}%
%       MM\textcolor{red}{/}YYYY}}
%     \cmdexamp{spdate}{\textnormal{YYYY\textcolor{red}{/}%
%       MM\textcolor{red}{/}DD}}
%   \end{minipage}\hfill
%   \begin{minipage}[t]{0.55\linewidth}
%     \cmdexamp{spdate}{\textnormal{DD\textcolor{red}{-}%
%       MM\textcolor{red}{-}YYYY}}
%     \cmdexamp{spdate}{\textnormal{YYYY\textcolor{red}{-}%
%       MM\textcolor{red}{-}DD}}
%   \end{minipage}
%
% \smallskip
%
%   El comando \ics*{spdate} trabaja en \emph{modo matemático}
%   y recibe el \emph{argumento} obligatorio
%   \mymarg{\textnormal{DD\textcolor{red}{/}MM\textcolor{red}{/}YYYY}},
%   \mymarg{\textnormal{DD\textcolor{red}{-}MM\textcolor{red}{-}YYYY}} o
%   \mymarg{\textnormal{YYYY\textcolor{red}{-}MM\textcolor{red}{-}DD}}.
%   Por ejemplo: |$\spdate{1981-01-02}$| y |$\spdate{02/01/1981}$|
%   imprimirán en \emph{modo texto} la salida \texttt{PDF} `1981-01-02'
%   y `02/01/1981' respectivamente y en ambos casos, \mynvmath{NVDA/MathCAT}
%   verbalizará \emph{«dos de enero de mil novecientos ochenta y uno»}.
% \end{function}
%
% \smallskip
%
% \begin{important}*
%   Si se pasa el argumento
%   \mymarg{\textnormal{YYYY\textcolor{red}{/}MM\textcolor{red}{/}DD}},
%   internamente se convertirá en
%   \mymarg{\textnormal{DD\textcolor{red}{/}MM\textcolor{red}{/}YYYY}}
%   y se imprimirá de esa forma por compatibilidad con \mynvmath{NVDA/MathCAT}.
% \end{important}
%
% \section{Utilidades para símbolos}
%
% El paquete \mypkg*[no-index]{spintent} provee tres comandos |\spread|, |\spdsym|
% y |\spmsym| para trabajar con \emph{«símbolos matemáticos»} descritos
% en esta sección. Están pensandos desde el punto de vista docente, para
% situaciones donde la lectura de un \emph{«símbolo»} (acorde al contexto)
% es relevante.
% \smallskip
%
% \begin{important}*
%  A diferencia de las otras utilidades implementados por
%  \mypkg*[no-index]{spintent} (que analizan y dan formato a la salida), estos no
%  procesan los \emph{«símbolos matemáticos»}, solo afectan la \emph{forma}
%  (semántica) en que serán verbalizados por \mynvmath{NVDA/MathCAT}.
% \end{important}
%
% \subsection*{El comando \cs[no-index]{spread}}\label{subsec:spread}
%
% \vspace*{-0.5\baselineskip}
%
% \begin{function}{\spread}
%   \begin{syntax}
%      \ics*{spread}\mymarg{verbalización \textnormal{NVDA}}\mymarg{símbolo matemático}
%   \end{syntax}
%   El comando \ics*{spread} trabaja en \emph{modo matemático} y recibe
%   dos \emph{argumentos obligatorios}: El primer \emph{argumento} \mymarg{verbalización \textnormal{NVDA}}
%   establece la \emph{«forma»} en la cual \mynvmath{NVDA/MathCAT}
%   verbalizará al segundo \emph{argumento} \mymarg{símbolo matemático}.
%   Por ejemplo: |$\spread{es-paralela}{\parallel}$| imprimirá `$\parallel$'
%   en la salida \texttt{PDF} y \mynvmath{NVDA/MathCAT} verbalizará \emph{«es paralela»};
%   |$\spread{es-paralelo}{\parallel}$| imprimirá `$\parallel$' en la
%   salida \texttt{PDF} y \mynvmath{NVDA/MathCAT} verbalizará \emph{«es paralelo»}.
% \end{function}
%
% \subsection*{El comando \cs[no-index]{spusep}}\label{subsec:spusep}
%
% \vspace*{-\baselineskip}
%
% \begin{function}{\spusep}
%   \begin{syntax}
%     \ics*{spusep}\mymarg{separador}
%     \ics*{spusep}\myoarg{key \textnormal{\textcolor{gray}{=}} val}\mymarg{separador}
%   \end{syntax}
%   El comando \ics*{spusep} trabaja en \emph{modo matemático} y recibe el
%   \emph{argumento obligatorio} \mymarg{separador}, el cual debe ser uno de
%   los \emph{«delimitadores»} soportados: `|(|', `|)|', `|[|', `|]|',
%   `|\{|' o `|\}|'. Imprime el \emph{delimitador} indicado con el tamaño
%   establecido por la llave \mykey[spusep]{size} y controla si
%   \mynvmath{NVDA/MathCAT} lo verbaliza o no mediante la llave
%   \mykey[spusep]{silent}.
%
%   \smallskip
%
%   Si no se proporciona el \emph{argumento opcional}
%   \myoarg{key \textnormal{\textcolor{gray}{=}} val}, |$\spusep{(}$|
%   imprimirá `$($' en la salida \texttt{PDF} y \mynvmath{NVDA/MathCAT}
%   verbalizará \emph{«paréntesis izquierdo»}.
% \end{function}
%
% \smallskip
%
% \begin{important}*
%   El motivo de existencia de |\spusep| es que los comandos |\left|,
%   |\right|, |\Uleft| y |\Uright| provistos por \hologo{LuaTeX} crean
%   un \emph{grupo} implícito que dificulta el código para \emph{tagged}
%   \texttt{PDF} cuando se necesita interrumpir algo entre ellos, por
%   ejemplo al construir expresiones complejas con |\spnZ|, |\spmQ| o
%   |\spdiv|. El comando |\spusep| imprime el \emph{delimitador} sin crear
%   ningún grupo.
% \end{important}
%
% \subsubsection*{Llaves para \cs[no-index]{spusep}}
%
% \keyexampcmd{size}{none \textnormal{\textcolor{lightgray}{\textbar}}
%   big \textnormal{\textcolor{lightgray}{\textbar}}
%   Big \textnormal{\textcolor{lightgray}{\textbar}}
%   bigg \textnormal{\textcolor{lightgray}{\textbar}}
%   Bigg}{none}[spusep]
% \smallskip
% Establece el \emph{tamaño} del \mymarg{delimitador} que se imprimirá en
% la salida \texttt{PDF}. Los valores corresponden a los cuatro tamaños
% fijos provistos por \hologo{LaTeX}: si se establece en |none|, el
% \mymarg{delimitador} se imprime en el tamaño por defecto del modo
% matemático en curso; si se establece en |big|, se usa |\bigl|/|\bigr|;
% si se establece en |Big|, se usa |\Bigl|/|\Bigr|; si se establece en
% |bigg|, se usa |\biggl|/|\biggr|; y si se establece en |Bigg|, se usa
% |\Biggl|/|\Biggr|. El tamaño se aplica automáticamente en el lado
% correcto (izquierdo o derecho) según el \mymarg{delimitador} pasado.
%
% \smallskip
%
% \keyexampcmd{silent}{true \textnormal{\textcolor{lightgray}{\textbar}}
%   false}{false}[spusep]
% \smallskip
% Establece si \mynvmath{NVDA/MathCAT} \emph{verbalizará} el
% \mymarg{delimitador} impreso por |\spusep|. Si se establece en |true|,
% el \mymarg{delimitador} se marcará con el \emph{intent} |:silent| y
% \mynvmath{NVDA/MathCAT} \emph{no} lo verbalizará; si se establece en
% |false|, \mynvmath{NVDA/MathCAT} verbalizará el \mymarg{delimitador}
% de forma normal.
%
% \subsection*{El comando \cs[no-index]{spdsym}}\label{subsec:spdsym}
%
% \vspace*{-0.5\baselineskip}
%
% \begin{function}{\spdsym}
%   \begin{syntax}
%    \ics*{spdsym}
%    \ics*{spdsym}\myoarg{key \textnormal{\textcolor{gray}{=}} val}
%   \end{syntax}
%   El comando \ics*{spdsym} trabaja en \emph{modo matemático} y recibe
%   el \emph{argumento opcional} \myoarg{key \textnormal{\textcolor{gray}{=}} val}
%   mediante el cual se establece la \emph{«forma»} en la que se imprimirá el
%   \emph{símbolo de división} y como se verbalizará en \mynvmath{NVDA/MathCAT}.
%
%   \smallskip
%
%   Si no se proporciona el \emph{argumento opcional} \myoarg{key \textnormal{\textcolor{gray}{=}} val},
%   |$\spdsym$| imprimirá `$\mathcolor{red}{:}$' en la salida \texttt{PDF}
%   y \mynvmath{NVDA/MathCAT} verbalizará \emph{«dividido»}.
% \end{function}
%
% \subsubsection*{Llaves para \cs[no-index]{spdsym}}\label{subsec:key-spdsym}
%
% \keyexampcmd{prefix}{prefijo}{no usado}[spdsym]
% \smallskip
% Establece el \emph{prefijo} que será \emph{verbalizado} por
% \mynvmath{NVDA/MathCAT} \emph{«antes»} de \emph{verbalizar} el \mymarg{valor} establecido
% por la llave \mykey[spdsym]{read-sym}. El valor de esta llave \emph{siempre}
% debe pasarse entre llaves `|{ }|'.
%
% \smallskip
%
% \keyexampcmd{suffix}{sufijo}{no usado}[spdsym]
% \smallskip
% Establece el \emph{sufijo} que será \emph{verbalizado} por
% \mynvmath{NVDA/MathCAT} \emph{«después»} de \emph{verbalizar} el \mymarg{valor} establecido
% por la llave \mykey[spdsym]{read-sym}. El valor de esta llave \emph{siempre}
% debe pasarse entre llaves `|{ }|'.
%
% \smallskip
%
% \keyexampcmd{set-sym}{old-style \textnormal{\textcolor{lightgray}{\textbar}} two-dots \textnormal{\textcolor{lightgray}{\textbar}} slash}{two-dots}[spdsym]
% \smallskip
% Establece el \emph{símbolo de división} que se imprimirá en la salida
% \texttt{PDF}. Si se establece en |old-style|, imprimirá
% `$\mathcolor{red}{\div}$'; si se establece en |two-dots|, imprimirá
% `$\mathcolor{red}{:}$' y si se establece |slash|, imprimirá
% `$\mathcolor{red}{/}$'.
%
% \smallskip
%
% \keyexampcmd{read-sym}{dividido \textnormal{\textcolor{lightgray}{\textbar}} dividido-por \textnormal{\textcolor{lightgray}{\textbar}} dividido-entre}{dividido}[spdsym]
% \smallskip
% Establece la \emph{forma} con la cual \mynvmath{NVDA/MathCAT} verbalizará el
% \emph{símbolo de división}. Si se establece en |dividido|, \mynvmath{NVDA/MathCAT}
% verbalizará \emph{«dividido»}; si se establece en |dividido-por|, \mynvmath{NVDA/MathCAT}
% verbalizará \emph{«dividido por»} y si se establece en |dividido-entre|,
% \mynvmath{NVDA/MathCAT} verbalizará \emph{«dividido entre»}.
%
% \subsection*{El comando \cs[no-index]{spmsym}}\label{subsec:spmsym}
%
% \vspace*{-0.5\baselineskip}
%
% \begin{function}{\spmsym}
%   \begin{syntax}
%     \ics*{spmsym}
%     \ics*{spmsym}\myoarg{key \textnormal{\textcolor{gray}{=}} val}
%   \end{syntax}
%   El comando \ics*{spmsym} trabaja en \emph{modo matemático} y recibe
%   el \emph{argumento opcional} \myoarg{key \textnormal{\textcolor{gray}{=}} val}
%   mediante el cual se establece la \emph{«forma»} en la que se imprimirá el
%   \emph{símbolo de multiplicación} y como se verbalizará en \mynvmath{NVDA/MathCAT}.
%
%   \smallskip
%
%   Si no se proporciona el \emph{argumento opcional} \myoarg{key \textnormal{\textcolor{gray}{=}} val},
%   |$\spmsym$| imprimirá `$\mathcolor{red}{\cdot}$' en la salida \texttt{PDF}
%   y \mynvmath{NVDA/MathCAT} verbalizará \emph{«por»}.
% \end{function}
%
% \subsubsection*{Llaves para \cs[no-index]{spmsym}}\label{subsec:key-spmsym}
%
% \keyexampcmd{prefix}{prefijo}{no usado}[spmsym]
% \smallskip
% Establece el \emph{prefijo} que será \emph{verbalizado} por \mynvmath{NVDA/MathCAT}
% \emph{«antes»} de \emph{verbalizar} el \mymarg{valor} establecido por
% la llave \mykey[spmsym]{read-sym}. El valor de esta llave \emph{siempre}
% debe pasarse entre llaves `|{ }|'.
%
% \smallskip
%
% \keyexampcmd{suffix}{sufijo}{no usado}[spmsym]
% \smallskip
% Establece el \emph{sufijo} que será \emph{verbalizado} por \mynvmath{NVDA/MathCAT}
% \emph{«después»} de \emph{verbalizar} el \mymarg{valor} establecido por
% la llave \mykey[spmsym]{read-sym}. El valor de esta llave \emph{siempre}
% debe pasarse entre llaves `|{ }|'.
%
% \smallskip
%
% \keyexampcmd{set-sym}{cross \textnormal{\textcolor{lightgray}{\textbar}} dot \textnormal{\textcolor{lightgray}{\textbar}} asterisk}{dot}[spmsym]
% \smallskip
% Establece el \emph{símbolo de multiplicación} que se imprimirá en la salida
% \texttt{PDF}. Si se establece en |cross|, imprimirá
% `$\mathcolor{red}{\times}$'; si se establece en |dot|, imprimirá
% `$\mathcolor{red}{\cdot}$' y si se establece en |asterisk|, imprimirá
% `$\mathcolor{red}{\ast}$'.
%
% \smallskip
%
% \keyexampcmd{read-sym}{por \textnormal{\textcolor{lightgray}{\textbar}} multiplicado-por \textnormal{\textcolor{lightgray}{\textbar}} veces}{por}[spmsym]
% \smallskip
% Establece la \emph{forma} con la cual \mynvmath{NVDA/MathCAT}
% verbalizará el \emph{símbolo de multiplicación}. Si se establece en |por|,
% \mynvmath{NVDA/MathCAT} verbalizará \emph{«por»}; si se establece en
% |multiplicado-por|, \mynvmath{NVDA/MathCAT} verbalizará \emph{«multiplicado por»}
% y si se establece en |veces|, \mynvmath{NVDA/MathCAT} verbalizará \emph{«veces»}.
%
% \keyexampcmd{not-print}{true \textnormal{\textcolor{lightgray}{\textbar}} false}{false}[spmsym]
% \smallskip
% Deshabilita la \emph{impresión} del \emph{símbolo de multiplicación}.
% Cuando se establece esta llave, \mynvmath{NVDA/MathCAT}
% verbalizará el \emph{símbolo de multiplicación} acorde al valor establecido
% por la llave \mykey[spmsym]{read-sym}, pero, no se imprimirá este en
% la salida \texttt{PDF}.
%
% \section{Utilidades para números}
%
% El paquete \mypkg*[no-index]{spintent} provee utilidades para trabajar con números
% descritas en esta sección.
%
% \subsection*{El comando \cs[no-index]{spdiv}}\label{subsec:spdiv}
%
% \vspace*{-0.5\baselineskip}
%
% \begin{function}{\spdiv}
%   \begin{syntax}
%     \ics*{spdiv}\mymarg{dividendo}\mymarg{divisor}
%     \ics*{spdiv}\myoarg{key \textnormal{\textcolor{gray}{=}} val}\mymarg{dividendo}\mymarg{divisor}
%   \end{syntax}
%   El comando \ics*{spdiv} trabaja en \emph{modo matemático} y recibe
%   dos \emph{argumentos obligatorios} \mymarg{dividendo} y \mymarg{divisor}.
%
%   \smallskip
%
%   Si no se proporciona el \emph{argumento opcional} \myoarg{key \textnormal{\textcolor{gray}{=}} val},
%   |$\spdiv{20}{26}$| imprimirá `$20:26$' en la salida \texttt{PDF}
%   y \mynvmath{NVDA/MathCAT} verbalizará \emph{«veinte dividido veinte y seis»}
% \end{function}
%
% \subsubsection*{Llaves para \cs[no-index]{spdiv}}
%
% El comando |\spdiv| posee las \mymeta{meta-keys} \mykey[spdiv]{cifras-sep},
% \mykey[spdiv]{read-sign}, \mykey[spdiv]{read-period-sep} y
% \mykey[spdiv]{notation-sym} pasadas al comando |\spnum| (\S\ref{subsec:key-spnum})
% junto a las \mymeta{meta-keys} \mykey[spdiv]{set-sym} y \mykey[spdiv]{read-sym}
% pasadas al comando |\spdsym| (pág.~\pageref{subsec:key-spdsym}).
%
% \smallskip
%
% \keyexampcmd{wrap-parens}{true \textnormal{\textcolor{lightgray}{\textbar}} false}{false}[spdiv]
% \smallskip
% Establece si se \emph{imprimen} paréntesis redondos `|(|\dots|)|'
% alrededor de los \emph{argumentos} \mymarg{dividendo} y \mymarg{divisor}
% pasados a |\spdiv|. Si se establece en |true|, se imprimirán paréntesis
% redondos alrededor de ambos \emph{argumentos}; si se establece en |false|,
% no se imprimirán los paréntesis.
%
% \keyexampcmd{read-parens}{true \textnormal{\textcolor{lightgray}{\textbar}} false}{true}[spdiv]
% \smallskip
% Establece si \mynvmath{NVDA/MathCAT} \emph{verbalizará} los paréntesis
% redondos `|(|\dots|)|' alrededor de los \emph{argumentos} \mymarg{dividendo}
% y \mymarg{divisor}  pasados a |\spdiv|. Si se establece en |true|,
% \mynvmath{NVDA/MathCAT} verbalizará la lectura de los paréntesis;
% si se establece en |false|, \mynvmath{NVDA/MathCAT} no los verbalizará.
%
% \subsection*{El comando \cs[no-index]{spmul}}\label{subsec:spmul}
%
% \vspace*{-0.5\baselineskip}
%
% \begin{function}{\spmul}
%   \begin{syntax}
%     \ics*{spmul}\mymarg{factor}\mymarg{factor}
%     \ics*{spmul}\myoarg{key \textnormal{\textcolor{gray}{=}} val}\mymarg{factor}\mymarg{factor}
%   \end{syntax}
%   El comando \ics*{spmul} trabaja en \emph{modo matemático} y recibe
%   dos \emph{argumentos obligatorios} \mymarg{factor} y \mymarg{factor}.
%
%   \smallskip
%
%   Si no se proporciona el \emph{argumento opcional} \myoarg{key \textnormal{\textcolor{gray}{=}} val},
%   |$\spmul{20}{26}$| imprimirá `$20 \cdot 26$' en la salida \texttt{PDF}
%   y \mynvmath{NVDA/MathCAT} verbalizará \emph{«veinte por veinte y seis»}.
% \end{function}
%
% \subsubsection*{Llaves para \cs[no-index]{spmul}}
%
% El comando |\spmul| posee las \mymeta{meta-keys} \mykey[spmul]{cifras-sep},
% \mykey[spmul]{read-sign}, \mykey[spmul]{read-period-sep} y
% \mykey[spmul]{notation-sym} pasadas al comando |\spnum| (\S\ref{subsec:key-spnum})
% junto a las \mymeta{meta-keys} \mykey[spmul]{set-sym} y \mykey[spmul]{read-sym}
% pasadas al comando |\spmsym| (pág.~\pageref{subsec:key-spmsym}).
%
% \smallskip
%
% \keyexampcmd{wrap-parens}{true \textnormal{\textcolor{lightgray}{\textbar}} false}{false}[spmul]
% \smallskip
% Establece si se \emph{imprimen} paréntesis redondos `|(|\dots|)|'
% alrededor de los \emph{argumentos} \mymarg{factor} y \mymarg{factor}
% pasados a |\spmul|. Si se establece en |true|, se imprimirán paréntesis
% redondos alrededor de ambos \emph{argumentos}; si se establece en |false|,
% no se imprimirán los paréntesis.
%
% \keyexampcmd{read-parens}{true \textnormal{\textcolor{lightgray}{\textbar}} false}{true}[spmul]
% \smallskip
% Establece si \mynvmath{NVDA/MathCAT} \emph{verbalizará} los paréntesis
% redondos `|(|\dots|)|' alrededor de los \emph{argumentos} \mymarg{factor}
% y \mymarg{factor}  pasados a |\spmul|. Si se establece en |true|,
% \mynvmath{NVDA/MathCAT} verbalizará la lectura de los paréntesis;
% si se establece en |false|, \mynvmath{NVDA/MathCAT} no los verbalizará.
%
% \subsection*{El comando \cs[no-index]{spnD}}\label{subsec:spnD}
%
% \vspace*{-\baselineskip}
%
% \begin{function}{\spnD}
%   \begin{syntax}
%     |\spnD|
%     |\spnD|\myparg{número natural}
%     |\spnD|\myoarg{key \textnormal{\textcolor{gray}{=}} val}\myparg{número natural}
%   \end{syntax}
%   El comando \ics*{spnD} trabaja en \emph{modo matemático} y recibe
%   el \emph{argumento delimitado por paréntesis} \myparg{número natural},
%   por ejemplo: `|(6)|'. Si se usa sin \myparg{argumento}, |$\spnD$|
%   imprimirá `$\mathrm{D}(n)$' en la salida \texttt{PDF} y \mynvmath{NVDA/MathCAT}
%   verbalizará \emph{«divisores de n»}.
%
%   \smallskip
%
%   Si se usa con \myparg{argumento}, por ejemplo, |$\spnD(6)$|, imprimirá
%   `$\mathrm{D}(6)$' en la salida \texttt{PDF} y \mynvmath{NVDA/MathCAT}
%   verbalizará \emph{«divisores de seis»}.
% \end{function}
%
% \subsubsection*{Llaves para \cs[no-index]{spnD}}
%
% \keyexampcmd{prefix}{prefijo}{no usado}[spnD]
% \smallskip
% Establece el \emph{prefijo} que será verbalizado por \mynvmath{NVDA/MathCAT}
% \emph{«antes»} de verbalizar |\spnD|\myparg{argumento}. El valor
% de esta llave \emph{siempre} debe pasarse entre llaves `|{ }|'.
%
% \smallskip
%
% \keyexampcmd{result}{true \textnormal{%
%   \textcolor{lightgray}{\textbar}} false}{false}[spnD]
% \smallskip
% Establece si se \emph{imprime} el resultado de calcular los
% \emph{divisores} del \myparg{número natural} pasado a |\spnD|. Si se
% establece en |true|, |$\spnD[result=true](6)$| imprimirá
% `$\mathrm{D}(6) = \{1, 2, 3, 6\}$' en la salida \texttt{PDF} y
% \mynvmath{NVDA/MathCAT} verbalizará \emph{«los divisores de seis son uno,
% dos, tres y seis»}; si se establece en |false|, no se imprimirá
% en la salida \texttt{PDF} el resultado de calcular los \emph{divisores}.
%
% \smallskip
%
% \keyexampcmd{result-num}{número natural}{5}[spnD]
% \smallskip
% Establece la \emph{cantidad} de números que se mostrarán en la salida
% \texttt{PDF} al activar |result=true|. Si se establece en |2|, solo
% se imprimirán dos valores en la salida \texttt{PDF}; si no se establece,
% se imprimirán como máximo |5| valores y luego se añadirá `\dots' en la
% salida \texttt{PDF}.
%
% \smallskip
%
% \keyexampcmd{sample-arg}{true \textnormal{\textcolor{lightgray}{\textbar}} false}{false}[spnD]
% \smallskip
% Desactiva la \emph{validación} interna del \myparg{argumento} pasado
% a |\spnD|, permitiendo usar ejemplos \emph{«no numéricos»} como entrada.
% Si se establece en |true|, |$\spnD[sample-arg=true](m)$| imprimirá `$\mathrm{D}(m)$'
% en la salida \texttt{PDF} y \mynvmath{NVDA/MathCAT} verbalizará
% \emph{«divisores de m»}; si se establece en |false|, el \myparg{argumento}
% debe ser un número natural, en caso contrario se devolverá un \enquote{error}.
%
% \smallskip
%
% \keyexampcmd{silent-cmd}{true \textnormal{\textcolor{lightgray}{\textbar}} false}{false}[spnD]
% \smallskip
% Establece si se \emph{silencia} la verbalización para \mynvmath{NVDA/MathCAT}
% de |\spnD|\myparg{argumento} cuando se establece |sample-arg=true|
% o cuando |\spnD| se usa sin \myparg{argumento}. Si se establece en |true|
% \mynvmath{NVDA/MathCAT} NO lo verbalizará; si se establece en |false|,
% \mynvmath{NVDA/MathCAT} SÍ lo verbalizará.
%
% \subsection*{El comando \cs[no-index]{spnM}}\label{subsec:spnM}
%
% \vspace*{-\baselineskip}
%
% \begin{function}{\spnM}
%   \begin{syntax}
%     |\spnM|
%     |\spnM|\myparg{número natural}
%     |\spnM|\myoarg{key \textnormal{\textcolor{gray}{=}} val}\myparg{número natural}
%   \end{syntax}
%   El comando \ics*{spnM} trabaja en \emph{modo matemático} y recibe
%   el \emph{argumento delimitado por paréntesis} \myparg{número natural},
%   por ejemplo: `|(6)|'. Si se usa sin \myparg{argumento}, |$\spnM$|
%   imprimirá `$\mathrm{M}(n)$' en la salida \texttt{PDF} y \mynvmath{NVDA/MathCAT}
%   verbalizará \emph{«múltiplos de n»}.
%
%   \smallskip
%
%   Si se usa con \myparg{argumento}, por ejemplo, |$\spnM(6)$|, imprimirá
%   `$\mathrm{M}(6)$' en la salida \texttt{PDF} y \mynvmath{NVDA/MathCAT}
%   verbalizará \emph{«múltiplos de seis»}.
% \end{function}
%
% \subsubsection*{Llaves para \cs[no-index]{spnM}}
%
% \keyexampcmd{prefix}{prefijo}{no usado}[spnM]
% \smallskip
% Establece el \emph{prefijo} que será verbalizado por \mynvmath{NVDA/MathCAT}
% \emph{«antes»} de verbalizar |\spnM|\myparg{argumento}. El valor
% de esta llave \emph{siempre} debe pasarse entre llaves `|{ }|'.
%
% \smallskip
%
% \keyexampcmd{result}{true \textnormal{%
%   \textcolor{lightgray}{\textbar}} false}{false}[spnM]
% \smallskip
% Establece si se \emph{imprime} el resultado de calcular los
% \emph{múltiplos} del \myparg{número natural} pasado a |\spnM|. Si se
% establece en |true|, |$\spnM[result=true](6)$| imprimirá
% `$\mathrm{M}(6) = \{6, 12, 18, 24, 30, \dots\}$' en la salida \texttt{PDF} y
% \mynvmath{NVDA/MathCAT} verbalizará \emph{«los múltiplos de seis son seis,
% doce, dieciocho, veinticuatro, treinta puntos suspencivos»}; si se
% establece en |false|, no se imprimirá en la salida \texttt{PDF} el
% resultado de calcular los \emph{múltiplos}.
%
% \smallskip
%
% \keyexampcmd{result-num}{número natural}{5}[spnM]
% \smallskip
% Establece la \emph{cantidad} de números que se mostrarán en la salida
% \texttt{PDF} al activar |result=true|. Si se establece en |2|, solo
% se imprimirán dos valores en la salida \texttt{PDF} seguidos de `\dots';
% si no se establece, se imprimirán por defecto los |5| primeros valores y
% luego se añadirá `\dots' en la salida \texttt{PDF}.
%
% \smallskip
%
% \keyexampcmd{sample-arg}{true \textnormal{\textcolor{lightgray}{\textbar}} false}{false}[spnM]
% \smallskip
% Desactiva la \emph{validación} interna del \myparg{argumento} pasado
% a |\spnM|, permitiendo usar ejemplos \emph{«no numéricos»} como entrada.
% Si se establece en |true|, |$\spnM[sample-arg=true](m)$| imprimirá `$\mathrm{M}(m)$'
% en la salida \texttt{PDF} y \mynvmath{NVDA/MathCAT} verbalizará
% \emph{«múltiplos de m»}; si se establece en |false|, el \myparg{argumento}
% debe ser un número natural, en caso contrario se devolverá un \enquote{error}.
%
% \smallskip
%
% \keyexampcmd{silent-cmd}{true \textnormal{\textcolor{lightgray}{\textbar}} false}{false}[spnM]
% \smallskip
% Establece si se \emph{silencia} la verbalización en \mynvmath{NVDA/MathCAT}
% de |\spnM|\myparg{argumento} cuando se establece |sample-arg=true|
% o cuando |\spnM| se usa sin \myparg{argumento}. Si se establece en |true|
% \mynvmath{NVDA/MathCAT} NO lo verbalizará; si se establece en |false|
% \mynvmath{NVDA/MathCAT} lo verbalizará.
%
% \subsection*{El comando \cs[no-index]{spmcd}}\label{subsec:spmcd}
%
% \vspace*{-\baselineskip}
%
% \begin{function}{\spmcd}
%   \begin{syntax}
%     |\spmcd|
%     |\spmcd|\myparg{números separados por coma}
%     |\spmcd|\myoarg{key \textnormal{\textcolor{gray}{=}} val}\myparg{números separados por coma}
%   \end{syntax}
%   El comando \ics*{spmcd} trabaja en \emph{modo matemático} y recibe el
%   \emph{argumento delimitado por paréntesis} \myparg{números separados por coma}, el cual
%   debe ser una lista de \emph{números naturales}, por ejemplo: `|(2, 3, 6)|'.
%   Si se usa sin \myparg{argumento} |$\spmcd$| imprimirá `$\mathrm{mcd}$' en
%   la salida \texttt{PDF}, y \mynvmath{NVDA/MathCAT} verbalizará \emph{«máximo
%   común divisor»} o \emph{«m c d»}, acorde al \mymarg{valor} establecido por
%   la llave |read-cmd|.
%
%   \smallskip
%
%   Si se usa con \myparg{argumento} |$\spmcd(2, 3, 6)$| imprimirá
%   `$\mathrm{mcd}(2, 3, 6)$' en la salida \texttt{PDF}, y
%   \mynvmath{NVDA/MathCAT} verbalizará \emph{«máximo común divisor entre dos,
%   tres y seis»} o \emph{«m c d entre dos, tres y seis»} acorde al
%   \mymarg{valor} establecido por la llave |read-cmd|.
% \end{function}
%
% \subsubsection*{Llaves para \cs[no-index]{spmcd}}
%
% \keyexampcmd{upper-case}{true \textnormal{%
%   \textcolor{lightgray}{\textbar}} false}{false}[spmcd]
% \smallskip
% Establece si se usarán \emph{mayúsculas} o \emph{minúsculas} en la salida
% \texttt{PDF} para el comando |\spmcd|. Si se establece en |true|, se
% imprimirá `$\mathrm{MCD}$' en la salida \texttt{PDF}; si se establece en
% |false|, se imprimirá `$\mathrm{mcd}$' en la salida \texttt{PDF}.
%
% \smallskip
%
% \keyexampcmd{read-cmd}{terse \textnormal{%
%   \textcolor{lightgray}{\textbar}} clear}{clear}[spmcd]
% \smallskip
% Establece la \emph{forma} en la que \mynvmath{NVDA/MathCAT} verbalizará el
% comando |\spmcd|. Si se establece en |terse|, \mynvmath{NVDA/MathCAT}
% verbalizará \emph{«m c d»}; si se establece en |clear|,
% \mynvmath{NVDA/MathCAT} verbalizará \emph{«máximo común divisor»}.
%
% \smallskip
%
% \keyexampcmd{prefix}{prefijo}{no usado}[spmcd]
% \smallskip
% Establece el \emph{prefijo} que será verbalizado por \mynvmath{NVDA/MathCAT}
% \emph{«antes»} de verbalizar el valor establecido por la llave \mykey[spmcd]{read-cmd}.
% El valor de esta llave \emph{siempre} debe pasarse entre llaves `|{ }|'.
%
% \smallskip
%
% \keyexampcmd{result}{true \textnormal{%
%   \textcolor{lightgray}{\textbar}} false}{false}[spmcd]
% \smallskip
% Establece si se \emph{imprime} el resultado de calcular el \emph{«máximo
% común divisor»} del \myparg{argumento} pasado a |\spmcd|. Si se
% establece en |true|: |$\spmcd[result=true](2, 3, 6)$| imprimirá
% `$\mathrm{mcd}(2, 3, 6) = 1$' en la salida \texttt{PDF} y
% \mynvmath{NVDA/MathCAT} verbalizará \emph{«máximo común divisor entre
% dos, tres y seis es igual a uno»}; si se establece en |false| no se
% imprimirá en la salida \texttt{PDF} el resultado de calcular el
% \emph{«máximo común divisor»}.
%
% \smallskip
%
% \keyexampcmd{sample-arg}{true \textnormal{%
%   \textcolor{lightgray}{\textbar}} false}{false}[spmcd]
% \smallskip
% Desactiva la \emph{validación} interna del \myparg{argumento} pasado
% a |\spmcd|, permitiendo usar ejemplos \emph{«no numéricos»} como entrada.
% Si se establece en |true|: |$\spmcd[sample-arg=true](m, n)$|, imprimirá
% `$\mathrm{mcd}(m, n)$' en la salida \texttt{PDF} y \mynvmath{NVDA/MathCAT}
% verbalizará \emph{«máximo común divisor entre m y n»}; si se establece
% en |false| el \myparg{argumento} debe ser una lista de \emph{números naturales},
% en caso contrario se devolverá un \enquote{error}.
%
% \smallskip
%
% \keyexampcmd{silent-cmd}{true \textnormal{%
%   \textcolor{lightgray}{\textbar}} false}{false}[spmcd]
% \smallskip
% Establece si se \emph{silencia} la verbalización en \mynvmath{NVDA/MathCAT}
% de |\spmcd|\myparg{argumento} cuando se establece |sample-arg=true|
% o cuando |\spmcd| se usa sin \myparg{argumento}. Si se establece en |true|
% \mynvmath{NVDA/MathCAT} NO lo verbalizará; si se establece en |false|,
% \mynvmath{NVDA/MathCAT} SÍ lo verbalizará.
%
% \subsection*{El comando \cs[no-index]{spmcm}}\label{subsec:spmcm}
%
% \vspace*{-\baselineskip}
%
% \begin{function}{\spmcm}
%   \begin{syntax}
%     |\spmcm|
%     |\spmcm|\myparg{números separados por coma}
%     |\spmcm|\myoarg{key \textnormal{\textcolor{gray}{=}} val}\myparg{números separados por coma}
%   \end{syntax}
%   El comando \ics*{spmcm} trabaja en \emph{modo matemático} y recibe el
%   \emph{argumento delimitado por paréntesis} \myparg{números separados
%   por coma}, el cual debe ser una lista de \emph{números naturales},
%   por ejemplo: `|(2, 3, 6)|'.
%   Si se usa sin \myparg{argumento} |$\spmcm$| imprimirá `$\mathrm{mcm}$' en
%   la salida \texttt{PDF}, y \mynvmath{NVDA/MathCAT} verbalizará \emph{«mínimo
%   común múltiplo»} o \emph{«m c m»}, acorde al \mymarg{valor} establecido por
%   la llave |read-cmd|.
%
%   \smallskip
%
%   Si se usa con \myparg{argumento} |$\spmcm(2, 3, 6)$| imprimirá
%   `$\mathrm{mcm}(2, 3, 6)$' en la salida \texttt{PDF}, y
%   \mynvmath{NVDA/MathCAT} verbalizará \emph{«mínimo común múltiplo entre dos,
%   tres y seis»} o \emph{«m c m entre dos, tres y seis»} acorde al
%   \mymarg{valor} establecido por la llave |read-cmd|.
% \end{function}
%
% \subsubsection*{Llaves para \cs[no-index]{spmcm}}
%
% \keyexampcmd{upper-case}{true \textnormal{%
%   \textcolor{lightgray}{\textbar}} false}{false}[spmcm]
% \smallskip
% Establece si se usarán \emph{mayúsculas} o \emph{minúsculas} en la salida
% \texttt{PDF} para el comando |\spmcm|. Si se establece en |true| se
% imprimirá `$\mathrm{MCM}$' en la salida \texttt{PDF}; si se establece en
% |false| se imprimirá `$\mathrm{mcm}$' en la salida \texttt{PDF}.
%
% \smallskip
%
% \keyexampcmd{read-cmd}{terse \textnormal{%
%   \textcolor{lightgray}{\textbar}} clear}{clear}[spmcm]
% \smallskip
% Establece la \emph{forma} en la que \mynvmath{NVDA/MathCAT} verbalizará el
% comando |\spmcm|. Si se establece en |terse|, \mynvmath{NVDA/MathCAT}
% verbalizará \emph{«m c m»}; si se establece en |clear|,
% \mynvmath{NVDA/MathCAT} verbalizará \emph{«mínimo común múltiplo»}.
%
% \smallskip
%
% \keyexampcmd{prefix}{prefijo}{no usado}[spmcm]
% \smallskip
% Establece el \emph{prefijo} que será verbalizado por \mynvmath{NVDA/MathCAT}
% \emph{«antes»} de verbalizar el valor establecido por la llave \mykey[spmcm]{read-cmd}.
% El valor de esta llave \emph{siempre} debe pasarse entre llaves `|{ }|'.
%
% \smallskip
%
% \keyexampcmd{result}{true \textnormal{%
%   \textcolor{lightgray}{\textbar}} false}{false}[spmcm]
% \smallskip
% Establece si se \emph{imprime} el resultado de calcular el \emph{«mínimo
% común múltiplo»} del \myparg{argumento} pasado a |\spmcm|. Si se
% establece en |true|: |$\spmcm[result=true](2, 3, 6)$| imprimirá
% `$\mathrm{mcm}(2, 3, 6) = 6$' en la salida \texttt{PDF} y
% \mynvmath{NVDA/MathCAT} verbalizará \emph{«mínimo común múltiplo entre
% dos, tres y seis es igual a seis»}; si se establece en |false| no se
% imprimirá en la salida \texttt{PDF} el resultado de calcular el
% \emph{«mínimo común múltiplo»}.
%
% \smallskip
%
% \keyexampcmd{sample-arg}{true \textnormal{%
%   \textcolor{lightgray}{\textbar}} false}{false}[spmcm]
% \smallskip
%
% Desactiva la \emph{validación} interna del \myparg{argumento} pasado
% a |\spmcm|, permitiendo usar ejemplos \emph{«no numéricos»} como entrada.
% Si se establece en |true|: |$\spmcm[sample-arg=true](m, n)$|, imprimirá
% `$\mathrm{mcm}(m, n)$' en la salida \texttt{PDF} y \mynvmath{NVDA/MathCAT}
% verbalizará \emph{«mínimo común múltiplo entre m y n»}; si se establece
% en |false| el \myparg{argumento} debe ser una lista de \emph{números
% naturales}, en caso contrario se devolverá un \enquote{error}.
%
% \smallskip
%
% \keyexampcmd{silent-cmd}{true \textnormal{%
%   \textcolor{lightgray}{\textbar}} false}{false}[spmcm]
% \smallskip
% Establece si se \emph{silencia} la verbalización en \mynvmath{NVDA/MathCAT}
% de |\spmcm|\myparg{argumento} cuando se establece |sample-arg=true|
% o cuando |\spmcm| se usa sin \myparg{argumento}. Si se establece en |true|
% \mynvmath{NVDA/MathCAT} NO lo verbalizará; si se establece en |false|
% \mynvmath{NVDA/MathCAT} SÍ lo verbalizará.
%
%
% \subsection*{El comando \cs[no-index]{spnZ}}\label{subsec:spnZ}
%
% \vspace*{-\baselineskip}
%
% \begin{function}{\spnZ}
%   \begin{syntax}
%     \ics*{spnZ}\mymarg{número entero}
%     \ics*{spnZ}\myoarg{key \textnormal{\textcolor{gray}{=}} val}\mymarg{número entero}
%   \end{syntax}
%   El comando \ics*{spnZ} trabaja en \emph{modo matemático} y recibe
%   el \emph{argumento} obligatorio \mymarg{número entero} el cual se procesa
%   internamente bajo el comando |\spnum|.
%   Este comando está dedicado exclusivamente al manejo de \emph{«números enteros»},
%   y añade algunas utilidades prácticas para el demarcado e
%   interacción con \mynvmath{NVDA/MathCAT}.
% \end{function}
%
% \subsubsection*{Llaves para \cs[no-index]{spnZ}}
%
% El comando |\spnZ| posee las \mymeta{meta-keys} \mykey[spnZ]{cifras-sep}
% y \mykey[spnZ]{read-sign} pasadas al comando |\spnum| (\S\ref{subsec:key-spnum}).
%
% \smallskip
%
% \keyexampcmd{wrap-parens}{true \textnormal{%
%   \textcolor{lightgray}{\textbar}} false}{false}[spnZ]
% \smallskip
% Establece si se \emph{imprimen} paréntesis redondos `|(|\dots|)|'
% alrededor del argumento \mymarg{número entero} pasado a |\spnZ|. Si se
% establece en |true|, se imprimirán paréntesis redondos alrededor del
% \mymarg{número entero}; si se establece en |false|, no se imprimirán
% los paréntesis.
%
% \keyexampcmd{read-parens}{true \textnormal{%
%   \textcolor{lightgray}{\textbar}} false}{false}[spnZ]
% \smallskip
% Establece si \mynvmath{NVDA/MathCAT} \emph{verbalizará} los paréntesis redondos
% `|(|\dots|)|' alrededor del \emph{argumento} \mymarg{número entero} pasado a
% |\spnZ|. Si se establece en |true|, \mynvmath{NVDA/MathCAT} verbalizará la lectura
% del paréntesis `|(|\dots|)|'; si se establece en |false|, \mynvmath{NVDA/MathCAT}
% no los verbalizará.
%
% \subsection*{El comando \cs[no-index]{spmQ}}\label{subsec:spmQ}
%
% \vspace*{-\baselineskip}
%
% \begin{function}{\spmQ}
%   \begin{syntax}
%     \ics*{spmQ}\mymarg{número entero}\mymarg{numerador}\mymarg{denominador}
%     \ics*{spmQ}\myoarg{key \textnormal{\textcolor{gray}{=}} val}\mymarg{número entero}\mymarg{numerador}\mymarg{denominador}
%   \end{syntax}
%   El comando \ics*{spmQ} trabaja en \emph{modo matemático} y recibe
%   \emph{tres argumentos obligatorios}: \mymarg{número entero},
%   \mymarg{numerador} y \mymarg{denominador}, los procesa internamente
%   y devuelve un \emph{«número mixto»}.
%
%   \smallskip
%
%   Si no se proporciona el \emph{argumento opcional}
%   \myoarg{key \textnormal{\textcolor{gray}{=}} val},
%   |$\spmQ{2}{5}{7}$| imprimirá `$\scalerel{2}{\frac{5}{7}}$'
%   en la salida \texttt{PDF} y \mynvmath{NVDA/MathCAT} verbalizará
%   \emph{«dos y cinco séptimos»}.
% \end{function}
%
% \smallskip
%
% \begin{important}*
%   En el caso de las fracciones, y en especial con los \emph{«números
%   mixtos»}, no se pueden manejar muchas cosas desde el lado de
%   \mypkg*[no-index]{spintent}. Por ejemplo, la entrada |$3\frac{5}{7}$|
%   con \emph{tagged} \texttt{PDF} se verbaliza por \mynvmath{NVDA/MathCAT}
%   como \emph{«tres más cinco séptimos»}, lo cual no está mal, pero
%   pedagógicamente hablando no es del todo correcto.
% \end{important}
%
% \subsubsection*{Llaves para \cs[no-index]{spmQ}}
%
% \keyexampcmd{join-word}{entero \textnormal{\textcolor{lightgray}{\textbar}}
%   enteros \textnormal{\textcolor{lightgray}{\textbar}} y}{y}[spmQ]
% \smallskip
% Establece la \emph{forma} en la que \mynvmath{NVDA/MathCAT} verbaliza la
% conexión entre el entero y la fracción que forman el \emph{número mixto}.
% Si se establece en |entero|, al usar |$\spmQ[join-word=entero]{1}{5}{7}$|
% imprimirá `$\scalerel{1}{\frac{5}{7}}$' en la salida \texttt{PDF}, y
% \mynvmath{NVDA/MathCAT} verbalizará \emph{«un entero cinco séptimos»};
% si se establece en |enteros|, al usar |$\spmQ[join-word=enteros]{2}{5}{7}$|
% imprimirá `$\scalerel{2}{\frac{5}{7}}$' en la salida \texttt{PDF}, y
% \mynvmath{NVDA/MathCAT} verbalizará \emph{«dos enteros cinco séptimos»};
% si se establece en |y|, al usar |$\spmQ[join-word=y]{2}{5}{7}$|
% imprimirá `$\scalerel{2}{\frac{5}{7}}$' en la salida \texttt{PDF}, y
% \mynvmath{NVDA/MathCAT} verbalizará \emph{«dos y cinco séptimos»}.
%
% \smallskip
%
% \keyexampcmd{use-spnZ}{true \textnormal{\textcolor{lightgray}{\textbar}}
%   false}{true}[spmQ]
% \smallskip
% Establece si el \mymarg{numerador} y el \mymarg{denominador} que forman
% la fracción del \emph{número mixto} se procesan internamente bajo el
% comando |\spnZ|, añadiendo así \emph{semántica de entero} en la
% estructura \mynvmath{MathML} de la salida \texttt{PDF}. Si se establece
% en |true|, \mymarg{numerador} y \mymarg{denominador} se pasarán a |\spnZ|
% y \mynvmath{NVDA/MathCAT} los verbalizará correctamente como
% \emph{números ordinales}; si se establece en |false|, los argumentos
% se tipografían directamente sin procesamiento semántico adicional.
%
% \smallskip
%
% \keyexampcmd{print-dfrac}{true \textnormal{\textcolor{lightgray}{\textbar}}
%   false}{false}[spmQ]
% \smallskip
% Establece si la fracción del \emph{número mixto} se imprime usando
% |\dfrac| en lugar de |\frac|. Si se establece en |true|, al usar
% |$\spmQ[print-dfrac=true]{2}{5}{7}$| imprimirá la fracción en tamaño
% \emph{display} `$\scalerel{2}{\dfrac{5}{7}}$' en la salida \texttt{PDF};
% si se establece en |false|, la fracción se imprime con |\frac| siguiendo
% el tamaño matemático del contexto en curso.
%
% \smallskip
%
% \keyexampcmd{wrap-parens}{true \textnormal{\textcolor{lightgray}{\textbar}}
%   false}{false}[spmQ]
% \smallskip
% Establece si se \emph{imprimen} paréntesis redondos `|(|\dots|)|'
% alrededor del \emph{número mixto} completo devuelto por |\spmQ|. Si se
% establece en |true|, se imprimirán paréntesis redondos cuyo tamaño se
% ajusta automáticamente al estilo matemático en curso; si se establece en
% |false|, no se imprimirán los paréntesis.
%
% \smallskip
%
% \keyexampcmd{read-parens}{true \textnormal{\textcolor{lightgray}{\textbar}}
%   false}{true}[spmQ]
% \smallskip
% Establece si \mynvmath{NVDA/MathCAT} \emph{verbalizará} los paréntesis
% redondos `|(|\dots|)|' alrededor del \emph{número mixto} cuando la llave
% \mykey[spmQ]{wrap-parens} está activa. Si se establece en |true|,
% \mynvmath{NVDA/MathCAT} verbalizará la lectura de los paréntesis; si se
% establece en |false|, \mynvmath{NVDA/MathCAT} no los verbalizará.
%
%^^A \section{Utilidades para geometría}
%
%^^A El paquete \mypkg*[no-index]{spintent} provee utilidades para trabajar con geometría
%^^A (plana), que describiremos en esta sección.
%
% \section{Utilidades para geometría plana}
%
% El paquete \mypkg*[no-index]{spintent} provee comandos para trabajar
% con los objetos y medidas de la \emph{geometría plana} descritos en
% esta sección. Todos los comandos trabajan en \emph{modo matemático} y
% comparten el mismo módulo \myluamod{} para construir la
% verbalización correcta.
%
% \smallskip
%
% \begin{important}*
%   Todos los comandos de esta sección aceptan como puntos solo
%   \emph{letras mayúsculas} como base, con modificadores opcionales de
%   subíndice `|_{...}|' o prima `|'|'. La única excepción es
%   |\sprecta| en \emph{notación de nombre}, que acepta también
%   \emph{letras minúsculas}. Los subíndices deben escribirse siempre
%   con la notación |A_{...}| y no con |A_n| sin llaves.
% \end{important}
%
% \subsection*{El comando \cs[no-index]{spPoint}}\label{subsec:spPoint}
%
% \vspace*{-\baselineskip}
%
% \begin{function}{\spPoint}
%   \begin{syntax}
%     \ics*{spPoint}\mymarg{punto}
%     \ics*{spPoint}\myoarg{key \textnormal{\textcolor{gray}{=}} val}\mymarg{punto}
%   \end{syntax}
%   El comando \ics*{spPoint} trabaja en \emph{modo matemático} y
%   recibe en \mymarg{punto} una \emph{única letra mayúscula} con
%   modificador opcional de subíndice `|_{...}|' o prima `|'|'.
%   Representa un punto geométrico como entidad abstracta, sin
%   coordenadas.
%
%   \smallskip
%
%   Si no se proporciona el \emph{argumento opcional}
%   \myoarg{key \textnormal{\textcolor{gray}{=}} val},
%   |$\spPoint{A}$| imprimirá `$A$' en la salida \texttt{PDF} y
%   \mynvmath{NVDA/MathCAT} verbalizará \emph{«punto A»};
%   |$\spPoint{A_{1}}$| imprimirá `$A_1$' y verbalizará
%   \emph{«punto A sub uno»}.
% \end{function}
%
% \smallskip
%
% \begin{important}*
%   El comando |\spPoint| solo acepta \emph{un punto}. Para pares o
%   listas de puntos (rectas, rayos, lados, ángulos, arcos, polígonos),
%   véanse los comandos correspondientes de esta sección.
% \end{important}
%
% \subsubsection*{Llaves para \cs[no-index]{spPoint}}
%
% \keyexampcmd{read-cmd}{school \textnormal{\textcolor{lightgray}{\textbar}}
%   mathcat}{school}[spPoint]
% \smallskip
% Establece la forma en la que \mynvmath{NVDA/MathCAT} verbalizará el
% punto. Si se establece en |school|, \mynvmath{NVDA/MathCAT} verbaliza
% la letra directamente sin anteponer «mayúscula»: |$\spPoint{A}$| se
% verbaliza \emph{«punto A»}; si se establece en |mathcat|,
% \mynvmath{NVDA/MathCAT} aplica sus propias reglas de habla, que
% pueden incluir el prefijo «mayúscula» según la configuración del
% usuario.
%
% \smallskip
%
% \keyexampcmd{prefix}{prefijo}{no usado}[spPoint]
% \smallskip
% Establece la palabra que \mynvmath{NVDA/MathCAT} verbalizará
% \emph{antes} de «punto». El valor de esta llave \emph{siempre} debe
% pasarse entre llaves `|{ }|'. Si se establece en |{el}|,
% \mynvmath{NVDA/MathCAT} verbalizará \emph{«el punto A»} en lugar de
% \emph{«punto A»}; si no se establece, no se antepone ninguna palabra.
%
% \smallskip
%
% \keyexampcmd{silent-cmd}{true \textnormal{\textcolor{lightgray}{\textbar}}
%   false}{false}[spPoint]
% \smallskip
% Establece si \mynvmath{NVDA/MathCAT} \emph{silenciará} la palabra
% «punto». Útil cuando el nombre del objeto ya aparece en el texto que
% rodea a la expresión matemática. Con |silent-cmd=true|,
% \mynvmath{NVDA/MathCAT} verbalizará solo \emph{«A»} en lugar de
% \emph{«punto A»}.
%
% \subsection*{El comando \cs[no-index]{sprecta}}\label{subsec:sprecta}
%
% \vspace*{-\baselineskip}
%
% \begin{function}{\sprecta}
%   \begin{syntax}
%     \ics*{sprecta}\mymarg{puntos}
%     \ics*{sprecta}\myoarg{key \textnormal{\textcolor{gray}{=}} val}\mymarg{puntos}
%   \end{syntax}
%   El comando \ics*{sprecta} trabaja en \emph{modo matemático} y
%   admite \emph{dos notaciones} detectadas automáticamente según el
%   contenido de \mymarg{puntos}:
%
%   \smallskip
%
%   \emph{Notación de dos puntos}: cuando \mymarg{puntos} contiene una
%   coma, se aplica \cs{overleftrightarrow} sobre los puntos. Por
%   ejemplo, |$\sprecta{A, B}$| imprimirá `$\overleftrightarrow{AB}$'
%   y \mynvmath{NVDA/MathCAT} verbalizará \emph{«recta A B»}.
%
%   \smallskip
%
%   \emph{Notación de nombre}: cuando \mymarg{puntos} no contiene
%   coma, se imprime el nombre sin ningún decorador visual. Por
%   ejemplo, |$\sprecta{m}$| imprimirá `$m$' y \mynvmath{NVDA/MathCAT}
%   verbalizará \emph{«recta m»}; |$\sprecta{L_{1}}$| imprimirá
%   `$L_{1}$' y verbalizará \emph{«recta L sub uno»}.
% \end{function}
%
% \smallskip
%
% \begin{important}*
%   En la \emph{notación de nombre} se aceptan tanto letras mayúsculas
%   como minúsculas ($m$, $r$, $l$, $s$, $L_1$, $L_2$). En esta
%   notación, con la llave |read-cmd=school| \mynvmath{NVDA/MathCAT}
%   siempre verbaliza sin «mayúscula» ya que el nombre forma parte del
%   texto verbalizado directamente.
% \end{important}
%
% \subsubsection*{Llaves para \cs[no-index]{sprecta}}
%
% \keyexampcmd{read-cmd}{school \textnormal{\textcolor{lightgray}{\textbar}}
%   mathcat}{school}[sprecta]
% \smallskip
% Establece la forma en la que \mynvmath{NVDA/MathCAT} verbalizará los
% puntos. Si se establece en |school|, \mynvmath{NVDA/MathCAT} verbaliza
% las letras directamente sin anteponer «mayúscula»: |$\sprecta{A, B}$|
% se verbaliza \emph{«recta A B»}; si se establece en |mathcat|,
% \mynvmath{NVDA/MathCAT} aplica sus propias reglas de habla, que
% pueden incluir el prefijo «mayúscula» según la configuración del
% usuario. Esta llave no afecta la \emph{notación de nombre}.
%
% \smallskip
%
% \keyexampcmd{prefix}{prefijo}{no usado}[sprecta]
% \smallskip
% Establece la palabra que \mynvmath{NVDA/MathCAT} verbalizará
% \emph{antes} de «recta». El valor de esta llave \emph{siempre} debe
% pasarse entre llaves `|{ }|'. Si se establece en |{la}|,
% \mynvmath{NVDA/MathCAT} verbalizará \emph{«la recta A B»} en lugar de
% \emph{«recta A B»}; si no se establece, no se antepone ninguna
% palabra.
%
% \smallskip
%
% \keyexampcmd{silent-cmd}{true \textnormal{\textcolor{lightgray}{\textbar}}
%   false}{false}[sprecta]
% \smallskip
% Establece si \mynvmath{NVDA/MathCAT} \emph{silenciará} la palabra
% «recta». Útil cuando el nombre del objeto ya aparece en el texto que
% rodea a la expresión matemática. Con |silent-cmd=true|,
% \mynvmath{NVDA/MathCAT} verbalizará solo los puntos o el nombre, sin
% anteponer «recta».
%
% \subsection*{El comando \cs[no-index]{sprayo}}\label{subsec:sprayo}
%
% \vspace*{-\baselineskip}
%
% \begin{function}{\sprayo}
%   \begin{syntax}
%     \ics*{sprayo}\mymarg{puntos}
%     \ics*{sprayo}\myoarg{key \textnormal{\textcolor{gray}{=}} val}\mymarg{puntos}
%   \end{syntax}
%   El comando \ics*{sprayo} trabaja en \emph{modo matemático} y
%   representa un \emph{rayo} con origen en el primer punto y dirección
%   hacia el segundo. La salida visual usa \cs{overrightarrow}.
%
%   \smallskip
%
%   Si no se proporciona el \emph{argumento opcional}
%   \myoarg{key \textnormal{\textcolor{gray}{=}} val},
%   |$\sprayo{A, B}$| imprimirá `$\overrightarrow{AB}$' en la salida
%   \texttt{PDF} y \mynvmath{NVDA/MathCAT} verbalizará
%   \emph{«rayo A B»}.
% \end{function}
%
% \smallskip
%
% \begin{important}*
%   El rayo $\overrightarrow{AB}$ y el rayo $\overrightarrow{BA}$ son
%   objetos distintos — el orden de los puntos determina la dirección.
% \end{important}
%
% \subsubsection*{Llaves para \cs[no-index]{sprayo}}
%
% \keyexampcmd{read-cmd}{school \textnormal{\textcolor{lightgray}{\textbar}}
%   mathcat}{school}[sprayo]
% \smallskip
% Establece la forma en la que \mynvmath{NVDA/MathCAT} verbalizará los
% puntos del rayo. Si se establece en |school|, \mynvmath{NVDA/MathCAT}
% verbaliza las letras directamente sin anteponer «mayúscula»: si se
% establece en |mathcat|, \mynvmath{NVDA/MathCAT} aplica sus propias
% reglas de habla, que pueden incluir el prefijo «mayúscula» según la
% configuración del usuario.
%
% \smallskip
%
% \keyexampcmd{prefix}{prefijo}{no usado}[sprayo]
% \smallskip
% Establece la palabra que \mynvmath{NVDA/MathCAT} verbalizará
% \emph{antes} del nombre del rayo. El valor de esta llave
% \emph{siempre} debe pasarse entre llaves `|{ }|'. Si se establece en
% |{el}|, \mynvmath{NVDA/MathCAT} verbalizará \emph{«el rayo A B»} en
% lugar de \emph{«rayo A B»}; si no se establece, no se antepone
% ninguna palabra.
%
% \smallskip
%
% \keyexampcmd{silent-cmd}{true \textnormal{\textcolor{lightgray}{\textbar}}
%   false}{false}[sprayo]
% \smallskip
% Establece si \mynvmath{NVDA/MathCAT} \emph{silenciará} el nombre del
% rayo. Útil cuando el objeto ya se ha mencionado en el texto. Con
% |silent-cmd=true|, \mynvmath{NVDA/MathCAT} verbalizará solo
% \emph{«A B»} en lugar de \emph{«rayo A B»}.
%
% \smallskip
%
% \keyexampcmd{read-sty}{rayo \textnormal{\textcolor{lightgray}{\textbar}}
%   semirrecta}{rayo}[sprayo]
% \smallskip
% Establece el \emph{término} que \mynvmath{NVDA/MathCAT} verbalizará
% para este objeto. Si se establece en |rayo| (valor por defecto),
% |$\sprayo{A, B}$| verbaliza \emph{«rayo A B»}; si se establece en
% |semirrecta|, |$\sprayo[read-sty=semirrecta]{A, B}$| verbaliza
% \emph{«semirrecta A B»}. La salida visual (`$\overrightarrow{AB}$')
% es idéntica en ambos casos. La elección depende de la terminología
% del libro de texto o del currículo nacional: en Chile y España se
% prefiere «semirrecta»; en México y otros países de América Latina se
% usa «rayo».
%
% \subsection*{El comando \cs[no-index]{spside}}\label{subsec:spside}
%
% \vspace*{-\baselineskip}
%
% \begin{function}{\spside}
%   \begin{syntax}
%     \ics*{spside}\mymarg{puntos}
%     \ics*{spside}\myoarg{key \textnormal{\textcolor{gray}{=}} val}\mymarg{puntos}
%   \end{syntax}
%   El comando \ics*{spside} trabaja en \emph{modo matemático} y
%   representa un \emph{lado}, \emph{segmento} o \emph{trazo} —según la
%   llave \mykey[spside]{read-sty}— determinado por dos puntos. La
%   salida visual usa \cs{overline}.
%
%   \smallskip
%
%   Si no se proporciona el \emph{argumento opcional}
%   \myoarg{key \textnormal{\textcolor{gray}{=}} val},
%   |$\spside{A, B}$| imprimirá `$\overline{AB}$' en la salida
%   \texttt{PDF} y \mynvmath{NVDA/MathCAT} verbalizará
%   \emph{«lado A B»}.
% \end{function}
%
% \subsubsection*{Llaves para \cs[no-index]{spside}}
%
% \keyexampcmd{read-cmd}{school \textnormal{\textcolor{lightgray}{\textbar}}
%   mathcat}{school}[spside]
% \smallskip
% Establece la forma en la que \mynvmath{NVDA/MathCAT} verbalizará los
% puntos. Si se establece en |school|, \mynvmath{NVDA/MathCAT} verbaliza
% las letras directamente sin anteponer «mayúscula»; si se establece en
% |mathcat|, \mynvmath{NVDA/MathCAT} aplica sus propias reglas de
% habla, que pueden incluir el prefijo «mayúscula» según la
% configuración del usuario.
%
% \smallskip
%
% \keyexampcmd{prefix}{prefijo}{no usado}[spside]
% \smallskip
% Establece la palabra que \mynvmath{NVDA/MathCAT} verbalizará
% \emph{antes} del nombre del objeto. El valor de esta llave
% \emph{siempre} debe pasarse entre llaves `|{ }|'. Si se establece en
% |{el}|, \mynvmath{NVDA/MathCAT} verbalizará \emph{«el lado A B»} en
% lugar de \emph{«lado A B»}; si no se establece, no se antepone
% ninguna palabra.
%
% \smallskip
%
% \keyexampcmd{silent-cmd}{true \textnormal{\textcolor{lightgray}{\textbar}}
%   false}{false}[spside]
% \smallskip
% Establece si \mynvmath{NVDA/MathCAT} \emph{silenciará} el nombre del
% objeto. Útil cuando el objeto ya se ha mencionado en el texto. Con
% |silent-cmd=true|, \mynvmath{NVDA/MathCAT} verbalizará solo
% \emph{«A B»} en lugar de \emph{«lado A B»}.
%
% \smallskip
%
% \keyexampcmd{read-sty}{lado \textnormal{\textcolor{lightgray}{\textbar}}
%   segmento \textnormal{\textcolor{lightgray}{\textbar}}
%   trazo}{lado}[spside]
% \smallskip
% Establece el \emph{término} que \mynvmath{NVDA/MathCAT} verbalizará
% para el objeto. Si se establece en |lado| (valor por defecto),
% |$\spside{A, B}$| verbaliza \emph{«lado A B»} — apropiado cuando el
% segmento es parte de una figura geométrica, como los lados de un
% triángulo o polígono. Si se establece en |segmento|,
% |$\spside[read-sty=segmento]{A, B}$| verbaliza \emph{«segmento A B»}
% — apropiado cuando el segmento se trata como objeto geométrico
% independiente. Si se establece en |trazo|,
% |$\spside[read-sty=trazo]{A, B}$| verbaliza \emph{«trazo A B»} —
% terminología habitual en construcciones con regla y compás. La salida
% visual (`$\overline{AB}$') es idéntica en los tres casos.
%
% \smallskip
%
% \keyexampcmd{read-arg}{nombre}{(según read-sty)}[spside]
% \smallskip
% Sobrescribe por completo el \emph{término} verbalizado, ignorando el
% valor de \mykey[spside]{read-sty}. El valor de esta llave
% \emph{siempre} debe pasarse entre llaves `|{ }|'. Por ejemplo,
% |$\spside[read-arg={hipotenusa}]{A, B}$| imprimirá `$\overline{AB}$'
% y \mynvmath{NVDA/MathCAT} verbalizará \emph{«hipotenusa A B»}.
%
% \subsection*{El comando \cs[no-index]{spmside}}\label{subsec:spmside}
%
% \vspace*{-\baselineskip}
%
% \begin{function}{\spmside}
%   \begin{syntax}
%     \ics*{spmside}\mymarg{puntos}
%     \ics*{spmside}\myoarg{key \textnormal{\textcolor{gray}{=}} val}\mymarg{puntos}
%   \end{syntax}
%   El comando \ics*{spmside} trabaja en \emph{modo matemático} y
%   representa la \emph{medida} del lado, segmento o trazo determinado
%   por dos puntos. La distinción entre |\spside| y |\spmside| es
%   semántica: |\spside| representa el \emph{objeto geométrico}
%   ($\overline{AB}$); |\spmside| representa su \emph{medida} (un
%   número real positivo).
%
%   \smallskip
%
%   Si no se proporciona el \emph{argumento opcional}
%   \myoarg{key \textnormal{\textcolor{gray}{=}} val},
%   |$\spmside{A, B}$| imprimirá `$AB$' en la salida \texttt{PDF} y
%   \mynvmath{NVDA/MathCAT} verbalizará \emph{«medida del lado A B»}.
% \end{function}
%
% \subsubsection*{Llaves para \cs[no-index]{spmside}}
%
% \keyexampcmd{read-cmd}{school \textnormal{\textcolor{lightgray}{\textbar}}
%   mathcat}{school}[spmside]
% \smallskip
% Establece la forma en la que \mynvmath{NVDA/MathCAT} verbalizará los
% puntos. Si se establece en |school|, \mynvmath{NVDA/MathCAT} verbaliza
% las letras directamente sin anteponer «mayúscula»; si se establece en
% |mathcat|, \mynvmath{NVDA/MathCAT} aplica sus propias reglas de
% habla, que pueden incluir el prefijo «mayúscula» según la
% configuración del usuario.
%
% \smallskip
%
% \keyexampcmd{prefix}{prefijo}{no usado}[spmside]
% \smallskip
% Establece la palabra que \mynvmath{NVDA/MathCAT} verbalizará
% \emph{antes} de «medida del ...». El valor de esta llave
% \emph{siempre} debe pasarse entre llaves `|{ }|'. Si se establece en
% |{la}|, \mynvmath{NVDA/MathCAT} verbalizará \emph{«la medida del
% lado A B»} en lugar de \emph{«medida del lado A B»}; si no se
% establece, no se antepone ninguna palabra.
%
% \smallskip
%
% \keyexampcmd{silent-cmd}{true \textnormal{\textcolor{lightgray}{\textbar}}
%   false}{false}[spmside]
% \smallskip
% Establece si \mynvmath{NVDA/MathCAT} \emph{silenciará} el nombre del
% objeto. Útil cuando el objeto ya se ha mencionado en el texto. Con
% |silent-cmd=true|, \mynvmath{NVDA/MathCAT} verbalizará solo
% \emph{«A B»} en lugar de \emph{«medida del lado A B»}.
%
% \smallskip
%
% \keyexampcmd{read-sty}{lado \textnormal{\textcolor{lightgray}{\textbar}}
%   segmento \textnormal{\textcolor{lightgray}{\textbar}}
%   trazo}{lado}[spmside]
% \smallskip
% Establece el \emph{término} verbalizado, con la misma semántica que
% en |\spside| (\S\ref{subsec:spside}). Si se establece en |segmento|,
% |$\spmside[read-sty=segmento]{A, B}$| verbaliza \emph{«medida del
% segmento A B»}; si se establece en |trazo|, verbaliza \emph{«medida
% del trazo A B»}.
%
% \smallskip
%
% \keyexampcmd{read-arg}{nombre}{(según read-sty)}[spmside]
% \smallskip
% Sobrescribe por completo el \emph{término} verbalizado, con la misma
% semántica que en |\spside| (\S\ref{subsec:spside}). El valor de esta
% llave \emph{siempre} debe pasarse entre llaves `|{ }|'.
%
% \smallskip
%
% \keyexampcmd{print-m}{true \textnormal{\textcolor{lightgray}{\textbar}}
%   false}{false}[spmside]
% \smallskip
% Establece la \emph{notación visual} de la medida. Si se establece en
% |true|, |$\spmside[print-m=true]{A, B}$| imprimirá
% `$\operatorname{m}\overline{AB}$' en la salida \texttt{PDF},
% notación habitual en libros de texto que prefieren hacer explícita la
% referencia al objeto; si se establece en |false| (valor por
% defecto), imprimirá simplemente `$AB$'. En ambos casos
% \mynvmath{NVDA/MathCAT} verbalizará \emph{«medida del lado A B»}.
%
% \subsection*{El comando \cs[no-index]{spangle}}\label{subsec:spangle}
%
% \vspace*{-\baselineskip}
%
% \begin{function}{\spangle}
%   \begin{syntax}
%     \ics*{spangle}\mymarg{arg}
%     \ics*{spangle}\myoarg{key \textnormal{\textcolor{gray}{=}} val}\mymarg{arg}
%   \end{syntax}
%   El comando \ics*{spangle} trabaja en \emph{modo matemático} y
%   admite \emph{tres notaciones} detectadas automáticamente:
%
%   \smallskip
%
%   \emph{Notación de tres puntos}: cuando \mymarg{arg} contiene una
%   coma (p.~ej.\ |A, O, B|), el segundo punto es el \emph{vértice}.
%   Por ejemplo, |$\spangle{A, O, B}$| imprimirá `$\angle AOB$' y
%   \mynvmath{NVDA/MathCAT} verbalizará \emph{«ángulo A O B»}.
%
%   \smallskip
%
%   \emph{Notación de un punto}: cuando \mymarg{arg} es una sola letra
%   mayúscula (p.~ej.\ |A| o |A_{1}|), representa el \emph{ángulo en
%   ese vértice}. Por ejemplo, |$\spangle{A}$| imprimirá `$\angle A$'
%   y \mynvmath{NVDA/MathCAT} verbalizará \emph{«ángulo en A»}.
%
%   \smallskip
%
%   \emph{Notación de nombre}: cuando \mymarg{arg} no contiene coma y
%   no es una sola letra mayúscula (p.~ej.\ |\alpha|, |1|, |ABC|). Por
%   ejemplo, |$\spangle{\alpha}$| imprimirá `$\angle\alpha$' y
%   \mynvmath{NVDA/MathCAT} verbalizará \emph{«ángulo alfa»}.
% \end{function}
%
% \subsubsection*{Llaves para \cs[no-index]{spangle}}
%
% \keyexampcmd{read-cmd}{school \textnormal{\textcolor{lightgray}{\textbar}}
%   mathcat}{school}[spangle]
% \smallskip
% Establece la forma en la que \mynvmath{NVDA/MathCAT} verbalizará los
% puntos del ángulo. Si se establece en |school|, \mynvmath{NVDA/MathCAT}
% verbaliza las letras directamente sin anteponer «mayúscula»; si se
% establece en |mathcat|, \mynvmath{NVDA/MathCAT} aplica sus propias
% reglas de habla, que pueden incluir el prefijo «mayúscula» según la
% configuración del usuario.
%
% \smallskip
%
% \keyexampcmd{prefix}{prefijo}{no usado}[spangle]
% \smallskip
% Establece la palabra que \mynvmath{NVDA/MathCAT} verbalizará
% \emph{antes} de «ángulo». El valor de esta llave \emph{siempre} debe
% pasarse entre llaves `|{ }|'. Si se establece en |{el}|,
% \mynvmath{NVDA/MathCAT} verbalizará \emph{«el ángulo A O B»} en lugar
% de \emph{«ángulo A O B»}; si no se establece, no se antepone ninguna
% palabra.
%
% \smallskip
%
% \keyexampcmd{silent-cmd}{true \textnormal{\textcolor{lightgray}{\textbar}}
%   false}{false}[spangle]
% \smallskip
% Establece si \mynvmath{NVDA/MathCAT} \emph{silenciará} la palabra
% «ángulo». Útil cuando el objeto ya se ha mencionado en el texto. Con
% |silent-cmd=true|, \mynvmath{NVDA/MathCAT} verbalizará solo
% \emph{«A O B»} en lugar de \emph{«ángulo A O B»}.
%
% \smallskip
%
% \keyexampcmd{recto}{true \textnormal{\textcolor{lightgray}{\textbar}}
%   false}{false}[spangle]
% \smallskip
% Establece si el ángulo es un \emph{ángulo recto}, cambiando el
% símbolo visual y la verbalización. Si se establece en |true|,
% |$\spangle[recto=true]{A}$| imprimirá `$\rightangle A$' en la salida
% \texttt{PDF} y \mynvmath{NVDA/MathCAT} verbalizará \emph{«ángulo
% recto en A»}; si se establece en |false|, imprimirá `$\angle A$' y
% verbalizará \emph{«ángulo en A»}. La misma llave aplica a la
% notación de tres puntos: |$\spangle[recto=true]{A, O, B}$| imprimirá
% `$\rightangle AOB$' y verbalizará \emph{«ángulo recto A O B»}.
%
% \subsection*{El comando \cs[no-index]{spmangle}}\label{subsec:spmangle}
%
% \vspace*{-\baselineskip}
%
% \begin{function}{\spmangle}
%   \begin{syntax}
%     \ics*{spmangle}\mymarg{arg}
%     \ics*{spmangle}\myoarg{key \textnormal{\textcolor{gray}{=}} val}\mymarg{arg}
%   \end{syntax}
%   El comando \ics*{spmangle} trabaja en \emph{modo matemático} y
%   representa la \emph{medida de un ángulo}. Admite las mismas tres
%   notaciones que |\spangle| (\S\ref{subsec:spangle}).
%
%   \smallskip
%
%   Si no se proporciona el \emph{argumento opcional}
%   \myoarg{key \textnormal{\textcolor{gray}{=}} val},
%   |$\spmangle{A, O, B}$| imprimirá `$\measuredangle AOB$' en la
%   salida \texttt{PDF} y \mynvmath{NVDA/MathCAT} verbalizará
%   \emph{«medida del ángulo A O B»}; |$\spmangle{A}$| imprimirá
%   `$\measuredangle A$' y verbalizará \emph{«medida del ángulo en
%   A»}.
% \end{function}
%
% \subsubsection*{Llaves para \cs[no-index]{spmangle}}
%
% \keyexampcmd{read-cmd}{school \textnormal{\textcolor{lightgray}{\textbar}}
%   mathcat}{school}[spmangle]
% \smallskip
% Establece la forma en la que \mynvmath{NVDA/MathCAT} verbalizará los
% puntos del ángulo. Si se establece en |school|, \mynvmath{NVDA/MathCAT}
% verbaliza las letras directamente sin anteponer «mayúscula»; si se
% establece en |mathcat|, \mynvmath{NVDA/MathCAT} aplica sus propias
% reglas de habla, que pueden incluir el prefijo «mayúscula» según la
% configuración del usuario.
%
% \smallskip
%
% \keyexampcmd{prefix}{prefijo}{no usado}[spmangle]
% \smallskip
% Establece la palabra que \mynvmath{NVDA/MathCAT} verbalizará
% \emph{antes} de «medida del ángulo». El valor de esta llave
% \emph{siempre} debe pasarse entre llaves `|{ }|'. Si se establece en
% |{la}|, \mynvmath{NVDA/MathCAT} verbalizará \emph{«la medida del
% ángulo en A»} en lugar de \emph{«medida del ángulo en A»}; si no se
% establece, no se antepone ninguna palabra.
%
% \smallskip
%
% \keyexampcmd{silent-cmd}{true \textnormal{\textcolor{lightgray}{\textbar}}
%   false}{false}[spmangle]
% \smallskip
% Establece si \mynvmath{NVDA/MathCAT} \emph{silenciará} la frase
% «medida del ángulo». Útil cuando el objeto ya se ha mencionado en el
% texto. Con |silent-cmd=true|, \mynvmath{NVDA/MathCAT} verbalizará
% solo \emph{«en A»} en lugar de \emph{«medida del ángulo en A»}.
%
% \smallskip
%
% \keyexampcmd{recto}{true \textnormal{\textcolor{lightgray}{\textbar}}
%   false}{false}[spmangle]
% \smallskip
% Establece si el ángulo cuya medida se representa es un \emph{ángulo
% recto}, con la misma semántica que en |\spangle|
% (\S\ref{subsec:spangle}).
%
% \smallskip
%
% \keyexampcmd{print-m}{true \textnormal{\textcolor{lightgray}{\textbar}}
%   false}{false}[spmangle]
% \smallskip
% Establece la \emph{notación visual} de la medida. Si se establece en
% |false| (valor por defecto), se usa el símbolo $\measuredangle$, que
% ya incorpora la noción de medida sin necesidad de anteponer una
% letra: |$\spmangle{A, O, B}$| imprimirá `$\measuredangle AOB$'. Si
% se establece en |true|, se antepone |\operatorname{m}| al símbolo
% $\angle$ simple: |$\spmangle[print-m=true]{A, O, B}$| imprimirá
% `$\operatorname{m}\angle AOB$'. En ambos casos \mynvmath{NVDA/MathCAT}
% verbalizará \emph{«medida del ángulo A O B»}. Esta llave funciona de
% forma \emph{independiente} de \mykey[spmangle]{recto}: si se
% establece en |true| junto con |recto=true|,
% |$\spmangle[print-m=true, recto=true]{A}$| imprimirá
% `$\operatorname{m}\rightangle A$' y verbalizará \emph{«medida del
% ángulo recto en A»}; si se establece en |false| junto con
% |recto=true|, |$\spmangle[recto=true]{A}$| imprimirá
% `$\measuredrightangle A$' con la misma verbalización.
%
% \subsection*{El comando \cs[no-index]{spang}}\label{subsec:spang}
%
% \vspace*{-\baselineskip}
%
% \begin{function}{\spang}
%   \begin{syntax}
%     \cmdexamp{spang}{grados}
%     \cmdexamp{spang}{grados \textnormal{\textcolor{red}{:}} minutos}
%     \cmdexamp{spang}{grados \textnormal{\textcolor{red}{:}} minutos \textnormal{\textcolor{red}{:}} segundos}
%   \end{syntax}
%   El comando \ics*{spang} trabaja en \emph{modo matemático} y recibe el
%   \emph{argumento obligatorio} \mymarg{grados}, \mymarg{grados \textnormal{\textcolor{red}{:}} minutos}
%   o \mymarg{grados \textnormal{\textcolor{red}{:}} minutos
%   \textnormal{\textcolor{red}{:}} segundos}. Los procesa internamente bajo
%   los comandos |\spnum| y |\spunit| para imprimir y verbalizar el
%   \emph{«ángulo sexagesimal»}. Por ejemplo, |$\spang{45:30:15}$|
%   imprimirá `$45°30′15″$' en la salida \texttt{PDF} y \mynvmath{NVDA/MathCAT}
%   verbalizará \emph{«cuarenta y cinco grados treinta minutos quince
%   segundos»}.
% \end{function}
%
% \subsection*{El comando \cs[no-index]{sparc}}\label{subsec:sparc}
%
% \vspace*{-\baselineskip}
%
% \begin{function}{\sparc}
%   \begin{syntax}
%     \ics*{sparc}\mymarg{puntos}
%     \ics*{sparc}\myoarg{key \textnormal{\textcolor{gray}{=}} val}\mymarg{puntos}
%   \end{syntax}
%   El comando \ics*{sparc} trabaja en \emph{modo matemático} y
%   representa un \emph{arco} geométrico. La salida visual usa un
%   símbolo de arco extensible, tomado de \cs{widearc} cuando está
%   disponible (provisto por los paquetes \mypkg{yhmath} o
%   \mypkg{MnSymbol}), o del glifo extensible U+23DC de la fuente
%   matemática cargada como alternativa.
%
%   \smallskip
%
%   Si no se proporciona el \emph{argumento opcional}
%   \myoarg{key \textnormal{\textcolor{gray}{=}} val},
%   |$\sparc{A, B}$| imprimirá `$\overset{\frown}{AB}$' en la salida
%   \texttt{PDF} y \mynvmath{NVDA/MathCAT} verbalizará
%   \emph{«arco A B»}.
% \end{function}
%
% \subsubsection*{Llaves para \cs[no-index]{sparc}}
%
% \keyexampcmd{read-cmd}{school \textnormal{\textcolor{lightgray}{\textbar}}
%   mathcat}{school}[sparc]
% \smallskip
% Establece la forma en la que \mynvmath{NVDA/MathCAT} verbalizará los
% puntos del arco. Si se establece en |school|, \mynvmath{NVDA/MathCAT}
% verbaliza las letras directamente sin anteponer «mayúscula»; si se
% establece en |mathcat|, \mynvmath{NVDA/MathCAT} aplica sus propias
% reglas de habla, que pueden incluir el prefijo «mayúscula» según la
% configuración del usuario.
%
% \smallskip
%
% \keyexampcmd{prefix}{prefijo}{no usado}[sparc]
% \smallskip
% Establece la palabra que \mynvmath{NVDA/MathCAT} verbalizará
% \emph{antes} de «arco». El valor de esta llave \emph{siempre} debe
% pasarse entre llaves `|{ }|'. Si se establece en |{el}|,
% \mynvmath{NVDA/MathCAT} verbalizará \emph{«el arco A B»} en lugar de
% \emph{«arco A B»}; si no se establece, no se antepone ninguna
% palabra.
%
% \smallskip
%
% \keyexampcmd{silent-cmd}{true \textnormal{\textcolor{lightgray}{\textbar}}
%   false}{false}[sparc]
% \smallskip
% Establece si \mynvmath{NVDA/MathCAT} \emph{silenciará} la palabra
% «arco». Útil cuando el objeto ya se ha mencionado en el texto. Con
% |silent-cmd=true|, \mynvmath{NVDA/MathCAT} verbalizará solo
% \emph{«A B»} en lugar de \emph{«arco A B»}.
%
% \subsection*{El comando \cs[no-index]{spmarc}}\label{subsec:spmarc}
%
% \vspace*{-\baselineskip}
%
% \begin{function}{\spmarc}
%   \begin{syntax}
%     \ics*{spmarc}\mymarg{puntos}
%     \ics*{spmarc}\myoarg{key \textnormal{\textcolor{gray}{=}} val}\mymarg{puntos}
%   \end{syntax}
%   El comando \ics*{spmarc} trabaja en \emph{modo matemático} y
%   representa la \emph{medida de un arco}. A diferencia de |\spmside|
%   y |\spmangle|, no existe la llave |print-m| porque la notación
%   $\operatorname{m}\overset{\frown}{AB}$ es la única forma estándar
%   para la medida de un arco: |\operatorname{m}| \emph{siempre} se
%   antepone al símbolo del arco.
%
%   \smallskip
%
%   Si no se proporciona el \emph{argumento opcional}
%   \myoarg{key \textnormal{\textcolor{gray}{=}} val},
%   |$\spmarc{A, B}$| imprimirá
%   `$\operatorname{m}\overset{\frown}{AB}$' en la salida \texttt{PDF}
%   y \mynvmath{NVDA/MathCAT} verbalizará \emph{«medida del arco
%   A B»}.
% \end{function}
%
% \subsubsection*{Llaves para \cs[no-index]{spmarc}}
%
% \keyexampcmd{read-cmd}{school \textnormal{\textcolor{lightgray}{\textbar}}
%   mathcat}{school}[spmarc]
% \smallskip
% Establece la forma en la que \mynvmath{NVDA/MathCAT} verbalizará los
% puntos del arco. Si se establece en |school|, \mynvmath{NVDA/MathCAT}
% verbaliza las letras directamente sin anteponer «mayúscula»; si se
% establece en |mathcat|, \mynvmath{NVDA/MathCAT} aplica sus propias
% reglas de habla, que pueden incluir el prefijo «mayúscula» según la
% configuración del usuario.
%
% \smallskip
%
% \keyexampcmd{prefix}{prefijo}{no usado}[spmarc]
% \smallskip
% Establece la palabra que \mynvmath{NVDA/MathCAT} verbalizará
% \emph{antes} de «medida del arco». El valor de esta llave
% \emph{siempre} debe pasarse entre llaves `|{ }|'. Si se establece en
% |{la}|, \mynvmath{NVDA/MathCAT} verbalizará \emph{«la medida del
% arco A B»} en lugar de \emph{«medida del arco A B»}; si no se
% establece, no se antepone ninguna palabra.
%
% \smallskip
%
% \keyexampcmd{silent-cmd}{true \textnormal{\textcolor{lightgray}{\textbar}}
%   false}{false}[spmarc]
% \smallskip
% Establece si \mynvmath{NVDA/MathCAT} \emph{silenciará} la frase
% «medida del arco». Útil cuando el objeto ya se ha mencionado en el
% texto. Con |silent-cmd=true|, \mynvmath{NVDA/MathCAT} verbalizará
% solo \emph{«A B»} en lugar de \emph{«medida del arco A B»}.
%
% \subsection*{El comando \cs[no-index]{spPoly}}\label{subsec:spPoly}
%
% \vspace*{-\baselineskip}
%
% \begin{function}{\spPoly}
%   \begin{syntax}
%     \ics*{spPoly}\mymarg{puntos}
%     \ics*{spPoly}\myoarg{key \textnormal{\textcolor{gray}{=}} val}\mymarg{puntos}
%   \end{syntax}
%   El comando \ics*{spPoly} trabaja en \emph{modo matemático} y
%   recibe en \mymarg{puntos} una lista de \emph{al menos tres letras
%   mayúsculas} separadas por comas. El tipo de figura y su símbolo
%   visual se determinan automáticamente según el número de puntos:
%   tres puntos producen un triángulo; cuatro puntos producen un
%   cuadrado; cinco o más puntos producen un polígono genérico sin
%   símbolo.
%
%   \smallskip
%
%   Con la llave \mykey[spPoly]{print-sym}|=true| (valor por
%   defecto), |$\spPoly{A, B, C}$| imprimirá `$\triangle ABC$' en la
%   salida \texttt{PDF} y \mynvmath{NVDA/MathCAT} verbalizará
%   \emph{«triángulo A B C»}; |$\spPoly{A, B, C, D}$| imprimirá
%   `$\square ABCD$' y verbalizará \emph{«cuadrado A B C D»};
%   |$\spPoly{A, B, C, D, E}$| imprimirá `$ABCDE$' y verbalizará
%   \emph{«A B C D E»}.
% \end{function}
%
% \smallskip
%
% \begin{important}*
%   Para figuras de cuatro lados que no son cuadrados (rectángulo,
%   paralelogramo, rombo, trapecio) o para polígonos de más de cuatro
%   lados, el nombre de la figura va en el texto que rodea la
%   expresión y |\spPoly| se usa con \mykey[spPoly]{print-sym}|=false|
%   para tipografiar solo los vértices. Por ejemplo: «el rectángulo
%   |$\spPoly[print-sym=false]{A, B, C, D}$|» imprimirá `$ABCD$' y
%   \mynvmath{NVDA/MathCAT} verbalizará \emph{«el rectángulo A B C
%   D»} sin repetir el nombre.
% \end{important}
%
% \subsubsection*{Llaves para \cs[no-index]{spPoly}}
%
% \keyexampcmd{read-cmd}{school \textnormal{\textcolor{lightgray}{\textbar}}
%   mathcat}{school}[spPoly]
% \smallskip
% Establece la forma en la que \mynvmath{NVDA/MathCAT} verbalizará los
% vértices. Si se establece en |school|, \mynvmath{NVDA/MathCAT}
% verbaliza las letras directamente sin anteponer «mayúscula»; si se
% establece en |mathcat|, \mynvmath{NVDA/MathCAT} aplica sus propias
% reglas de habla, que pueden incluir el prefijo «mayúscula» según la
% configuración del usuario.
%
% \smallskip
%
% \keyexampcmd{prefix}{prefijo}{no usado}[spPoly]
% \smallskip
% Establece la palabra que \mynvmath{NVDA/MathCAT} verbalizará
% \emph{antes} del nombre de la figura. El valor de esta llave
% \emph{siempre} debe pasarse entre llaves `|{ }|'. Si se establece en
% |{el}|, \mynvmath{NVDA/MathCAT} verbalizará \emph{«el triángulo
% A B C»} en lugar de \emph{«triángulo A B C»}; si no se establece,
% no se antepone ninguna palabra.
%
% \smallskip
%
% \keyexampcmd{print-sym}{true \textnormal{\textcolor{lightgray}{\textbar}}
%   false}{true}[spPoly]
% \smallskip
% Establece si se imprime el símbolo visual de la figura y
% \mynvmath{NVDA/MathCAT} verbaliza su nombre. Si se establece en
% |false|, |$\spPoly[print-sym=false]{A, B, C}$| imprimirá `$ABC$' sin
% símbolo y \mynvmath{NVDA/MathCAT} verbalizará solo \emph{«A B C»};
% si se establece en |true|, imprimirá `$\triangle ABC$' y verbalizará
% \emph{«triángulo A B C»}. Para polígonos de cinco o más puntos esta
% llave no tiene efecto — siempre se imprimen solo los vértices.
%
% \smallskip
%
% \keyexampcmd{silent-cmd}{true \textnormal{\textcolor{lightgray}{\textbar}}
%   false}{false}[spPoly]
% \smallskip
% Alias de \mykey[spPoly]{print-sym}|=false|, incluido por consistencia
% con el resto de los comandos de geometría plana. Escribir
% |silent-cmd=true| equivale a escribir |print-sym=false|.
%
% \subsection*{El comando \cs[no-index]{spAfig}}\label{subsec:spAfig}
%
% \vspace*{-\baselineskip}
%
% \begin{function}{\spAfig}
%   \begin{syntax}
%     \ics*{spAfig}
%     \ics*{spAfig}(\mymeta{puntos})
%     \ics*{spAfig}\myoarg{key \textnormal{\textcolor{gray}{=}} val}(\mymeta{puntos})
%   \end{syntax}
%   El comando \ics*{spAfig} trabaja en \emph{modo matemático} y
%   representa el \emph{área} de una figura poligonal. A diferencia
%   del resto de los comandos de geometría plana, el argumento
%   \mymeta{puntos} se delimita con \emph{paréntesis} en lugar de
%   llaves, y es \emph{completamente opcional}.
%
%   \smallskip
%
%   Si se usa sin paréntesis, |$\spAfig$| imprimirá `$\mathcal{A}$'
%   en la salida \texttt{PDF} y \mynvmath{NVDA/MathCAT} verbalizará
%   \emph{«área»}. Si se proporcionan tres o cuatro puntos, el tipo de
%   figura se determina automáticamente por el número de vértices y se
%   imprime como subíndice: |$\spAfig(A, B, C)$| imprimirá
%   `$\mathcal{A}_{\triangle ABC}$' y \mynvmath{NVDA/MathCAT}
%   verbalizará \emph{«área del triángulo A B C»};
%   |$\spAfig(A, B, C, D)$| imprimirá `$\mathcal{A}_{\square ABCD}$'
%   y verbalizará \emph{«área del cuadrado A B C D»}. Con cinco o más
%   puntos, el subíndice solo contiene los vértices, y
%   \mynvmath{NVDA/MathCAT} verbaliza el nombre del polígono
%   correspondiente sin símbolo visual.
% \end{function}
%
% \smallskip
%
% \begin{important}*
%   El subíndice de |\spAfig| solo identifica la figura por sus
%   vértices — nunca contiene medidas (base, altura, lados). Para
%   expresar el cálculo mismo, la medida se escribe aparte en la
%   ecuación: por ejemplo, `$\mathcal{A}_{\triangle ABC} =
%   \dfrac{b \cdot h}{2}$'.
% \end{important}
%
% \subsubsection*{Llaves para \cs[no-index]{spAfig}}
%
% \keyexampcmd{read-cmd}{school \textnormal{\textcolor{lightgray}{\textbar}}
%   mathcat}{school}[spAfig]
% \smallskip
% Establece la forma en la que \mynvmath{NVDA/MathCAT} verbalizará los
% vértices. Si se establece en |school|, \mynvmath{NVDA/MathCAT}
% verbaliza las letras directamente sin anteponer «mayúscula»; si se
% establece en |mathcat|, \mynvmath{NVDA/MathCAT} aplica sus propias
% reglas de habla, que pueden incluir el prefijo «mayúscula» según la
% configuración del usuario.
%
% \smallskip
%
% \keyexampcmd{prefix}{prefijo}{no usado}[spAfig]
% \smallskip
% Establece la palabra que \mynvmath{NVDA/MathCAT} verbalizará
% \emph{antes} de «área». El valor de esta llave \emph{siempre} debe
% pasarse entre llaves `|{ }|'. Si se establece en |{el}|,
% \mynvmath{NVDA/MathCAT} verbalizará \emph{«el área del triángulo
% A B C»} en lugar de \emph{«área del triángulo A B C»}; si no se
% establece, no se antepone ninguna palabra.
%
% \smallskip
%
% \keyexampcmd{print-sym}{true \textnormal{\textcolor{lightgray}{\textbar}}
%   false}{true}[spAfig]
% \smallskip
% Establece si se imprime el \emph{símbolo visual} de la figura
% (\cs{triangle} o \cs{mdlgwhtsquare}) dentro del subíndice. Esta llave
% solo tiene efecto cuando el número de vértices es tres o cuatro. Si
% se establece en |false|, |$\spAfig[print-sym=false](A, B, C)$|
% imprimirá `$\mathcal{A}_{ABC}$' sin el símbolo del triángulo,
% mientras que \mynvmath{NVDA/MathCAT} sigue verbalizando \emph{«área
% del triángulo A B C»} — la llave afecta solo lo visual, no la
% verbalización.
%
% \smallskip
%
% \keyexampcmd{read-arg}{nombre}{(automático)}[spAfig]
% \smallskip
% Sobrescribe el \emph{nombre de la figura} que \mynvmath{NVDA/MathCAT}
% verbaliza, sin afectar el símbolo visual ni el número de vértices.
% Es especialmente útil con cuatro vértices, donde la cardinalidad por
% sí sola no puede distinguir un cuadrado de un rectángulo, un rombo,
% un paralelogramo o un trapecio. El valor de esta llave \emph{siempre}
% debe pasarse entre llaves `|{ }|'. Por ejemplo,
% |$\spAfig[read-arg={trapecio}](A, B, C, D)$| imprimirá
% `$\mathcal{A}_{\square ABCD}$' (el símbolo por defecto de cuatro
% vértices) pero \mynvmath{NVDA/MathCAT} verbalizará \emph{«área del
% trapecio A B C D»}. Para además ocultar el símbolo, que en este caso
% no representa correctamente un trapecio, combínese con
% \mykey[spAfig]{print-sym}|=false|:
% |$\spAfig[read-arg={trapecio}, print-sym=false](A, B, C, D)$|
% imprimirá `$\mathcal{A}_{ABCD}$' con la misma verbalización.
%
% \subsection*{El comando \cs[no-index]{spPfig}}\label{subsec:spPfig}
%
% \vspace*{-\baselineskip}
%
% \begin{function}{\spPfig}
%   \begin{syntax}
%     \ics*{spPfig}
%     \ics*{spPfig}(\mymeta{puntos})
%     \ics*{spPfig}\myoarg{key \textnormal{\textcolor{gray}{=}} val}(\mymeta{puntos})
%   \end{syntax}
%   El comando \ics*{spPfig} es el homólogo de |\spAfig| para el
%   \emph{perímetro} de una figura poligonal: misma sintaxis, mismo
%   comportamiento, usando $\mathcal{P}$ como símbolo.
%
%   \smallskip
%
%   Si se usa sin paréntesis, |$\spPfig$| imprimirá `$\mathcal{P}$' en
%   la salida \texttt{PDF} y \mynvmath{NVDA/MathCAT} verbalizará
%   \emph{«perímetro»}. Con vértices, |$\spPfig(A, B, C)$| imprimirá
%   `$\mathcal{P}_{\triangle ABC}$' y verbalizará \emph{«perímetro del
%   triángulo A B C»}; |$\spPfig(A, B, C, D)$| imprimirá
%   `$\mathcal{P}_{\square ABCD}$' y verbalizará \emph{«perímetro del
%   cuadrado A B C D»}.
% \end{function}
%
% \smallskip
%
% \begin{important}*
%   Para el área y el perímetro de un \emph{círculo} o una
%   \emph{circunferencia}, |\spAfig| y |\spPfig| no son los comandos
%   adecuados — véase |\spCirc|, que provee la notación propia de
%   estos objetos.
% \end{important}
%
% \subsubsection*{Llaves para \cs[no-index]{spPfig}}
%
% \keyexampcmd{read-cmd}{school \textnormal{\textcolor{lightgray}{\textbar}}
%   mathcat}{school}[spPfig]
% \smallskip
% Establece la forma en la que \mynvmath{NVDA/MathCAT} verbalizará los
% vértices. Si se establece en |school|, \mynvmath{NVDA/MathCAT}
% verbaliza las letras directamente sin anteponer «mayúscula»; si se
% establece en |mathcat|, \mynvmath{NVDA/MathCAT} aplica sus propias
% reglas de habla, que pueden incluir el prefijo «mayúscula» según la
% configuración del usuario.
%
% \smallskip
%
% \keyexampcmd{prefix}{prefijo}{no usado}[spPfig]
% \smallskip
% Establece la palabra que \mynvmath{NVDA/MathCAT} verbalizará
% \emph{antes} de «perímetro». El valor de esta llave \emph{siempre}
% debe pasarse entre llaves `|{ }|'. Si se establece en |{el}|,
% \mynvmath{NVDA/MathCAT} verbalizará \emph{«el perímetro del
% triángulo A B C»} en lugar de \emph{«perímetro del triángulo
% A B C»}; si no se establece, no se antepone ninguna palabra.
%
% \smallskip
%
% \keyexampcmd{print-sym}{true \textnormal{\textcolor{lightgray}{\textbar}}
%   false}{true}[spPfig]
% \smallskip
% Establece si se imprime el \emph{símbolo visual} de la figura dentro
% del subíndice, con la misma semántica que en |\spAfig|
% (\S\ref{subsec:spAfig}). Si se establece en |false|,
% |$\spPfig[print-sym=false](A, B, C, D)$| imprimirá
% `$\mathcal{P}_{ABCD}$' y \mynvmath{NVDA/MathCAT} sigue verbalizando
% \emph{«perímetro del cuadrado A B C D»}.
%
% \smallskip
%
% \keyexampcmd{read-arg}{nombre}{(automático)}[spPfig]
% \smallskip
% Sobrescribe el \emph{nombre de la figura} verbalizado, con la misma
% semántica que en |\spAfig| (\S\ref{subsec:spAfig}). El valor de esta
% llave \emph{siempre} debe pasarse entre llaves `|{ }|'.
%
%
% \newpage
%
% \let\stdsection\section
% \def\section*#1{\stdsection{#1}}
%
% \begin{thebibliography}{10}
%
% \bibitem{unicode-math} \textsc{Robertson, Will }. \enquote{unicode-math
% - Unicode mathematics support for \hologo{XeTeX} and \hologo{LuaTeX}}. Available from
% \textsc{ctan}, \url{https://www.ctan.org/pkg/unicode-math}, 2023.
%
% \bibitem{luamml} \textsc{Krueger, Marcel}. \enquote{LuaMML%
% - Automatically generate MathML from \hologo{LuaTeX} math mode material}.
% Available from \textsc{ctan}, \url{https://ctan.org/pkg/luamml}, 2026.
%
% \bibitem{fontspec} The \hologo{LaTeX} Project. \enquote{fontspec%
% - Advanced font selection for \hologo{XeLaTeX} and \hologo{LuaTeX}}.
% Available from \textsc{ctan}, \url{https://www.ctan.org/pkg/fontspec},%
% 2025.
%
% \bibitem{libertinus-otf} \textsc{Voß, Herbert}. \enquote{libertinus-otf%
% - Support for Libertinus OpenType}. Available from \textsc{ctan},%
% \url{https://www.ctan.org/pkg/libertinus-otf}, 2025.
%
% \bibitem{erewhon-math} \textsc{Flipo, Daniel}. \enquote{erewhon-math%
% - Utopia based OpenType Math font}. Available from%
% \textsc{ctan}, \url{https://ctan.org/pkg/erewhon-math}, 2026.
%
% \bibitem{sourcecodepro} \textsc{Hofstra, Silke}. \enquote{sourcecodepro%
% - Use SourceCodePro with \hologo{TeX}(-alike) systems}. Available from%
% \textsc{ctan}, \url{https://ctan.org/pkg/sourcecodepro}, 2025.
%
% \bibitem{scontents} \textsc{González, Pablo}. \enquote{\textsf{scontents} - Stores
% \hologo{LaTeX} contents in memory or files}. Available from
% \textsc{ctan}, \url{https://www.ctan.org/pkg/scontents}, 2026.
%
% \bibitem{enumext} \textsc{González, Pablo}. \enquote{\textsf{enumext} -
% Enumerate exercise sheets}. Available from \textsc{ctan},
% \url{https://www.ctan.org/pkg/enumext}, 2026.
%
% \bibitem{babel} \textsc{Braams, Johannes and Bezos, Javier}.
% \enquote{\textsf{babel} - Multilingual support for \hologo{LaTeX},
% \hologo{LuaLaTeX}, \hologo{XeLaTeX}, and Plain \TeX}. Available from
% \textsc{ctan}, \url{https://www.ctan.org/pkg/babel}, 2026.
%
% \bibitem{babel-spanish} \textsc{Bezos, Javier}.
% \enquote{\textsf{babel-spanish} - Babel support for Spanish}.
% Available from \textsc{ctan}, \url{https://www.ctan.org/pkg/babel-spanish},
% 2026.
%
% \bibitem{nvda} \textsc{NV Access}. \enquote{\textsf{NVDA} - NonVisual
% Desktop Access screen reader}. Available from \textsc{nv access},
% \url{https://www.nvaccess.org/}, 2026.
%
% \bibitem{mathcat} \textsc{Soiffer, Neil}. \enquote{\textsf{MathCAT} -
% Math to speech and braille translation}. Available from \textsc{github},
% \url{https://nsoiffer.github.io/MathCAT/}, 2026.
%
% \bibitem{acrobat} \textsc{Adobe Inc.} \enquote{\textsf{Adobe Acrobat
% Reader} - Reliable \texttt{PDF} viewer}. Available from \textsc{adobe},
% \url{https://get.adobe.com/reader/}, 2026.
%
% \bibitem{pdf20} \textsc{ISO}. \enquote{\textsf{ISO 32000-2:2020} - Portable
% Document Format (\texttt{PDF} 2.0)}. Available from \textsc{\texttt{PDF} association},
% \url{https://pdfa.org/resource/iso-32000-2/}, 2020.
%
% \bibitem{wtpdf} \textsc{\texttt{PDF} Association}. \enquote{\textsf{WTPDF} -
% Well-Tagged \texttt{PDF} specification}. Available from \textsc{pdfa},
% \url{https://pdfa.org/wtpdf/}, 2024.
%
% \bibitem{pdf20-errata} \textsc{\texttt{PDF} Association}. \enquote{\textsf{\texttt{PDF} 2.0
% Errata Collection 3}}. Available from \textsc{pdfa},
% \url{https://pdfa.org/PDF-2-0-errata-collection-3-now-available/}, 2026.
%
% \bibitem{bpg-math} \textsc{LaTeX Project LWG}. \enquote{\textsf{Best
% Practice Guide: Math in PDF}}. Available from \textsc{pdfa},
% \url{https://pdfa.org/resource/best-practice-guide-math-in-PDF/},
% 2026.
%
% \bibitem{pdf20-impact} \textsc{Johnson, Duff}. \enquote{\textsf{\texttt{PDF} 2.0:
% New Features, Real-World Impact}}. Available from \textsc{pdfa},
% \url{https://pdfa.org/PDF-2-0-new-features-real-world-impact/}, 2026.
%
% \bibitem{tagged-project} The \hologo{LaTeX} Project. \enquote{\textsf{LaTeX
% Tagged \texttt{PDF} Project} - Documentation and resources}. Available from
% \textsc{latex project}, \url{https://latex3.github.io/tagging-project/},
% 2026.
%
% \bibitem{verapdf} \textsc{veraPDF Consortium}. \enquote{\textsf{veraPDF} -
% Industry supported \texttt{PDF}/A and \texttt{PDF/UA} validator}. Available from
% \textsc{verapdf}, \url{https://verapdf.org/}, 2026.
%
% \bibitem{ngpdf} \textsc{Dual Lab}. \enquote{\textsf{ngPDF} -
% Next Generation \texttt{PDF} demonstrator and validator}. Available from
% \textsc{ngpdf}, \url{https://ngpdf.com/}, 2026.
%
% \bibitem{latexref} \textsc{Berry, Karl}. \enquote{\hologo{LaTeX2e}: An
% Unofficial Reference Manual}. Available from \textsc{ctan},
% \url{https://ctan.org/pkg/latex2e-help-texinfo}, 2026.
%
% \bibitem{hyperref} The \hologo{LaTeX} Project. \enquote{Extensive support
% for hypertext in \hologo{LaTeX}}. Available from
% \textsc{ctan}, \url{https://www.ctan.org/pkg/hyperref}, 2026.
%
% \bibitem{expl3} The \hologo{LaTeX} Project. \enquote{The \textsf{expl3}
% package}. Available from
% \textsc{ctan}, \url{https://www.ctan.org/pkg/l3kernel}, 2026.
%
% \bibitem{inter3} The \hologo{LaTeX} Project. \enquote{The \hologo{LaTeX}3
% Interfaces}. Available from
% \textsc{ctan}, \url{https://www.ctan.org/pkg/l3kernel}, 2026.
%
% \bibitem{latex2e} The \hologo{LaTeX} Project. \enquote{The \hologo{LaTeX2e}
% sources}. Available from
% \textsc{ctan}, \url{https://ctan.org/tex-archive/macros/latex/base}, 2026.
%
% \bibitem{ltnews} The \hologo{LaTeX} Project. \enquote{\hologo{LaTeX} News,
% Issue 41, June 2025}. Available from
% \textsc{ctan}, \url{https://ctan.org/tex-archive/macros/latex/base}, 2025.
%
% \bibitem{userguide} The \hologo{LaTeX} Project. \enquote{\hologo{LaTeX} for authors
% current version}. Available from
% \textsc{ctan}, \url{https://ctan.org/pkg/latex-base}, 2026.
%
% \bibitem{tagpdf} \textsc{Fischer, Ulrike}. \enquote{\textsf{tagpdf} – \hologo{LaTeX}
% kernel code for \texttt{PDF} tagging}. Available from
% \textsc{ctan}, \url{https://www.ctan.org/pkg/tagpdf}, 2026.
%
% \bibitem{latexlab} The \hologo{LaTeX} Project. \enquote{\textsf{latex-lab} – \hologo{LaTeX}
% laboratory}. Available from
% \textsc{ctan}, \url{https://www.ctan.org/pkg/latex-lab}, 2026.
%
% \end{thebibliography}
%
% \let\section\stdsection
%
% \section{Change history}
% \label{sec:changes}
%
% \setlist[itemize,1]{label=\textendash,wide=0.5em,nosep,noitemsep,leftmargin=10pt}
% \newlength\descrwidth
% \settowidth{\descrwidth}{\textsf{v1.0, (ctan), 2026-06-20} }
% \begin{description}[font=\small\sffamily,wide=0pt,style=multiline,leftmargin=1.1\descrwidth,nosep,noitemsep]
%  \item [\fileversion{} (ctan), \filedate]
%    \begin{itemize}
%      \item First public release.
%    \end{itemize}
% \end{description}
%
% \newpage
%
% \indexprologue{
% Los números en cursiva indican las páginas donde se describe la entrada
% correspondiente.}
%
% \printindex[userdoc]
%
% \newpage
%
% \selectlanguage{english}
%
% \StartImplementation
%
% \StopEventually{^^A
% \newgeometry{top=0.5in,bottom=0.3in,left=1.0in,right=0.5in,footskip=0.2in,headheight=1cm,headsep=0.27cm}
% \addtocontents{toc}{\protect\setcounter{tocdepth}{2}}
% \cleardoublepage
% \phantomsection
% \indexprologue{
% The italic numbers denote the pages where the corresponding entry is
% described, the numbers underlined and all others indicate the line on
% which they are implemented in the package code.
% }
% \printindex
% }
%
% \newgeometry{top=0.5in,bottom=0.3in,left=1.95in,right=0.5in,footskip=0.2in,headheight=1cm,headsep=0.27cm}
%
% \section{Implementation}
% \label{sec:Implementation}
% \addtocontents{toc}{\protect\setcounter{tocdepth}{0}}
%
% The most recent publicly released version of \mypkg*[no-index]{spintent} is
% available at \textsc{ctan}: \url{https://www.ctan.org/pkg/spintent}.
% While general feedback via email is welcomed, specific bugs or
% feature requests should be reported through the issue
% tracker: \textcolor{gray}{\scriptsize\faIcon[regular]{github}} \url{https://github.com/pablgonz/spintent/issues}.
%
% \smallskip
%
% All variables and functions defined in this package are \enquote{private}
% and are NOT intended to work or be used by another package or module.
%
% \smallskip
%
% \begin{important}*
% The documentation presented here is far from professional, it contains
% a lot of obvious information that to the eye of a \hologo{TeX}pert are
% superfluous, but, after so many years developing this project is the
% only way to remember what does what.
% \end{important}
%
% \subsection{Initial set up}
%
% Start the \pkg{DocStrip} guards.
%    \begin{macrocode}
%<*package>
%    \end{macrocode}
% Identify the internal prefix (\LaTeX3 \pkg{DocStrip} convention) for
% \mypkg{l3doc} class.
%    \begin{macrocode}
%<@@=spintent>
%    \end{macrocode}
%
% \smallskip
%
% \begin{important}*
% Functions of the form \myhlc|\@@_luafun_some:nn| or variables of the
% form \myhlc|\l_@@_luaset_some_tl| or \myhlc|\l_@@_luaset_some_str| are
% defined in the Lua module \myluamod{} and only declared in \mypkg{expl3}.
% \end{important}
%
% \subsection{Declaration of the package}
%
% First we will make sure we have a minimum (super updated) version of
% \hologo{LaTeX} to work correctly.
%    \begin{macrocode}
\NeedsTeXFormat{LaTeX2e}[2026-11-01]
%    \end{macrocode}
% Now declare the \mypkg*[no-index]{spintent} package.
%    \begin{macrocode}
\ProvidesExplPackage {spintent} {2026-08-20} {0.98} {Spanish parse intents}
%    \end{macrocode}
%
% We check if the requirements for running the \mypkg*[no-index]{spintent} package
% are met; that is: |\DocumentMetadata| must be active, \hologo{LuaTeX}
% must be running, and the \mypkg{unicode-math} or \mypkg{lua-unicode-math}
% package must be loaded. If neither is present, we will load \mypkg{unicode-math}.
% If the conditions cannot be met, we will return a \enquote{fatal error}.
%    \begin{macrocode}
\IfDocumentMetadataTF
  {
    \sys_if_engine_luatex:TF
      {
        \msg_new:nnn { spintent } { package-not-load }
          {
            The~'#1'~package~will~be~loaded~as~a~dependency.
          }
        \IfPackageLoadedF { unicode-math }
          {
            \IfPackageLoadedF { lua-unicode-math }
              {
                \msg_info:nnn { spintent } { package-not-load } { unicode-math }
                \RequirePackage{unicode-math}
              }
          }
        \msg_new:nnn { spintent } { env-success }
          { Everything~looks~correct,~spintent~will~make~the~intent! }
        \msg_info:nn { spintent } { env-success }
      }
      {
        \msg_new:nnnn { spintent } { fatal-engine-requires-luatex }
          { The~spintent~package~requires~LuaLaTeX. }
          { This~package~relies~on~Lua~module.~Compile~with~lualatex. }
        \msg_fatal:nn { spintent } { fatal-engine-requires-luatex }
      }
  }
  {
    \msg_new:nnnn { spintent } { fatal-metadata-missing }
      { \token_to_str:N \DocumentMetadata \c_space_tl is~missing. }
      {
        The~spintent~package~requires~tagged~PDF~active.\\
        You~must~declare~\token_to_str:N \DocumentMetadata { tagging=on } ~
      }
    \msg_fatal:nn { spintent } { fatal-metadata-missing }
  }
%    \end{macrocode}
%
% \subsection{Some kernel variants}
%
% \begin{macro}[int]{\str_if_eq:nVTF, \str_if_eq:nVF,
%                    \str_uppercase:V, \str_lowercase:V,
%                    \seq_set_from_clist:Ne}
%    Non-standard kernel variants.
%    \begin{macrocode}
\cs_generate_variant:Nn \str_if_eq:nnTF { V }
\cs_generate_variant:Nn \str_if_eq:nnF { V }
\cs_generate_variant:Nn \str_uppercase:n { V }
\cs_generate_variant:Nn \str_lowercase:n { V }
\cs_generate_variant:Nn \seq_set_from_clist:Nn { Ne }
%    \end{macrocode}
% \end{macro}
%
% \subsection{Loading the Lua module and functions}
%
% \begin{macro}[int]
%  {
%    \@@_luafun_spunit_lookup_alias:V,
%    \@@_luafun_spmoney_lookup_data:V,
%    \@@_luafun_spmQ_scale_int:Vn,
%    \@@_luafun_clean_split_and_set:V,
%  }
% Now we load the Lua module \myluamod{} and generate the variants for the
% functions \myhlc|\@@_luafun_spunit_lookup_alias:n| used by |\spunit| (\S\ref{cmd:spunit}),
% \myhlc|\@@_luafun_spmoney_lookup_data:V| used by |\spmoney| (\S\ref{cmd:spmoney}),
% \myhlc|\@@_luafun_spmQ_scale_int:Vn| used by |\spmQ| (\S\ref{cmd:spmQ}) and
% \myhlc|\@@_luafun_clean_split_and_set:V| used by |\spnum| (\S\ref{cmd:spnum})
% defined in it.
% \iffalse
%% Load Lua module
% \fi
%    \begin{macrocode}
\lua_load_module:n { spintent.lua }
\cs_generate_variant:Nn \@@_luafun_spunit_lookup_alias:n { V }
\cs_generate_variant:Nn \@@_luafun_spmoney_lookup_data:n { V }
\cs_generate_variant:Nn \@@_luafun_spmQ_scale_int:nn { VV, Vn }
\cs_generate_variant:Nn \@@_luafun_clean_split_and_set:n { V }
%    \end{macrocode}
% \end{macro}
%
% \subsection{Internal copy of the unicode invisible separator}
%
% \begin{macro}[int]{\@@_invisible_sep:, \c_@@_invisible_sep_tl}
%   The function \myhlc|\@@_invisible_sep:| is an internal copy of the
%   Unicode invisible separator |U-2063| that uses |\luamml_annotate:en|
%   from the package \mypkg{luamml}\cite{luamml} and let \lstinline+<mo>...</mo>+ in
%   the \mynvmath{MathML} structure embedded in the \texttt{PDF} output.
% \iffalse
%% Internal copy of invisible separator U-2063 => <mo>...</mo>.
% \fi
%    \begin{macrocode}
\cs_new_protected:Nn \@@_invisible_sep:
  {
    \luamml_annotate:en
      {
        core = {[0] = 'mo', intent=':silent', "^^^^2063"}
      }
      {
        \latelua{}
      }
  }
%    \end{macrocode}
%   The token list constant \myhlc|\c_@@_invisible_sep_tl| is an internal
%   copy of the Unicode invisible separator |U-2063| using the primitive
%   |\Umathchardef| from \hologo{LuaTeX} which leaves \lstinline+<mi>...</mi>+ in the
%   structure \mynvmath{MathML} embedded in the output \texttt{PDF}.
% \iffalse
%% Internal copy of invisible separator U-2063 => <mi>...</mi>.
% \fi
%    \begin{macrocode}
\hook_gput_code:nnn { begindocument } { spintent }
  {
    \Umathchardef \c_@@_invisible_sep_tl = "0 "0 "2063
  }
%    \end{macrocode}
% \begin{important}*
%   Both copies are used for different purposes, the function
%   \myhlc|\@@_invisible_sep:| will be used to force the start of a \lstinline+<mrow>+
%   in the structure \mynvmath{MathML}, the token list constant
%   \myhlc|\c_@@_invisible_sep_tl| will be used to inject text into the
%   structure \mynvmath{MathML} when we use |\MathMLintent|\mymarg{text}.
% \end{important}
% \end{macro}
%
% \subsection{Normalize argument and affixes}
%
% We normalize the \mymarg{argument} to support affixes (prefixes/suffixes)
% by trimming leading and trailing spaces, converting internal spaces to
% hyphens, collapsing consecutive hyphens, and removing leading or
% trailing hyphens.
%
% \begin{macro}[int]{\@@_collapse_hyphens:N}
%   A recursive helper to collapse multiple hyphens into a single one.
%    \begin{macrocode}
\cs_new_protected:Nn \@@_collapse_hyphens:N
  {
    \tl_if_in:NnT #1 { -- }
      {
        \tl_replace_all:Nnn #1 { -- } { - }
        \@@_collapse_hyphens:N #1
      }
  }
%    \end{macrocode}
% \end{macro}
%
% \begin{macro}[int]{\@@_sanitize_affix:Nn, \@@_sanitize_affix:cn}
%   The function \myhlc|\@@_sanitize_affix:Nn| checks if the \mymarg{argument}
%   passed in `|#2|' is \emph{blank}. If so, it clears the target token list
%   variable passed in `|#1|'. Otherwise, trimming bounds, converting spaces `\verb*+ +' to
%   hyphens `|-|', collapsing multiple hyphens, and safely stripping
%   leading/trailing hyphens. For example, if you pass a target variable and
%   raw text:
%
%   \begin{center}
%     \myhlc|\@@_sanitize_affix:Nn| \myhlc|\l_@@_spsomesym_prefix_tl||{ foo bar baz }|
%   \end{center}
%
%   The function will \emph{\enquote{store}} the normalized string
%   |foo-bar-baz| inside \myhlc|\l_@@_spsomesym_prefix_tl|, ready to be
%   conditionally appended.
%    \begin{macrocode}
\cs_new_protected:Nn \@@_sanitize_affix:Nn
  {
    \tl_if_blank:nTF {#2}
      { \tl_clear:N #1 }
      {
        \tl_set:Ne #1 {#2}
        \tl_trim_spaces:N #1
        \tl_replace_all:Nnn #1 { ~ } { - }
        \@@_collapse_hyphens:N #1
        \tl_put_left:Nn  #1 { \q_mark }
        \tl_put_right:Nn #1 { \q_stop }
        \tl_replace_all:Nnn #1 { \q_mark - } { \q_mark }
        \tl_replace_all:Nnn #1 { - \q_stop } { \q_stop }
        \tl_remove_all:Nn #1 { \q_mark }
        \tl_remove_all:Nn #1 { \q_stop }
      }
  }
\cs_generate_variant:Nn \@@_sanitize_affix:Nn { cn }
%    \end{macrocode}
% \end{macro}
%
% \subsection{The command \cs{setspintent}}\label{cmd:setspintent}
%
% The \ics*{setspintent} command is used to configure \mymeta{keys}
% \emph{\enquote{globally}} for the utilities provided by \mypkg*[no-index]{spintent}.
%
% \subsubsection{Internal variables and error messages for \cs[no-index]{setspintent}}
%
% \begin{macro}[int]{\l_@@_setspintent_input_arg_cmd_tl, unknown-command-setup}
%   The variable \myhlc|\l_@@_setspintent_input_arg_cmd_tl| will \emph{store} the
%   command or command name passed to the \emph{\enquote{first}} \mymarg{argument} `|#1|'
%   of the command |\setspintent|.
%    \begin{macrocode}
\tl_new:N \l_@@_setspintent_input_arg_cmd_tl
%    \end{macrocode}
%   Error message if the \mymarg{argument} `|#1|' passed to |\setspintent|
%   is a \emph{\enquote{non-existent}} command.
%    \begin{macrocode}
\msg_new:nnnn { spintent } { unknown-command-setup }
  {
    Cannot~configure~unknown~command~'\exp_not:c{#1}'.
  }
  {
    You~attempted~to~use~\token_to_str:N \setspintent{#1}{...},~but~the~
    command~\exp_not:c{#1}~is~not~defined.~Please~check~for~typos.
  }
%    \end{macrocode}
% \end{macro}
%
% \begin{macro}[int]{\@@_setspintent_set_keys:nn, \@@_setspintent_set_keys:en}
%   The function \myhlc|\@@_setspintent_set_keys:nn| will check if the
%   command passed as \mymarg{argument} exists in `|#1|' and set the
%   \mymeta{keys} passed as \mymarg{argument} in `|#2|' to the correct
%   path; otherwise, it will return an error.
%    \begin{macrocode}
\cs_new_protected:Nn \@@_setspintent_set_keys:nn
  {
    \cs_if_exist:cTF { #1 }
      {
        \keys_set:nn { spintent / #1 } { #2 }
      }
      {
        \msg_error:nnn { spintent } { unknown-command-setup } { #1 }
      }
  }
\cs_generate_variant:Nn \@@_setspintent_set_keys:nn { en }
%    \end{macrocode}
% \end{macro}
% \begin{macro}[int]{\@@_setspintent_parse_args:nn, \@@_setspintent_parse_args:Vn}
%   The function \myhlc|\@@_setspintent_parse_args:nn| will check if the
%   argument passed in |#1| begins with `|\|' or is just a \emph{command name}. If it
%   begins with `\myhlc|\|', it will remove it using \myhlc|\cs_to_str:N| and then
%   execute \myhlc|\@@_setspintent_set_keys:en|. Otherwise, i.e., if `|\|'
%   was not passed, it will directly execute \myhlc|\@@_setspintent_set_keys:nn|.
%    \begin{macrocode}
\cs_new_protected:Nn \@@_setspintent_parse_args:nn
  {
    \tl_if_single:nTF { #1 }
      {
        \token_if_cs:NTF #1
          {
            \@@_setspintent_set_keys:en { \cs_to_str:N #1 } { #2 }
          }
          {
            \@@_setspintent_set_keys:nn { #1 } { #2 }
          }
      }
      {
        \@@_setspintent_set_keys:nn { #1 } { #2 }
      }
  }
\cs_generate_variant:Nn \@@_setspintent_parse_args:nn { Vn }
%    \end{macrocode}
% \end{macro}
% \begin{macro}{\setspintent}
%   Now we define the user command |\setspintent|.
%    \begin{macrocode}
\NewDocumentCommand \setspintent { m m }
  {
    \tl_if_blank:nTF { #1 }
      {
        \msg_error:nnn { spintent } { unknown-command-setup } { <empty> }
      }
      {
        \tl_set:Nn \l_@@_setspintent_input_arg_cmd_tl { #1 }
        \tl_trim_spaces:N \l_@@_setspintent_input_arg_cmd_tl
        \@@_setspintent_parse_args:Vn \l_@@_setspintent_input_arg_cmd_tl { #2 }
      }
  }
%    \end{macrocode}
% \end{macro}
%
% \subsection{Handling unknown keys}\label{key:unknown-keys}
%
% \begin{macro}[int]{\l_@@_command_name_str, unknown-choice, cmd-key-unknown,
%                    cmd-key-value-unknown,}
%   We define the variable \myhlc|\l_@@_command_name_str| which will \emph{store}
%   the \emph{\enquote{name}} of each command within the implementation along with the
%   messages |unknown-choice| used by \mymeta{keys} that use |.choice:n| in
%   their implementation and the messages |cmd-key-unknown|, |cmd-key-value-unknown|
%   used by the \mymeta{key} |unknown| in commands that use \myoarg{key \textnormal{\textcolor{gray}{=}} val}.
% \iffalse
%% Common messages for command in the keys |unknown| and |unknown-choice|.
% \fi
%    \begin{macrocode}
\str_new:N  \l_@@_command_name_str
\msg_new:nnn { spintent } { unknown-choice }
  {
    The~value~'#3'~for~'#1'~key~is~invalid~use~('#2').
  }
\msg_new:nnnn { spintent } { cmd-key-unknown }
  {
    The~key~'#1'~is~unknown~by~command~
    \c_backslash_str \l_@@_command_name_str \c_space_tl
    and~is~being~ignored~\msg_line_context:.
  }
  {
    The~command~\c_backslash_str \l_@@_command_name_str \c_space_tl
    does~not~have~a~key~called~'#1'.\\
    Check~that~you~have~spelled~the~key~name~correctly.
  }
\msg_new:nnnn { spintent } { cmd-key-value-unknown }
  {
    The~key~'#1=#2'~is~unknown~by~command~
    \c_backslash_str \l_@@_command_name_str \c_space_tl
    and~is~being~ignored~\msg_line_context:.
  }
  {
    The~command~\c_backslash_str \l_@@_command_name_str \c_space_tl
    does~not~have~a~key~called~'#1'.\\
    Check~that~you~have~spelled~the~key~name~correctly.
  }
%    \end{macrocode}
% \end{macro}
% \begin{macro}[int]{\@@_unknown_keys:n, \@@_unknown_keys:nn,}
%   The internal functions \myhlc|\@@_unknown_keys:n| and \myhlc|\@@_unknown_keys:nn|
%   used by the \mymeta{key} |unknown|.
% \iffalse
%% Internal fucnctions for the key |unknown|.
% \fi
%    \begin{macrocode}
\cs_new_protected:Nn \@@_unknown_keys_cmd:n
  {
    \exp_args:NV \@@_unknown_keys_cmd:nn \l_keys_key_str {#1}
  }
\cs_new_protected:Nn \@@_unknown_keys_cmd:nn
  {
    \tl_if_blank:nTF {#2}
      { \msg_error:nnn { spintent } { cmd-key-unknown } {#1} }
      { \msg_error:nnnn { spintent } { cmd-key-value-unknown } {#1} {#2} }
  }
%    \end{macrocode}
% \end{macro}
%
% \subsection{Definition of core base variables}\label{sec:core-base-vars}
%
% \begin{macro}[int]{\@@_luafun_clean_split_and_set:n, \@@_luafun_clean_split_and_set:V}
%   The function \myhlc|\@@_luafun_clean_split_and_set:n| defined in the
%   Lua module \myluamod{}, is common to several commands in this
%   implementation. It analyzes, processes, cleans and splits the
%   \mymarg{argument} passed to commands using internal functions defined
%   in the Lua module \myluamod. It isolates signs, integer components, decimal
%   separators, decimal parte, decimal periodic part, exponents and final physical
%   \emph{\enquote{units}}, assigning them directly to \mypkg{expl3} variables.
%
%   \smallskip
%
%   This function assigns the states of the variables `|_str|' or `|_tl|'
%   described below to |true|, |false|, |error| or \emph{stores}
%   information according to the role it fulfills within the
%   implementation.
% \end{macro}
%
% \begin{variable}[int]{
%   \l_@@_luaset_sign_tl,
%   \l_@@_luaset_part_int_tl,
%   \l_@@_luaset_dec_sep_tl,
%   \l_@@_luaset_part_dec_tl,
%   \l_@@_luaset_part_period_tl,
%   \l_@@_luaset_has_integer_str,
%   \l_@@_luaset_has_natural_str,
% }
%   The variable \myhlc|\l_@@_luaset_sign_tl| \emph{stores} the explicit
%   `$\mathcolor{red}{+}$' or `$\mathcolor{red}{-}$' sign if present.
%   The variable \myhlc|\l_@@_luaset_part_int_tl| \emph{stores} the raw digit
%   sequence of the \emph{\enquote{integer part}}. The variable
%   \myhlc|\l_@@_luaset_dec_sep_tl| \emph{stores} the matched
%   `\textcolor{red}{\texttt{,}}' or `\textcolor{red}{\texttt{.}}' used as
%   the decimal separator. The variable \myhlc|\l_@@_luaset_part_dec_tl|
%   \emph{stores} the raw digits of the \emph{\enquote{finite decimal part}}
%   that follow the separator, and the variable \myhlc|\l_@@_luaset_part_period_tl|
%   \emph{stores} the digits of the \emph{\enquote{infinite periodic decimal part}}
%   that follow a semicolon `\textcolor{red}{\texttt{;}}'.
%
%   \smallskip
%
%    The variable \myhlc|\l_@@_luaset_has_integer_str| \emph{store} the state
%   |true| when it is \emph{only} an \emph{\enquote{integer}}. The variable
%   \myhlc|\l_@@_luaset_has_natural_str| \emph{store} the state  |true| when
%   it is \emph{only} an \emph{\enquote{natural number}}.
% \iffalse
%% Core base variables set in Lua module spintent.lua (\__spintent_luafun_clean_split_and_set:n)
% \fi
%    \begin{macrocode}
\tl_new:N  \l_@@_luaset_sign_tl
\tl_new:N  \l_@@_luaset_part_int_tl
\tl_new:N  \l_@@_luaset_dec_sep_tl
\tl_new:N  \l_@@_luaset_part_dec_tl
\tl_new:N  \l_@@_luaset_part_period_tl
\str_new:N \l_@@_luaset_has_integer_str
\str_new:N \l_@@_luaset_has_natural_str
%    \end{macrocode}
% \end{variable}
%
% \begin{variable}[int]{
%   \l_@@_luaset_only_part_int_str,
%   \l_@@_luaset_only_part_dec_str,
%   \l_@@_luaset_only_part_period_str,
%   \l_@@_luaset_format_part_int_str,
%   \l_@@_luaset_format_part_dec_str,
% }
%   The literal entries of the numbers that we will pass to the
%   \emph{intent function} |:number| of |\MathMLintent| are \emph{stored}
%   in \myhlc|\l_@@_luaset_only_part_int_str| for the \emph{\enquote{integer
%   part}}, in \myhlc|\l_@@_luaset_only_part_dec_str| for the
%   \emph{\enquote{finite decimal part}}, and in \myhlc|\l_@@_luaset_only_part_period_str|
%   for the \emph{\enquote{infinite periodic decimal part}}.
%
%   \smallskip
%
%   Formatted versions (\mykey[spnum]{cifras-sep}|=true|) are \emph{stored}
%   in \myhlc|\l_@@_luaset_format_part_int_str| for the
%   \emph{\enquote{integer part}} and \myhlc|\l_@@_luaset_format_part_dec_str| for
%   the \emph{\enquote{finite decimal part}}.
%    \begin{macrocode}
\str_new:N \l_@@_luaset_only_part_int_str
\str_new:N \l_@@_luaset_only_part_dec_str
\str_new:N \l_@@_luaset_only_part_period_str
\str_new:N \l_@@_luaset_format_part_int_str
\str_new:N \l_@@_luaset_format_part_dec_str
%    \end{macrocode}
% \end{variable}
%
% \begin{variable}[int]{
%   \l_@@_luaset_millons_str,
%   \l_@@_luaset_decimal_period_str,
%   \l_@@_luaset_not_decimal_period_str,
%   \l_@@_luaset_not_decimal_not_period_str,
%   \l_@@_luaset_dec_is_all_zeros_str,
%   \l_@@_luaset_has_exponent_str,
%   \l_@@_luaset_exponent_tl,
% }
%
%   \smallskip
%
%   The flag variable \myhlc|\l_@@_luaset_millons_str| indicates whether the
%   integer part translates exactly to a multiple of one million when
%   passed to |\MathMLintent|. To properly structure the reading and writing of numbers in
%   \mynvmath{MathML}, we use the flag variables \myhlc|\l_@@_luaset_decimal_period_str|,
%   \myhlc|\l_@@_luaset_not_decimal_period_str|, and \myhlc|\l_@@_luaset_not_decimal_not_period_str|.
%   The flag variable \myhlc|\l_@@_luaset_dec_is_all_zeros_str| monitors whether the \emph{\enquote{finite
%   decimal part}} is completely zero (useful for \emph{\enquote{pure periodic decimals}}).
%
%   \smallskip
%
%   The flag variable \myhlc|\l_@@_luaset_has_exponent_str| indicates
%   whether active \emph{\enquote{scientific notation}} was detected at the end of the
%   entered number, passed with `\textcolor{red}{\texttt{e}}' or
%   `\textcolor{red}{\texttt{E}}'. If true, it will store the digits of the
%   exponent in \myhlc|\l_@@_luaset_exponent_tl|.
%    \begin{macrocode}
\str_new:N \l_@@_luaset_millons_str
\str_new:N \l_@@_luaset_decimal_period_str
\str_new:N \l_@@_luaset_not_decimal_period_str
\str_new:N \l_@@_luaset_not_decimal_not_period_str
\str_new:N \l_@@_luaset_dec_is_all_zeros_str
\str_new:N \l_@@_luaset_has_exponent_str
\tl_new:N  \l_@@_luaset_exponent_tl
%    \end{macrocode}
% \end{variable}
%
% \begin{variable}[int]{
%   \l_@@_luaset_has_units_str,
%   \l_@@_luaset_status_str,
%   \l_@@_luaset_arg_above_tl,
%   \l_@@_luaset_arg_below_tl,
%   \l_@@_luaset_denom_has_number_str,
% }
%
% Since the function \myhlc|\@@_luafun_clean_split_and_set:n| is common to
% |\spnum| (\S\ref{cmd:spnum}) and |\spunit| (\S\ref{cmd:spunit}) (as well as others), it determines the following
% unit-related variables after processing the \mymarg{arguments} via the Lua module \myluamod.
%
% \smallskip
%
% First, it checks if there is any residual text \emph{after} the processed number
% that is a candidate for physical \emph{\enquote{units}}. It sets the flag variable
% \myhlc|\l_@@_luaset_has_units_str| accordingly. Following this, it
% verifies if this text contains more than one \enquote{/}. It then sets
% the variable \myhlc|\l_@@_luaset_status_str| to |valid| or |multislash|.
% If the structure is |valid|, it splits the unit into
% \myhlc|\l_@@_luaset_arg_above_tl| and \myhlc|\l_@@_luaset_arg_below_tl|.
% Finally, it analyzes \myhlc|\l_@@_luaset_arg_below_tl|; if the denominator
% starts with a number, it sets \myhlc|\l_@@_luaset_denom_has_number_str| to
% |true|, otherwise to |false|.
%    \begin{macrocode}
\str_new:N \l_@@_luaset_has_units_str
\str_new:N \l_@@_luaset_status_str
\tl_new:N  \l_@@_luaset_arg_above_tl
\tl_new:N  \l_@@_luaset_arg_below_tl
\str_new:N \l_@@_luaset_denom_has_number_str
%    \end{macrocode}
% \end{variable}
%
% \subsection{Common messages for text mode and math mode}
%
% Messages shared by the different commands for text mode or math mode.
% \iffalse
%% Common msm for command in math and text mode
% \fi
%    \begin{macrocode}
\msg_new:nnnn { spintent } { need-math-mode }
  {
    Math~mode~required~for~'#1'~\msg_line_context:.
  }
  {
    Use~inside~$...$~or~in~math~environment.
  }
\msg_new:nnnn { spintent } { need-text-mode }
  {
    Text~mode~required~for~'#1'~\msg_line_context:.
  }
  {
    Use~outside~of~$...$~or~math~environment.
  }
%    \end{macrocode}
%
% \subsection{The command \cs{spnum}}\label{cmd:spnum}
%
% The user command \ics*{spnum} is \emph{first} of \emph{\enquote{the big four}}
% provided by the \mypkg*[no-index]{spintent} package and is fundamental to the
% implementation of other commands. It handles the processing of numerical
% input, scientific notation, and units of measurement, using the Lua module \myluamod,
% which is responsible for separating and assigning variables defined in
% \mypkg{expl3}.
%
% \smallskip
%
% From the \mypkg{expl3} side, we handle the printing logic and the
% |intent| \emph{functions} passed to |\MathMLintent| to achieve a valid
% \mynvmath{MathML} structure that is accepted by \mynvmath{MathCAT}.
%
% \subsubsection{Definition of keys and variables for \cs{spnum}}
%
% \begin{variable}[int]{
%   \l_@@_spnum_read_terse_bool,
%   \l_@@_spnum_read_period_sep_str,
%   \l_@@_spnum_notation_sym_tl,
%   \l_@@_spnum_intent_read_sign_str,
%   \l_@@_spnum_intent_read_dec_sep_str,
% }
%
%   The variable \cs{l_@@_spnum_read_terse_bool} handled by the key \mykey[spnum]{read-sign}, controls the
%   |intent| passed to |\MathMLintent| to spoken the `$\mathcolor{red}{+}$'
%   or `$\mathcolor{red}{-}$'. If set to |clear|, it will say
%   \emph{\enquote{positivo}} or \emph{\enquote{negativo}} \emph{after} reading the
%   number; if set to |terse|, it will say \emph{\enquote{más}} or \emph{\enquote{menos}}
%   \emph{before} reading the number.
%
%  \smallskip
%
%   The variable \cs{l_@@_spnum_read_period_sep_str} handled by the key \mykey[spnum]{read-period-sep}, controls how the
%   decimal separator `\textcolor{red}{\texttt{;}}' (periodic part) is printed
%   and spoken when the \mymarg{argument} is a \emph{pure decimal periodic} and
%   is passed the |intent| to |\MathMLintent|. If set to |coma|, it will say
%   \emph{\enquote{coma}} and print `\textcolor{red}{,}', if set to |punto| it
%   will say \emph{\enquote{punto}} and print `\textcolor{red}{.}'.
%
%   \smallskip
%
%   The variable \cs{l_@@_spnum__notation_sym_tl}, handled by the key
%   \mykey[spnum]{notation-sym}, determines whether `$\mathcolor{red}{\times}$'
%   will be printed when set to |times| or `$\mathcolor{red}{\cdot}$'
%   when set to |cdot|.
%
%   \smallskip
%
%   The variable \cs{l_@@_spnum_intent_read_sign_str} will \emph{store} the
%   spoken translation used in the |intent| function passed to
%   |\MathMLintent| for the `$\mathcolor{red}{+}$' or
%   `$\mathcolor{red}{-}$' according to the value passed to the key
%   \mykey[spnum]{read-sign}.
%
%   \smallskip
%
%   The variable \cs{l_@@_spnum_intent_read_dec_sep_str} \emph{stores} the
%   spoken translation that will be used in the |intent| function passed to
%   |\MathMLintent| for reading the decimal separation as a
%   \emph{\enquote{coma}} or a \emph{\enquote{punto}} according to the
%   value found in the variable \myhlc|\l_@@_luaset_dec_sep_tl|.
% \iffalse
%% Internal vars for \spnum
% \fi
%    \begin{macrocode}
\bool_new:N \l_@@_spnum_read_terse_bool
\str_new:N  \l_@@_spnum_read_period_sep_str
\tl_new:N   \l_@@_spnum_notation_sym_tl
\str_new:N  \l_@@_spnum_intent_read_sign_str
\str_new:N  \l_@@_spnum_intent_read_dec_sep_str
%    \end{macrocode}
% \end{variable}
%
% \begin{macro}[int]{cifras-sep, read-sign, read-period-sep, notation-sym}
%   Now we add the keys \mykey[spnum]{cifras-sep}, \mykey[spnum]{read-sign},
%   \mykey[spnum]{read-period-sep} and \mykey[spnum]{notation-sym}.
% \iffalse
%% Keys for \spnum
% \fi
%    \begin{macrocode}
\keys_define:nn { spintent/spnum }
  {
    cifras-sep                .bool_set:N = \l_@@_spnum_cifras_sep_bool,
    cifras-sep                .initial:n = true,
    cifras-sep                .value_required:n = true,
    read-sign                 .choices:nn =
      { terse , clear }
      {
        \int_case:nn { \l_keys_choice_int }
          {
            {1} { \bool_set_true:N  \l_@@_spnum_read_terse_bool }
            {2} { \bool_set_false:N \l_@@_spnum_read_terse_bool }
          }
      },
    read-sign                 .initial:n  = terse,
    read-sign                 .value_required:n = true,
    read-sign / unknown       .code:n =
      \msg_error:nneee { spintent } { unknown-choice }
        { read-sign } { terse , ~clear } { \exp_not:n {#1} },
    read-period-sep           .choices:nn =
      { coma , punto }
      {
        \int_case:nn { \l_keys_choice_int }
          {
            {1} { \str_set:Nn \l_@@_spnum_read_period_sep_str { coma } }
            {2} { \str_set:Nn \l_@@_spnum_read_period_sep_str { punto } }
          }
      },
    read-period-sep           .initial:n  = coma,
    read-period-sep           .value_required:n = true,
    read-period-sep / unknown .code:n =
      \msg_error:nneee { spintent } { unknown-choice }
        { read-period-sep } { coma, ~punto } { \exp_not:n {#1} },
    notation-sym              .choices:nn =
      { times , cdot }
      {
        \int_case:nn { \l_keys_choice_int }
          {
            {1} { \tl_set:Nn \l_@@_spnum_notation_sym_tl { \times } }
            {2} { \tl_set:Nn \l_@@_spnum_notation_sym_tl { \cdot } }
          }
      },
    notation-sym              .initial:n  = cdot,
    notation-sym              .value_required:n = true,
    notation-sym / unknown    .code:n =
      \msg_error:nneee { spintent } { unknown-choice }
        { notation-sym } { times , ~ cdot } { \exp_not:n {#1} },
    unknown                   .code:n =
      {
        \str_set:Nn \l_@@_command_name_str { spnum }
        \@@_unknown_keys_cmd:n {#1}
      },
  }
%    \end{macrocode}
% \end{macro}
%
% \subsubsection{Internal functions for \cs{spnum}}
%
% \begin{macro}[int]{\@@_spnum_render_part:n, \@@_spnum_render_int:, \@@_spnum_render_dec:}
%   The function \myhlc|\@@_spnum_render_part:n| processes the
%   \mymarg{argument} passed in `|#1|' used by the functions \myhlc|\@@_spnum_render_int:| and
%   \myhlc|\@@_spnum_render_dec:| to define the \emph{intent function} |:number| for
%   the integer and decimal part passed to |\MathMLintent|. It evaluates
%   whether the key \mykey[spnum]{cifras-sep} is active or not.
% \iffalse
%% Internal functions for \spnum
% \fi
%    \begin{macrocode}
\cs_new_protected:Nn \@@_spnum_render_part:n
  {
    \MathMLintent { \str_use:c { l_@@_luaset_only_part_#1_str } :number }
      {
        \bool_if:NTF \l_@@_spnum_cifras_sep_bool
          { \str_use:c { l_@@_luaset_format_part_#1_str } }
          { \tl_use:c { l_@@_luaset_part_#1_tl } }
      }
  }
\cs_new_protected:Nn \@@_spnum_render_int:
  {
    \@@_spnum_render_part:n { int }
  }
\cs_new_protected:Nn \_@@_spnum_render_dec:
  {
    \@@_spnum_render_part:n { dec }
  }
%    \end{macrocode}
% \end{macro}
%
% \begin{macro}[int]{\@@_spnum_render_only_period:, \@@_spnum_print_period_bar:,\@@_spnum_render_period_sep:}
%   The function \myhlc|\@@_spnum_render_only_period:| isolates the
%   \emph{\enquote{infinite periodic decimal part}} from the
%   \emph{\enquote{finite decimal part}}, assigning the \emph{intent
%   function} |:number| passed to |\MathMLintent| which will be used
%   by |\MathMLarg| in the function \myhlc|\@@_spnum_print_period_bar:|.
%    \begin{macrocode}
\cs_new_protected:Nn \@@_spnum_render_only_period:
  {
    \MathMLintent { \l_@@_luaset_only_part_period_str :number }
      {
        \l_@@_luaset_part_period_tl
      }
  }
%    \end{macrocode}
%   The function \myhlc|\@@_spnum_print_period_bar:| inserts the function
%   \myhlc|\@@_spnum_render_only_period:| as an argument within the \emph{intent
%   function} |_periódico:postfix($dp)| passed to |\MathMLintent|, generating
%   the correct spoken for the \emph{\enquote{infinite periodic decimal
%   part}}.
%    \begin{macrocode}
\cs_new_protected:Nn \@@_spnum_print_period_bar:
  {
    \MathMLintent { _periódico:postfix($dp) }
      {
        {
          \overline
            {
              \MathMLarg { dp } { \@@_spnum_render_only_period: }
            }
        }
      }
  }
%    \end{macrocode}
%   The function \myhlc|\@@_spnum_render_period_sep:| sets how the decimal
%   separator handled by the key \mykey[spnum]{read-period-sep} will be
%   printed and spoken when a \emph{\enquote{pure periodic decimal}} is passed.
%    \begin{macrocode}
\cs_new_protected:Nn \@@_spnum_render_period_sep:
  {
    \str_if_eq:VnTF \l_@@_spnum_read_period_sep_str { coma }
      {
        \MathMLintent { _coma:prefix($pp) }
          {
            \mathord{,}
            \MathMLarg { pp } { \@@_spnum_print_period_bar: }
          }
       }
       {
         \MathMLintent { _punto:prefix($pp) }
           {
             \mathord{.}
             \MathMLarg { pp } { \@@_spnum_print_period_bar: }
           }
       }
  }
%    \end{macrocode}
% \end{macro}
%
% \begin{macro}[int]{\@@_spnum_render_decimal_sep:}
%  The function \myhlc|\@@_spnum_render_decimal_sep:| parses the value of
%  the variable \myhlc|\l_@@_luaset_dec_sep_tl| and sets the variable
%  \cs{l_@@_spnum_intent_read_dec_sep_str} with the function
%  |_coma:prefix($dp)| for reading a \emph{\enquote{coma}} or |_punto:prefix($dp)| for
%  reading a \emph{\enquote{punto}} to |\MathMLintent|. We print the decimal separator
%  stored in \myhlc|\l_@@_luaset_dec_sep_tl| using |\mathord| to remove
%  spaces around it and then use the command |\MathMLarg| where we insert
%  \myhlc|\@@_invisible_sep:| for compatibility with \mynvmath{MathML} and
%  finally call the function \myhlc|\@@_spnum_render_dec:| to print the
%  finite decimal part.
%    \begin{macrocode}
\cs_new_protected:Nn \@@_spnum_render_decimal_sep:
  {
    \str_if_eq:VnTF \l_@@_luaset_dec_sep_tl { , }
      { \str_set:Nn \l_@@_spnum_intent_read_dec_sep_str { _coma } }
      { \str_set:Nn \l_@@_spnum_intent_read_dec_sep_str { _punto } }
    \MathMLintent { \l_@@_spnum_intent_read_dec_sep_str :prefix($dp) }
      {
        \mathord { \l_@@_luaset_dec_sep_tl }
        \MathMLarg { dp }
          {
            \@@_invisible_sep:
            \@@_spnum_render_dec:
          }
      }
    % period part
    \str_if_eq:VnT \l_@@_luaset_decimal_period_str { true }
      {
        \str_if_eq:VnTF \l_@@_luaset_dec_is_all_zeros_str { true }
          {
            \@@_invisible_sep:
            \@@_spnum_print_period_bar:
          }
          {
            \MathMLintent { _con:prefix($pnum) }
              {
                \@@_invisible_sep:
                \MathMLarg { pnum }
                  {
                    \@@_spnum_print_period_bar:
                  }
              }
          }
      }
  }
%    \end{macrocode}
% \end{macro}
%
%
% \begin{macro}[int]{\@@_spnum_print_part_dec:, \@@_spnum_print_part_period:}
%   The function \myhlc|\@@_spnum_print_part_dec:| checks if a decimal part exists.
%   If it does, it determines whether the decimal separator is a comma or a period
%   to set the correct speech intent (\emph{\enquote{coma}} or \emph{\enquote{punto}}).
%   It then prints the separator, renders the decimal digits, and if a periodic
%   sequence follows, it decides how to attach the periodic bar.
%
%   \smallskip
%
%   The function \myhlc|\@@_spnum_print_part_period:| handles cases where a
%   periodic part exists. It evaluates if it should render the periodic separator
%   directly, particularly when the number has a periodic part but lacks a standard
%   decimal part.
%
%   \smallskip
%
%    \begin{macrocode}
\cs_new_protected:Nn \@@_spnum_print_part_dec:
  {
    \tl_if_empty:NF \l_@@_luaset_part_dec_tl
      {
        \@@_spnum_render_decimal_sep:
      }
  }
\cs_new_protected:Nn \@@_spnum_print_part_period:
  {
    \tl_if_empty:NF \l_@@_luaset_part_period_tl
      {
        \str_if_eq:VnTF \l_@@_luaset_not_decimal_period_str { true }
          { \@@_spnum_render_period_sep: }
          {
            \tl_if_empty:NT \l_@@_luaset_part_dec_tl
              {
                \@@_spnum_render_period_sep:
              }
          }
      }
  }
%    \end{macrocode}
% \end{macro}
%
% \begin{macro}[int]{\@@_spnum_print_integer:, \@@_spnum_print_num_start:, \@@_spnum_print_number_core:}
%   The function \myhlc|\@@_spnum_print_integer:| checks if the number consists
%   only of an integer. If so, it simply renders it. If a decimal or periodic
%   part exists, it sets up a \mynvmath{MathML} intent to link the integer part (respecting
%   digit separators if enabled) to the subsequent decimal or periodic functions.
%
%   \smallskip
%
%   The function \myhlc|\@@_spnum_print_num_start:| checks if an explicit sign
%   is present. If it finds one, it evaluates the terse reading flag to decide
%   whether the sign should be read as a prefix (\emph{\enquote{más}} or
%   \emph{\enquote{menos}}) or as a postfix (\emph{\enquote{positivo}} or
%   \emph{\enquote{negativo}}), and wraps the intent around the rest of the number.
%
%   \smallskip
%
%   The function \myhlc|\@@_spnum_print_number_core:| acts as the main wrapper.
%   It triggers the sign and base number evaluation, and then checks if a scientific
%   exponent is present. If it is, it appends the correct base-10 notation symbol
%   and exponent to complete the sequence.
%
%   \smallskip
%
% \begin{important}*
%   Note for \myhlc|\@@_invisible_sep:|.
% \end{important}
%    \begin{macrocode}
\cs_new_protected:Nn \@@_spnum_print_integer:
  {
    \str_if_eq:VnTF \l_@@_luaset_not_decimal_not_period_str { true }
      { \@@_spnum_render_int: }
      {
        \MathMLintent { \l_@@_luaset_only_part_int_str :number:prefix($next) }
          {
            \bool_if:NTF \l_@@_spnum_cifras_sep_bool
              { \l_@@_luaset_format_part_int_str }
              { \l_@@_luaset_part_int_tl }
            \MathMLarg { next }
              {
                \tl_if_empty:NTF \l_@@_luaset_part_dec_tl
                  { \@@_spnum_print_part_period: }
                  { \@@_spnum_print_part_dec: }
              }
          }
      }
  }
\cs_new_protected:Nn \@@_spnum_print_num_start:
  {
    \tl_if_empty:NTF \l_@@_luaset_sign_tl
      {
        \@@_invisible_sep: % need
        \@@_spnum_print_integer:
      }
      {
        \bool_if:NTF \l_@@_spnum_read_terse_bool
          {
            \str_if_eq:VnTF \l_@@_luaset_sign_tl { + }
              { \str_set:Nn \l_@@_spnum_intent_read_sign_str { _más:prefix($n)} }
              { \str_set:Nn \l_@@_spnum_intent_read_sign_str { _menos:prefix($n)} }
          }
          {
            \str_if_eq:VnTF \l_@@_luaset_sign_tl { + }
              { \str_set:Nn \l_@@_spnum_intent_read_sign_str { _positivo:postfix($n) } }
              { \str_set:Nn \l_@@_spnum_intent_read_sign_str { _negativo:postfix($n) } }
          }
        \@@_invisible_sep: % need
        \MathMLintent { \l_@@_spnum_intent_read_sign_str }
          {
            \l_@@_luaset_sign_tl
            \MathMLarg { n } { \@@_spnum_print_integer: }
          }
      }
  }
\cs_new_protected:Nn \@@_spnum_print_number_core:
  {
    \@@_spnum_print_num_start:
    \str_if_eq:VnT \l_@@_luaset_has_exponent_str { true }
      {
        \MathMLintent { _por } { \l_@@_spnum_notation_sym_tl }
        10^{ \l_@@_luaset_exponent_tl }
      }
  }
%    \end{macrocode}
% \end{macro}
%
% \subsubsection{Definition of error messages}
%
% \begin{macro}[int]{spnum-has-units, spnum-no-digits}
%   We define two error messages for the |\spnum| command. The first
%   \enquote{error} triggers if the user includes \emph{\enquote{units}} in the
%   \mymarg{argument} passed in `|#1|', advising them to use the |\spunit|
%   command instead. The second \enquote{error} is raised if the provided
%   \mymarg{argument} contains no numerical digits.
% \iffalse
%% messages errors for \spnum
% \fi
%    \begin{macrocode}
\msg_new:nnnn { spintent } { spnum-has-units }
  {
    The~\token_to_str:N \spnum \c_space_tl command~does~not~accept~units~('#1').
  }
  {
    Use~\token_to_str:N \spunit \c_space_tl if~you~need~to~print~numbers~with~units.
  }
\msg_new:nnnn { spintent } { spnum-no-digits }
  {
    The~\token_to_str:N \spnum \c_space_tl command~requires~a~valid~numerical~value.
  }
  {
    You~provided~'#1',~which~contains~no~digits.
  }
%    \end{macrocode}
% \end{macro}
%
% \subsubsection{Implementation of the \cs[no-index]{spum} command}
%
% \begin{macro}{\spnum}
%   The public user command |\spnum| accepts an optional \myoarg{key
%   \textnormal{\textcolor{gray}{=}} val} argument in `|#1|' and a mandatory
%   \mymarg{rational number} in `|#2|'. It first checks if it is executed inside
%   math mode, raising an error if it is not. If correct, it opens a local group
%   to apply the options and passes the input to Lua for parsing via
%   \myhlc|\@@_luafun_clean_split_and_set:n|.
%
%   \smallskip
%
%   Before printing, it validates the extracted data. It throws an error if it
%   finds trailing units in the input or if the integer part is completely empty
%   (no valid digits). If all checks pass, it calls \myhlc|\@@_spnum_print_number_core:|
%   to render the final number, and finally closes the group.
% \iffalse
%% Public command \spnum
% \fi
%    \begin{macrocode}
\NewDocumentCommand \spnum { O{} m }
  {
    \mode_if_math:TF
      {
        \group_begin:
          \keys_set:nn { spintent/spnum } { #1 }
          \@@_luafun_clean_split_and_set:n { \exp_not:n { #2 } }
          \str_if_eq:VnTF \l_@@_luaset_has_units_str { true }
            { \msg_error:nnn { spintent } { spnum-has-units } { #2 } }
            {
              \tl_if_empty:NTF \l_@@_luaset_part_int_tl
                { \msg_error:nnn { spintent } { spnum-no-digits } { #2 } }
                { \@@_spnum_print_number_core: }
            }
        \group_end:
      }
      { \msg_error:nnn { spintent } { need-math-mode } { \spnum } }
  }
%    \end{macrocode}
% \end{macro}
%
% \subsection{The command \cs{spunit}}\label{cmd:spunit}
%
% The user command \ics*{spunit} is the \emph{second} of \emph{\enquote{the big four}}
% provided by the \mypkg*[no-index]{spintent} package and is used to parse and print physical units
% and angular systems. It sends the \mymarg{argument} to Lua module \myluamod{} via
% \myhlc|\@@_luafun_clean_split_and_set:n| to separate and clean the \mymarg{argument}.
%
% \smallskip
%
% This ensures that all \emph{\enquote{units}} (such as SÍ fractions, products, or sexagesimal
% angles) are processed consistently by the \myluamod{}.
%
% \subsubsection{Definition of variables}
%
% \begin{variable}[int]{
%   \l_@@_spunit_luaset_canonical_str,
%   \l_@@_spunit_luaset_lookup_status_str,
%   \l_@@_spunit_luaset_read_str,
%   \l_@@_spunit_luaset_is_compact_str,
%   \l_@@_spunit_luaset_compact_exp_str,
%   \l_@@_spunit_luaset_is_sexagesimal_str,
%   \l_@@_spunit_luaset_status_str
% }
% These string variables store the data returned by the Lua module \myluamod{}.
% They hold the canonical \emph{\enquote{unit symbol}}, the dictionary lookup
% status, the correct speech text for \mynvmath{NVDA/MathCAT}, and specific flags that
% indicate if the \emph{\enquote{unit}} is compact or represents a sexagesimal angle.
% \iffalse
%% Variables for \spunit
% \fi
%    \begin{macrocode}
\str_new:N  \l_@@_spunit_luaset_canonical_str
\str_new:N  \l_@@_spunit_luaset_lookup_status_str
\str_new:N  \l_@@_spunit_luaset_read_str
\str_new:N  \l_@@_spunit_luaset_is_compact_str
\str_new:N  \l_@@_spunit_luaset_compact_exp_str
\str_new:N  \l_@@_spunit_luaset_is_sexagesimal_str
\str_new:N  \l_@@_spunit_luaset_status_str
%    \end{macrocode}
% \end{variable}
%
% \begin{variable}[int]{
%   \l_@@_spunit_exponent_tl,
%   \l_@@_spunit_unit_name_tl,
%   \l_@@_spunit_read_slash_str,
%   \l_@@_spunit_exp_seq,
%   \l_@@_spunit_mult_seq,
%   \l_@@_spunit_is_mult_bool,
%   \l_@@_spunit_read_cdot_bool
% }
% These variables are used to process and format the \emph{\enquote{unit}}. The
% sequences \myhlc|\l_@@_spunit_mult_seq| and \myhlc|\l_@@_spunit_exp_seq| split
% compound \emph{\enquote{units}} at multiplication `|*|' or exponent `|^|' boundaries.
%
% \smallskip
%
% The boolean variables control specific rendering options, such as whether
% the multiplication dot `|\cdot|' is explicitly read aloud or silently skipped.
%    \begin{macrocode}
\tl_new:N   \l_@@_spunit_exponent_tl
\tl_new:N   \l_@@_spunit_unit_name_tl
\str_new:N  \l_@@_spunit_read_slash_str
\seq_new:N  \l_@@_spunit_exp_seq
\seq_new:N  \l_@@_spunit_mult_seq
\bool_new:N \l_@@_spunit_is_mult_bool
\bool_new:N \l_@@_spunit_read_cdot_bool
%    \end{macrocode}
% \end{variable}
%
% \subsubsection{Definition of error messages}
%
% \begin{macro}[int]{spunit-invalid-unit, spunit-no-units, inline-numeric-denominator}
%   We define three error messages for the |\spunit| command. The first
%   \enquote{error} triggers if the \emph{\enquote{unit expression}} provided
%   in `|#1|' is unknown or malformed. The second \enquote{error} occurs
%   if no \emph{\enquote{units}} are found in the \mymarg{argument},
%   advising the user to use |\spnum| for pure numbers. The third \enquote{error}
%   is raised if a numeric coefficient is illegally placed in the denominator
%   of an inline fraction (after a slash).
% \iffalse
%% messages errors for \spunit
% \fi
%    \begin{macrocode}
\msg_new:nnnn { spintent } { spunit-invalid-unit }
  {
    The~expression~'#1'~contains~an~unknown~or~malformed~unit~structure.
  }
  {
    Check~spelling~or~verify~the~unit~is~defined~in~the~dictionary.
  }
\msg_new:nnnn { spintent } { spunit-no-units }
  {
    The~\token_to_str:N \spunit \c_space_tl command~requires~at~least~one~unit~('#1').
  }
  {
    Use~\token_to_str:N \spnum \c_space_tl if~you~only~need~to~format~a~pure~number.
  }
\msg_new:nnnn { spintent } { inline-numeric-denominator }
  {
    Illegal~numeric~coefficient~in~inline~denominator~'#1'.
  }
  {
    The~inline~mode~of~\token_to_str:N \spunit \c_space_tl
    does~not~allow~numbers~after~the~slash~(/).
  }
%    \end{macrocode}
% \end{macro}
%
% \subsubsection{Definition of keys for \cs{spunit}}
%
% \begin{macro}[int]{read-cdot, read-slash, cifras-sep, read-period-sep, notation-sym,}
%   Now we add the keys \mykey[spunit]{print-cdot}, \mykey[spunit]{read-cdot}, \mykey[spunit]{read-slash},
%   \mykey[spunit]{cifras-sep}, \mykey[spunit]{read-period-sep} and \mykey[spunit]{notation-sym}.
%    \begin{macrocode}
\keys_define:nn { spintent/spunit }
  {
    print-cdot           .bool_set:N = \l_@@_spunit_print_cdot_bool,
    print-cdot           .initial:n  = false,
    print-cdot           .value_required:n = true,
    read-cdot            .bool_set:N = \l_@@_spunit_read_cdot_bool,
    read-cdot            .initial:n  = false,
    read-cdot            .value_required:n = true,
    read-slash           .choices:nn =
      { por , partido-por }
      {
        \int_case:nn { \l_keys_choice_int }
          {
            {1} { \str_set:Nn \l_@@_spunit_read_slash_str { por } }
            {2} { \str_set:Nn \l_@@_spunit_read_slash_str { partido-por } }
          }
      },
    read-slash           .initial:n  = por,
    read-slash           .value_required:n = true,
    read-slash / unknown .code:n =
      \msg_error:nneee { spintent } { unknown-choice }
        { read-slash } { por, ~partido-por } { \exp_not:n {#1} },
    cifras-sep           .meta:nn = { spintent/spnum } { cifras-sep = {#1} },
    cifras-sep           .value_required:n = true,
    read-period-sep      .meta:nn = { spintent/spnum } { read-period-sep = {#1} },
    read-period-sep      .value_required:n = true,
    notation-sym         .meta:nn = { spintent/spnum } { notation-sym = {#1} },
    notation-sym         .value_required:n = true,
    unknown              .code:n =
      {
        \str_set:Nn \l_@@_command_name_str { spunit }
        \@@_unknown_keys_cmd:n {#1}
      },
  }
%    \end{macrocode}
% \end{macro}
%
% \subsubsection{Internal functions for \cs[no-index]{spunit}}
%
% \begin{macro}[int]{\@@_spunit_format_canonical:, \@@_spunit_process_base_exp:n}
%   The function \myhlc|\@@_spunit_format_canonical:| handles the visual formatting
%   of the \emph{\enquote{unit symbol}}. It specifically checks for the \emph{liter symbol}
%   `ℓ' to print it as |\ell|, and formats all other valid \emph{\enquote{units}}
%   in standard upright text using |\mathrm|.
%
%   \smallskip
%
%   The function \myhlc|\@@_spunit_process_base_exp:n| evaluates a \emph{\enquote{single unit}}
%   along with its exponent. It first looks up the unit `|#1|` in the dictionary.
%   If the \emph{\enquote{unit}} is not found, it flags it as invalid. If found, it handles the
%   spacing and multiplication dot `|\cdot|' intents, and then renders the base
%   unit and its exponent (either using a pre-defined compact exponent from Lua
%   or by manually splitting the string at the `|^|' symbol).
%    \begin{macrocode}
\cs_new_protected:Nn \@@_spunit_format_canonical:
  {
    \str_case:VnF \l_@@_spunit_luaset_canonical_str
      {
        { ℓ } { \ell }
      }
      { \mathrm { \l_@@_spunit_luaset_canonical_str } }
  }
\cs_new_protected:Nn \@@_spunit_process_base_exp:n
  {
    \tl_clear:N \l_@@_spunit_exponent_tl
    \@@_luafun_spunit_lookup_alias:n { #1 }
    \str_if_eq:VnTF \l_@@_spunit_luaset_lookup_status_str { notfound }
      {
        \str_set:Nn \l_@@_luaset_status_str { invalid }
      }
      {
        \bool_if:NF \l_@@_spunit_is_mult_bool
          {
            \bool_if:NTF \l_@@_spunit_print_cdot_bool
              {
                \bool_if:NTF \l_@@_spunit_read_cdot_bool
                  { \MathMLintent { por } { \cdot } }
                  { \MathMLintent { :silent } { \cdot } }
              }
              {
                \bool_if:NTF \l_@@_spunit_read_cdot_bool
                  { \MathMLintent { por } { \, } }
                  { \, }
              }
          }
        \str_if_eq:VnTF \l_@@_spunit_luaset_is_compact_str { true }
          {
            {
              \MathMLintent { \l_@@_spunit_luaset_read_str }
                {
                  \@@_spunit_format_canonical:
                }
            }
            ^ { \l_@@_spunit_luaset_compact_exp_str }
          }
          {
            \seq_set_split:Nnn \l_@@_spunit_exp_seq { ^ } { #1 }
            \seq_pop_left:NN \l_@@_spunit_exp_seq \l_@@_spunit_unit_name_tl
            \seq_if_empty:NF \l_@@_spunit_exp_seq
              {
                \seq_pop_left:NN \l_@@_spunit_exp_seq \l_@@_spunit_exponent_tl
              }
            \MathMLintent { \l_@@_spunit_luaset_read_str }
              {
                \@@_spunit_format_canonical:
              }
            \tl_if_empty:NF \l_@@_spunit_exponent_tl
              {
                ^ { \l_@@_spunit_exponent_tl }
              }
          }
        \bool_set_false:N \l_@@_spunit_is_mult_bool
      }
  }
%    \end{macrocode}
% \end{macro}
%
% \begin{macro}[int]{\@@_spunit_process_mult:n, \@@_spunit_process_mult:V}
%   The function \myhlc|\@@_spunit_process_mult:n| handles compound \emph{\enquote{units}}
%   joined by an explicit multiplication star `|*|'. It splits the \mymarg{argument} at
%   each \emph{\enquote{star symbol}} and applies \myhlc|\@@_spunit_process_base_exp:n|
%   to evaluate and render each resulting unit individually.
%    \begin{macrocode}
\cs_new_protected:Nn \@@_spunit_process_mult:n
  {
    \bool_set_true:N \l_@@_spunit_is_mult_bool
    \seq_set_split:Nnn \l_@@_spunit_mult_seq { * } { #1 }
    \seq_map_function:NN \l_@@_spunit_mult_seq \@@_spunit_process_base_exp:n
  }
\cs_generate_variant:Nn \@@_spunit_process_mult:n { V }
%    \end{macrocode}
% \end{macro}
%
% \begin{macro}[int]{\@@_print_spunit_hierarchy:}
%   The function \myhlc|\@@_print_spunit_hierarchy:| handles the layout of inline
%   \emph{\enquote{unit}} fractions. It first processes the \emph{\enquote{units}}
%   before the slash (the numerator). If a denominator exists, it prints the slash |/| `$/$'
%   with its corresponding speech intent, and then processes the remaining \emph{\enquote{units}}.
%    \begin{macrocode}
\cs_new_protected:Nn \@@_print_spunit_hierarchy:
  {
    \tl_if_empty:NF \l_@@_luaset_arg_above_tl
      {
        \@@_spunit_process_mult:V \l_@@_luaset_arg_above_tl
        \tl_if_empty:NF \l_@@_luaset_arg_below_tl
          {
            \MathMLintent { \l_@@_spunit_read_slash_str } { / }
            \@@_spunit_process_mult:V \l_@@_luaset_arg_below_tl
          }
      }
  }
%    \end{macrocode}
% \end{macro}
%
% \begin{macro}[int]{\@@_spunit_safe_exec:n}
%   The function \myhlc|\@@_spunit_safe_exec:n| performs the main validation
%   and layout for \emph{\enquote{units}}. It first checks if the \mymarg{argument}
%   has units, is a valid expression, and does not contain illegal numbers
%   in the denominator, throwing the appropriate \enquote{errors} if any check fails.
%
%   \smallskip
%
%   If the unit is valid, it checks if it is a sexagesimal angle (degrees,
%   minutes, or seconds). For sexagesimal \emph{\enquote{units}}, it prints
%   the number and the symbol \emph{\enquote{without any space in between}},
%   applying a slight negative space `|\!|' for |arcseconds|. For standard units,
%   it prints the number followed by a thin space `|\,|' and then processes
%   the unit sequence.
%    \begin{macrocode}
\cs_new_protected:Nn \@@_spunit_safe_exec:n
  {
    \str_if_eq:VnTF \l_@@_luaset_has_units_str { false }
      { \msg_error:nnn { spintent } { spunit-no-units } { #1 } }
      {
        \str_if_eq:VnTF \l_@@_luaset_status_str { valid }
          {
            \str_if_eq:VnTF
              \l_@@_luaset_denom_has_number_str { true }
              {
                \msg_error:nnn { spintent }
                  { inline-numeric-denominator } { #1 }
              }
              {
                \@@_luafun_spunit_lookup_alias:V \l_@@_luaset_arg_above_tl
                \str_if_eq:VnTF \l_@@_spunit_luaset_is_sexagesimal_str { true }
                  {
                    \tl_if_empty:NF \l_@@_luaset_part_int_tl
                      {
                        \@@_spnum_print_number_core:
                      }
                    \MathMLintent { \l_@@_spunit_luaset_read_str }
                      {
                        \mathord
                          {
                            \str_case:Vn \l_@@_spunit_luaset_canonical_str
                              {
                                { ″ } { \prime \! \prime }
                                { ′ } { \prime           }
                                { ° } { °                }
                              }
                          }
                      }
                  }
                  {
                    \tl_if_empty:NF \l_@@_luaset_part_int_tl
                      {
                        \@@_spnum_print_number_core: \,
                      }
                    \@@_print_spunit_hierarchy:
                  }
              }
          }
          { \msg_error:nnn { spintent } { spunit-invalid-unit } { #1 } }
      }
  }
%    \end{macrocode}
% \end{macro}
%
% \subsubsection{Implementation of the \cs{spunit} command}
%
% \begin{macro}{\spunit}
%   The public user command |\spunit| accepts an optional \myoarg{key
%   \textnormal{\textcolor{gray}{=}} val} argument in `|#1|' and a mandatory
%   \mymarg{unit expression} in `|#2|'. It first checks if it is executed inside
%   \emph{math mode}, raising an \enquote{error} if it is not. If correct, it opens a local group
%   to apply the options and passes the \mymarg{argument} to Lua module for parsing via
%   \myhlc|\@@_luafun_clean_split_and_set:n| function.
%
%   \smallskip
%
%   Finally, it delegates the execution to the function \myhlc|\@@_spunit_safe_exec:n|
%   to perform the safety checks and render the parsed data, before closing the group.
%    \begin{macrocode}
\NewDocumentCommand \spunit { O{} m }
  {
    \mode_if_math:TF
      {
        \group_begin:
          \keys_set:nn { spintent/spunit } { #1 }
          \@@_luafun_clean_split_and_set:n { \exp_not:n { #2 } }
          \@@_spunit_safe_exec:n { #2 }
        \group_end:
      }
      { \msg_error:nnn { spintent } { need-math-mode } { \spunit } }
  }
%    \end{macrocode}
% \end{macro}
%
% \subsection{The commands \cs{spnewunit} and \cs{spunitalias}}\label{cmd:spnewunit}
%
% The user commands |\spnewunit| and |\spunitalias| allow users to add
% new custom \emph{\enquote{units}} or \emph{\enquote{aliases}} to the
% package's dictionary. Similar to |\spunit|, these commands delegate
% the processing and validation to the \myluamod{} via
% \myhlc|\@@_luafun_define_custom_unit:nn| and \myhlc|\@@_luafun_define_spunit_alias:nn|.
%
% \smallskip
%
% During the definition, the Lua module checks if the new \emph{\enquote{unit}}
% or \emph{\enquote{alias}} is valid and ensures it does not conflict with existing \emph{\enquote{units}}.
% It returns the result of this operation to \TeX{} using the \myhlc|\l_@@_spunit_luaset_status_str| variable.
% If a collision is detected or the \mymarg{argument} is invalid, the package throws the
% corresponding \enquote{error} message.
%
% \subsubsection{Definition of error messages}
%
% \begin{macro}[int]{spnewunit-duplicate, spunitalias-duplicate, spunitalias-invalid-expr}
%   We define three error messages for the creation of \emph{\enquote{units}}
%   and \emph{\enquote{aliases}}. The first two \enquote{errors} trigger if
%   the user attempts to define a \emph{\enquote{new unit}} or alias passed in `|#1|'
%   that already exists in the dictionary. The third \enquote{error}
%   is raised by |\spunitalias| if the target expression provided in `|#2|'
%   is invalid (for example, if it contains unknown \emph{\enquote{units}}
%   or incorrect syntax).
% \iffalse
%% messages errors for \spnewunit and \spunitalias
% \fi
%    \begin{macrocode}
\msg_new:nnnn { spintent } { spnewunit-duplicate }
  {
    The~unit~symbol~or~alias~'#1'~is~already~defined!
  }
  {
    You~cannot~redefine~with~\token_to_str:N \spnewunit.
  }
\msg_new:nnnn { spintent } { spunitalias-duplicate }
  {
    The~alias~'#1'~is~already~defined~or~reserved!
  }
  {
    You~cannot~redefine~an~existing~unit~or~another~alias.
  }
\msg_new:nnnn { spintent } { spunitalias-invalid-expr }
  {
    The~expression~'#2'~for~alias~'#1'~is~invalid.
  }
  {
    It~must~contain~only~defined~units,~exponents~and~one~'/'.
  }
%    \end{macrocode}
% \end{macro}
%
% \begin{macro}{\spnewunit}
%   The public command |\spnewunit| defines a new custom unit. It passes
%   the new unit symbol `|#1|' and its speech text `|#2|' to Lua via
%   \myhlc|\@@_luafun_define_custom_unit:nn|. If the module reports that the
%   unit already exists, it throws the corresponding duplicate error.
%    \begin{macrocode}
\NewDocumentCommand \spnewunit { m m }
  {
    \mode_if_math:TF
      {
        \msg_error:nnn { spintent } { need-text-mode } { \spnewunit }
      }
      {
        \@@_luafun_define_custom_unit:nn { #1 } { #2 }
        \str_if_eq:VnT \l_@@_spunit_luaset_status_str { duplicate }
          {
            \msg_error:nnn { spintent } { spnewunit-duplicate } { #1 }
          }
      }
  }
%    \end{macrocode}
% \end{macro}
%
% \begin{macro}{\spunitalias}
%   The public command |\spunitalias| creates a shortcut for a complex
%   unit expression. It passes the new alias name `|#1|' and the target
%   expression `|#2|' to Lua via \myhlc|\@@_luafun_define_spunit_alias:nn|.
%   It throws an error if the alias already exists, or if the provided target
%   expression contains invalid syntax or unknown units.
%    \begin{macrocode}
\NewDocumentCommand \spunitalias { m m }
  {
    \mode_if_math:TF
      {
        \msg_error:nnn { spintent } { need-text-mode } { \spunitalias }
      }
      {
        \@@_luafun_define_spunit_alias:nn { #1 } { #2 }
        \str_if_eq:VnT \l_@@_spunit_luaset_status_str { duplicate }
          {
            \msg_error:nnn { spintent } { spunitalias-duplicate } { #1 }
          }
        \str_if_eq:VnT \l_@@_spunit_luaset_status_str { invalid-expr }
          {
            \msg_error:nnnn { spintent } { spunitalias-invalid-expr } { #1 } { #2 }
          }
      }
  }
%    \end{macrocode}
% \end{macro}
%
% \subsection{The command \cs{spmoney}}\label{cmd:spmoney}
%
% The user command \ics*{spmoney} is the \emph{third} of
% \emph{\enquote{the big four}} provided by the \mypkg*[no-index]{spintent} package,
% designed to format monetary values and currencies. It delegates
% \mymarg{argument} parsing to the Lua module \myluamod{} via
% \myhlc|\@@_luafun_clean_split_and_set:n|. The module then retrieves the
% correct \emph{\enquote{currency symbol}}, pluralization rules, and layout
% structures from the database.
%
% \smallskip
%
% The command accepts \emph{optional arguments} to override the default
% \emph{\enquote{currency symbol}} or enable \emph{\enquote{subunit}} rendering.
% Finally, it constructs the \mynvmath{MathML} intent for \mynvmath{MathCAT}.
%
% \begin{variable}[int]{
%   \l_@@_spmoney_luaset_symbol_str,
%   \l_@@_spmoney_luaset_position_str,
%   \l_@@_spmoney_luaset_sing_str,
%   \l_@@_spmoney_luaset_plur_str,
%   \l_@@_spmoney_luaset_conde_str,
%   \l_@@_spmoney_luaset_subsing_str,
%   \l_@@_spmoney_luaset_subplur_str,
%   \l_@@_spmoney_luaset_print_iso_str,
%   \l_@@_spmoney_luaset_currency_arg_str,
%   \l_@@_spmoney_luaset_status_str
% }
% These string variables store the currency data returned by the Lua module
% \myluamod{}. Upon a successful database lookup, the module populates these
% variables with the \emph{\enquote{currency symbol}}, its layout position
% (|pre| or |post|), the appropriate singular and plural names for both the
% main unit and subunits, and the execution status flags. \TeX{} then uses
% this information to render the visual layout and construct the correct
% \mynvmath{MathML} speech intents.
%    \begin{macrocode}
\str_new:N  \l_@@_spmoney_luaset_symbol_str
\str_new:N  \l_@@_spmoney_luaset_position_str
\str_new:N  \l_@@_spmoney_luaset_sing_str
\str_new:N  \l_@@_spmoney_luaset_plur_str
\str_new:N  \l_@@_spmoney_luaset_conde_str
\str_new:N  \l_@@_spmoney_luaset_subsing_str
\str_new:N  \l_@@_spmoney_luaset_subplur_str
\str_new:N  \l_@@_spmoney_luaset_print_iso_str
\str_new:N  \l_@@_spmoney_luaset_currency_arg_str
\str_new:N  \l_@@_spmoney_luaset_status_str
%    \end{macrocode}
% \end{variable}
%
% \begin{variable}[int]{
%   \l_@@_spmoney_sub_units_bool,
%   \l_@@_spmoney_intent_money_str,
%   \l_@@_spmoney_intent_sign_str,
%   \l_@@_spmoney_intent_sub_unit_str,
%   \l_@@_spmoney_res_main_voice_str,
%   \l_@@_spmoney_subnum_int,
%   \l_@@_spmoney_extracted_curr_str,
%   \l_@@_spmoney_extracted_num_tl
% }
% The boolean variable tracks whether subunit rendering
% is enabled. The string variables assemble the specific speech text
% (intents) for \mynvmath{NVDA/MathCAT}, including the sign, the main currency,
% and the subunits. Finally, the remaining variables temporarily store
% the numeric value and currency code extracted from the user's input.
%    \begin{macrocode}
\bool_new:N \l_@@_spmoney_sub_units_bool
\str_new:N  \l_@@_spmoney_intent_money_str
\str_new:N  \l_@@_spmoney_intent_sign_str
\str_new:N  \l_@@_spmoney_intent_sub_unit_str
\str_new:N  \l_@@_spmoney_res_main_voice_str
\int_new:N  \l_@@_spmoney_subnum_int
\str_new:N  \l_@@_spmoney_extracted_curr_str
\tl_new:N   \l_@@_spmoney_extracted_num_tl
%    \end{macrocode}
% \end{variable}
%
% \subsubsection{Definition of error messages}
%
% \begin{macro}[int]{currency-no-subunits, invalid-currency}
%   We define two messages to handle invalid currency configurations.
%   The first \enquote{error} warns the user if subunit rendering is requested
%   for a currency that does not support fractional divisions (such as
%   JPY or CLP), safely ignoring the flag. The second \enquote{error} is triggered
%   if the provided currency identifier is not registered in the Lua database.
% \iffalse
%% messages errors for \spmoney
% \fi
%    \begin{macrocode}
\msg_new:nnnn { spintent } { currency-no-subunits }
  {
    The~currency~'#1'~does~not~support~subunits.~
    The~option~'sub-units=true'~is~ignored.
  }
  {
    Monies~like~the~Japanese~Yen~(JPY)~or~Chilean~Peso~(CLP)~
    do~not~have~fractional~subunits~registered~in~Lua~module.~
  }
\msg_new:nnnn { spintent } { invalid-currency }
  {
    Invalid~currency~identifier~'#1'~in~\token_to_str:N \spmoney .
  }
  {
    The~provided~symbol,~alias,~or~ISO~code~does~not~exist~in~
    the~Lua~module.
  }
%    \end{macrocode}
% \end{macro}
%
% \subsubsection{Auxiliary functions for \cs[no-index]{spmoney}}
%
% \begin{macro}[int]{\@@_spmoney_set_currency:n, \@@_spmoney_normalize_key:n, \@@_spmoney_normalize_key:V}
%   The function \myhlc|\@@_spmoney_set_currency:n| validates the currency
%   identifier provided in `|#1|' against the database. It delegates the
%   lookup to the helper function \myhlc|\@@_spmoney_normalize_key:n|, which
%   interfaces with the Lua module \myluamod{}. If the module returns a
%   |notfound| status, an \enquote{error} is triggered. Otherwise, it stores the valid
%   currency identifier in \myhlc|\l_@@_spmoney_luaset_currency_arg_str| for
%   subsequent processing.
%    \begin{macrocode}
\cs_new_protected:Nn \@@_spmoney_set_currency:n
  {
    \@@_spmoney_normalize_key:n {#1}
    \str_if_eq:VnTF \l_@@_spmoney_luaset_status_str { notfound }
      {
        \msg_error:nnn { spintent } { invalid-currency } {#1}
      }
      { \str_set:Nn \l_@@_spmoney_luaset_currency_arg_str {#1} }
  }
\cs_new_protected:Nn \@@_spmoney_normalize_key:n
  {
    \@@_luafun_spmoney_normalize_key:n { #1 }
  }
\cs_generate_variant:Nn \@@_spmoney_normalize_key:n { V }
%    \end{macrocode}
% \end{macro}
%
% \subsubsection{Definition of keys for \cs[no-index]{spmoney}}
%
% \begin{macro}[int]{currency, sub-units, dolar-math, cifras-sep }
%   Now we add the keys \mykey[spmoney]{currency}, \mykey[spmoney]{sub-units},
%   \mykey[spmoney]{dolar-math} and \mykey[spmoney]{cifras-sep}.
%    \begin{macrocode}
\keys_define:nn { spintent / spmoney }
  {
    currency     .code:n      = { \@@_spmoney_set_currency:n {#1} },
    currency     .initial:n   = { pesos },
    currency     .value_required:n = true,
    sub-units    .bool_set:N  = \l_@@_spmoney_sub_units_bool,
    sub-units    .value_required:n = true,
    dolar-math   .bool_set:N  = \l_@@_spmoney_dolar_math_bool,
    dolar-math   .initial:n   = { true },
    dolar-math   .value_required:n = true,
    cifras-sep   .meta:nn     = { spintent/spnum } { cifras-sep = {#1} },
    cifras-sep   .value_required:n = true,
    unknown      .code:n =
      {
        \str_set:Nn \l_@@_command_name_str { spmoney }
        \@@_unknown_keys_cmd:n {#1}
      },
  }
%    \end{macrocode}
% \end{macro}
%
% \subsubsection{Internal functions for \cs[no-index]{spmoney}}
%
% \begin{macro}[int]{\@@_spmoney_validate_money:}
%   The function \myhlc|\@@_spmoney_validate_money:| verifies if the parsed
%   input contains an embedded currency identifier by checking
%   \myhlc|\l_@@_luaset_has_units_str|. If a currency is detected, it extracts
%   and normalizes the key from \myhlc|\l_@@_luaset_arg_above_tl|. It throws
%   an \enquote{error} if the currency is not found in the database. If valid, it
%   updates \myhlc|\l_@@_spmoney_luaset_currency_arg_str| and retrieves the
%   full currency metadata from the \myluamod{} via \myhlc|\@@_luafun_spmoney_lookup_data:V|.
%    \begin{macrocode}
\cs_new_protected:Nn \@@_spmoney_validate_money:
  {
    \str_if_eq:VnT \l_@@_luaset_has_units_str { true }
      {
        \@@_spmoney_normalize_key:V \l_@@_luaset_arg_above_tl
        \str_if_eq:VnTF \l_@@_spmoney_luaset_status_str { notfound }
          { \msg_error:nne { spintent } { invalid-currency } { \l_@@_luaset_arg_above_tl } }
          {
            \str_set:NV
              \l_@@_spmoney_luaset_currency_arg_str
              \l_@@_luaset_arg_above_tl
          }
      }
    \str_set:Ne \l_@@_spmoney_luaset_status_str { found }
    \@@_luafun_spmoney_lookup_data:V \l_@@_spmoney_luaset_currency_arg_str
  }
%    \end{macrocode}
% \end{macro}
%
% \begin{macro}[int]{\@@_spmoney_print_symbol:}
%   The function \myhlc|\@@_spmoney_print_symbol:| renders the \emph{\enquote{currency symbol}}.
%   If \myhlc|\l_@@_spmoney_luaset_print_iso_str| is |true|, it outputs the ISO
%   code using |\mathrm|. Otherwise, it prints the \emph{\enquote{standard symbol}}.
%   For dollar signs `\$', it checks \myhlc|\l_@@_spmoney_dolar_math_bool| to
%   determine whether to render it directly as a math character `\$' or
%   wrap it inside |\text|.
%^^A   All other \emph{\enquote{currency symbols}}
%^^A   are safely wrapped in |\text| to ensure proper font rendering.
%    \begin{macrocode}
\cs_new_protected:Nn \@@_spmoney_print_symbol:
  {
    \str_if_eq:VnTF \l_@@_spmoney_luaset_print_iso_str { true }
      { \mathrm { \l_@@_spmoney_luaset_currency_arg_str } }
      {
        \str_if_eq:VnTF \l_@@_spmoney_luaset_symbol_str { $ }
          {
            \bool_if:NTF \l_@@_spmoney_dolar_math_bool
              { \$ } { \text { \$ } }
          }
          { \l_@@_spmoney_luaset_symbol_str }
      }
  }
%    \end{macrocode}
% \end{macro}
%
% \begin{macro}[int]{\@@_spmoney_render_layout:}
%   The function \myhlc|\@@_spmoney_render_layout:| determines the layout
%   order of the \emph{\enquote{currency symbol}} relative to the numeric
%   value. It evaluates \myhlc|\l_@@_spmoney_luaset_position_str|; if set to |pre|,
%   it renders the \emph{\enquote{currency symbol}} \emph{before} the number
%   node. Otherwise, it places the number first, followed by the
%   \emph{\enquote{currency symbol}}.
%    \begin{macrocode}
\cs_new_protected:Nn \@@_spmoney_render_layout:
  {
    \str_if_eq:VnTF \l_@@_spmoney_luaset_position_str { pre }
      {
        \@@_spmoney_print_symbol:
        \MathMLarg { n } { \@@_spmoney_render_number_node: }
      }
      {
        \MathMLarg { n } { \@@_spmoney_render_number_node: }
        \@@_spmoney_print_symbol:
      }
  }
%    \end{macrocode}
% \end{macro}
%
% \begin{macro}[int]{\@@_spmoney_render_number_node:, \@@_spmoney_render_only_integer_node:}
%   The function \myhlc|\@@_spmoney_render_number_node:| renders the full
%   numeric value, encompassing both the integer and decimal components.
%   Conversely, the function \myhlc|\@@_spmoney_render_only_integer_node:|
%   outputs exclusively the integer portion, omitting any decimal data.
%    \begin{macrocode}
\cs_new_protected:Nn \@@_spmoney_render_number_node:
  {
    \@@_spnum_render_int:
    \@@_spnum_print_part_dec:
  }
\cs_new_protected:Nn \@@_spmoney_render_only_integer_node:
  {
    \@@_spnum_render_int:
  }
%    \end{macrocode}
% \end{macro}
%
% \begin{macro}[int]{\@@_spmoney_print_int:, \@@_spmoney_intent_int:}
%   The function \myhlc|\@@_spmoney_print_int:| evaluates the presence of an
%   explicit sign in \myhlc|\l_@@_luaset_sign_tl|. If a sign is detected, it
%   assigns a prefix speech intent `|más|' or `|menos|' and embeds the sign
%   and subsequent layout inside a |\MathMLintent|. If no sign is present,
%   it routes execution based on whether subunits are enabled.
%
%   \smallskip
%
%   The function \myhlc|\@@_spmoney_intent_int:| determines the correct
%   grammatical inflection for the currency string (singular, plural, or the
%   prepositional `|de|' variant for exact millions). It evaluates the integer
%   component, decimal boundaries, and zero-fills to select the appropriate
%   string register, applying it as a postfix \mynvmath{MathML} intent over
%   the layout rendered by \myhlc|\@@_spmoney_render_layout:|.
%    \begin{macrocode}
\cs_new_protected:Nn \@@_spmoney_print_int:
  {
    \tl_if_empty:NTF \l_@@_luaset_sign_tl
      {
        \bool_if:NTF \l_@@_spmoney_sub_units_bool
          { \@@_spmoney_set_and_print_sub_units: }
          { \@@_spmoney_intent_int: }
      }
      {
        \str_if_eq:VnTF \l_@@_luaset_sign_tl { + }
          {
            \str_set:Nn \l_@@_spmoney_intent_sign_str { _más:prefix ($n) }
          }
          {
            \str_set:Nn \l_@@_spmoney_intent_sign_str { _menos:prefix ($n) }
          }
        \@@_invisible_sep:
        \MathMLintent { \l_@@_spmoney_intent_sign_str }
          {
            \l_@@_luaset_sign_tl
            \MathMLarg { n }
              {
                \bool_if:NTF \l_@@_spmoney_sub_units_bool
                  { \@@_spmoney_set_and_print_sub_units: }
                  { \@@_spmoney_intent_int: }
              }
          }
      }
  }
\cs_new_protected:Nn \@@_spmoney_intent_int:
  {
    \tl_if_empty:NTF \l_@@_luaset_dec_sep_tl
      {
        \str_if_eq:VnTF \l_@@_luaset_only_part_int_str { 1 }
          {
            \str_set:NV \l_@@_spmoney_intent_money_str \l_@@_spmoney_luaset_sing_str
          }
          {
            \str_set:NV \l_@@_spmoney_intent_money_str \l_@@_spmoney_luaset_plur_str
          }
        \str_if_eq:VnT \l_@@_luaset_millons_str { true }
          {
            \str_set:NV \l_@@_spmoney_intent_money_str \l_@@_spmoney_luaset_conde_str
          }
      }
      {
        \str_set:NV \l_@@_spmoney_intent_money_str \l_@@_spmoney_luaset_plur_str
        \str_if_eq:VnT \l_@@_luaset_millons_str { true }
          {
            \str_if_eq:VnTF \l_@@_luaset_dec_is_all_zeros_str { true }
              {
                \str_set:NV \l_@@_spmoney_intent_money_str \l_@@_spmoney_luaset_conde_str
              }
              {
                \str_set:NV \l_@@_spmoney_intent_money_str \l_@@_spmoney_luaset_plur_str
              }
          }
      }
    \@@_invisible_sep: % need for MathCAT
    \MathMLintent { \l_@@_spmoney_intent_money_str :postfix($n) }
      {
        \@@_spmoney_render_layout:
      }
  }
%    \end{macrocode}
% \end{macro}
%
% \begin{macro}[int]{\@@_spmoney_check_sub_units:, \@@_spmoney_set_and_print_sub_units:}
%   The function \myhlc|\@@_spmoney_check_sub_units:| verifies whether the
%   selected currency supports \emph{\enquote{subunits}}. If the database returns empty
%   subunit data, it triggers a warning and automatically disables the
%   \emph{\enquote{subunit}} rendering flag.
%
%   \smallskip
%
%   The function \myhlc|\@@_spmoney_set_and_print_sub_units:| processes
%   and formats the speech intents when \emph{\enquote{subunits}} are active. It normalizes
%   the decimal component (for instance, padding a single-digit decimal
%   like `.5` to `50`), evaluates the correct singular or plural inflections
%   for both the main \emph{\enquote{currency}} and the \emph{\enquote{subunits}},
%   and delegates the final \mynvmath{MathML} tree assembly to \myhlc|\@@_spmoney_build_tree:|.
%    \begin{macrocode}
\cs_new_protected:Nn \@@_spmoney_check_sub_units:
  {
    \str_if_empty:NT \l_@@_spmoney_luaset_subsing_str
      {
        \msg_warning:nne { spintent } { currency-no-subunits }
          { \l_@@_spmoney_luaset_currency_arg_str }
        \bool_set_false:N \l_@@_spmoney_sub_units_bool
      }
  }
\cs_new_protected:Nn \@@_spmoney_set_and_print_sub_units:
  {
    \int_compare:nTF { \str_count:N \l_@@_luaset_only_part_dec_str == 1 }
      {
        \int_set:Nn \l_@@_spmoney_subnum_int
          { \l_@@_luaset_only_part_dec_str * 10 }
        \str_set:Ne \l_@@_luaset_only_part_dec_str
          { \l_@@_luaset_only_part_dec_str 0 }
      }
      {
        \int_set:Nn \l_@@_spmoney_subnum_int
          { \l_@@_luaset_only_part_dec_str }
      }
    \int_compare:nTF { \l_@@_spmoney_subnum_int == 1 }
      {
        \str_set:Ne \l_@@_spmoney_intent_sub_unit_str
          { \l_@@_spmoney_luaset_subsing_str :postfix($sub) }
      }
      {
        \str_set:Ne \l_@@_spmoney_intent_sub_unit_str
          { \l_@@_spmoney_luaset_subplur_str :postfix($sub) }
      }
    \str_if_eq:VnTF \l_@@_luaset_only_part_int_str { 1 }
      {
        \str_set:NV \l_@@_spmoney_res_main_voice_str
          \l_@@_spmoney_luaset_sing_str
      }
      {
        \str_set:NV \l_@@_spmoney_res_main_voice_str
          \l_@@_spmoney_luaset_plur_str
      }
    \str_if_eq:VnT \l_@@_luaset_millons_str { true }
      {
        \str_set:NV \l_@@_spmoney_res_main_voice_str
          \l_@@_spmoney_luaset_conde_str
      }
    \@@_spmoney_build_tree:
  }
%    \end{macrocode}
% \end{macro}
%
% \begin{macro}[int]{\@@_spmoney_build_tree:, \@@_spmoney_print_integer_continue:,
%                    \@@_spmoney_print_sep_con:, \@@_spmoney_print_dec_sub_unit:}
%   The function \myhlc|\@@_spmoney_build_tree:| structures the
%   final \mynvmath{MathML} tree assembly for fractional \emph{\enquote{currencies}}. It arranges the
%   integer component, the decimal separator, and the \emph{\enquote{subunit}} component
%   based on the layout position (|pre| or |post|) of the \emph{\enquote{currency symbol}}.
%
%   \smallskip
%
%   The function \myhlc|\@@_spmoney_print_integer_continue:| isolates the rendering
%   of the integer node with its associated main \emph{\enquote{currency}} intent.
%   Furthermore, \myhlc|\@@_spmoney_print_sep_con:| applies the specific
%   speech intent `|con|' to the decimal separator, while
%   \myhlc|\@@_spmoney_print_dec_sub_unit:| binds the computed \emph{\enquote{subunit}}
%   intent to the decimal digits, ensuring accurate accessibility output.
%    \begin{macrocode}
\cs_new_protected:Nn \@@_spmoney_build_tree:
  {
    \str_if_eq:VnTF \l_@@_spmoney_luaset_position_str { pre }
      {
        \MathMLintent { \l_@@_spmoney_res_main_voice_str :postfix($n) }
          {
            \@@_spmoney_print_symbol:
            \MathMLarg { n } { \@@_spmoney_render_only_integer_node: }
          }
        \@@_spmoney_print_sep_con:
        \@@_spmoney_print_dec_sub_unit:
      }
      {
        \MathMLintent { \l_@@_spmoney_res_main_voice_str :postfix($n) }
          {
            \@@_invisible_sep:
            \MathMLarg { n } { \@@_spmoney_render_only_integer_node: }
          }
        \@@_spmoney_print_sep_con:
        \@@_spmoney_print_dec_sub_unit:
        \MathMLintent {:silent } { \@@_spmoney_print_symbol: }
      }
  }
\cs_new_protected:Nn \@@_spmoney_print_integer_continue:
  {
    \MathMLintent { \l_@@_spmoney_res_main_voice_str :postfix($n) }
      {
        \MathMLarg { n } { \@@_spmoney_render_only_integer_node: }
      }
  }
\cs_new_protected:Nn \@@_spmoney_print_sep_con:
  {
    \MathMLintent { con }
      {
        \mathord { \l_@@_luaset_dec_sep_tl }
      }
  }
\cs_new_protected:Nn \@@_spmoney_print_dec_sub_unit:
  {
    \MathMLintent { \l_@@_spmoney_intent_sub_unit_str }
      {
        \@@_invisible_sep:
        \MathMLarg { sub }
          {
            \@@_spnum_render_dec:
          }
      }
  }
%    \end{macrocode}
% \end{macro}
%
% \begin{macro}[int]{\@@_spmoney_parse_and_route:n}
%   The internal function \myhlc|\@@_spmoney_parse_and_route:n| processes
%   the \mymarg{argument} passed in `|#1|` using Lua module to extract
%   the numeric value and any \emph{\enquote{currency}} identifier. If a
%   \emph{\enquote{currency}} is found (\myhlc|\l_@@_spmoney_extracted_curr_str|),
%   it sets the corresponding variables. It then calls \myhlc|\@@_spmoney_validate_money:|
%   to verify the \emph{\enquote{currency}}. If \emph{\enquote{subunit}}
%   rendering is enabled (\myhlc|\l_@@_spmoney_sub_units_bool|),
%   it calls \myhlc|\@@_spmoney_check_sub_units:|. Finally, it calls
%   \myhlc|\@@_spmoney_print_int:| to format the output.
%    \begin{macrocode}
\cs_new_protected:Nn \@@_spmoney_parse_and_route:n
  {
    \@@_luafun_spmoney_prepare_input:n { #1 }
    \@@_luafun_clean_split_and_set:V \l_@@_spmoney_extracted_num_tl
    \str_if_empty:NF \l_@@_spmoney_extracted_curr_str
      {
        \str_set:Ne \l_@@_luaset_has_units_str { true }
        \str_set:NV \l_@@_luaset_arg_above_tl \l_@@_spmoney_extracted_curr_str
      }
    \@@_spmoney_validate_money:
    \bool_if:NT \l_@@_spmoney_sub_units_bool
      { \@@_spmoney_check_sub_units: }
    \@@_spmoney_print_int:
  }
%    \end{macrocode}
% \end{macro}
%
% \begin{macro}{\spmoney}
%   The public document command |\spmoney| formats monetary values
%   with accessible speech intents. It accepts an optional argument `|#1|'
%   for local key-value configurations and a mandatory argument `|#2|'
%   containing the numeric value and optional embedded currency identifier.
%   It first evaluates the typesetting context, triggering a fatal error if
%   executed outside of math mode. Upon validation, it establishes a local
%   \TeX{} group to safely isolate the key assignments, delegates the raw
%   input string to \myhlc|\@@_spmoney_parse_and_route:n| for parsing and
%   rendering, and subsequently closes the group to prevent scope leakage.
%    \begin{macrocode}
\NewDocumentCommand \spmoney { O{} m }
  {
    \mode_if_math:TF
      {
        \group_begin:
          \keys_set:nn { spintent / spmoney } { #1 }
          \@@_spmoney_parse_and_route:n { #2 }
        \group_end:
      }
      { \msg_error:nnn { spintent } { need-math-mode } { \spmoney } }
  }
%    \end{macrocode}
% \end{macro}
%
% \subsection{The command \cs{spshort}}\label{cmd:spshort}
%
% The user command |\spshort| is the \emph{fourth} and \emph{final} of
% \emph{\enquote{the big four}}. It manages the semantic injection and
% typographical formatting of \emph{text mode} abbreviations, ordinals, and
% acronyms. Unlike \emph{math mode} commands, which rely on \mynvmath{MathML}
% attributes, this command uses the \texttt{PDF} \emph{tagging tools}
% provided by the \mypkg{tagpdf}\cite{tagpdf} package. Specifically, it uses
% the \mykey{E} (expanded) key to ensure that \mynvmath{NVDA} interpret the full
% form of the abbreviation, ordinal, or acronym in \emph{text mode}.
%
% \begin{variable}[int]{
%   \l_@@_spshort_luaset_status_str,
%   \l_@@_spshort_luaset_layout_str,
%   \l_@@_spshort_luaset_expanded_str,
%   \l_@@_spshort_luaset_output_str,
%   \l_@@_spshort_luaset_base_str,
%   \l_@@_spshort_luaset_suffix_str
% }
% These internal string variables serve as the primary data interface with
% the Lua module \myluamod{}. During processing, the Lua module
% \emph{\enquote{stores}} within these registers the execution status, the
% targeted typographical layout parameters, the phonetic replacement text
% for the \mykey{E} key intended for \mynvmath{NVDA}, the final
% sanitized output string, and the isolated structural components for
% ordinals (base and suffix).
%    \begin{macrocode}
\str_new:N  \l_@@_spshort_luaset_status_str
\str_new:N  \l_@@_spshort_luaset_layout_str
\str_new:N  \l_@@_spshort_luaset_expanded_str
\str_new:N  \l_@@_spshort_luaset_output_str
\str_new:N  \l_@@_spshort_luaset_base_str
\str_new:N  \l_@@_spshort_luaset_suffix_str
%    \end{macrocode}
% \end{variable}
%
% \subsubsection{Definition of error messages for \cs[no-index]{spshort}}
%
% \begin{macro}[int]{unknown-abbreviation}
%   The internal \enquote{error} message \texttt{unknown-abbreviation} is
%   triggered when the command |\spshort| encounters an invalid lookup
%   state from the Lua module. It acts as a strict fallback exception when
%   the provided \mymarg{argument} cannot be mapped to any registered
%   database entry in \myluamod{} nor successfully resolved through the
%   grammatical ordinal expansion rules.
% \iffalse
% % messages errors for \spshort
% \fi
%    \begin{macrocode}
\msg_new:nnnn { spintent } { unknown-abbreviation }
  { The~abbreviation~or~ordinal~'#1'~is~invalid~or~not~recognized. }
  { Verify~the~spelling~or~check~the~grammar~rules~for~ordinals. }
%    \end{macrocode}
% \end{macro}
%
% \subsubsection{Accessibility interface}
%
% \begin{macro}[int]{\@@_spshort_tag_span:nn, \@@_spshort_tag_span:en}
%   The internal function \myhlc|\@@_spshort_tag_span:nn| processes the
%   \mymarg{argument} passed in `|#2|' using the tools provided by \mypkg{tagpdf}.
%   By utilizing the \mykey{tag} and \mykey{E} keys, it creates a
%   native \texttt{PDF} |/E (#1)| structure and expands the \mymarg{argument}
%   passed in `|#1|' for proper reading with \mynvmath{NVDA}.
%    \begin{macrocode}
\cs_new_protected:Nn \@@_spshort_tag_span:nn
  {
    \tag_mc_end_push:
      \tag_struct_begin:n { tag=Span, E={ #1 } }
        \tag_mc_begin:n { tag=Span }
          \tag_suspend:n { \spshort }
            #2
          \tag_resume:n { \spshort }
        \tag_mc_end:
      \tag_struct_end:
    \tag_mc_begin_pop:n {}
  }
\cs_generate_variant:Nn \@@_spshort_tag_span:nn { en }
%    \end{macrocode}
% \end{macro}
%
% \subsubsection{Definition of keys for \cs[no-index]{spshort}}
%
% \begin{macro}[int]{small-caps, read-arg}
%   Now we add the keys \mykey[spshort]{small-caps} and \mykey[spshort]{read-arg}.
%    \begin{macrocode}
\keys_define:nn { spintent / spshort }
  {
    small-caps .bool_set:N = \l_@@_spshort_smallcaps_bool,
    small-caps .initial:n  = false,
    small-caps .value_required:n = true,
    read-arg   .str_set:N  = \l_@@_spshort_read_arg_str,
    read-arg   .initial:n  = {},
    read-arg   .value_required:n = true,
    unknown    .code:n =
      {
        \str_set:Nn \l_@@_command_name_str { spshort }
        \@@_unknown_keys_cmd:n {#1}
      },
  }
%    \end{macrocode}
% \end{macro}
%
% \subsubsection{Internal functions for \cs[no-index]{spshort}}
%
% \begin{macro}[int]{\@@_spshort_print:nn}
%   The internal function \myhlc|\@@_spshort_print:nn| acts as the primary
%   formatter for text abbreviations. It clears the string variable
%   \myhlc|\l_@@_spshort_read_arg_str| and evaluates the
%   optional \mymeta{keys} passed in `|#1|'. It then checks if a manual
%   reading override was provided via those keys. If that string variable
%   remains empty, it calls \myhlc|\@@_spshort_process_and_render:n| to
%   process the \mymarg{argument} passed in `|#2|'. Otherwise, if an
%   override is detected, it directly delegates the execution to
%   \myhlc|\@@_spshort_render_bypass:n|.
%    \begin{macrocode}
\cs_new_protected:Nn \@@_spshort_print:nn
  {
    \str_clear:N \l_@@_spshort_read_arg_str
    \keys_set:nn { spintent / spshort } {#1}
    \str_if_empty:NTF \l_@@_spshort_read_arg_str
      {
        \@@_spshort_process_and_render:n {#2}
      }
      {
        \@@_spshort_render_bypass:n {#2}
      }
  }
%    \end{macrocode}
% \end{macro}
%
% \begin{macro}[int]{\@@_spshort_process_and_render:n}
%   The internal function \myhlc|\@@_spshort_process_and_render:n| interfaces
%   with the Lua module \myluamod{} function \myhlc|\@@_luafun_spshort_process:n|
%   to resolve the abbreviation or ordinal \mymarg{argument} passed in `|#1|'.
%
%   \smallskip
%
%   It subsequently evaluates the execution state stored within the string
%   variable \myhlc|\l_@@_spshort_luaset_status_str| and dynamically routes
%   execution to the corresponding rendering subroutine |success|,
%   |fallback| or |error|. Any unrecognized status automatically defaults
%   to the \enquote{error} handler.
%    \begin{macrocode}
\cs_new_protected:Nn \@@_spshort_process_and_render:n
  {
    \@@_luafun_spshort_process:n {#1}
    \str_case:VnF \l_@@_spshort_luaset_status_str
      {
        { success }  { \@@_spshort_render_success: }
        { fallback } { \@@_spshort_render_fallback: }
        { error }    { \@@_spshort_render_error:n {#1} }
      }
      { \@@_spshort_render_error:n {#1} }
  }
%    \end{macrocode}
% \end{macro}
%
% \begin{macro}[int]{\@@_spshort_render_success:, \@@_spshort_render_fallback:,
%                    \@@_spshort_render_bypass:n, \@@_spshort_render_error:n}
%   The function \myhlc|\@@_spshort_render_success:| injects the key
%   \mykey{E} and executes the targeted typographical layout.
%   The \myhlc|\@@_spshort_render_fallback:| function handles default cases,
%   applying optional formatting such as \textsc{small caps} to unrecognized
%   or partially resolved terms.
%
%   \smallskip
%
%   Conversely, the function \myhlc|\@@_spshort_render_bypass:n| directly
%   maps the user-defined \emph{\enquote{speech}} to the \mymarg{argument},
%   bypassing the automated lookup Lua module. Finally, the function
%   \myhlc|\@@_spshort_render_error:n| triggers a structural exception while
%   yielding the unformatted \mymarg{argument} to prevent fatal compilation.
%    \begin{macrocode}
\cs_new_protected:Nn \@@_spshort_render_success:
  {
    \@@_spshort_tag_span:en
      { \l_@@_spshort_luaset_expanded_str }
      { \@@_spshort_execute_layout: }
  }
\cs_new_protected:Nn \@@_spshort_render_fallback:
  {
    \@@_spshort_tag_span:en
      { \l_@@_spshort_luaset_expanded_str }
      {
        \bool_if:NTF \l_@@_spshort_smallcaps_bool
          { \textsc{ \l_@@_spshort_luaset_output_str } }
          { \l_@@_spshort_luaset_output_str }
      }
  }
\cs_new_protected:Nn \@@_spshort_render_bypass:n
  {
    \@@_spshort_tag_span:en { \l_@@_spshort_read_arg_str } { #1 }
  }
\cs_new_protected:Nn \@@_spshort_render_error:n
  {
    \msg_error:nnn { spintent } { unknown-abbreviation } {#1}
  }
%    \end{macrocode}
% \end{macro}
%
% \begin{macro}[int]{\@@_spshort_execute_layout:}
%   The internal function \myhlc|\@@_spshort_execute_layout:| executes the
%   typographical layout. It evaluates the descriptor \emph{\enquote{stored}}
%   within the string variable \myhlc|\l_@@_spshort_luaset_layout_str|
%   (such as |linear_regular|, |superscript|, or |small_caps_pure|) and
%   maps it to the corresponding \LaTeX{} formatting. Furthermore, it
%   manages complex nested configurations, dynamically applying
%   combinations such as superscripts integrated with
%   conditional \textsc{small caps} based on the active boolean variables.
%    \begin{macrocode}
\cs_new_protected:Nn \@@_spshort_execute_layout:
  {
    \str_case:VnF \l_@@_spshort_luaset_layout_str
      {
        { linear_regular }
          {
            \bool_if:NTF \l_@@_spshort_smallcaps_bool
              { \textsc{ \l_@@_spshort_luaset_output_str } }
              { \l_@@_spshort_luaset_output_str }
          }
        { linear_caps }
          { \l_@@_spshort_luaset_output_str }
        { linear_mixed }
          { \l_@@_spshort_luaset_output_str }
        { superscript }
          {
            \l_@@_spshort_luaset_base_str
            \bool_if:NTF \l_@@_spshort_smallcaps_bool
              {
                \textsuperscript{
                  \textsc{
                    \str_lowercase:V \l_@@_spshort_luaset_suffix_str
                  }
                }
              }
              {
                \textsuperscript{ \l_@@_spshort_luaset_suffix_str }
              }
          }
        { small_caps_pure }
          {
            \bool_if:NTF \l_@@_spshort_smallcaps_bool
              {
                \textsc{
                  \str_lowercase:V \l_@@_spshort_luaset_output_str
                }
              }
              { \l_@@_spshort_luaset_output_str }
          }
        { acronym_long }
          { \l_@@_spshort_luaset_output_str }
      }
      { \l_@@_spshort_luaset_output_str }
  }
%    \end{macrocode}
% \end{macro}
%
% \begin{macro}{\spshort}
%   The public document command |\spshort| manages the accessible
%   formatting of \emph{text mode} abbreviations, acronyms, and ordinals.
%   It first triggers an error if invoked within a \emph{math mode},
%   upon validation, puts |\mode_leave_vertical:| and establishes a local
%   group to call the function \myhlc|\@@_spshort_print:nn| for processing
%   and printing.
%    \begin{macrocode}
\NewDocumentCommand \spshort { O{} m }
  {
    \mode_if_math:TF
      {
        \msg_error:nnn { spintent } { need-text-mode } { \spshort }
      }
      {
        \mode_leave_vertical:
        \group_begin:
          \@@_spshort_print:nn {#1} {#2}
        \group_end:
      }
  }
%    \end{macrocode}
% \end{macro}
%
% \subsection{The command \cs{spsiglo}}\label{cmd:spsiglo}
%
% The public document command \ics*{spsiglo} provides a semantic interface
% for typesetting century designations (e.g., siglo XXI). It accepts both
% Arabic and Roman numerals as \mymarg{argument}, interfacing with the Lua
% module \myluamod{} to validate, parse, and normalize chronological values.
%
% \smallskip
%
% This command work in dual-mode context awareness, operating uniformly across
% both \emph{text mode} and \emph{math mode}. When executed within \emph{math
% mode}, it constructs an accessible \mynvmath{MathML} \texttt{intent}
% expression tree. Conversely, when invoked in standard \emph{text mode},
% it encapsulates the rendered output within a native \texttt{PDF} tagging
% structure (such as |/Alt (#1)|) containing an \mykey{alt} key.
%
% \smallskip
%
% \begin{variable}[int]{
%   \l_@@_spsiglo_luaset_error_str,
%   \l_@@_spsiglo_luaset_arabic_str,
%   \l_@@_spsiglo_luaset_roman_min_str,
%   \l_@@_spsiglo_print_era_str
% }
%   These internal string variables handle data transfer and synchronization
%   with the Lua module \myluamod{}. They isolate the exception evaluation
%   states, the Arabic numeral required for speech synthesis, the lowercase
%   Roman numeral designated for small caps typesetting via |\textsc|, and
%   the localized chronological era modifier.
%    \begin{macrocode}
\str_new:N \l_@@_spsiglo_luaset_error_str
\str_new:N \l_@@_spsiglo_luaset_arabic_str
\str_new:N \l_@@_spsiglo_luaset_roman_min_str
\str_new:N \l_@@_spsiglo_print_era_str
%    \end{macrocode}
% \end{variable}
%
% \subsubsection{Definition of error messages for \cs[no-index]{spsiglo}}
%
% \begin{macro}[int]{invalid-century-format}
%   This internal error message is issued when the provided century
%   identifier fails both Arabic and Roman numeral validation. It is
%   triggered when the Lua module \myluamod{} cannot map the \mymarg{argument}
%   to an integer within the valid range of |1| to |4000|, preventing
%   incorrect speech synthesis.
%    \begin{macrocode}
\msg_new:nnnn { spintent } { invalid-century-format }
  {
    Invalid~century~format~in~'#1'.
  }
  {
    The~argument~'#2'~is~not~a~valid~Arabic~or~Roman~numeral.~
    Ensure~the~value~is~between~1~and~4000.
  }
%    \end{macrocode}
% \end{macro}
%
% \subsubsection{Definition of keys for \cs[no-index]{spsiglo}}
%
% \begin{macro}[int]{print-era}
% Now we add the key \mykey[spsiglo]{print-era}. User configuration
% determining if an era designation is appended, available values are |none|,
% |ac| (antes de cristo) and |dc| (después de cristo).
%    \begin{macrocode}
\keys_define:nn { spintent / spsiglo }
  {
    print-era           .choices:nn =
      { none , ac , dc }
      {
        \int_case:nn { \l_keys_choice_int }
          {
            {1} { \str_set:Nn \l_@@_spsiglo_print_era_str { none } }
            {2} { \str_set:Nn \l_@@_spsiglo_print_era_str { ac } }
            {3} { \str_set:Nn \l_@@_spsiglo_print_era_str { dc } }
          }
      },
    print-era           .initial:n  = none,
    print-era           .value_required:n = true,
    print-era / unknown .code:n =
      \msg_error:nneee { spintent } { unknown-choice }
        { print-era } { none, ~ac, ~dc } { \exp_not:n {#1} },
    unknown             .code:n =
      {
        \str_set:Nn \l_@@_command_name_str { spsiglo }
        \@@_unknown_keys_cmd:n {#1}
      },
  }
%    \end{macrocode}
% \end{macro}
%
% \subsubsection{Internal functions for \cs[no-index]{spsiglo}}
%
% \begin{macro}[int]{\@@_spsiglo_render_simple:nn, \@@_spsiglo_render_simple:VV,
%                    \@@_spsiglo_render_with_era:nnn, \@@_spsiglo_render_with_era:VVV}
%   The functions \myhlc|\@@_spsiglo_render_simple:nn| and
%   \myhlc|\@@_spsiglo_render_with_era:nnn| generate \mynvmath{MathML} structures
%   when century designations are evaluated within \emph{math mode}. They
%   construct \mynvmath{MathML} \texttt{intent} structures, mapping the Arabic
%   numerals to their phonetic equivalents. Furthermore,
%   \myhlc|\@@_spsiglo_render_with_era:nnn| integrates the era modifiers
%   \mykey{ac} or \mykey{dc}. It adds invisible spacing and postfix
%   intents to ensure correct \mynvmath{NVDA/MathCAT} pronunciation.
%    \begin{macrocode}
\cs_new_protected:Nn \@@_spsiglo_render_simple:nn
  {
    \MathMLintent { #1:number } { \text { \textsc {#2} } }
  }
\cs_generate_variant:Nn \@@_spsiglo_render_simple:nn { VV }
\cs_new_protected:Nn \@@_spsiglo_render_with_era:nnn
  {
    \str_if_eq:nnTF { #1 } { ac }
      {
        \MathMLintent { antes-de-cristo:postfix($z) }
          {
            \@@_invisible_sep:
            \MathMLarg { z }
              {
                \MathMLintent { #2:number }
                  {
                    \text { \textsc {#3}~a.~C. }
                  }
              }
          }
      }
      {
        \MathMLintent { después-de-cristo:postfix($z) }
          {
            \@@_invisible_sep:
            \MathMLarg { z }
              {
                \MathMLintent { #2:number }
                  {
                    \text { \textsc {#3}~d.~C. }
                  }
              }
          }
      }
  }
\cs_generate_variant:Nn \@@_spsiglo_render_with_era:nnn { VVV }
%    \end{macrocode}
% \end{macro}
%
% \begin{macro}[int]{\@@_spsiglo_tag_span:nn, \@@_spsiglo_tag_span:en}
%   The function \myhlc|\@@_spsiglo_tag_span:nn| handles century formatting
%   within \emph{text mode}. It utilizes the \mypkg{tagpdf} package to
%   enclose the content inside a \texttt{PDF} |/Alt (#1)| structure. By passing
%   the \mymarg{argument} to the \mykey{alt} key, it embeds the
%   chronological value directly into the \texttt{PDF}, ensuring correct
%   \mynvmath{NVDA/MathCAT} pronunciation without relying on \mynvmath{MathML} intents.
%    \begin{macrocode}
\cs_new_protected:Nn \@@_spsiglo_tag_span:nn
  {
    \tag_mc_end_push:
      \tag_struct_begin:n { tag=Span, alt={ #1 } }
        \tag_mc_begin:n { tag=Span }
          \tag_suspend:n { \spsiglo }
            #2
          \tag_resume:n { \spsiglo }
        \tag_mc_end:
      \tag_struct_end:
    \tag_mc_begin_pop:n {}
  }
\cs_generate_variant:Nn \@@_spsiglo_tag_span:nn { en }
%    \end{macrocode}
% \end{macro}
%
% \begin{macro}[int]{\@@_spsiglo_print_math:n, \@@_spsiglo_print_text:n}
%   The functions \myhlc|\@@_spsiglo_print_math:n| and
%   \myhlc|\@@_spsiglo_print_text:n| format century designations based on the
%   active mode. Both first check the \enquote{error} status in the variable
%   \myhlc|\l_@@_spsiglo_luaset_error_str| returned by the Lua module
%   \myluamod{} and issue an \enquote{error} if validation fails.
%
%   \smallskip
%
%   If the \mymarg{argument} is valid, the function \myhlc|\@@_spsiglo_print_math:n|
%   operates in \emph{math mode}, applying \mynvmath{MathML} structures
%   depending on whether an era modifier is present. Conversely, the function
%   \myhlc|\@@_spsiglo_print_text:n| operates in \emph{text mode},
%   utilizing the \mykey{tag} and \mykey{alt} keys to expand the value
%   of the \mykey[spsiglo]{print-era} option into full text (e.g.,
%   \enquote{antes de cristo}) for \mynvmath{NVDA/MathCAT}.
%    \begin{macrocode}
\cs_new_protected:Nn \@@_spsiglo_print_math:n
  {
    \str_if_eq:VnTF \l_@@_spsiglo_luaset_error_str { true }
      {
        \msg_error:nnnn { spintent } { invalid-century-format }
          { \spsiglo } { #1 }
      }
      {
        \str_if_eq:VnTF \l_@@_spsiglo_print_era_str { none }
          {
            \@@_spsiglo_render_simple:VV
              \l_@@_spsiglo_luaset_arabic_str
              \l_@@_spsiglo_luaset_roman_min_str
          }
          {
            \@@_spsiglo_render_with_era:VVV
              \l_@@_spsiglo_print_era_str
              \l_@@_spsiglo_luaset_arabic_str
              \l_@@_spsiglo_luaset_roman_min_str
          }
      }
  }
\cs_new_protected:Nn \@@_spsiglo_print_text:n
  {
    \str_if_eq:VnTF \l_@@_spsiglo_luaset_error_str { true }
      {
        \msg_error:nnnn { spintent } { invalid-century-format }
          { \spsiglo } { #1 }
      }
      {
         \str_case:Vn \l_@@_spsiglo_print_era_str
          {
            { none }
            {
              \@@_spsiglo_tag_span:en
                { \l_@@_spsiglo_luaset_arabic_str }
                { \textsc { \l_@@_spsiglo_luaset_roman_min_str } }
            }
            { ac }
            {
              \@@_spsiglo_tag_span:en
                { \l_@@_spsiglo_luaset_arabic_str ~ antes ~ de ~ cristo }
                { \textsc { \l_@@_spsiglo_luaset_roman_min_str } ~ a.~C. }
            }
            { dc }
            {
              \@@_spsiglo_tag_span:en
                { \l_@@_spsiglo_luaset_arabic_str ~ después ~ de ~ cristo }
                { \textsc { \l_@@_spsiglo_luaset_roman_min_str } ~ d.~C. }
            }
          }
      }
  }
%    \end{macrocode}
% \end{macro}
%
% \begin{macro}{\spsiglo}
%   The public command \ics*{spsiglo} is the main interface for century
%   typesetting. It opens a local group and evaluates any optional
%   \mymeta{keys} passed in `|#1|', then executes the function
%   \myhlc|\@@_luafun_spsiglo_parse:n| defined in the Lua module \myluamod{}
%   passed in `|#2|' and formats the output for \emph{math mode} or \emph{text mode}.
%    \begin{macrocode}
\NewDocumentCommand \spsiglo { O{} m }
  {
    \group_begin:
      \keys_set:nn { spintent / spsiglo } { #1 }
      \@@_luafun_spsiglo_parse:n { #2 }
      \mode_if_math:TF
        { \@@_spsiglo_print_math:n { #2 } }
        {
          \mode_leave_vertical:
          \@@_spsiglo_print_text:n { #2 }
        }
    \group_end:
  }
%    \end{macrocode}
% \end{macro}
%
% \subsection{The command \cs{sptime}}\label{cmd:sptime}
%
% The command |\sptime| is the main public interface for typesetting and
% tagging time expressions. It parses time string formats (such as |HH:MM|
% or |HH:MM:SS|) using the Lua function \myhlc|\@@_luafun_sptime_parse:n|.
% This separates the hours, minutes, seconds, and fractions to construct
% a \mynvmath{MathML} |intent| tree for \mynvmath{NVDA/MathCAT}.
% \smallskip
%
% \begin{variable}[int]{
%   \l_@@_sptime_luaset_seconds_str,
%   \l_@@_sptime_luaset_has_seconds_str,
%   \l_@@_sptime_luaset_error_str,
%   \l_@@_sptime_luaset_is_fraction_str,
%   \l_@@_sptime_luaset_base_str
% }
% These internal string variables store the values returned from the Lua
% module \myluamod{}. They hold the extracted time units, error flags, and
% boolean indicators (such as the presence of seconds or decimal fractions)
% used during formatting.
%    \begin{macrocode}
\str_new:N \l_@@_sptime_luaset_seconds_str
\str_new:N \l_@@_sptime_luaset_has_seconds_str
\str_new:N \l_@@_sptime_luaset_error_str
\str_new:N \l_@@_sptime_luaset_is_fraction_str
\str_new:N \l_@@_sptime_luaset_base_str
%    \end{macrocode}
% \end{variable}
%
% \subsubsection{Definition of error messages}
%
% \begin{macro}[int]{invalid-time-format}
%   The error message |invalid-time-format| defines the text displayed when
%   an invalid time input is detected. It is triggered when an argument
%   fails the format checks or contains values outside normal ranges,
%   such as hours exceeding 23 or minutes and seconds exceeding 59.
%    \begin{macrocode}
\msg_new:nnnn { spintent } { invalid-time-format }
  {
    Invalid~time~format~in~#1. \\
    The~argument~'#2'~does~not~match~any~supported~time~format.
  }
  {
    Allowed~formats:~HH:MM,~HH:MM:SS,~or~HH:MM:SS.ss. \\
    Please~verify~that~hours~are~0-23~and~minutes/seconds~are~0-59.
  }
%    \end{macrocode}
% \end{macro}
%
% \subsubsection{Internal functions for \cs[no-index]{sptime}}
%
% \begin{macro}[int]{\@@_sptime_build_seconds:nn}
%   The internal function \myhlc|\@@_sptime_build_seconds:nn| builds the \mynvmath{MathML}
%   structure for the seconds component of a time expression. It wraps the
%   numeric value `|#2|' in speech intent tags, applying a connector prefix
%   (passed in `|#1|') and the "segundos" postfix so it is read correctly by
%   \mynvmath{NVDA/MathCAT}.
%    \begin{macrocode}
\cs_new_protected:Nn \@@_sptime_build_seconds:nn
  {
    \MathMLintent { #1:prefix($z) }
      {
        \text { : }
        \MathMLarg { z }
          {
            \MathMLintent { segundos:postfix($m) }
              {
                \@@_invisible_sep:
                \MathMLarg { m }
                  {
                    \text { #2 }
                  }
              }
          }
      }
  }
%    \end{macrocode}
% \end{macro}
%
% \begin{macro}[int]{\@@_sptime_build_tree:}
%   The internal function \myhlc|\@@_sptime_build_tree:| builds the complete
%   \mynvmath{MathML} structure for the time expression. It checks the string variables
%   \myhlc|\l_@@_sptime_luaset_has_seconds_str| and
%   \myhlc|\l_@@_sptime_luaset_is_fraction_str| to determine if the time includes
%   seconds or decimal fractions. Based on these values, it inserts the base
%   time string and either calls \myhlc|\@@_sptime_build_seconds:nn| with the
%   appropriate prefix ('con' or 'minutos-con'), or simply appends the 'minutos'
%   intent tag if no seconds are present.
%    \begin{macrocode}
\cs_new_protected:Nn \@@_sptime_build_tree:
  {
    \str_if_eq:VnTF \l_@@_sptime_luaset_has_seconds_str { true }
      {
        \MathMLintent { :time }
          {
            \text { \l_@@_sptime_luaset_base_str }
          }
        \str_if_eq:VnTF \l_@@_sptime_luaset_is_fraction_str { true }
          {
            \@@_sptime_build_seconds:nn { con }
              { \l_@@_sptime_luaset_seconds_str }
          }
          {
            \@@_sptime_build_seconds:nn { minutos-con }
              { \l_@@_sptime_luaset_seconds_str }
          }
      }
      {
        \MathMLintent { :time }
          {
            \text { \l_@@_sptime_luaset_base_str }
          }
        \str_if_eq:VnT \l_@@_sptime_luaset_is_fraction_str { false }
          {
            \MathMLintent { minutos }
              {
                \@@_invisible_sep:
              }
          }
      }
  }
%    \end{macrocode}
% \end{macro}
%
% \begin{macro}{\sptime}
%   The command |\sptime| is the public interface for formatting time
%   expressions. It first checks if it is being used in math mode, issuing
%   an error if not. Inside a local group, it parses the input using
%   \cs{@@_luafun_sptime_parse:n} and checks the error status in
%   \myhlc|\l_@@_sptime_luaset_error_str|. If an invalid format is detected,
%   it issues the |invalid-time-format| error; otherwise, it calls
%   \cs{@@_sptime_build_tree:} to assemble the \mynvmath{MathML} output.
%    \begin{macrocode}
\NewDocumentCommand \sptime { m }
  {
    \mode_if_math:TF
      {
        \group_begin:
          \@@_luafun_sptime_parse:n { #1 }
          \str_if_eq:VnTF \l_@@_sptime_luaset_error_str { true }
            {
               \msg_error:nnnn { spintent } { invalid-time-format }
                 { \sptime } { #1 }
            }
            {
               \@@_sptime_build_tree:
            }
        \group_end:
      }
      { \msg_error:nnn { spintent } { need-math-mode } { \sptime } }
  }
%    \end{macrocode}
% \end{macro}
%
% \subsection{The command \cs{spdate}}\label{cmd:spdate}
%
% The command |\spdate| is the main public command for typesetting and
% tagging date expressions. It parses date string formats (such as |DD/MM/YYYY|
% or |YYYY-MM-DD|) using the Lua function \myhlc|\@@_luafun_spdate_parse:n|.
% This verifies the input and constructs a \mynvmath{MathML} |intent| tree for
% \mynvmath{NVDA/MathCAT}.
%
% \smallskip
%
% \begin{variable}[int]{
%   \l_@@_spdate_luaset_output_str,
%   \l_@@_spdate_luaset_error_str
% }
% These internal string variables store the values returned from the Lua
% module \myluamod{}. They hold the formatted date string and the error
% status used during formatting.
%    \begin{macrocode}
\str_new:N \l_@@_spdate_luaset_output_str
\str_new:N \l_@@_spdate_luaset_error_str
%    \end{macrocode}
% \end{variable}
%
% \subsubsection{Definition of error messages}
%
% \begin{macro}[int]{invalid-date-format}
%   The error message |invalid-date-format| defines the text displayed when
%   an invalid date input is detected. It is triggered when an argument
%   fails the format checks or represents an invalid calendar date,
%   such as |31/02/2024| or using an unsupported pattern like |YYYY/MM/DD|.
%    \begin{macrocode}
\msg_new:nnnn { spintent } { invalid-date-format }
  {
    Invalid~date~or~unsupported~format~in~#1. \\
    The~argument~'#2'~does~not~match~a~supported~format.
  }
  {
    Allowed~formats:~DD\slash MM\slash YYYY,~DD-MM-YYYY,~or~YYYY-MM-DD. \\
    Ensure~the~day~and~month~are~valid~(e.g.,~avoid~31\slash 02\slash 2024).
  }
%    \end{macrocode}
% \end{macro}
%
% \subsubsection{Internal functions for \cs[no-index]{spdate}}
%
% \begin{macro}[int]{\@@_spdate_process:n}
%   The internal function \myhlc|\@@_spdate_process:n| parses the input using
%   \cs{@@_luafun_spdate_parse:n} and checks the error status in
%   \myhlc|\l_@@_spdate_luaset_error_str|. If an invalid format is detected,
%   it issues the |invalid-date-format| error. Otherwise, it builds the
%   \mynvmath{MathML} |intent| structure using the formatted date stored in
%   \myhlc|\l_@@_spdate_luaset_output_str|.
%    \begin{macrocode}
\cs_new_protected:Nn \@@_spdate_process:n
  {
    \@@_luafun_spdate_parse:n { #1 }
    \str_if_eq:VnTF \l_@@_spdate_luaset_error_str { true }
      {
        \msg_error:nnnn { spintent } { invalid-date-format }
          { \spdate } { #1 }
      }
      {
        \MathMLintent { :date }
          {
            \text { \l_@@_spdate_luaset_output_str }
          }
      }
  }
%    \end{macrocode}
% \end{macro}
%
% \begin{macro}{\spdate}
%   The command |\spdate| is the public interface for formatting date
%   expressions. It first checks if it is being used in math mode, issuing
%   an error if not. Inside a local group, it calls \cs{@@_spdate_process:n}
%   to process the input.
%    \begin{macrocode}
\NewDocumentCommand \spdate { m }
  {
    \mode_if_math:TF
      {
        \group_begin:
          \@@_spdate_process:n { #1 }
        \group_end:
      }
      { \msg_error:nnn { spintent } { need-math-mode } { \spdate } }
  }
%    \end{macrocode}
% \end{macro}
%
%
% \subsection{Semantic Operators}
%
% This section provides purely semantic operators. Unlike their formatting
% counterparts (which parse and format numbers), these commands do not
% process mandatory arguments. Instead, they solely inject a mathematical
% glyph wrapped in its corresponding regional \mynvmath{MathCAT} intent. This
% grants teachers absolute control when building complex algebraic
% expressions, mixed numbers, or fraction divisions.
%
% \subsection{The command \cs{spusep}}\label{cmd:spusep}
%
% For the use of parentheses, we will provide the command |\spusep| with
% only two \mymeta{keys} |size| and |silent|. Most of the time we can use
% |\left| and |\right| or |\Uleft| and |\Uright|, but both forms create a
% group, and this hinders the code for \emph{tagged} \texttt{PDF} when we
% want to interrupt something between them.
%
% \begin{macro}[int]{spusep-empty-arg, spusep-invalid-arg}
%   The first \enquote{error} message |spusep-empty-arg| is raised when |\spusep|
%   receives an empty or blank \mymarg{argument}, the second this is when
%   an input is invalid.
%    \begin{macrocode}
\msg_new:nnnn { spintent } { spusep-empty-arg }
  { Empty~argument~in~\token_to_str:N \spusep. }
  { \token_to_str:N \spusep~requires~a~non-empty~argument. }
\msg_new:nnn { spintent } { spusep-invalid-arg }
  {
    Invalid~argument~'#1'~for~\token_to_str:N \spusep.
  }
%    \end{macrocode}
% \end{macro}
%
% \begin{variable}[int]{\l_@@_spusep_arg_tl, \l_@@_spusep_left_tl, \l_@@_spusep_right_tl}
%   The variable \cs{l_@@_spusep_arg_tl} stores the \mymarg{argument},
%   the variable \cs{l_@@_spusep_left_tl} and \cs{l_@@_spusep_right_tl}
%   stores the parentheses.
%    \begin{macrocode}
\tl_new:N \l_@@_spusep_arg_tl
\tl_new:N \l_@@_spusep_left_tl
\tl_new:N \l_@@_spusep_right_tl
%    \end{macrocode}
% \end{variable}
%
% \begin{macro}[int]{size, silent, unknown}
%   We define the \mymeta{keys} \mykey[spusep]{size} and \mykey[spusep]{silent}.
%    \begin{macrocode}
\keys_define:nn { spintent / spusep }
  {
    size           .choices:nn =
      { none , big , Big, bigg, Bigg }
      {
        \int_case:nn { \l_keys_choice_int }
          {
            { 1 }
              {
                \tl_clear:N \l_@@_spusep_left_tl
                \tl_clear:N \l_@@_spusep_right_tl
              }
            { 2 }
              {
                \tl_set:Nn \l_@@_spusep_left_tl  { \bigl }
                \tl_set:Nn \l_@@_spusep_right_tl { \bigr }
              }
            { 3 }
              {
                \tl_set:Nn \l_@@_spusep_left_tl  { \Bigl }
                \tl_set:Nn \l_@@_spusep_right_tl { \Bigr }
              }
            { 4 }
              {
                \tl_set:Nn \l_@@_spusep_left_tl  { \biggl }
                \tl_set:Nn \l_@@_spusep_right_tl { \biggr }
              }
            { 5 }
              {
                \tl_set:Nn \l_@@_spusep_left_tl  { \Biggl }
                \tl_set:Nn \l_@@_spusep_right_tl { \Biggr }
              }
          }
      },
    size           .initial:n  = none,
    size           .value_required:n = true,
    size / unknown .code:n =
      \msg_error:nneee { spintent } { unknown-choice }
        { size } { none,~big,~Big,~bigg,~Bigg } { \exp_not:n {#1} },
    silent         .bool_set:N = \l_@@_spusep_silent_bool,
    silent         .initial:n = false,
    silent         .value_required:n = true,
    unknown        .code:n =
      {
        \str_set:Nn \l_@@_command_name_str { spusep }
        \@@_unknown_keys_cmd:n {#1}
      },
  }
%    \end{macrocode}
% \end{macro}
%
%
% \begin{macro}[int]{\@@_spusep_set_parens:nn,}
%   The function \myhlc|\@@_spusep_set_parens:nn| will place the argument
%   passed in `|#2|' to the right of the variables \cs{l_@@_spusep_left_tl}
%   and \cs{l_@@_spusep_right_tl} and then check the state of the variable
%   \cs{l_@@_spusep_silent_bool} handled by the key \mykey{silent} to
%   print.
%    \begin{macrocode}
\cs_new_protected:Nn \@@_spusep_set_parens:nn
  {
    \tl_put_right:cn { l_@@_spusep_ #1 _tl } { #2 }
    \bool_if:NTF \l_@@_spusep_silent_bool
      {
         \MathMLintent { :silent } { \tl_use:c { l_@@_spusep_ #1 _tl } }
      }
      {
        \tl_use:c { l_@@_spusep_ #1 _tl }
      }
  }
%    \end{macrocode}
% \end{macro}
%
% \begin{macro}[int]{\@@_spusep_convert:n,}
%   The function \myhlc|\@@_spusep_convert:n| captures and stores the
%   \mymarg{argument} passed in `|#1|` in the variable \cs{l_@@_spusep_arg_tl},
%   validates the cases by converting them to the macros provided by
%   \mypkg{unicode-math}, and passes them to the function \myhlc|\@@_spusep_set_parens:nn|
%   for printing. If the input is not supported or is empty, we return an \enquote{error}.
%    \begin{macrocode}
\cs_new_protected:Nn \@@_spusep_convert:n
  {
    \tl_set:Ne \l_@@_spusep_arg_tl { \tl_trim_spaces:n {#1} }
    \tl_if_empty:NT \l_@@_spusep_arg_tl
      {
        \msg_error:nnn { spintent } { spusep-invalid-arg } {#1}
      }
    \str_case:VnF \l_@@_spusep_arg_tl
      {
        { (  }
          {
            \@@_spusep_set_parens:nn { left } { \lparen }
          }
        { )  }
          {
            \@@_spusep_set_parens:nn { right } { \rparen }
          }
        { [  }
          {
            \@@_spusep_set_parens:nn { left } { \lbrack }
          }
        { ]  }
          {
            \@@_spusep_set_parens:nn { right } { \rbrack }
          }
        { \{ }
          {
            \@@_spusep_set_parens:nn { left } { \lbrace }
          }
        { \} }
          {
            \@@_spusep_set_parens:nn { right } { \rbrace }
          }
      }
      {
        \msg_error:nnn { spintent } { spusep-invalid-arg } {#1}
      }
  }
%    \end{macrocode}
% \end{macro}
%
% \begin{macro}{\spusep}
%   The command |\spusep| checks that it is running in math mode, opens a
%   local group, sets the keys passed in `|#1|' and processes the
%   \mymarg{argument} passed in `|#2|' by means of the function
%   \myhlc|\@@_spusep_convert:n|.
%    \begin{macrocode}
\NewDocumentCommand \spusep { O{} m }
  {
    \mode_if_math:TF
      {
        \group_begin:
          \keys_set:nn { spintent / spusep } {#1}
          \@@_spusep_convert:n {#2}
        \group_end:
      }
      { \msg_error:nnn { spintent } { need-math-mode } { \spusep } }
  }
%    \end{macrocode}
% \end{macro}
%
% \subsubsection{The command \cs{spread}}\label{cmd:spread}
%
% The command |\spread| provides a semantic read for \mynvmath{MathCAT}
% of mathematical symbols for situations in which it is necessary to
% specify gender or plurality in Spanish.
%
% \begin{macro}[int]{spread-empty-args}
%   The error message |spread-empty-args| is raised when |\spread|
%   receives an empty or blank \mymarg{argument}.
%    \begin{macrocode}
\msg_new:nnnn { spintent } { spread-empty-args }
  { Empty~argument~in~\token_to_str:N \spread. }
  { \token_to_str:N \spread~requires~two~non-empty~arguments. }
%    \end{macrocode}
% \end{macro}
%
% \begin{variable}[int]{\l_@@_spread_intent_tl}
%   The variable \myhlc|\l_@@_spread_intent_tl| stores the sanitized semantic
%   reading \mymarg{argument} for symbol to be injected into the |\MathMLintent|.
%    \begin{macrocode}
\tl_new:N \l_@@_spread_intent_tl
%    \end{macrocode}
% \end{variable}
%
% \begin{macro}[int]{\@@_spread_exec:nn}
%   The function \myhlc|\@@_spread_exec:nn| implements the core logic for |\spread|.
%   It verifies that neither \mymarg{argument} is empty or blank. If valid,
%   it evaluates the \emph{reading} \mymarg{argument} `|#1|' using
%   \myhlc|\@@_sanitize_affix:Nn| and embeds it via |\MathMLintent| into the
%   \emph{math symbol} `|#2|'.
%    \begin{macrocode}
\cs_new_protected:Nn \@@_spread_exec:nn
  {
    \bool_lazy_or:nnTF { \tl_if_blank_p:n { #1 } } { \tl_if_blank_p:n { #2 } }
      { \msg_error:nn { spintent } { spread-empty-args } }
      {
        \@@_sanitize_affix:Nn \l_@@_spread_intent_tl { #1 }
        \MathMLintent { \l_@@_spread_intent_tl } { #2 }
      }
  }
%    \end{macrocode}
% \end{macro}
%
% \begin{macro}{\spread}
%   The command |\spread| assigns a custom semantic \emph{reading} `|#1|' to
%   any mathematical \emph{symbol} `|#2|'. It enforces math mode and passes execution
%   to \myhlc|\@@_spread_exec:nn|.
%    \begin{macrocode}
\NewDocumentCommand \spread { m m }
  {
    \mode_if_math:TF
      {
        \group_begin:
          \@@_spread_exec:nn { #1 } { #2 }
        \group_end:
      }
      { \msg_error:nnn { spintent } { need-math-mode } { \spread } }
  }
%    \end{macrocode}
% \end{macro}
%
% \subsubsection{The command \cs{spdsym}}
%
% \begin{variable}[int]{\l_@@_spdsym_read_sym_tl,
%                       \l_@@_spdsym_symbol_tl,
%                       \l_@@_spdsym_prefix_tl,
%                       \l_@@_spdsym_suffix_tl}
%   Variables to store the regional reading string, the visual operator,
%   and the optional semantic affixes specifically for |\spdsym|.
%    \begin{macrocode}
\tl_new:N \l_@@_spdsym_read_sym_tl
\tl_new:N \l_@@_spdsym_symbol_tl
\tl_new:N \l_@@_spdsym_prefix_tl
\tl_new:N \l_@@_spdsym_suffix_tl
%    \end{macrocode}
% \end{variable}
%
% \begin{macro}[int]{prefix, suffix, set-sym, read-sym, unknown}
%   We define the \mymeta{keys} \mykey[spdsym]{prefix}, \mykey[spdsym]{suffix},
%   \mykey[spdsym]{set-sym} and \mykey[spdsym]{read-sym} for |\spdsym|.
%    \begin{macrocode}
\keys_define:nn { spintent / spdsym }
  {
    prefix             .code:n =
      { \@@_sanitize_affix:Nn \l_@@_spdsym_prefix_tl {#1} },
    prefix             .initial:n = ,
    prefix             .value_required:n = true,
    suffix             .code:n =
      { \@@_sanitize_affix:Nn \l_@@_spdsym_suffix_tl {#1} },
    suffix             .initial:n = ,
    suffix             .value_required:n = true,
    set-sym            .choices:nn =
      { old-style , two-dots , slash }
      {
        \int_case:nn { \l_keys_choice_int }
          {
            {1} { \tl_set:Nn \l_@@_spdsym_symbol_tl { \div } }
            {2} { \tl_set:Nn \l_@@_spdsym_symbol_tl {   :  } }
            {3} { \tl_set:Nn \l_@@_spdsym_symbol_tl {   /  } }
          }
      },
    set-sym            .initial:n  = two-dots,
    set-sym            .value_required:n = true,
    set-sym / unknown  .code:n =
      \msg_error:nneee { spintent } { unknown-choice }
        { set-sym } { old-style, ~two-dots, ~slash } { \exp_not:n {#1} },
    read-sym           .choices:nn =
      { dividido , dividido-por , dividido-entre }
      {
        \int_case:nn { \l_keys_choice_int }
          {
            {1} { \tl_set:Nn \l_@@_spdsym_read_sym_tl { dividido       } }
            {2} { \tl_set:Nn \l_@@_spdsym_read_sym_tl { dividido-por   } }
            {3} { \tl_set:Nn \l_@@_spdsym_read_sym_tl { dividido-entre } }
          }
      },
    read-sym           .initial:n  = dividido,
    read-sym           .value_required:n = true,
    read-sym / unknown .code:n =
      \msg_error:nneee { spintent } { unknown-choice }
        { read-sym } { dividido, ~dividido-por, ~dividido-entre }
        { \exp_not:n {#1} },
    unknown            .code:n =
      {
        \str_set:Nn \l_@@_command_name_str { spdsym }
        \@@_unknown_keys_cmd:n {#1}
      },
  }
%    \end{macrocode}
% \end{macro}
%
% \begin{macro}{\spdsym}
%   Defines the |\spdsym| operator. It takes a single optional argument
%   for local setup and outputs the configured intent structure, safely
%   attaching prefixes and suffixes if present.
%    \begin{macrocode}
\NewDocumentCommand \spdsym { O{} }
  {
    \mode_if_math:TF
      {
        \group_begin:
          \keys_set:nn { spintent / spdsym } { #1 }
          \MathMLintent
            {
              \tl_if_empty:NF \l_@@_spdsym_prefix_tl
                { \l_@@_spdsym_prefix_tl - }
              \l_@@_spdsym_read_sym_tl
              \tl_if_empty:NF \l_@@_spdsym_suffix_tl
                { - \l_@@_spdsym_suffix_tl }
            }
            { \l_@@_spdsym_symbol_tl }
        \group_end:
      }
      { \msg_error:nnn { spintent } { need-math-mode } { \spdsym } }
  }
%    \end{macrocode}
% \end{macro}
%
% \subsubsection{The command \cs{spmsym}}
%
% \begin{variable}[int]{\l_@@_spmsym_read_sym_tl,
%                       \l_@@_spmsym_symbol_tl,
%                       \l_@@_spmsym_prefix_tl,
%                       \l_@@_spmsym_suffix_tl}
%   Variables to store the regional reading string, the visual operator,
%   and the optional semantic affixes specifically for |\spmsym|.
%    \begin{macrocode}
\tl_new:N \l_@@_spmsym_read_sym_tl
\tl_new:N \l_@@_spmsym_symbol_tl
\tl_new:N \l_@@_spmsym_prefix_tl
\tl_new:N \l_@@_spmsym_suffix_tl
%    \end{macrocode}
% \end{variable}
%
% \begin{macro}[int]{prefix, suffix, set-sym, read-sym, not-print, unknown}
%   We define the \mymeta{keys} \mykey[spmsym]{prefix}, \mykey[spmsym]{suffix},
%   \mykey[spmsym]{set-sym}, \mykey[spmsym]{read-sym} and \mykey[spmsym]{not-print}
%   for |\spmsym|.
%    \begin{macrocode}
\keys_define:nn { spintent / spmsym }
  {
    prefix             .code:n =
      { \@@_sanitize_affix:Nn \l_@@_spmsym_prefix_tl {#1} },
    prefix             .initial:n = ,
    prefix             .value_required:n = true,
    suffix             .code:n =
      { \@@_sanitize_affix:Nn \l_@@_spmsym_suffix_tl {#1} },
    suffix             .initial:n = ,
    suffix             .value_required:n = true,
    set-sym            .choices:nn =
      { cross , dot , asterisk }
      {
        \int_case:nn { \l_keys_choice_int }
          {
            {1} { \tl_set:Nn \l_@@_spmsym_symbol_tl { \times } }
            {2} { \tl_set:Nn \l_@@_spmsym_symbol_tl { \cdot  } }
            {3} { \tl_set:Nn \l_@@_spmsym_symbol_tl { \ast   } }
          }
      },
    set-sym            .initial:n  = dot,
    set-sym            .value_required:n = true,
    set-sym / unknown  .code:n =
      \msg_error:nneee { spintent } { unknown-choice }
        { set-sym } { cross, ~dot, ~asterisk } { \exp_not:n {#1} },
    read-sym           .choices:nn =
      { por , multiplicado-por , veces }
      {
        \int_case:nn { \l_keys_choice_int }
          {
            {1} { \tl_set:Nn \l_@@_spmsym_read_sym_tl { por } }
            {2} { \tl_set:Nn \l_@@_spmsym_read_sym_tl { multiplicado-por } }
            {3} { \tl_set:Nn \l_@@_spmsym_read_sym_tl { veces } }
          }
      },
    read-sym           .initial:n  = por,
    read-sym           .value_required:n = true,
    read-sym / unknown .code:n =
      \msg_error:nneee { spintent } { unknown-choice }
        { read-sym } { por, ~multiplicado-por, ~veces } { \exp_not:n {#1} },
    not-print          .bool_set:N = \l_@@_spmsym_not_print_bool,
    not-print          .initial:n  = false,
    not-print          .value_required:n = true,
    unknown            .code:n =
      {
        \str_set:Nn \l_@@_command_name_str { spmsym }
        \@@_unknown_keys_cmd:n {#1}
      },
  }
%    \end{macrocode}
% \end{macro}
%
% \begin{macro}{\spmsym}
%   Defines the |\spmsym| operator. Similar to |\spdsym|, it
%   outputs the correct mathematical glyph based on the local \mymeta{keys},
%   safely attaching prefixes and suffixes if present.
%    \begin{macrocode}
\NewDocumentCommand \spmsym { O{} }
  {
    \mode_if_math:TF
      {
        \group_begin:
          \keys_set:nn { spintent / spmsym } { #1 }
          \MathMLintent
            {
              \tl_if_empty:NF \l_@@_spmsym_prefix_tl
                { \l_@@_spmsym_prefix_tl - }
              \l_@@_spmsym_read_sym_tl
              \tl_if_empty:NF \l_@@_spmsym_suffix_tl
                { - \l_@@_spmsym_suffix_tl }
            }
            {
              \bool_if:NTF \l_@@_spmsym_not_print_bool
                {
                  \invisibletimes
                }
                { \l_@@_spmsym_symbol_tl }
            }
        \group_end:
      }
      { \msg_error:nnn { spintent } { need-math-mode } { \spmsym } }
  }
%    \end{macrocode}
% \end{macro}
%
% \subsection{The command \cs{spdiv}}\label{cmd:spdiv}
%
% The user-level command |\spdiv| formats division representations. It
% takes two mandatory arguments: the \mymarg{dividend} and the
% \mymarg{divisor}. Both \mymarg{arguments} are internally processed by the
% |\spnum| command to ensure proper digit separation. This command enforces
% strict validation, ensuring that neither the \mymarg{dividend} nor
% the \mymarg{divisor} is \emph{blank}.
%
% \subsubsection{Definition of keys for \cs[no-index]{spdiv}}
%
% \begin{macro}[int]{cifras-sep, read-sign, read-period-sep, notation-sym,
%                    set-sym, read-sym, wrap-parens, read-parens}
%   We define the \mymeta{keys} \mykey[spdiv]{wrap-parens} and
%   \mykey[spdiv]{read-parens} for |\spdiv| command. Number-related
%   \mymeta{keys} \mykey[spdiv]{cifras-sep}, \mykey[spdiv]{read-sign},
%   \mykey[spdiv]{read-period-sep} and \mykey[spdiv]{notation-sym}
%   are forwarded to |\spnum| command using |.meta:nn|, while semantic
%   operator \mymeta{keys} \mykey[spdiv]{set-sym} and \mykey[spdiv]{read-sym}
%   are forwarded to the |\spdsym| command using |.meta:nn|.
%    \begin{macrocode}
\keys_define:nn { spintent / spdiv }
  {
    cifras-sep      .meta:nn = { spintent / spnum  } { cifras-sep      = {#1} },
    read-sign       .meta:nn = { spintent / spnum  } { read-sign       = {#1} },
    read-period-sep .meta:nn = { spintent / spnum  } { read-period-sep = {#1} },
    notation-sym    .meta:nn = { spintent / spnum  } { notation-sym    = {#1} },
    set-sym         .meta:nn = { spintent / spdsym } { set-sym         = {#1} },
    read-sym        .meta:nn = { spintent / spdsym } { read-sym        = {#1} },
    wrap-parens     .bool_set:N = \l_@@_spdiv_wrap_parens_bool,
    wrap-parens     .initial:n  = false,
    wrap-parens     .value_required:n = true,
    read-parens     .bool_set:N = \l_@@_spdiv_read_parens_bool,
    read-parens     .initial:n  = true,
    read-parens     .value_required:n = true,
    unknown         .code:n =
      {
        \str_set:Nn \l_@@_command_name_str { spdiv }
        \@@_unknown_keys_cmd:n {#1}
      },
  }
%    \end{macrocode}
% \end{macro}
%
% \subsubsection{Definition of error messages}
%
% \begin{macro}[int]{spdiv-empty-argument}
%   Raised when either the \mymarg{dividend} or the \mymarg{divisor} is
%   missing or \emph{blank}.
%    \begin{macrocode}
\msg_new:nnnn { spintent } { spdiv-empty-argument }
  { The~arguments~for~\token_to_str:N \spdiv~cannot~be~empty. }
  {
    You~must~provide~both~a~dividend~and~a~divisor.\\
  }
%    \end{macrocode}
% \end{macro}
%
% \subsubsection{Internal functions for \cs[no-index]{spdiv}}
%
% \begin{macro}[int]{\@@_spdiv_render_op:nn}
%   The function \myhlc|\@@_spdiv_render_op:nn| parses both \mymarg{arguments}
%   through the |\spnum| command and embeds the semantic operator from
%   |\spdsym|.
%    \begin{macrocode}
\cs_new_protected:Nn \@@_spdiv_render_op:nn
  {
    \spnum {#1}
    \MathMLintent { \l_@@_spdsym_read_sym_tl }
      {
        \l_@@_spdsym_symbol_tl
      }
    \spnum {#2}
  }
%    \end{macrocode}
% \end{macro}
%
% \begin{macro}[int]{\@@_spdiv_render_parens:nn}
%   The function \myhlc|\@@_spdiv_render_parens:nn| handles the structural
%   parentheses around the division.
%    \begin{macrocode}
\cs_new_protected:Nn \@@_spdiv_render_parens:nn
  {
    \bool_if:NTF \l_@@_spdiv_read_parens_bool
      {
        ( \@@_spdiv_render_op:nn {#1} {#2} )
      }
      {
        \MathMLintent { :silent } { ( }
        \@@_spdiv_render_op:nn {#1} {#2}
        \MathMLintent { :silent } { ) }
      }
  }
%    \end{macrocode}
% \end{macro}
%
% \begin{macro}[int]{\@@_spdiv_process:nnn}
%   The function \myhlc|\@@_spdiv_process:nnn| sets \mymeta{keys}, validates
%   \mymarg{arguments}, and decides on parenthesis wrapping.
%    \begin{macrocode}
\cs_new_protected:Nn \@@_spdiv_process:nnn
  {
    \keys_set:nn { spintent / spdiv } { #1 }
    \tl_if_blank:nTF { #2 }
      { \msg_error:nnn { spintent } { spdiv-empty-argument } }
      {
        \tl_if_blank:nTF { #3 }
          { \msg_error:nnn { spintent } { spdiv-empty-argument } }
          {
            \bool_if:NTF \l_@@_spdiv_wrap_parens_bool
              { \@@_spdiv_render_parens:nn {#2} {#3} }
              { \@@_spdiv_render_op:nn {#2} {#3} }
          }
      }
  }
%    \end{macrocode}
% \end{macro}
%
% \begin{macro}{\spdiv}
%   Defines the user command |\spdiv|. It strictly relies on \emph{math mode}.
%    \begin{macrocode}
\NewDocumentCommand \spdiv { O{} m m }
  {
    \mode_if_math:TF
      {
        \group_begin:
          \@@_spdiv_process:nnn { #1 } { #2 } { #3 }
        \group_end:
      }
      { \msg_error:nnn { spintent } { need-math-mode } { \spdiv } }
  }
%    \end{macrocode}
% \end{macro}
%
% \subsection{The command \cs{spmul}}\label{cmd:spmul}
%
% The user-level command |\spmul| formats multiplication representations. It
% takes two mandatory arguments: the \mymarg{first factor} and the
% \mymarg{second factor}. Both \mymarg{arguments} are internally processed by the
% |\spnum| command to ensure proper digit separation. This command enforces
% strict validation, ensuring that neither the \mymarg{first factor} nor
% the \mymarg{second factor} is \emph{blank}.
%
% \subsubsection{Definition of keys for \cs[no-index]{spmul}}
%
% \begin{macro}[int]{cifras-sep, read-sign, read-period-sep, notation-sym,
%                    set-sym, read-sym, wrap-parens, read-parens}
%   We define the \mymeta{keys} \mykey[spmul]{wrap-parens} and
%   \mykey[spmul]{read-parens} for |\spmul| command. Number-related
%   \mymeta{keys} \mykey[spmul]{cifras-sep}, \mykey[spmul]{read-sign},
%   \mykey[spmul]{read-period-sep} and \mykey[spmul]{notation-sym}
%   are forwarded to |\spnum| command using |.meta:nn|, while semantic
%   operator \mymeta{keys} \mykey[spmul]{set-sym} and \mykey[spmul]{read-sym}
%   are forwarded to the |\spmsym| command using |.meta:nn|.
%    \begin{macrocode}
\keys_define:nn { spintent / spmul }
  {
    cifras-sep      .meta:nn = { spintent / spnum  } { cifras-sep       = {#1} },
    read-sign       .meta:nn = { spintent / spnum  } { read-sign        = {#1} },
    read-period-sep .meta:nn = { spintent / spnum  } { read-period-sep  = {#1} },
    notation-sym    .meta:nn = { spintent / spnum  } { notation-sym     = {#1} },
    set-sym         .meta:nn = { spintent / spmsym } { set-sym          = {#1} },
    read-sym        .meta:nn = { spintent / spmsym } { read-sym         = {#1} },
    wrap-parens     .bool_set:N = \l_@@_spmul_wrap_parens_bool,
    wrap-parens     .initial:n  = false,
    wrap-parens     .value_required:n = true,
    read-parens     .bool_set:N = \l_@@_spmul_read_parens_bool,
    read-parens     .initial:n  = true,
    read-parens     .value_required:n = true,
    unknown         .code:n =
      {
        \str_set:Nn \l_@@_command_name_str { spmul }
        \@@_unknown_keys_cmd:n {#1}
      },
  }
%    \end{macrocode}
% \end{macro}
%
% \subsubsection{Definition of error messages}
%
% \begin{macro}[int]{spmul-empty-argument}
%   Raised when either the \mymarg{first factor} or the \mymarg{second factor} is
%   missing or \emph{blank}.
%    \begin{macrocode}
\msg_new:nnnn { spintent } { spmul-empty-argument }
  { The~arguments~for~\token_to_str:N \spmul~cannot~be~empty. }
  {
    You~must~provide~both~first~factor~and~second~factor.\\
  }
%    \end{macrocode}
% \end{macro}
%
% \subsubsection{Internal functions for \cs[no-index]{spmul}}
%
% \begin{macro}[int]{\@@_spmul_render_op:nn}
%   The function \myhlc|\@@_spmul_render_op:nn| parses both \mymarg{arguments}
%   through the |\spnum| command and embeds the semantic operator from
%   |\spmsym|.
%    \begin{macrocode}
\cs_new_protected:Nn \@@_spmul_render_op:nn
  {
    \spnum {#1}
    \MathMLintent { \l_@@_spmsym_read_sym_tl }
      {
        \l_@@_spmsym_symbol_tl
      }
    \spnum {#2}
  }
%    \end{macrocode}
% \end{macro}
%
% \begin{macro}[int]{\@@_spmul_render_parens:nn}
%   The function \myhlc|\@@_spmul_render_parens:nn| handles the structural
%   parentheses around the multiplication.
%    \begin{macrocode}
\cs_new_protected:Nn \@@_spmul_render_parens:nn
  {
    \bool_if:NTF \l_@@_spmul_read_parens_bool
      {
        ( \@@_spmul_render_op:nn {#1} {#2} )
      }
      {
        \MathMLintent { :silent } { ( }
        \@@_spmul_render_op:nn {#1} {#2}
        \MathMLintent { :silent } { ) }
      }
  }
%    \end{macrocode}
% \end{macro}
%
% \begin{macro}[int]{\@@_spmul_process:nnn}
%   The function \myhlc|\@@_spmul_process:nnn| sets \mymeta{keys}, validates
%   \mymarg{arguments}, and decides on parenthesis wrapping.
%    \begin{macrocode}
\cs_new_protected:Nn \@@_spmul_process:nnn
  {
    \keys_set:nn { spintent / spmul } { #1 }
    \tl_if_blank:nTF { #2 }
      { \msg_error:nnn { spintent } { spmul-empty-argument } }
      {
        \tl_if_blank:nTF { #3 }
          { \msg_error:nnn { spintent } { spmul-empty-argument } }
          {
            \bool_if:NTF \l_@@_spmul_wrap_parens_bool
              { \@@_spmul_render_parens:nn {#2} {#3} }
              { \@@_spmul_render_op:nn {#2} {#3} }
          }
      }
  }
%    \end{macrocode}
% \end{macro}
%
% \begin{macro}{\spmul}
%   Defines the user command |\spmul|. It strictly relies on \emph{math mode}.
%    \begin{macrocode}
\NewDocumentCommand \spmul { O{} m m }
  {
    \mode_if_math:TF
      {
        \group_begin:
          \@@_spmul_process:nnn { #1 } { #2 } { #3 }
        \group_end:
      }
      { \msg_error:nnn { spintent } { need-math-mode } { \spmul } }
  }
%    \end{macrocode}
% \end{macro}
%
% \subsection{The commands \cs{spmcd} and \cs{spmcm}}\label{cmd:spmcd}
%
% The user-level commands |\spmcd| and |\spmcm| provide a semantic interface
% for the Greatest Common Divisor (MCD in Spanish) and the Least Common
% Multiple (MCM in Spanish). They delegate argument validation and calculation
% to the Lua module \myluamod{}, constructing a valid \mynvmath{MathCAT} intent.
%
% \subsubsection{Variables for \cs{spmcd} and \cs{spmcm}}
% \begin{variable}[int]{
%   \l_@@_spmcm_spmcd_luaset_error_str, \l_@@_mc_args_seq,
%   \l_@@_spmcd_print_result_bool,   \l_@@_spmcm_print_result_bool,
%   \l_@@_spmcd_read_terse_bool,     \l_@@_spmcm_read_terse_bool,
%   \l_@@_spmcd_print_upper_bool,    \l_@@_spmcm_print_upper_bool,
%   \l_@@_spmcd_prefix_tl,           \l_@@_spmcm_prefix_tl,
%   \l_@@_spmcd_suffix_tl,           \l_@@_spmcm_suffix_tl,
% }
% Internal variables generated dynamically via a comma-separated list
% to avoid code duplication for both symmetric commands.
%    \begin{macrocode}
\str_new:N \l_@@_spmcm_spmcd_luaset_error_str
\seq_new:N \l_@@_mc_args_seq
\clist_map_inline:nn { mcd , mcm }
  {
    \tl_new:c   { l_@@_sp #1 _luaset_ #1 _value_tl }
    \bool_new:c { l_@@_sp #1 _print_result_bool }
    \bool_new:c { l_@@_sp #1 _read_terse_bool }
    \bool_new:c { l_@@_sp #1 _print_upper_bool }
    \tl_new:c   { l_@@_sp #1 _prefix_tl }
    \tl_new:c   { l_@@_sp #1 _suffix_tl }
  }
%    \end{macrocode}
% \end{variable}
%
% \subsubsection{Definition of error messages}
%
% \begin{macro}[int]{mc-invalid-arguments}
%   Exception triggered when the comma-separated sequence contains negative
%   integers, floats, or alphabetical characters.
%    \begin{macrocode}
\msg_new:nnnn { spintent } { mc-invalid-arguments }
  {
    Invalid~arguments~provided~to~#1.
  }
  {
    The~arguments~must~be~a~comma-separated~list. \\
  }
%    \end{macrocode}
% \end{macro}
%
% \subsubsection{Keys for \cs[no-index]{spmcd} and \cs[no-index]{spmcm}}
%
% \begin{macro}[int]{result,prefix,read-cmd, upper-case, sample-arg, silent-cmd}
%   We define the shared \mymeta{keys} \mykey[spmcd,spmcm]{result},
%   \mykey[spmcd,spmcm]{prefix}, \mykey[spmcd,spmcm]{read-cmd}, \mykey[spmcd,spmcm]{upper-case},
%   \mykey[spmcd,spmcm]{sample-arg} and \mykey[spmcd,spmcm]{silent-cmd}.
%    \begin{macrocode}
\clist_map_inline:nn { mcd , mcm }
  {
    \keys_define:nn { spintent / sp #1 }
      {
        result             .bool_set:c = { l_@@_sp #1 _print_result_bool },
        result             .initial:n  = false,
        result             .value_required:n = true,
        prefix             .code:n =
          { \@@_sanitize_affix:cn { l_@@_sp #1 _prefix_tl } {##1} },
        prefix             .initial:n = ,
        prefix             .value_required:n = true,
        read-cmd           .choices:nn =
          { terse , clear }
          {
            \int_case:nn { \l_keys_choice_int }
              {
                {1} { \bool_set_true:c  { l_@@_sp #1 _read_terse_bool } }
                {2} { \bool_set_false:c { l_@@_sp #1 _read_terse_bool } }
              }
          },
        read-cmd           .initial:n = clear,
        read-cmd           .value_required:n = true,
        read-cmd / unknown .code:n =
          \msg_error:nneee { spintent } { unknown-choice }
            { read-cmd } { terse , ~clear } { \exp_not:n {##1} },
        upper-case         .bool_set:c = { l_@@_sp #1 _print_upper_bool },
        upper-case         .initial:n  = false,
        upper-case         .value_required:n = true,
        sample-arg         .bool_set:c = { l_@@_sp #1 _sample_arg_bool },
        sample-arg         .initial:n  = false,
        sample-arg         .value_required:n = true,
        silent-cmd         .bool_set:c = { l_@@_sp #1 _silent_cmd_bool },
        silent-cmd         .initial:n  = false,
        silent-cmd         .value_required:n = true,
        unknown            .code:n =
          {
            \str_set:Nn \l_@@_command_name_str { sp #1 }
            \@@_unknown_keys_cmd:n {##1}
          },
      }
  }
%    \end{macrocode}
% \end{macro}
%
% \subsubsection{Internal functions for \cs[no-index]{spmcd} and \cs[no-index]{spmcm}}
%
% \begin{macro}[int]{\@@_mc_process_args:nn}
%   The function \myhlc|\@@_mc_process_args:nn| parses the operation `|#1|'
%   and its arguments `|#2|', injecting \mynvmath{MathML} intents.
%    \begin{macrocode}
\cs_new_protected:Nn \@@_mc_process_args:nn
  {
    \seq_set_from_clist:Nn \l_@@_mc_args_seq {#2}
    \bool_if:cTF { l_@@_sp #1 _silent_cmd_bool }
      {
        \seq_set_map:NNn \l_@@_mc_args_seq \l_@@_mc_args_seq
          { \MathMLintent { :silent } { ##1 } }
        \MathMLintent { entre:silent } { ( }
        \seq_use:Nnnn \l_@@_mc_args_seq
          { \MathMLintent { _y:silent } { , } }
          { \MathMLintent { :silent }   { , } }
          { \MathMLintent { _y:silent } { , } }
      }
      {
        \MathMLintent { entre } { ( }
        \seq_use:Nnnn \l_@@_mc_args_seq
          { \MathMLintent { _y }      { , } }
          { \MathMLintent { :silent } { , } }
          { \MathMLintent { _y }      { , } }
      }
    \MathMLintent { :silent } { ) }
  }
%    \end{macrocode}
% \end{macro}
%
% \begin{macro}[int]{\@@_spmc_process:nnnn}
%   The internal function \myhlc|\@@_spmc_process:nnnn| serves as a shared
%   central handler for both the |\spmcd| and |\spmcm| commands. It takes
%   the command identifier `|#1|', its full reading string `|#2|', the local
%   \mymeta{keys} `|#3|', and the optional values `|#4|'. It dynamically resolves
%   variables using the identifier to build the \mynvmath{MathML} |intent| string
%   (including any prefixes or suffixes) and prints the math operator.
%   If values are provided in `|#4|', it calculates the result using the Lua
%   backend. If an error occurs, it throws the |mc-invalid-arguments| error;
%   otherwise, it processes the arguments and conditionally prints the calculated
%   result.
%    \begin{macrocode}
\cs_new_protected:Nn \@@_spmc_process:nnnn
  {
    \keys_set:nn { spintent / sp #1 } { #3 }
    \MathMLintent
      {
        \tl_if_empty:cF { l_@@_sp #1 _prefix_tl }
          { \tl_use:c { l_@@_sp #1 _prefix_tl } - }
        \bool_if:cTF { l_@@_sp #1 _read_terse_bool }
          { #1 } { #2 }
        \bool_if:cT { l_@@_sp #1 _silent_cmd_bool }
          {
            \tl_if_novalue:nTF {#4}
              { :silent }
              { \bool_if:cT { l_@@_sp #1 _sample_arg_bool } { :silent } }
          }
      }
      {
        \operatorname
          {
            \bool_if:cTF { l_@@_sp #1 _print_upper_bool }
              { \text_uppercase:n {#1} } { #1 }
          }
      }
    \tl_if_novalue:nF { #4 }
      {
        \bool_if:cTF { l_@@_sp #1 _sample_arg_bool }
          {
            \@@_mc_process_args:nn {#1} {#4}
          }
          {
            \use:c { @@_luafun_calculate_ #1 :n } { #4 }
            \str_if_eq:VnTF \l_@@_spmcm_spmcd_luaset_error_str { true }
              {
                \msg_error:nne { spintent } { mc-invalid-arguments } { \exp_not:c { sp #1 } }
              }
              {
                \@@_mc_process_args:nn {#1} {#4}
                \bool_if:cT { l_@@_sp #1 _print_result_bool }
                  {
                    \MathMLintent { es-igual-a } { = }
                    \tl_use:c { l_@@_sp #1 _luaset_ #1 _value_tl }
                  }
              }
          }
      }
  }
%    \end{macrocode}
% \end{macro}
%
% \begin{macro}{\spmcd, \spmcm}
%   The commands |\spmcd| and |\spmcm| act as the public interface
%   wrappers for typesetting the greatest common divisor and least common
%   multiple. They enforce math mode, open a local group, and defer
%   execution to the main processing macro \myhlc|\@@_spmc_process:nnnn|,
%   passing along the appropriate reading strings and arguments.
%    \begin{macrocode}
\NewDocumentCommand \spmcd { O{} d() }
  {
    \mode_if_math:TF
      {
        \group_begin:
          \@@_spmc_process:nnnn { mcd } { máximo-común-divisor } {#1} {#2}
        \group_end:
      }
      { \msg_error:nnn { spintent } { need-math-mode } { \spmcd } }
  }
\NewDocumentCommand \spmcm { O{} d() }
  {
    \mode_if_math:TF
      {
        \group_begin:
          \@@_spmc_process:nnnn { mcm } { mínimo-común-múltiplo } {#1} {#2}
        \group_end:
      }
      { \msg_error:nnn { spintent } { need-math-mode } { \spmcm } }
  }
%    \end{macrocode}
% \end{macro}
%
% \subsection{The commands \cs{spnM} and \cs{spnD}}\label{cmd:spnM-spnD}
%
% The user-level commands |\spnM| and |\spnD| provide a semantic interface
% for the Multiples of a number (\emph{\enquote{múltiplos}} in Spanish)
% and the Divisors of a number (\emph{\enquote{divisores}} in Spanish).
% They delegate argument validation and calculation to the Lua module \myluamod{},
% constructing a valid \mynvmath{MathCAT} intent.
%
% \subsubsection{Variables for \cs{spnM} and \cs{spnD}}
%
% \begin{variable}[int]{
%   \l_@@_md_out_seq,
%   \l_@@_spnM_luaset_result_tl,        \l_@@_spnD_luaset_result_tl,
%   \l_@@_spnM_luaset_is_truncated_str, \l_@@_spnD_luaset_is_truncated_str,
%   \l_@@_spnM_print_result_bool,       \l_@@_spnD_print_result_bool,
%   \l_@@_spnM_silent_cmd_bool,         \l_@@_spnD_silent_cmd_bool,
%   \l_@@_spnM_prefix_tl,               \l_@@_spnD_prefix_tl,
%   \l_@@_spnM_print_result_num_tl,        \l_@@_spnD_print_result_num_tl
% }
%   Internal variables generated dynamically for multiples and divisors.
%    \begin{macrocode}
\seq_new:N \l_@@_md_out_seq
\clist_map_inline:nn { nM , nD }
  {
    \tl_new:c   { l_@@_sp #1 _luaset_result_tl        }
    \str_new:c  { l_@@_sp #1 _luaset_is_truncated_str }
    \bool_new:c { l_@@_sp #1 _print_result_bool       }
    \bool_new:c { l_@@_sp #1 _silent_cmd_bool         }
    \tl_new:c   { l_@@_sp #1 _prefix_tl               }
    \tl_new:c   { l_@@_sp #1 _print_result_num_tl     }
  }
%    \end{macrocode}
% \end{variable}
%
% \subsubsection{Keys for \cs[no-index]{spnM} and \cs[no-index]{spnD}}
%
% \begin{macro}[int]{result, silent-cmd, result-num, prefix, sample-arg,}
% We define the shared \mymeta{keys} \mykey[spnM,spnD]{result},
% \mykey[spnM,spnD]{silent-cmd}, \mykey[spmcd,spmcm]{result-num} and
% \mykey[spmcd,spmcm]{prefix}.
%    \begin{macrocode}
\clist_map_inline:nn { nM , nD }
  {
    \keys_define:nn { spintent / sp #1 }
      {
        result     .bool_set:c = { l_@@_sp #1 _print_result_bool },
        result     .initial:n  = false,
        result     .value_required:n = true,
        silent-cmd .bool_set:c = { l_@@_sp #1 _silent_cmd_bool },
        silent-cmd .initial:n  = false,
        silent-cmd .value_required:n = true,
        result-num .tl_set:c   = { l_@@_sp #1 _print_result_num_tl },
        result-num .initial:n  = ,
        prefix     .code:n =
          { \@@_sanitize_affix:cn { l_@@_sp #1 _prefix_tl } {##1} },
        prefix     .initial:n = ,
        prefix     .value_required:n = true,
        sample-arg .bool_set:c = { l_@@_sp #1 _sample_arg_bool },
        sample-arg .initial:n  = false,
        sample-arg .value_required:n = true,
        unknown    .code:n =
          {
            \str_set:Nn \l_@@_command_name_str { sp #1 }
            \@@_unknown_keys_cmd:n {##1}
          },
      }
  }
%    \end{macrocode}
% \end{macro}
%
% \subsubsection{Definition of error messages for \cs[no-index]{spnM} and \cs[no-index]{spnD}}
%
% \begin{macro}[int]{md-invalid-argument}
% Exception triggered when the argument for multiples or divisors is invalid.
%    \begin{macrocode}
\msg_new:nnnn { spintent } { md-invalid-argument }
  { Invalid~argument~provided~to~#1. }
  {
    The~argument~must~be~a~single~positive~integer. \\
  }
%    \end{macrocode}
% \end{macro}
%
% \subsubsection{Internal functions for \cs[no-index]{spnM} and \cs[no-index]{spnD}}
%
% \begin{macro}[int]{\@@_spnM_spnD_process:nnnnn}
%   The internal function \myhlc|\@@_spnM_spnD_process:nnnnn| acts as the
%   primary entry point for the multiples and divisors commands. It takes
%   the command identifier `|#1|', its printed mathematical symbol `|#2|',
%   its base reading string `|#3|', the local key-value options `|#4|', and
%   the optional argument `|#5|'. Its sole purpose is to evaluate and set
%   the provided keys within the current scope before delegating the core
%   branching and typesetting logic to \myhlc|\@@_spn_M_and_D_exec:nnnn|.
%    \begin{macrocode}
\cs_new_protected:Nn \@@_spnM_spnD_process:nnnnn
  {
    \keys_set:nn { spintent / sp #1 } { #4 }
    \@@_spnM_spnD_exec:nnnn { #1 } { #2 } { #3 } { #5 }
  }
%    \end{macrocode}
% \end{macro}
%
% \begin{macro}[int]{\@@_spnM_spnD_intent:nn}
%   The internal function \myhlc|\@@_spnM_spnD_intent:nn| resolves the main
%   reading intent string. It takes the command identifier `|#1|' (such as
%   "nM" or "nD") and the base reading string `|#2|'. It automatically
%   prepends the localized prefix if defined and handles specific string
%   normalizations, such as ensuring the correct accentuation for "múltiplos".
%    \begin{macrocode}
\cs_new:Nn \@@_spnM_spnD_intent:nn
  {
    \tl_if_empty:cF { l_@@_sp #1 _prefix_tl }
      { \tl_use:c { l_@@_sp #1 _prefix_tl } - }
    \str_case:nnF { #2 }
      { { multiplos } { múltiplos } }
      { #2 }
  }
%    \end{macrocode}
% \end{macro}
%
% \begin{macro}[int]{\@@_spnM_spnD_silent:nn}
%   The internal function \myhlc|\@@_spn_M_and_D_silent:nn| prints the atomized
%   notation structure in full silent mode. It takes the mathematical symbol
%   `|#1|' (such as "D" or "M") and the variable or explicit argument `|#2|'.
%   Every component is wrapped in a silent intent, and an invisible separator
%   is injected between the parentheses to prevent assistive technologies from
%   applying default contextual reading rules like "de".
%    \begin{macrocode}
\cs_new_protected:Nn \@@_spnM_spnD_silent:nn
  {
    \MathMLintent { :silent } { \operatorname { #1 } }
    \MathMLintent { :silent } { ( }
    \MathMLintent { :silent } { \@@_invisible_sep: }
    \MathMLintent { :silent } { #2 }
    \MathMLintent { :silent } { ) }
  }
%    \end{macrocode}
% \end{macro}
%
% \begin{macro}[int]{\@@_spnM_spnD_result:n}
%   The internal function \myhlc|\@@_spn_M_and_D_result:n| handles the
%   calculation and accessible formatting of the computed mathematical set.
%   It takes the command identifier `|#1|'. It invokes the corresponding Lua
%   backend calculation, maps the resulting clist into a sequence, and formats
%   the final printed output utilizing custom accessibility intents for both
%   separators and conjunctions. It also safely manages truncation indicators.
%    \begin{macrocode}
\cs_new_protected:Nn \@@_spnM_spnD_result:n
  {
    \use:c { @@_luafun_calcular_ #1 :nn }
      { \l_@@_luaset_part_int_tl }
      { \tl_use:c { l_@@_sp #1 _print_result_num_tl } }
    \MathMLintent { son } { = \{ }
    \seq_set_from_clist:Ne
      \l_@@_md_out_seq { \tl_use:c { l_@@_sp #1 _luaset_result_tl } }
    \seq_use:Nnnn \l_@@_md_out_seq
      { \MathMLintent { _y }      { , } }
      { \MathMLintent { :silent } { , } }
      { \MathMLintent { _y }      { , } }
    \str_if_eq:vnT
      { l_@@_sp #1 _luaset_is_truncated_str } { true }
        {
          \MathMLintent { :silent } { , } ~ \dots
        }
    \MathMLintent { :silent } { \} }
  }
%    \end{macrocode}
% \end{macro}
%
% \begin{macro}[int]{\@@_spnM_spnD_exec:nnnn}
%   The function \myhlc|\@@_spn_M_and_D_exec:nnnn| implements the multiples
%   and divisors commands. It takes the \mymarg{argument} `|#1|' (the command name),
%   the \mymarg{argument} `|#2|' (the mathematical symbol), the base reading
%   \mymarg{argument} `|#3|', and the optional \mymarg{argument} `|#4|'. If no
%   \mymarg{argument} is provided, it evaluates the key \mykey{silent-cmd}
%   using the default token `|n|'. If an \mymarg{argument} is given, it evaluates
%   the key \mykey{sample-arg}. If true, it processes the \mymarg{argument}
%   directly and evaluates the key \mykey{silent-cmd}. If false, it
%   bypasses the \emph{silent mode}, validates the \mymarg{argument} as a natural
%   number, and evaluates the key \mykey{result} optionally
%   calculate and append the resulting set.
%    \begin{macrocode}
\cs_new_protected:Nn \@@_spnM_spnD_exec:nnnn
  {
    \tl_if_novalue:nTF { #4 }
      {
        \bool_if:cTF { l_@@_sp #1 _silent_cmd_bool }
          { \@@_spnM_spnD_silent:nn { #2 } { n } }
          {
            \MathMLintent
              { \@@_spnM_spnD_intent:nn { #1 } { #3 } }
              { \operatorname { #2 } ( n ) }
          }
      }
      {
        \bool_if:cTF { l_@@_sp #1 _sample_arg_bool }
          {
            \bool_if:cTF { l_@@_sp #1 _silent_cmd_bool }
              { \@@_spnM_spnD_silent:nn { #2 } { #4 } }
              {
                \MathMLintent
                  { \@@_spnM_spnD_intent:nn { #1 } { #3 } }
                  { \operatorname { #2 } } ( #4 )
              }
          }
          {
            \@@_luafun_clean_split_and_set:n { \exp_not:n { #4 } }
            \str_if_eq:VnTF \l_@@_luaset_has_natural_str { true }
              {
                \MathMLintent
                  { \@@_spnM_spnD_intent:nn { #1 } { #3 } }
                  { \operatorname { #2 } } ( #4 )
                \bool_if:cT { l_@@_sp #1 _print_result_bool }
                  { \@@_spnM_spnD_result:n { #1 } }
              }
              {
                \msg_error:nne { spintent } { md-invalid-argument } { \exp_not:c { sp #1 } }
              }
          }
      }
  }
%    \end{macrocode}
% \end{macro}
%
% \begin{macro}{\spnM, \spnD}
%   The commands |\spnM| and |\spnD| act as the public interface
%   wrappers for typesetting multiples and divisors, respectively. They
%   enforce math mode, open a local group, and defer execution to the
%   function \myhlc|\@@_spnM_spnD_process:nnnnn|, passing along
%   the optional \mymeta{keys} `|#1|' and the numeric \mymarg{argument} `|#2|'.
%    \begin{macrocode}
\NewDocumentCommand \spnM { O{} d() }
  {
    \mode_if_math:TF
      {
        \group_begin:
          \@@_spnM_spnD_process:nnnnn { nM } { M } { multiplos } {#1} {#2}
        \group_end:
      }
      { \msg_error:nnn { spintent } { need-math-mode } { \spnM } }
  }
\NewDocumentCommand \spnD { O{} d() }
  {
    \mode_if_math:TF
      {
        \group_begin:
          \@@_spnM_spnD_process:nnnnn { nD } { D } { divisores } {#1} {#2}
        \group_end:
      }
      { \msg_error:nnn { spintent } { need-math-mode } { \spnD } }
  }
%    \end{macrocode}
% \end{macro}
%
% \subsection{The command \cs{spang}}\label{cmd:spang}
%
% The command |\spang| is the main public command for typesetting and
% tagging sexagesimal angle expressions (degrees, minutes, and seconds).
% It parses the input using the Lua function \myhlc|\@@_luafun_spang_parse:n|.
% This separates the degrees, minutes, and seconds to construct a \mynvmath{MathML}
% |intent| tree with the correct postfixes for \mynvmath{NVDA/MathCAT}.
% \smallskip
%
% \begin{variable}[int]{
%   \l_@@_spang_luaset_grado_str,
%   \l_@@_spang_luaset_minuto_str,
%   \l_@@_spang_luaset_segundo_str,
%   \l_@@_spang_luaset_error_str
% }
% These internal string variables store the values returned from the Lua
% module \myluamod{}. They hold the extracted numerical values for degrees,
% minutes, and seconds, as well as the error status used during formatting.
%    \begin{macrocode}
\str_new:N  \l_@@_spang_luaset_grado_str
\str_new:N  \l_@@_spang_luaset_minuto_str
\str_new:N  \l_@@_spang_luaset_segundo_str
\str_new:N  \l_@@_spang_luaset_error_str
%    \end{macrocode}
% \end{variable}
%
% \subsubsection{Definition of error messages for \cs[no-index]{spang}}
%
% \begin{macro}[int]{invalid-angle-format}
%   The error message |invalid-angle-format| defines the text displayed when
%   an invalid angle input is detected. It is triggered when an argument
%   fails the format checks, does not match the expected sexagesimal pattern,
%   or contains invalid measurement symbols.
%    \begin{macrocode}
\msg_new:nnnn { spintent } { invalid-angle-format }
  { Invalid~angle~format~in~#1. }
  {
    The~argument~'#2'~could~not~be~parsed. \\
  }
%    \end{macrocode}
% \end{macro}
%
% \subsubsection{Internal functions for \cs[no-index]{spang}}
%
% \begin{macro}[int]{\@@_spang_process:n}
%   The internal function \myhlc|\@@_spang_process:n| parses the angle input using
%   \cs{@@_luafun_spang_parse:n} after resetting the error flag in
%   \myhlc|\l_@@_spang_luaset_error_str|. If an invalid format is detected, it
%   issues the |invalid-angle-format| error. Otherwise, it checks the extracted
%   string variables \myhlc|\l_@@_spang_luaset_grado_str|,
%   \myhlc|\l_@@_spang_luaset_minuto_str|, and \myhlc|\l_@@_spang_luaset_segundo_str|.
%   For each non-empty variable, it constructs the corresponding \mynvmath{MathML} |intent|
%   structure with the correct measurement symbols and postfixes for \mynvmath{NVDA/MathCAT}.
%    \begin{macrocode}
\cs_new_protected:Nn \@@_spang_process:n
  {
    \str_set:Nn \l_@@_spang_luaset_error_str { false }
    \@@_luafun_spang_parse:n { #1 }
    \str_if_eq:VnTF \l_@@_spang_luaset_error_str { true }
      {
        \msg_error:nnnn { spintent } { invalid-angle-format } { \spang } { #1 }
      }
      {
        \str_if_empty:NF \l_@@_spang_luaset_grado_str
          {
            \spunit{ \l_@@_spang_luaset_grado_str °}
          }
        \str_if_empty:NF \l_@@_spang_luaset_minuto_str
          {
            \MathMLintent { minutos : postfix($m) }
              {
                \MathMLarg { m }
                  {
                    \spunit{ \l_@@_spang_luaset_minuto_str ′}
                  }
              }
          }
        \str_if_empty:NF \l_@@_spang_luaset_segundo_str
          {
            \MathMLintent { segundos : postfix($s) }
              {
                \MathMLarg { s }
                  {
                    \spunit { \l_@@_spang_luaset_segundo_str ″}
                  }
              }
          }
      }
  }
%    \end{macrocode}
% \end{macro}
%
% \begin{macro}{\spang}
%   The command |\spang| is the public interface for formatting
%   sexagesimal angle expressions. It first checks if it is being used in
%   math mode, issuing an error if not. Inside a local group, it calls
%   \myhlc|\@@_spang_process:n| to process the \mymarg{argument}.
%    \begin{macrocode}
\NewDocumentCommand \spang { m }
  {
    \mode_if_math:TF
      {
        \group_begin:
          \@@_spang_process:n { #1 }
        \group_end:
      }
      { \msg_error:nnn { spintent } { need-math-mode } { \spang } }
  }
%    \end{macrocode}
% \end{macro}
%
% \subsection{The command \cs{spnZ}}\label{cmd:spnz}
%
% The user-level command |\spnZ| formats integers (the set
% $\mathbb{Z}$). It includes specific options to automatically enclose
% numbers (particularly negative ones) in visual parentheses and allows
% you to control whether these parentheses are read by \mynvmath{MathCAT}.
%
% \smallskip
%
% This command verifies that the \mymarg{argument} is indeed an \emph{integer},
% checking for the complete absence of decimal separators or attached
% units of measurement.
%
% \subsubsection{Definition of keys for \cs[no-index]{spnZ}}
%
% \begin{macro}[int]{cifras-sep, read-sign, wrap-parens, read-parens,}
%   Now we define the \mymeta{keys} \mykey[spnZ]{cifras-sep} and
%   \mykey[spnZ]{read-sign}. These are |.meta:nn| \mymeta{keys} passed directly to
%   the |\spnum| command. We also introduce \mykey[spnZ]{wrap-parens},
%   which determines whether the number is enclosed in parentheses, and
%   \mykey[spnZ]{read-parens}, which controls if these parentheses are
%   read by \mynvmath{MathCAT}.
%    \begin{macrocode}
\keys_define:nn { spintent / spnZ }
  {
    cifras-sep  .meta:nn = { spintent / spnum } { cifras-sep = {#1} },
    read-sign   .meta:nn = { spintent / spnum } { read-sign  = {#1} },
    wrap-parens .bool_set:N = \l_@@_spnz_wrap_parens_bool,
    wrap-parens .initial:n  = false,
    wrap-parens .value_required:n = true,
    read-parens .bool_set:N = \l_@@_spnz_read_parens_bool,
    read-parens .initial:n  = true,
    read-parens .value_required:n = true,
    unknown     .code:n =
      {
        \str_set:Nn \l_@@_command_name_str { spnZ }
        \@@_unknown_keys_cmd:n {#1}
      },
  }
%    \end{macrocode}
% \end{macro}
%
% \subsubsection{Definition of error messages}
%
% \begin{macro}[int]{spnz-not-integer}
%   Raised when the \mymarg{argument} contains decimal separators or attached
%   units, which conflict with the constraints of a pure integer
%   representation.
%    \begin{macrocode}
\msg_new:nnnn { spintent } { spnz-not-integer }
  { The~value~'#1'~is~not~a~valid~integer. }
  {
    Command~\token_to_str:N \spnZ~only~accepts~integers.\\
   }
%    \end{macrocode}
% \end{macro}
%
% \subsubsection{Internal functions for \cs[no-index]{spnZ}}
%
% \begin{macro}[int]{\@@_spnz_render_parens:}
%   The function \myhlc|\@@_spnz_render_parens:| first checks the state of
%   the boolean variable \myhlc|\l_@@_spnz_read_parens_bool| (controlled by
%   the \mykey[spnZ]{read-parens} key). If it is |true|, it prints the
%   number surrounded by parentheses. Otherwise, it applies the |:silent|
%   \emph{function intent} to the parentheses while printing the number.
%    \begin{macrocode}
\cs_new_protected:Nn \@@_spnz_render_parens:
  {
    \bool_if:NTF \l_@@_spnz_read_parens_bool
      {
        ( \@@_spnum_print_num_start: )
      }
      {
        \MathMLintent { :silent } { ( }
        \@@_spnum_print_num_start:
        \MathMLintent { :silent } { ) }
      }
  }
%    \end{macrocode}
% \end{macro}
%
% \begin{macro}[int]{\@@_spnz_process:nn}
%   The function \myhlc|\@@_spnz_process:nn| first processes the
%   \mymeta{keys} passed in `|#1|' and then calls
%   \myhlc|\@@_luafun_clean_split_and_set:n| on `|#2|'. It verifies
%   that the input is a valid integer by checking if the variable
%   \myhlc|\l_@@_luaset_has_integer_str| is |true|. If not |true|,
%   an \enquote{error} is raised.
%
%   \smallskip
%
%   If it is |true|, it checks whether \myhlc|\l_@@_luaset_part_int_tl|
%   is \emph{empty} and raises an \enquote{error} if so. Otherwise, it checks
%   \myhlc|\l_@@_spnz_wrap_parens_bool| set by \mykey[spnZ]{wrap-parens} and
%   calls either the function \myhlc|\@@_spnz_render_parens:| or
%   \myhlc|\@@_spnum_print_num_start:|.
%    \begin{macrocode}
\cs_new_protected:Nn \@@_spnz_process:nn
  {
    \keys_set:nn { spintent / spnZ } {#1}
    \@@_luafun_clean_split_and_set:n { \exp_not:n {#2} }
    \str_if_eq:VnTF \l_@@_luaset_has_integer_str { true }
      {
        \tl_if_empty:NTF \l_@@_luaset_part_int_tl
          { \msg_error:nnn { spintent } { spnum-no-digits } {#2} }
          {
            \bool_if:NTF \l_@@_spnz_wrap_parens_bool
              { \@@_spnz_render_parens: }
              { \@@_spnum_print_num_start: }
          }
      }
      { \msg_error:nnn { spintent } { spnz-not-integer } {#2} }
  }
%    \end{macrocode}
% \end{macro}
%
% \begin{macro}{\spnZ}
%   Defines the user command |\spnZ|, which is executed within a
%   \emph{group} and strictly enforces \emph{math mode}.
%    \begin{macrocode}
\NewDocumentCommand \spnZ { O{} m }
  {
    \mode_if_math:TF
      {
        \group_begin:
          \@@_spnz_process:nn { #1 } { #2 }
        \group_end:
      }
      { \msg_error:nnn { spintent } { need-math-mode } { \spnZ } }
  }
%    \end{macrocode}
% \end{macro}
%
% \subsection{The command \cs{spmQ}}\label{cmd:spmQ}
%
% The command |\spmQ| typesets \emph{\enquote{mixed numbers}} within Spanish
% mathematical contexts (e.g., $3\frac{1}{2}$). It delegates the scaling
% of the integer part to the Lua module \myluamod{} via
% \myhlc|\@@_luafun_spmQ_scale_int:Vn|, which measures both the integer
% glyph and the fraction box at the node level and computes a scaling factor
% to align their heights visually. Finally, it constructs the correct
% \mynvmath{MathML} intent for \mynvmath{MathCAT}.
%
% \subsubsection{Definition of variables for \cs[no-index]{spmQ}}
%
% \begin{variable}[int]{
%   \l_@@_spmQ_join_word_str,
%   \l_@@_spmQ_math_style_tl,
%   \l_@@_spnQ_intent_read_sign_str,
%   \l_@@_spnQ_wrap_parens_l_tl,
%   \l_@@_spnQ_wrap_parens_r_tl,
%   \l_@@_saved_luamml_flag_int}
%   The variable \myhlc|\l_@@_spmQ_join_word_str| stores the value set by
%   the key \mykey[spmQ]{join-word} which will be verbalized by
%   \mynvmath{MathCAT} as a grammatical connector between the integer and
%   fractional parts.
%
%   \smallskip
%
%   The variable \myhlc|\l_@@_spmQ_math_style_tl| stores the current
%   mathematical style used by the function \myhlc|\@@_spmQ_set_wrap_parens_size:|.
%
%   \smallskip
%
%   The variable \myhlc|\l_@@_spnQ_intent_read_sign_str| stores the
%   |:intent| passed to |\MathMLintent| when the signs
%   `$\mathcolor{red}{+}$' or `$\mathcolor{red}{-}$' are present in the
%   integer part, used by the function \myhlc|\@@_spmQ_print_sacle_int:n|.
%
%   \smallskip
%
%   The variables \myhlc|\l_@@_spnQ_wrap_parens_l_tl| and
%   \myhlc|\l_@@_spnQ_wrap_parens_r_tl| store the sized parentheses set
%   by the function \myhlc|\@@_spmQ_set_wrap_parens_size:| and used by the
%   key \mykey[spmQ]{wrap-parens}.
%
%   \smallskip
%
%   The integer variable \myhlc|\l_@@_saved_luamml_flag_int| temporarily
%   stores the current value of |\l__luamml_flag_int| before disabling it
%   during the injection of the scaled integer node, and restores it
%   afterwards. This prevents \mypkg{luamml} from generating a spurious
%   \lstinline+<math>+ element around the |sub_box| processed in isolation,
%   before |\frac| has expanded.
%    \begin{macrocode}
\str_new:N \l_@@_spmQ_join_word_str
\tl_new:N  \l_@@_spmQ_math_style_tl
\str_new:N \l_@@_spnQ_intent_read_sign_str
\tl_new:N  \l_@@_spnQ_wrap_parens_l_tl
\tl_new:N  \l_@@_spnQ_wrap_parens_r_tl
\int_new:N \l_@@_saved_luamml_flag_int
%    \end{macrocode}
% \end{variable}
%
% \subsubsection{Definition of error messages for \cs[no-index]{spmQ}}
%
% \begin{macro}[int]{spmQ-zero-denominator, spmQ-not-integer}
%   The error message |spmQ-zero-denominator| is raised when a zero
%   denominator is passed to |\spmQ|. The second error |spmQ-not-integer|
%   is returned when the integer part passed in \mymarg{argument} is not
%   actually an integer.
% \iffalse
%% Error messages for \spmQ.
% \fi
%    \begin{macrocode}
\msg_new:nnnn { spintent } { spmQ-zero-denominator }
  { Zero~denominator~in~\token_to_str:N \spmQ. }
  {
    The~denominator~of~a~fraction~cannot~be~zero.\\
  }
\msg_new:nnnn { spintent } { spmQ-not-integer }
  { The~value~'#1'~is~not~a~valid~integer. }
  {
    Command~\token_to_str:N \spmQ \c_space_tl only~accepts~integers.\\
  }
%    \end{macrocode}
% \end{macro}
%
% \subsubsection{Definition of keys for \cs[no-index]{spmQ}}
%
% \begin{macro}[int]{join-word, use-spnZ, wrap-parens, read-parens,
%                    print-dfrac, unknown}
%   Now we add the keys \mykey[spmQ]{join-word}, \mykey[spmQ]{use-spnZ},
%   \mykey[spmQ]{wrap-parens}, \mykey[spmQ]{read-parens} and
%   \mykey[spmQ]{print-dfrac}.
%    \begin{macrocode}
\keys_define:nn { spintent / spmQ }
  {
    join-word         .choices:nn =
      { entero , enteros , y }
      {
        \int_case:nn { \l_keys_choice_int }
          {
            {1} { \str_set:Nn \l_@@_spmQ_join_word_str { entero } }
            {2} { \str_set:Nn \l_@@_spmQ_join_word_str { enteros } }
            {3} { \str_set:Nn \l_@@_spmQ_join_word_str { y } }
          }
      },
    join-word         .initial:n  = y,
    join-word         .value_required:n = true,
    join-word/unknown .code:n =
      \msg_error:nneee { spintent } { unknown-choice }
        { join-word } { entero, ~enteros, ~y } { \exp_not:n {#1} },
    use-spnZ          .bool_set:N = \l_@@_spmQ_use_spnZ_bool,
    use-spnZ          .initial:n = true,
    use-spnZ          .value_required:n = true,
    print-dfrac       .bool_set:N = \l_@@_spmQ_print_dfrac_bool,
    print-dfrac       .initial:n = false,
    print-dfrac       .value_required:n = true,
    wrap-parens       .bool_set:N = \l_@@_spnQ_wrap_parens_bool,
    wrap-parens       .initial:n  = false,
    wrap-parens       .value_required:n = true,
    read-parens       .bool_set:N = \l_@@_spnQ_read_parens_bool,
    read-parens       .initial:n  = true,
    read-parens       .value_required:n = true,
    unknown           .code:n =
      {
        \str_set:Nn \l_@@_command_name_str { spmQ }
        \@@_unknown_keys_cmd:n {#1}
      },
  }
%    \end{macrocode}
% \end{macro}
%
% \subsubsection{Internal functions for \cs[no-index]{spmQ}}
%
% \begin{macro}[int]{\@@_spmQ_set_wrap_parens_size:}
%   The function \myhlc|\@@_spmQ_set_wrap_parens_size:| captures the
%   current mathematical style using the primitive |\mathstyle| from
%   \hologo{LuaTeX} and sets the variables \cs{l_@@_spnQ_wrap_parens_l_tl}
%   and \cs{l_@@_spnQ_wrap_parens_r_tl} with the appropriate size for the
%   parentheses. Display styles (|0| and |1|) receive |\Bigl\lparen| and
%   |\Bigr\rparen|; text styles (|2| and |3|) receive |\bigl\lparen| and
%   |\bigr\rparen|; any other style falls back to unsized |\lparen| and
%   |\rparen|.
%
%   \smallskip
%
%   If the key \mykey[spmQ]{print-dfrac} is active, the size is always
%   forced to |\Bigl|/|\Bigr| regardless of the current mathematical style,
%   since |\dfrac| always produces a display-height fraction.
%    \begin{macrocode}
\cs_new_protected:Nn \@@_spmQ_set_wrap_parens_size:
  {
    \int_case:nnF { \mathstyle }
      {
        { 0 }
          {
            \tl_set:Nn \l_@@_spnQ_wrap_parens_l_tl { \Bigl \lparen }
            \tl_set:Nn \l_@@_spnQ_wrap_parens_r_tl { \Bigr \rparen }
          }
        { 1 }
          {
            \tl_set:Nn \l_@@_spnQ_wrap_parens_l_tl { \Bigl \lparen }
            \tl_set:Nn \l_@@_spnQ_wrap_parens_r_tl { \Bigr \rparen }
          }
        { 2 }
          {
            \tl_set:Nn \l_@@_spnQ_wrap_parens_l_tl { \bigl \lparen }
            \tl_set:Nn \l_@@_spnQ_wrap_parens_r_tl { \bigr \rparen }
          }
        { 3 }
          {
            \tl_set:Nn \l_@@_spnQ_wrap_parens_l_tl { \bigl \lparen }
            \tl_set:Nn \l_@@_spnQ_wrap_parens_r_tl { \bigr \rparen }
          }
      }
      {
        \tl_set:Nn \l_@@_spnQ_wrap_parens_l_tl { \lparen }
        \tl_set:Nn \l_@@_spnQ_wrap_parens_r_tl { \rparen }
      }
%    \end{macrocode}
%   If the key \mykey[spmQ]{print-dfrac} is active, we must reset the
%   variables \cs{l_@@_spnQ_wrap_parens_l_tl} and
%   \cs{l_@@_spnQ_wrap_parens_r_tl} to |\Bigl|/|\Bigr| unconditionally.
%    \begin{macrocode}
    \bool_if:NT \l_@@_spmQ_print_dfrac_bool
      {
        \tl_set:Nn \l_@@_spnQ_wrap_parens_l_tl { \Bigl \lparen }
        \tl_set:Nn \l_@@_spnQ_wrap_parens_r_tl { \Bigr \rparen }
      }
  }
%    \end{macrocode}
% \end{macro}
%
% \begin{macro}[int]{\@@_spmQ_start_wrap_parens:, \@@_spmQ_stop_wrap_parens:}
%   The functions \myhlc|\@@_spmQ_start_wrap_parens:| and
%   \myhlc|\@@_spmQ_stop_wrap_parens:| open and close the parentheses
%   around the \emph{\enquote{mixed number}} when the key
%   \mykey[spmQ]{wrap-parens} is active. They first call
%   \myhlc|\@@_spmQ_set_wrap_parens_size:| to determine the correct size
%   for the current mathematical style, then print the appropriate
%   parenthesis stored in \cs{l_@@_spnQ_wrap_parens_l_tl} or
%   \cs{l_@@_spnQ_wrap_parens_r_tl}.
%
%   \smallskip
%
%   If the variable \cs{l_@@_spnQ_read_parens_bool}, controlled by the
%   key \mykey[spmQ]{read-parens}, is \emph{false}, each parenthesis is
%   wrapped inside |\MathMLintent{:silent}| so that
%   \mynvmath{NVDA/MathCAT} does not verbalize them.
%    \begin{macrocode}
\cs_new_protected:Nn \@@_spmQ_start_wrap_parens:
  {
    \bool_if:NT \l_@@_spnQ_wrap_parens_bool
      {
        \@@_spmQ_set_wrap_parens_size:
        \bool_if:NTF \l_@@_spnQ_read_parens_bool
          { \l_@@_spnQ_wrap_parens_l_tl }
          { \MathMLintent { :silent } { \l_@@_spnQ_wrap_parens_l_tl } }
      }
  }
\cs_new_protected:Nn \@@_spmQ_stop_wrap_parens:
  {
    \bool_if:NT \l_@@_spnQ_wrap_parens_bool
      {
        \@@_spmQ_set_wrap_parens_size:
        \bool_if:NTF \l_@@_spnQ_read_parens_bool
          { \l_@@_spnQ_wrap_parens_r_tl }
          { \MathMLintent { :silent } { \l_@@_spnQ_wrap_parens_r_tl } }
      }
  }
%    \end{macrocode}
% \end{macro}
%
% \begin{macro}[int]{\@@_spmQ_set_join_word:}
%   The function \myhlc|\@@_spmQ_set_join_word:| injects the grammatical
%   connector between the integer part and the fraction by dispatching on
%   the value stored in \cs{l_@@_spmQ_join_word_str}, set by the key
%   \mykey[spmQ]{join-word}. In all three cases the visual output is the
%   Unicode invisible separator |U+2063| via \cs{c_@@_invisible_sep_tl},
%   which produces a \lstinline+<mi>+ node in the \mynvmath{MathML}
%   structure; what differs is the |:intent| text attached to it via
%   |\MathMLintent|, which \mynvmath{NVDA/MathCAT} will verbalize as
%   \emph{«entero»}, \emph{«enteros»} or \emph{«y»} respectively.
%    \begin{macrocode}
\cs_new_protected:Nn \@@_spmQ_set_join_word:
  {
    \str_case:Vn \l_@@_spmQ_join_word_str
      {
        { entero  } { \MathMLintent { _entero  } { \c_@@_invisible_sep_tl } }
        { enteros } { \MathMLintent { _enteros } { \c_@@_invisible_sep_tl } }
        { y       } { \MathMLintent { _y-      } { \c_@@_invisible_sep_tl } }
      }
  }
%    \end{macrocode}
% \end{macro}
%
% \begin{macro}[int]{\@@_spmQ_print_scale_int:}
%   The function \myhlc|\@@_spmQ_print_scale_int:| prints the integer part
%   of the \emph{\enquote{mixed number}} scaled to match the visual height
%   of the accompanying fraction. Before calling the Lua function
%   \myhlc|\@@_luafun_spmQ_scale_int:Vn|, it temporarily disables
%   \mypkg{luamml} by saving the current value of |\l__luamml_flag_int|
%   into \cs{l_@@_saved_luamml_flag_int} and setting it to |0| (skip
%   completely). This prevents \mypkg{luamml} from generating a spurious
%   \lstinline+<math>+ wrapper around the |sub_box| node that holds the
%   scaled integer, which is processed in isolation before |\frac| has
%   expanded. Once the Lua function returns, the original flag value is
%   restored via \cs{l_@@_saved_luamml_flag_int}.
%
%   \smallskip
%
%   The Lua function receives the integer digits from
%   \cs{l_@@_luaset_part_int_tl} and the fraction command name
%   (|\frac| or |\dfrac|, determined by the key \mykey[spmQ]{print-dfrac})
%   as its two arguments, and inserts the appropriately scaled |sub_box|
%   node directly into the current horizontal list.
%    \begin{macrocode}
\cs_new_protected:Nn \@@_spmQ_print_scale_int:
  {
    \int_set_eq:Nc \l_@@_saved_luamml_flag_int { l__luamml_flag_int }
    \int_set:Nn \l__luamml_flag_int { 0 }
    \bool_if:NTF \l_@@_spmQ_print_dfrac_bool
      { \@@_luafun_spmQ_scale_int:Vn \l_@@_luaset_part_int_tl { dfrac } }
      { \@@_luafun_spmQ_scale_int:Vn \l_@@_luaset_part_int_tl { frac  } }
    \int_set_eq:cN { l__luamml_flag_int } \l_@@_saved_luamml_flag_int
  }
%    \end{macrocode}
% \end{macro}
%
% \begin{macro}[int]{\@@_spmQ_print_sacle_int:n}
%   The function \myhlc|\@@_spmQ_print_sacle_int:n| takes the integer part
%   of the \emph{\enquote{mixed number}} passed in `|#1|'. It first calls
%   \myhlc|\@@_luafun_clean_split_and_set:n| on `|#1|' to parse and
%   populate the shared Lua variables, then checks the status of
%   \cs{l_@@_luaset_has_integer_str}.
%
%   \smallskip
%
%   If the status is |true|, it inspects \cs{l_@@_luaset_sign_tl}. When
%   the sign is \emph{empty}, it delegates directly to
%   \myhlc|\@@_spmQ_print_scale_int:| to render the scaled integer. When
%   a sign is present, it sets \myhlc|\l_@@_spnQ_intent_read_sign_str| to
%   either |_más:prefix($e)| or |_menos:prefix($e)| and wraps the sign
%   glyph together with a \mynvmath{MathML} argument |($e)| containing the
%   scaled integer inside a |\MathMLintent|, so that
%   \mynvmath{NVDA/MathCAT} verbalizes the sign as a prefix before the
%   integer value. An \myhlc|\@@_invisible_sep:| call between the sign and
%   the |MathMLarg| forces the opening of a new \lstinline+<mrow>+ in the
%   \mynvmath{MathML} structure, which is required for the intent tree to
%   be well-formed.
%
%   \smallskip
%
%   If the status is |false|, the |spmQ-not-integer| \enquote{error} is
%   raised immediately.
%    \begin{macrocode}
\cs_new_protected:Nn \@@_spmQ_print_sacle_int:n
  {
    \@@_luafun_clean_split_and_set:n { \exp_not:n {#1} }
    \str_if_eq:VnTF \l_@@_luaset_has_integer_str { true }
      {
        \tl_if_empty:NTF \l_@@_luaset_sign_tl
          {
            \@@_spmQ_print_scale_int:
          }
          {
            \str_if_eq:VnTF \l_@@_luaset_sign_tl { + }
              {
                \str_set:Nn
                  \l_@@_spnQ_intent_read_sign_str { _más:prefix($e) }
              }
              {
                \str_set:Nn
                  \l_@@_spnQ_intent_read_sign_str { _menos:prefix($e) }
              }
            \MathMLintent { \l_@@_spnQ_intent_read_sign_str }
              {
                \l_@@_luaset_sign_tl \@@_invisible_sep:
                \MathMLarg { e } { \@@_spmQ_print_scale_int: }
              }
          }
      }
      { \msg_error:nn { spintent } { spmQ-not-integer } }
  }
%    \end{macrocode}
% \end{macro}
%
% \begin{macro}[int]{\@@_spmQ_print_spnZ_dfrac:nn}
%   The function \myhlc|\@@_spmQ_print_spnZ_dfrac:nn| typesets the
%   fractional part of the \emph{\enquote{mixed number}}. It takes the
%   numerator in `|#1|' and the denominator in `|#2|'. It first calls
%   \myhlc|\@@_luafun_clean_split_and_set:n| on `|#1|' to validate that
%   the numerator is an integer; if \cs{l_@@_luaset_has_integer_str} is
%   |false|, the |spmQ-not-integer| \enquote{error} is raised immediately.
%
%   \smallskip
%
%   If the validation passes, the function dispatches on the two boolean
%   variables \cs{l_@@_spmQ_use_spnZ_bool} and
%   \cs{l_@@_spmQ_print_dfrac_bool}, controlled respectively by the keys
%   \mykey[spmQ]{use-spnZ} and \mykey[spmQ]{print-dfrac}. When
%   \mykey[spmQ]{use-spnZ} is |true|, both the numerator and denominator
%   are passed through |\spnZ| to gain \mynvmath{MathML} integer semantics;
%   otherwise, the raw token lists `|#1|' and `|#2|' are used directly.
%   In both branches, the fraction command is selected dynamically via
%   |\use:c|: if \mykey[spmQ]{print-dfrac} is |true|, |dfrac| is
%   assembled; otherwise, |frac| is used.
%    \begin{macrocode}
\cs_new_protected:Nn \@@_spmQ_print_spnZ_dfrac:nn
  {
    \@@_luafun_clean_split_and_set:n { \exp_not:n {#1} }
    \str_if_eq:VnT \l_@@_luaset_has_integer_str { false }
      { \msg_error:nn { spintent } { spmQ-not-integer } }
    \bool_if:NTF \l_@@_spmQ_use_spnZ_bool
      {
        \use:c { \bool_if:NT \l_@@_spmQ_print_dfrac_bool { d } frac }
          { \spnZ { #1 } } { \spnZ { #2 } }
      }
      {
        \use:c { \bool_if:NT \l_@@_spmQ_print_dfrac_bool { d } frac }
          { #1 } { #2 }
      }
  }
%    \end{macrocode}
% \end{macro}
%
% \begin{macro}[int]{\@@_spmQ_print:nnn}
%   The function \myhlc|\@@_spmQ_print:nnn| is the main rendering
%   orchestrator for |\spmQ|. It takes three arguments: the integer part
%   `|#1|', the numerator `|#2|', and the denominator `|#3|'. It opens
%   the optional parentheses via \myhlc|\@@_spmQ_start_wrap_parens:|,
%   then calls \myhlc|\@@_spmQ_print_sacle_int:n| to scale and render the
%   integer part. Next, it injects the grammatical connector between the
%   integer and the fraction via \myhlc|\@@_spmQ_set_join_word:|. Then it
%   calls \myhlc|\@@_spmQ_print_spnZ_dfrac:nn| to typeset the fraction
%   using the keys \mykey[spmQ]{print-dfrac} and \mykey[spmQ]{use-spnZ}.
%   Finally, it closes the optional parentheses via
%   \myhlc|\@@_spmQ_stop_wrap_parens:|.
%    \begin{macrocode}
\cs_new_protected:Nn \@@_spmQ_print:nnn
  {
    \@@_spmQ_start_wrap_parens:
    \@@_spmQ_print_sacle_int:n { #1 }
    \@@_spmQ_set_join_word:
    \@@_spmQ_print_spnZ_dfrac:nn { #2 } { #3 }
    \@@_spmQ_stop_wrap_parens:
  }
%    \end{macrocode}
% \end{macro}
%
% \begin{macro}[int]{\@@_spmQ_exec:nnnn}
%   The internal function \myhlc|\@@_spmQ_exec:nnnn| takes four arguments:
%   the optional \mymeta{keys} `|#1|', the integer `|#2|', the numerator
%   `|#3|', and the denominator `|#4|'. First sets the optional \mymeta{keys}
%   from `|#1|', then processes the denominator input in `|#4|' using
%   \myhlc|\@@_luafun_clean_split_and_set:n| to validate it is a
%   non-zero integer.
%
%   \smallskip
%
%   It checks the status of \cs{l_@@_luaset_has_integer_str}. If |true|
%   and the integer part is exactly zero, it issues the
%   |spmQ-zero-denominator| \enquote{error}. Otherwise, whether the status
%   is |false| or the integer part is \emph{non-zero}, it calls
%   \myhlc|\@@_spmQ_print:nnn| to format the \emph{\enquote{mixed number}}.
%    \begin{macrocode}
\cs_new_protected:Nn \@@_spmQ_exec:nnnn
  {
    \keys_set:nn { spintent / spmQ } { #1 }
    \@@_luafun_clean_split_and_set:n { \exp_not:n { #4 } }
    \str_if_eq:VnTF \l_@@_luaset_has_integer_str { true }
      {
        \int_compare:nNnTF { \l_@@_luaset_part_int_tl } = { 0 }
          { \msg_error:nn { spintent } { spmQ-zero-denominator } }
          { \@@_spmQ_print:nnn { #2 } { #3 } { #4 } }
      }
      { \@@_spmQ_print:nnn { #2 } { #3 } { #4 } }
  }
%    \end{macrocode}
% \end{macro}
%
% \begin{macro}{\spmQ}
%   The command |\spmQ| is the public interface for formatting
%   \emph{\enquote{mixed numbers}}. It first checks if it is being used in
%   \emph{math mode}, issuing an \enquote{error} if not. Inside a local
%   group, it calls \myhlc|\@@_spmQ_exec:nnnn| to process the optional
%   \mymeta{keys} and the \mymarg{arguments}.
%    \begin{macrocode}
\NewDocumentCommand \spmQ { O{} m m m }
  {
    \mode_if_math:TF
      {
        \group_begin:
          \@@_spmQ_exec:nnnn { #1 } { #2 } { #3 } { #4 }
        \group_end:
      }
      { \msg_error:nnn { spintent } { need-math-mode } { \spmQ } }
  }
%    \end{macrocode}
% \end{macro}
%
% \subsection{The commands of the \texttt{geo} family}\label{cmd:geo}
%
% The following commands provide a semantic interface for typesetting
% geometric objects in Spanish mathematical contexts. They all delegate
% argument parsing to the Lua module \myluamod{}, which builds the
% \mynvmath{MathML} |intent| string and the visual token list, sharing
% the variables \cs{l__spintent_geo_error_str},
% \cs{l__spintent_geo_intent_str}, \cs{l__spintent_geo_visual_tl} and,
% where applicable, \cs{l__spintent_geo_is_nombre_str},
% \cs{l__spintent_geo_poly_sym_str}, \cs{l__spintent_geo_intent_pts_str}
% and \cs{l__spintent_geo_angulo_caso_str}.
%
% Several commands in this family are implemented in pairs that share
% almost all of their structure — a plain object and its measure
% (|\spside|/|\spmside|, |\spangle|/|\spmangle|, |\sparc|/|\spmarc|),
% or area and perimeter (|\spAfig|/|\spPfig|). For these pairs, variable
% declarations, error messages and keys are generated once via
% \cs{clist_map_inline:nn}, following the exact pattern already used by
% |\spmcd|/|\spmcm|: a short identifier (e.g.\ |side|/|mside|) is
% substituted into csnames built with the |:c| variants of the standard
% variable and key-setting functions, so that each command in the pair
% still has its own fully independent set of variables and keys, without
% duplicating the surrounding code.
%
% \begin{important}*
%   Inside every \cs{clist_map_inline:nn} body, the mapper's own
%   parameter is |#1| (the short identifier). Any placeholder meant for
%   a function defined \emph{inside} that body — such as the |#1|/|#2|
%   arguments of a \cs{msg_new:nnnn} message text, or the |##1| already
%   required by a key's |.code:n| property — must be written one level
%   deeper than usual: message placeholders become |##1|/|##2|, and key
%   |.code:n| placeholders remain |##1| as they already are outside a
%   map. Forgetting this produces \texttt{Illegal parameter number}
%   errors from \hologo{LaTeX}.
% \end{important}
%
% \subsubsection{Definition of shared variables for the \texttt{geo} family}
%
% \begin{variable}[int]{
%   \l__spintent_geo_error_str,
%   \l__spintent_geo_intent_str,
%   \l__spintent_geo_visual_tl,
%   \l__spintent_geo_is_nombre_str,
%   \l__spintent_geo_poly_sym_str,
%   \l__spintent_geo_intent_pts_str,
%   \l__spintent_geo_angulo_caso_str,
%   \l__spintent_geo_side_cmd_tl,
%   \l__spintent_geo_fig_op_tl}
%   The variable \myhlc|\l__spintent_geo_error_str| stores the error
%   status returned by the Lua parser; it is set to |true| when the
%   input contains invalid tokens, and to |false| otherwise.
%
%   \smallskip
%
%   The variable \myhlc|\l__spintent_geo_intent_str| stores the
%   \mynvmath{MathML} |intent| string constructed by the Lua module for
%   the current geometric command.
%
%   \smallskip
%
%   The variable \myhlc|\l__spintent_geo_visual_tl| stores the
%   concatenated token list of the geometric points as they will be
%   typeset inside the visual decorator. It is shared across all
%   commands of the \texttt{geo} family.
%
%   \smallskip
%
%   The variable \myhlc|\l__spintent_geo_is_nombre_str| is set by the
%   Lua parser to |true| when the argument contains no comma and is
%   therefore interpreted as a proper name (used only by |\sprecta|),
%   and to |false| otherwise.
%
%   \smallskip
%
%   The variable \myhlc|\l__spintent_geo_poly_sym_str| is set by the
%   Lua parser to |triangle| (three points), |square| (four points), or
%   an empty string (five or more points, or when a name override
%   applies). It is read by the |\sym:| helper functions of |\spPoly|
%   and of the |\spAfig|/|\spPfig| pair.
%
%   \smallskip
%
%   The variable \myhlc|\l__spintent_geo_intent_pts_str| stores an
%   |intent| string containing only the vertex labels, with no figure
%   name prepended (e.g.\ |A-B-C| rather than |triángulo-A-B-C|); it is
%   always set alongside \cs{l__spintent_geo_intent_str} by the |\spPoly|
%   parser, allowing the expl3 side to pick the appropriate one
%   depending on \mykey[spPoly]{print-sym} without a further call into
%   Lua.
%
%   \smallskip
%
%   The variable \myhlc|\l__spintent_geo_angulo_caso_str| is set by the
%   Lua angle parser to |tres-puntos|, |un-punto|, or |nombre| to
%   indicate which of the three angle notations was detected.
%
%   \smallskip
%
%   The scratch variables \myhlc|\l__spintent_geo_side_cmd_tl| and
%   \myhlc|\l__spintent_geo_fig_op_tl| are used internally by the
%   generic processing functions of |\spside|/|\spmside|,
%   |\spangle|/|\spmangle|, |\sparc|/|\spmarc| and |\spAfig|/|\spPfig|
%   to assemble the concept name passed to Lua before it is known
%   whether the command is the plain or the measure variant of the pair.
%    \begin{macrocode}
\str_new:N \l__spintent_geo_error_str
\str_new:N \l__spintent_geo_intent_str
\tl_new:N  \l__spintent_geo_visual_tl
\str_new:N \l__spintent_geo_is_nombre_str
\str_new:N \l__spintent_geo_poly_sym_str
\str_new:N \l__spintent_geo_intent_pts_str
\str_new:N \l__spintent_geo_angulo_caso_str
\tl_new:N  \l__spintent_geo_side_cmd_tl
\tl_new:N  \l__spintent_geo_fig_op_tl
%    \end{macrocode}
% \end{variable}
%
% \subsection{The command \cs{spPoint}}\label{cmd:spPoint}
%
% The user-level command |\spPoint| provides a semantic interface for
% typesetting a single geometric point as an abstract entity, without
% coordinates. It reuses the point grammar and the helper functions
% already used by the name-notation branch of |\sprecta|
% (\S\ref{cmd:sprecta}), anchored so that any trailing content beyond a
% single point is rejected as invalid. It does not take part in the
% generic |:c| pattern, since it has no measure counterpart.
%
% \subsubsection{Definition of variables for \cs[no-index]{spPoint}}
%
% \begin{variable}[int]{
%   \l__spintent_spPoint_mode_str,
%   \l__spintent_spPoint_prefix_tl,
%   \l__spintent_spPoint_silent_cmd_bool}
%   All three variables follow the same semantics as the corresponding
%   variables of |\spside| (\S\ref{cmd:spside}).
%    \begin{macrocode}
\str_new:N  \l__spintent_spPoint_mode_str
\tl_new:N   \l__spintent_spPoint_prefix_tl
\bool_new:N \l__spintent_spPoint_silent_cmd_bool
%    \end{macrocode}
% \end{variable}
%
% \subsubsection{Definition of error messages for \cs[no-index]{spPoint}}
%
% \begin{macro}[int]{spPoint-invalid-point}
%   Exception triggered when the argument is not a single uppercase
%   letter with an optional subscript or prime modifier.
%    \begin{macrocode}
\msg_new:nnnn { spintent } { spPoint-invalid-point }
  { Invalid~point~format~in~\token_to_str:N \spPoint:~'#1'. }
  {
    The~argument~must~be~a~single~uppercase~letter~with~an~optional~
    subscript~(\_{...})~or~prime~(').\\
    Valid~examples:~'A',~'A_\{1\}',~'A'''.
  }
%    \end{macrocode}
% \end{macro}
%
% \subsubsection{Definition of keys for \cs[no-index]{spPoint}}
%
% \begin{macro}[int]{read-cmd, prefix, silent-cmd, unknown}
%    \begin{macrocode}
\keys_define:nn { spintent / spPoint }
  {
    read-cmd            .choices:nn = { school , mathcat }
      {
        \int_case:nn { \l_keys_choice_int }
          {
            {1} { \str_set:Nn \l__spintent_spPoint_mode_str { school  } }
            {2} { \str_set:Nn \l__spintent_spPoint_mode_str { mathcat } }
          }
      },
    read-cmd            .initial:n  = school,
    read-cmd            .value_required:n = true,
    read-cmd / unknown  .code:n =
      \msg_error:nneee { spintent } { unknown-choice }
        { read-cmd } { school,~mathcat } { \exp_not:n {#1} },
    prefix              .code:n =
      { \__spintent_sanitize_affix:Nn \l__spintent_spPoint_prefix_tl {#1} },
    prefix              .initial:n  = ,
    prefix              .value_required:n = true,
    silent-cmd          .bool_set:N = \l__spintent_spPoint_silent_cmd_bool,
    silent-cmd          .initial:n  = false,
    silent-cmd          .value_required:n = true,
    unknown             .code:n =
      {
        \str_set:Nn \l__spintent_command_name_str { spPoint }
        \__spintent_unknown_keys_cmd:n {#1}
      },
  }
%    \end{macrocode}
% \end{macro}
%
% \subsubsection{Internal functions for \cs[no-index]{spPoint}}
%
% \begin{macro}[int]{\__spintent_spPoint_print:n}
%   Calls \myhlc|\__spintent_luafun_geo_point_parse_and_set:nnn| with
%   the point label in `|#1|', the current mode, and the silent flag.
%   Unlike the rest of the family, no visual decorator is applied — the
%   content of \cs{MathMLintent} is simply
%   \cs{l__spintent_geo_visual_tl}, typeset directly, or wrapped in
%   \cs{MathMLarg}\marg{pt} in |mathcat| mode.
%    \begin{macrocode}
\cs_new_protected:Nn \__spintent_spPoint_print:n
  {
    \__spintent_luafun_geo_point_parse_and_set:nnn
      { \exp_not:n {#1} }
      { \l__spintent_spPoint_mode_str }
      { \bool_if:NTF \l__spintent_spPoint_silent_cmd_bool { true } { false } }
    \str_if_eq:VnT \l__spintent_geo_error_str { true }
      {
        \msg_error:nne { spintent } { spPoint-invalid-point }
          { \exp_not:n {#1} }
      }
    \str_if_eq:VnTF \l__spintent_spPoint_mode_str { school }
      {
        \MathMLintent
          {
            \tl_if_empty:NF \l__spintent_spPoint_prefix_tl
              { \tl_use:N \l__spintent_spPoint_prefix_tl - }
            \l__spintent_geo_intent_str
          }
          { \c__spintent_invisible_sep_tl \l__spintent_geo_visual_tl }
        \__spintent_invisible_sep:
      }
      {
        \MathMLintent
          {
            \tl_if_empty:NF \l__spintent_spPoint_prefix_tl
              { \tl_use:N \l__spintent_spPoint_prefix_tl - }
            \l__spintent_geo_intent_str
          }
          {
            \c__spintent_invisible_sep_tl
            \MathMLarg { pt } { \l__spintent_geo_visual_tl }
          }
        \__spintent_invisible_sep:
      }
  }
%    \end{macrocode}
% \end{macro}
%
% \begin{macro}[int]{\__spintent_spPoint_exec:nn}
%    \begin{macrocode}
\cs_new_protected:Nn \__spintent_spPoint_exec:nn
  {
    \keys_set:nn { spintent / spPoint } { #1 }
    \__spintent_spPoint_print:n { #2 }
  }
%    \end{macrocode}
% \end{macro}
%
% \begin{macro}{\spPoint}
%    \begin{macrocode}
\NewDocumentCommand \spPoint { O{} m }
  {
    \mode_if_math:TF
      {
        \group_begin:
          \__spintent_spPoint_exec:nn { #1 } { #2 }
        \group_end:
      }
      { \msg_error:nnn { spintent } { need-math-mode } { \spPoint } }
  }
%    \end{macrocode}
% \end{macro}
%
% \subsection{The command \cs{sprecta}}\label{cmd:sprecta}
%
% The user-level command |\sprecta| supports two notations detected
% automatically by the Lua parser: \emph{two-point notation} (comma
% present) renders \cs{overleftrightarrow}; \emph{name notation} (no
% comma) renders the proper name without any visual decorator. The
% shared variable \cs{l__spintent_geo_is_nombre_str} carries the
% detection result from Lua to the rendering function. It does not take
% part in the generic |:c| pattern.
%
% \subsubsection{Definition of variables for \cs[no-index]{sprecta}}
%
% \begin{variable}[int]{
%   \l__spintent_sprecta_mode_str,
%   \l__spintent_sprecta_prefix_tl,
%   \l__spintent_sprecta_silent_cmd_bool}
%   All three variables follow the same semantics as the corresponding
%   variables of |\spside| (\S\ref{cmd:spside}).
%    \begin{macrocode}
\str_new:N  \l__spintent_sprecta_mode_str
\tl_new:N   \l__spintent_sprecta_prefix_tl
\bool_new:N \l__spintent_sprecta_silent_cmd_bool
%    \end{macrocode}
% \end{variable}
%
% \subsubsection{Definition of error messages for \cs[no-index]{sprecta}}
%
% \begin{macro}[int]{sprecta-invalid-points}
%    \begin{macrocode}
\msg_new:nnnn { spintent } { sprecta-invalid-points }
  { Invalid~format~in~\token_to_str:N \sprecta:~'#1'. }
  {
    The~argument~must~be~either~a~comma-separated~pair~of~uppercase~
    letters~(two-point~notation)~or~a~single~letter~(uppercase~or~
    lowercase)~with~optional~subscript~or~prime~(name~notation).\\
    Valid~examples:~'A,~B',~'A_\{1\},~B_\{1\}',~'L_\{1\}',~'m',~'r'.
  }
%    \end{macrocode}
% \end{macro}
%
% \subsubsection{Definition of keys for \cs[no-index]{sprecta}}
%
% \begin{macro}[int]{read-cmd, prefix, silent-cmd, unknown}
%    \begin{macrocode}
\keys_define:nn { spintent / sprecta }
  {
    read-cmd            .choices:nn = { school , mathcat }
      {
        \int_case:nn { \l_keys_choice_int }
          {
            {1} { \str_set:Nn \l__spintent_sprecta_mode_str { school  } }
            {2} { \str_set:Nn \l__spintent_sprecta_mode_str { mathcat } }
          }
      },
    read-cmd            .initial:n  = school,
    read-cmd            .value_required:n = true,
    read-cmd / unknown  .code:n =
      \msg_error:nneee { spintent } { unknown-choice }
        { read-cmd } { school,~mathcat } { \exp_not:n {#1} },
    prefix              .code:n =
      { \__spintent_sanitize_affix:Nn \l__spintent_sprecta_prefix_tl {#1} },
    prefix              .initial:n  = ,
    prefix              .value_required:n = true,
    silent-cmd          .bool_set:N = \l__spintent_sprecta_silent_cmd_bool,
    silent-cmd          .initial:n  = false,
    silent-cmd          .value_required:n = true,
    unknown             .code:n =
      {
        \str_set:Nn \l__spintent_command_name_str { sprecta }
        \__spintent_unknown_keys_cmd:n {#1}
      },
  }
%    \end{macrocode}
% \end{macro}
%
% \subsubsection{Internal functions for \cs[no-index]{sprecta}}
%
% \begin{macro}[int]{\__spintent_sprecta_print:n}
%   Calls \myhlc|\__spintent_luafun_geo_parse_and_set:nnnn| with the
%   concept name |recta|, then reads
%   \cs{l__spintent_geo_is_nombre_str} to dispatch between the
%   name-notation branch (no visual decorator) and the two-point branch
%   (\cs{overleftrightarrow}).
%    \begin{macrocode}
\cs_new_protected:Nn \__spintent_sprecta_print:n
  {
    \__spintent_luafun_geo_parse_and_set:nnnn
      { recta }
      { \exp_not:n {#1} }
      { \l__spintent_sprecta_mode_str }
      { \bool_if:NTF \l__spintent_sprecta_silent_cmd_bool { true } { false } }
    \str_if_eq:VnT \l__spintent_geo_error_str { true }
      {
        \msg_error:nne { spintent } { sprecta-invalid-points }
          { \exp_not:n {#1} }
      }
    \str_if_eq:VnTF \l__spintent_geo_is_nombre_str { true }
      {
        \str_if_eq:VnTF \l__spintent_sprecta_mode_str { school }
          {
            \MathMLintent
              {
                \tl_if_empty:NF \l__spintent_sprecta_prefix_tl
                  { \tl_use:N \l__spintent_sprecta_prefix_tl - }
                \l__spintent_geo_intent_str
              }
              { \c__spintent_invisible_sep_tl \l__spintent_geo_visual_tl }
            \__spintent_invisible_sep:
          }
          {
            \MathMLintent
              {
                \tl_if_empty:NF \l__spintent_sprecta_prefix_tl
                  { \tl_use:N \l__spintent_sprecta_prefix_tl - }
                \l__spintent_geo_intent_str
              }
              {
                \c__spintent_invisible_sep_tl
                \MathMLarg { pts } { \l__spintent_geo_visual_tl }
              }
            \__spintent_invisible_sep:
          }
      }
      {
        \str_if_eq:VnTF \l__spintent_sprecta_mode_str { school }
          {
            \MathMLintent
              {
                \tl_if_empty:NF \l__spintent_sprecta_prefix_tl
                  { \tl_use:N \l__spintent_sprecta_prefix_tl - }
                \l__spintent_geo_intent_str
              }
              {
                \c__spintent_invisible_sep_tl
                \overleftrightarrow { \l__spintent_geo_visual_tl }
              }
            \__spintent_invisible_sep:
          }
          {
            \MathMLintent
              {
                \tl_if_empty:NF \l__spintent_sprecta_prefix_tl
                  { \tl_use:N \l__spintent_sprecta_prefix_tl - }
                \l__spintent_geo_intent_str
              }
              {
                \c__spintent_invisible_sep_tl
                \overleftrightarrow
                  { \MathMLarg { pts } { \l__spintent_geo_visual_tl } }
              }
            \__spintent_invisible_sep:
          }
      }
  }
%    \end{macrocode}
% \end{macro}
%
% \begin{macro}[int]{\__spintent_sprecta_exec:nn}
%    \begin{macrocode}
\cs_new_protected:Nn \__spintent_sprecta_exec:nn
  {
    \keys_set:nn { spintent / sprecta } { #1 }
    \__spintent_sprecta_print:n { #2 }
  }
%    \end{macrocode}
% \end{macro}
%
% \begin{macro}{\sprecta}
%    \begin{macrocode}
\NewDocumentCommand \sprecta { O{} m }
  {
    \mode_if_math:TF
      {
        \group_begin:
          \__spintent_sprecta_exec:nn { #1 } { #2 }
        \group_end:
      }
      { \msg_error:nnn { spintent } { need-math-mode } { \sprecta } }
  }
%    \end{macrocode}
% \end{macro}
%
% \subsection{The command \cs{sprayo}}\label{cmd:sprayo}
%
% The user-level command |\sprayo| provides a semantic interface for
% typesetting a ray. It absorbs what used to be the separate command
% |\spsemirecta|: the key \mykey[sprayo]{read-sty} selects between the
% terms |rayo| and |semirrecta|, both rendered identically with
% \cs{overrightarrow}. Since there is no measure counterpart for a ray,
% it does not take part in the generic |:c| pattern.
%
% \subsubsection{Definition of variables for \cs[no-index]{sprayo}}
%
% \begin{variable}[int]{
%   \l__spintent_sprayo_mode_str,
%   \l__spintent_sprayo_prefix_tl,
%   \l__spintent_sprayo_silent_cmd_bool,
%   \l__spintent_sprayo_read_sty_str}
%   The variables \myhlc|\l__spintent_sprayo_mode_str|,
%   \myhlc|\l__spintent_sprayo_prefix_tl| and
%   \myhlc|\l__spintent_sprayo_silent_cmd_bool| follow the same
%   semantics as the corresponding variables of |\spside|
%   (\S\ref{cmd:spside}). The variable
%   \myhlc|\l__spintent_sprayo_read_sty_str| stores the term set by the
%   key \mykey[sprayo]{read-sty}, either |rayo| or |semirrecta|, and is
%   passed directly to the Lua parser as the concept name.
%    \begin{macrocode}
\str_new:N  \l__spintent_sprayo_mode_str
\tl_new:N   \l__spintent_sprayo_prefix_tl
\bool_new:N \l__spintent_sprayo_silent_cmd_bool
\str_new:N  \l__spintent_sprayo_read_sty_str
%    \end{macrocode}
% \end{variable}
%
% \subsubsection{Definition of error messages for \cs[no-index]{sprayo}}
%
% \begin{macro}[int]{sprayo-invalid-points}
%    \begin{macrocode}
\msg_new:nnnn { spintent } { sprayo-invalid-points }
  { Invalid~point~format~in~\token_to_str:N \sprayo:~'#1'. }
  {
    The~argument~must~be~a~comma-separated~list~of~exactly~two~
    uppercase~letters~with~optional~subscripts~(\_{...})~or~primes~(').\\
    Valid~examples:~'A,~B',~'A_\{1\},~B_\{1\}',~'A'',~B'''.
  }
%    \end{macrocode}
% \end{macro}
%
% \subsubsection{Definition of keys for \cs[no-index]{sprayo}}
%
% \begin{macro}[int]{read-cmd, prefix, silent-cmd, read-sty, unknown}
%   The key \mykey[sprayo]{read-sty} replaces the former separate
%   command |\spsemirecta|: it defaults to |rayo| and accepts
%   |semirrecta| as the alternative term.
%    \begin{macrocode}
\keys_define:nn { spintent / sprayo }
  {
    read-cmd            .choices:nn = { school , mathcat }
      {
        \int_case:nn { \l_keys_choice_int }
          {
            {1} { \str_set:Nn \l__spintent_sprayo_mode_str { school  } }
            {2} { \str_set:Nn \l__spintent_sprayo_mode_str { mathcat } }
          }
      },
    read-cmd            .initial:n  = school,
    read-cmd            .value_required:n = true,
    read-cmd / unknown  .code:n =
      \msg_error:nneee { spintent } { unknown-choice }
        { read-cmd } { school,~mathcat } { \exp_not:n {#1} },
    prefix              .code:n =
      { \__spintent_sanitize_affix:Nn \l__spintent_sprayo_prefix_tl {#1} },
    prefix              .initial:n  = ,
    prefix              .value_required:n = true,
    silent-cmd          .bool_set:N = \l__spintent_sprayo_silent_cmd_bool,
    silent-cmd          .initial:n  = false,
    silent-cmd          .value_required:n = true,
    read-sty            .choices:nn = { rayo , semirrecta }
      {
        \int_case:nn { \l_keys_choice_int }
          {
            {1} { \str_set:Nn \l__spintent_sprayo_read_sty_str { rayo       } }
            {2} { \str_set:Nn \l__spintent_sprayo_read_sty_str { semirrecta } }
          }
      },
    read-sty            .initial:n  = rayo,
    read-sty            .value_required:n = true,
    read-sty / unknown  .code:n =
      \msg_error:nneee { spintent } { unknown-choice }
        { read-sty } { rayo,~semirrecta } { \exp_not:n {#1} },
    unknown             .code:n =
      {
        \str_set:Nn \l__spintent_command_name_str { sprayo }
        \__spintent_unknown_keys_cmd:n {#1}
      },
  }
%    \end{macrocode}
% \end{macro}
%
% \subsubsection{Internal functions for \cs[no-index]{sprayo}}
%
% \begin{macro}[int]{\__spintent_sprayo_print:n}
%   Passes \cs{l__spintent_sprayo_read_sty_str} — |rayo| or
%   |semirrecta| — as the concept name to
%   \myhlc|\__spintent_luafun_geo_parse_and_set:nnnn|, and uses
%   \cs{overrightarrow} as the visual decorator in both cases.
%    \begin{macrocode}
\cs_new_protected:Nn \__spintent_sprayo_print:n
  {
    \__spintent_luafun_geo_parse_and_set:nnnn
      { \l__spintent_sprayo_read_sty_str }
      { \exp_not:n {#1} }
      { \l__spintent_sprayo_mode_str }
      { \bool_if:NTF \l__spintent_sprayo_silent_cmd_bool { true } { false } }
    \str_if_eq:VnT \l__spintent_geo_error_str { true }
      {
        \msg_error:nne { spintent } { sprayo-invalid-points }
          { \exp_not:n {#1} }
      }
    \str_if_eq:VnTF \l__spintent_sprayo_mode_str { school }
      {
        \MathMLintent
          {
            \tl_if_empty:NF \l__spintent_sprayo_prefix_tl
              { \tl_use:N \l__spintent_sprayo_prefix_tl - }
            \l__spintent_geo_intent_str
          }
          {
            \c__spintent_invisible_sep_tl
            \overrightarrow { \l__spintent_geo_visual_tl }
          }
        \__spintent_invisible_sep:
      }
      {
        \MathMLintent
          {
            \tl_if_empty:NF \l__spintent_sprayo_prefix_tl
              { \tl_use:N \l__spintent_sprayo_prefix_tl - }
            \l__spintent_geo_intent_str
          }
          {
            \c__spintent_invisible_sep_tl
            \overrightarrow
              { \MathMLarg { pts } { \l__spintent_geo_visual_tl } }
          }
        \__spintent_invisible_sep:
      }
  }
%    \end{macrocode}
% \end{macro}
%
% \begin{macro}[int]{\__spintent_sprayo_exec:nn}
%    \begin{macrocode}
\cs_new_protected:Nn \__spintent_sprayo_exec:nn
  {
    \keys_set:nn { spintent / sprayo } { #1 }
    \__spintent_sprayo_print:n { #2 }
  }
%    \end{macrocode}
% \end{macro}
%
% \begin{macro}{\sprayo}
%    \begin{macrocode}
\NewDocumentCommand \sprayo { O{} m }
  {
    \mode_if_math:TF
      {
        \group_begin:
          \__spintent_sprayo_exec:nn { #1 } { #2 }
        \group_end:
      }
      { \msg_error:nnn { spintent } { need-math-mode } { \sprayo } }
  }
%    \end{macrocode}
% \end{macro}
%
% \subsection{The commands \cs{spside} and \cs{spmside}}\label{cmd:spside}
%
% The user-level commands |\spside| and |\spmside| absorb what used to
% be four separate commands (|\splado|, |\spsegmento|, |\spmlado|,
% |\spmsegmento|): the key \mykey[spside]{read-sty} selects the
% verbalized term (|lado|, |segmento|, or |trazo|, the latter used in
% straightedge-and-compass constructions), and the key
% \mykey[spside]{read-arg} allows overriding that term entirely with a
% free-form string, independently of \mykey[spside]{read-sty}.
% |\spmside| additionally provides \mykey[spmside]{print-m}, which
% controls whether |\operatorname{m}| is typeset before the segment.
%
% Both commands are generated by a single generic backend, following
% the |:c|-based pattern already established by |\spmcd|/|\spmcm|: all
% variables and keys are declared once via \cs{clist_map_inline:nn}
% over the short identifiers |side| and |mside|, and the processing
% function \myhlc|\__spintent_spside_process:nn| takes the identifier as
% its first argument to build the correct csnames at run time via the
% |:c| variants of the standard variable-access functions.
%
% \subsubsection{Definition of variables for \cs[no-index]{spside} and
%   \cs[no-index]{spmside}}
%
% \begin{variable}[int]{
%   \l__spintent_spside_mode_str,       \l__spintent_spmside_mode_str,
%   \l__spintent_spside_prefix_tl,      \l__spintent_spmside_prefix_tl,
%   \l__spintent_spside_silent_cmd_bool,\l__spintent_spmside_silent_cmd_bool,
%   \l__spintent_spside_read_sty_str,   \l__spintent_spmside_read_sty_str,
%   \l__spintent_spside_read_arg_str,   \l__spintent_spmside_read_arg_str,
%   \l__spintent_spside_print_m_bool,   \l__spintent_spmside_print_m_bool}
%   Each pair of variables is declared once for both |side| and |mside|
%   via \cs{clist_map_inline:nn}. The |_mode_str|, |_prefix_tl| and
%   |_silent_cmd_bool| variables follow the same semantics as the
%   corresponding variables of |\spPoint| (\S\ref{cmd:spPoint}). The
%   |_read_sty_str| variable stores the value set by
%   \mykey[spside]{read-sty} (|lado|, |segmento|, or |trazo|); the
%   |_read_arg_str| variable stores the optional override set by
%   \mykey[spside]{read-arg}. The |_print_m_bool| variable is only
%   meaningful for the |mside| identifier, controlling whether
%   |\operatorname{m}| is typeset — it is declared for |side| as well
%   for uniformity, but is never read by |\spside|.
%    \begin{macrocode}
\clist_map_inline:nn { side , mside }
  {
    \str_new:c  { l__spintent_sp #1 _mode_str        }
    \tl_new:c   { l__spintent_sp #1 _prefix_tl       }
    \bool_new:c { l__spintent_sp #1 _silent_cmd_bool }
    \str_new:c  { l__spintent_sp #1 _read_sty_str    }
    \str_new:c  { l__spintent_sp #1 _read_arg_str    }
    \bool_new:c { l__spintent_sp #1 _print_m_bool    }
  }
%    \end{macrocode}
% \end{variable}
%
% \subsubsection{Definition of error messages for \cs[no-index]{spside}
%   and \cs[no-index]{spmside}}
%
% \begin{macro}[int]{spside-invalid-points, spmside-invalid-points}
%   Both messages are generated by the same \cs{clist_map_inline:nn}
%   loop. The message text uses |##1|/|##2| — doubled because the text
%   is written one level inside the map's own |#1| substitution — for
%   the command name and the offending argument, both supplied later at
%   the point where \cs{msg_error:nnee} is called.
%    \begin{macrocode}
\clist_map_inline:nn { side , mside }
  {
    \msg_new:nnnn { spintent } { sp #1 -invalid-points }
      { Invalid~point~format~in~##1:~'##2'. }
      {
        The~argument~must~be~a~comma-separated~list~of~uppercase~
        letters~with~optional~subscripts~(\_{...})~or~primes~(').\\
        Valid~examples:~'A,~B',~'A_\{1\},~B_\{1\}',~'A'',~B'''.
      }
  }
%    \end{macrocode}
% \end{macro}
%
% \subsubsection{Definition of keys for \cs[no-index]{spside} and
%   \cs[no-index]{spmside}}
%
% \begin{macro}[int]{read-cmd, prefix, silent-cmd, read-sty, read-arg,
%                    print-m, unknown}
%   The keyspaces |spintent/spside| and |spintent/spmside| are both
%   generated by the same loop, each with its own independent set of
%   keys pointing at its own |:c|-named variables. The key
%   \mykey[spside]{read-sty} defaults to |lado|. The key
%   \mykey[spside]{read-arg}, when non-empty, overrides
%   \mykey[spside]{read-sty} entirely — this precedence is resolved
%   later in \myhlc|\__spintent_spside_process:nn|, not here. The key
%   \mykey[spmside]{print-m} defaults to |false|, matching the |\angle|
%   symbol convention documented for |\spmangle| (\S\ref{cmd:spmangle}).
%    \begin{macrocode}
\clist_map_inline:nn { side , mside }
  {
    \keys_define:nn { spintent / sp #1 }
      {
        read-cmd            .choices:nn = { school , mathcat }
          {
            \int_case:nn { \l_keys_choice_int }
              {
                {1} { \str_set:cn { l__spintent_sp #1 _mode_str } { school  } }
                {2} { \str_set:cn { l__spintent_sp #1 _mode_str } { mathcat } }
              }
          },
        read-cmd            .initial:n  = school,
        read-cmd            .value_required:n = true,
        read-cmd / unknown  .code:n =
          \msg_error:nneee { spintent } { unknown-choice }
            { read-cmd } { school,~mathcat } { \exp_not:n {##1} },
        prefix              .code:n =
          { \__spintent_sanitize_affix:cn { l__spintent_sp #1 _prefix_tl } {##1} },
        prefix              .initial:n  = ,
        prefix              .value_required:n = true,
        silent-cmd          .bool_set:c = { l__spintent_sp #1 _silent_cmd_bool },
        silent-cmd          .initial:n  = false,
        silent-cmd          .value_required:n = true,
        read-sty            .choices:nn = { lado , segmento , trazo }
          {
            \int_case:nn { \l_keys_choice_int }
              {
                {1} { \str_set:cn { l__spintent_sp #1 _read_sty_str } { lado      } }
                {2} { \str_set:cn { l__spintent_sp #1 _read_sty_str } { segmento  } }
                {3} { \str_set:cn { l__spintent_sp #1 _read_sty_str } { trazo     } }
              }
          },
        read-sty            .initial:n  = lado,
        read-sty            .value_required:n = true,
        read-sty / unknown  .code:n =
          \msg_error:nneee { spintent } { unknown-choice }
            { read-sty } { lado,~segmento,~trazo } { \exp_not:n {##1} },
        read-arg            .code:n =
          { \str_set:cn { l__spintent_sp #1 _read_arg_str } {##1} },
        read-arg            .initial:n  = ,
        read-arg            .value_required:n = true,
        print-m             .bool_set:c = { l__spintent_sp #1 _print_m_bool },
        print-m             .initial:n  = false,
        print-m             .value_required:n = true,
        unknown             .code:n =
          {
            \str_set:Nn \l__spintent_command_name_str { sp #1 }
            \__spintent_unknown_keys_cmd:n {##1}
          },
      }
  }
%    \end{macrocode}
% \end{macro}
%
% \subsubsection{Internal functions for \cs[no-index]{spside} and
%   \cs[no-index]{spmside}}
%
% \begin{macro}[int]{\__spintent_spside_process:nn}
%   The function \myhlc|\__spintent_spside_process:nn| takes the short
%   identifier |side|/|mside| in `|#1|' and the point list in `|#2|'.
%   It first resolves the concept name into
%   \myhlc|\l_@@_geo_side_cmd_tl|: \myhlc|l_@@_sp#1_read_arg_str|
%   when non-empty, otherwise \myhlc|l_@@_sp#1_read_sty_str|. When
%   |#1| is |mside|, |medida-del-| is prepended to that concept name.
%   The resulting string is passed as the first argument to
%   \myhlc|\_@@_luafun_geo_parse_and_set:nnnn|.
%
%   \smallskip
%
%   After checking \myhlc|\l_@@_geo_error_str|, the function builds
%   the \myhlc|\MathMLintent|. When |#1| is |mside| and
%   \myhlc|l_@@_sp#1_print_m_bool| is |true|, |\operatorname{m}| is
%   typeset wrapped in \myhlc|\MathMLintent|\marg{:silent} before
%   \cs{overline}, following the same pause-avoiding pattern documented
%   for |\spmarc| (\S\ref{cmd:sparc}).
%    \begin{macrocode}
\cs_new_protected:Nn \__spintent_spside_process:nn
  {
    \tl_set:Nn \l__spintent_geo_side_cmd_tl
      {
        \str_if_empty:cTF { l__spintent_sp #1 _read_arg_str }
          { \tl_use:c { l__spintent_sp #1 _read_sty_str } }
          { \tl_use:c { l__spintent_sp #1 _read_arg_str } }
      }
    \str_if_eq:eeTF {#1} { mside }
      { \tl_put_left:Nn \l__spintent_geo_side_cmd_tl { medida-del- } }
      { }
    \__spintent_luafun_geo_parse_and_set:nnnn
      { \l__spintent_geo_side_cmd_tl }
      { \exp_not:n {#2} }
      { \tl_use:c { l__spintent_sp #1 _mode_str } }
      { \bool_if:cTF { l__spintent_sp #1 _silent_cmd_bool } { true } { false } }
    \str_if_eq:VnT \l__spintent_geo_error_str { true }
      {
        \msg_error:nnee { spintent } { sp #1 -invalid-points }
          { \c_backslash_str sp #1 } { \exp_not:n {#2} }
      }
    \str_if_eq:eeTF { \tl_use:c { l__spintent_sp #1 _mode_str } } { school }
      {
        \MathMLintent
          {
            \tl_if_empty:cF { l__spintent_sp #1 _prefix_tl }
              { \tl_use:c { l__spintent_sp #1 _prefix_tl } - }
            \l__spintent_geo_intent_str
          }
          {
            \c__spintent_invisible_sep_tl
            \str_if_eq:eeT {#1} { mside }
              {
                \bool_if:cT { l__spintent_sp #1 _print_m_bool }
                  { \MathMLintent { :silent } { \operatorname { m } } }
              }
            \overline { \l__spintent_geo_visual_tl }
          }
        \__spintent_invisible_sep:
      }
      {
        \MathMLintent
          {
            \tl_if_empty:cF { l__spintent_sp #1 _prefix_tl }
              { \tl_use:c { l__spintent_sp #1 _prefix_tl } - }
            \l__spintent_geo_intent_str
          }
          {
            \c__spintent_invisible_sep_tl
            \str_if_eq:eeT {#1} { mside }
              {
                \bool_if:cT { l__spintent_sp #1 _print_m_bool }
                  { \MathMLintent { :silent } { \operatorname { m } } }
              }
            \overline
              { \MathMLarg { pts } { \l__spintent_geo_visual_tl } }
          }
        \__spintent_invisible_sep:
      }
  }
%    \end{macrocode}
% \end{macro}
%
% \begin{macro}[int]{\__spintent_spside_exec:nnn}
%   Receives the short identifier in `|#1|', the optional \mymeta{keys}
%   in `|#2|', and the point list in `|#3|'. Sets the keys in the
%   correct keyspace (|spintent/spside| or |spintent/spmside|) via
%   \cs{keys_set:nn} — the keyspace name is plain text substituted from
%   `|#1|', not a csname, so no |:c| variant is required here.
%    \begin{macrocode}
\cs_new_protected:Nn \__spintent_spside_exec:nnn
  {
    \keys_set:nn { spintent / sp #1 } { #2 }
    \__spintent_spside_process:nn { #1 } { #3 }
  }
%    \end{macrocode}
% \end{macro}
%
% \begin{macro}{\spside, \spmside}
%   Both public commands are thin wrappers around
%   \myhlc|\__spintent_spside_exec:nnn|, passing their own short
%   identifier as the first argument.
%    \begin{macrocode}
\NewDocumentCommand \spside { O{} m }
  {
    \mode_if_math:TF
      {
        \group_begin:
          \__spintent_spside_exec:nnn { side } { #1 } { #2 }
        \group_end:
      }
      { \msg_error:nnn { spintent } { need-math-mode } { \spside } }
  }
\NewDocumentCommand \spmside { O{} m }
  {
    \mode_if_math:TF
      {
        \group_begin:
          \__spintent_spside_exec:nnn { mside } { #1 } { #2 }
        \group_end:
      }
      { \msg_error:nnn { spintent } { need-math-mode } { \spmside } }
  }
%    \end{macrocode}
% \end{macro}
%
% \subsection{The commands \cs{spangle} and \cs{spmangle}}\label{cmd:spmangle}
%
% The user-level commands |\spangle| and |\spmangle| — renamed from the
% earlier |\spangulo|/|\spmangulo| — provide a semantic interface for
% angles, following the same generic |:c| pattern as |\spside|/
% |\spmside|. The key \mykey[spmangle]{print-m} chooses between the
% |\angle|/|\rightangle| pair (typeset together with an explicit
% |\operatorname{m}|) and the |\measuredangle|/|\measuredrightangle|
% pair (which already conveys the notion of measure on its own); the
% key \mykey[spangle]{recto} independently chooses the right-angle
% variant of whichever pair is active.
%
% \subsubsection{Definition of variables for \cs[no-index]{spangle} and
%   \cs[no-index]{spmangle}}
%
% \begin{variable}[int]{
%   \l__spintent_spangle_mode_str,        \l__spintent_spmangle_mode_str,
%   \l__spintent_spangle_prefix_tl,       \l__spintent_spmangle_prefix_tl,
%   \l__spintent_spangle_silent_cmd_bool, \l__spintent_spmangle_silent_cmd_bool,
%   \l__spintent_spangle_recto_bool,      \l__spintent_spmangle_recto_bool,
%   \l__spintent_spangle_print_m_bool,    \l__spintent_spmangle_print_m_bool}
%   Declared once for both |angle| and |mangle| via
%   \cs{clist_map_inline:nn}, following the same pattern as
%   |\spside|/|\spmside| (\S\ref{cmd:spside}). The |_recto_bool|
%   variable stores the value set by \mykey[spangle]{recto}; the
%   |_print_m_bool| variable stores the value set by
%   \mykey[spmangle]{print-m}, declared for both identifiers for
%   uniformity but only read for |mangle|.
%    \begin{macrocode}
\clist_map_inline:nn { angle , mangle }
  {
    \str_new:c  { l__spintent_sp #1 _mode_str        }
    \tl_new:c   { l__spintent_sp #1 _prefix_tl       }
    \bool_new:c { l__spintent_sp #1 _silent_cmd_bool }
    \bool_new:c { l__spintent_sp #1 _recto_bool      }
    \bool_new:c { l__spintent_sp #1 _print_m_bool    }
  }
%    \end{macrocode}
% \end{variable}
%
% \subsubsection{Definition of error messages for \cs[no-index]{spangle}
%   and \cs[no-index]{spmangle}}
%
% \begin{macro}[int]{spangle-invalid-arg, spmangle-invalid-arg}
%    \begin{macrocode}
\clist_map_inline:nn { angle , mangle }
  {
    \msg_new:nnnn { spintent } { sp #1 -invalid-arg }
      { Invalid~format~in~##1:~'##2'. }
      {
        The~argument~must~be~one~of:\\
        - three~uppercase~points~separated~by~commas~(e.g.\ 'A,~O,~B');\\
        - a~single~uppercase~point~with~optional~subscript~or~prime~
          (e.g.\ 'A',~'A_\{1\}');\\
        - an~angle~name~(e.g.\ '\string\alpha',~'1',~'ABC').
      }
  }
%    \end{macrocode}
% \end{macro}
%
% \subsubsection{Definition of keys for \cs[no-index]{spangle} and
%   \cs[no-index]{spmangle}}
%
% \begin{macro}[int]{read-cmd, prefix, silent-cmd, recto, print-m,
%                    unknown}
%    \begin{macrocode}
\clist_map_inline:nn { angle , mangle }
  {
    \keys_define:nn { spintent / sp #1 }
      {
        read-cmd            .choices:nn = { school , mathcat }
          {
            \int_case:nn { \l_keys_choice_int }
              {
                {1} { \str_set:cn { l__spintent_sp #1 _mode_str } { school  } }
                {2} { \str_set:cn { l__spintent_sp #1 _mode_str } { mathcat } }
              }
          },
        read-cmd            .initial:n  = school,
        read-cmd            .value_required:n = true,
        read-cmd / unknown  .code:n =
          \msg_error:nneee { spintent } { unknown-choice }
            { read-cmd } { school,~mathcat } { \exp_not:n {##1} },
        prefix              .code:n =
          { \__spintent_sanitize_affix:cn { l__spintent_sp #1 _prefix_tl } {##1} },
        prefix              .initial:n  = ,
        prefix              .value_required:n = true,
        silent-cmd          .bool_set:c = { l__spintent_sp #1 _silent_cmd_bool },
        silent-cmd          .initial:n  = false,
        silent-cmd          .value_required:n = true,
        recto               .bool_set:c = { l__spintent_sp #1 _recto_bool },
        recto               .initial:n  = false,
        recto               .value_required:n = true,
        print-m             .bool_set:c = { l__spintent_sp #1 _print_m_bool },
        print-m             .initial:n  = false,
        print-m             .value_required:n = true,
        unknown             .code:n =
          {
            \str_set:Nn \l__spintent_command_name_str { sp #1 }
            \__spintent_unknown_keys_cmd:n {##1}
          },
      }
  }
%    \end{macrocode}
% \end{macro}
%
% \subsubsection{Internal functions for \cs[no-index]{spangle} and
%   \cs[no-index]{spmangle}}
%
% \begin{macro}[int]{\__spintent_spangle_sym:n}
%   Takes the short identifier in `|#1|' and selects one of four visual
%   symbols. When |#1| is |mangle| and |print-m| is |true|:
%   \cs{angle} (plain) or \cs{rightangle} (|recto|); when |#1| is
%   |mangle| and |print-m| is |false|: \cs{measuredangle} or
%   \cs{measuredrightangle}. When |#1| is |angle| (no measure
%   involved), only \cs{recto} matters: \cs{angle} or \cs{rightangle}.
%    \begin{macrocode}
\cs_new_protected:Nn \__spintent_spangle_sym:n
  {
    \str_if_eq:eeTF {#1} { mangle }
      {
        \bool_if:cTF { l__spintent_sp #1 _print_m_bool }
          {
            \bool_if:cTF { l__spintent_sp #1 _recto_bool }
              { \rightangle } { \angle }
          }
          {
            \bool_if:cTF { l__spintent_sp #1 _recto_bool }
              { \measuredrightangle } { \measuredangle }
          }
      }
      {
        \bool_if:cTF { l__spintent_sp #1 _recto_bool }
          { \rightangle } { \angle }
      }
  }
%    \end{macrocode}
% \end{macro}
%
% \begin{macro}[int]{\__spintent_spangle_process:nn}
%   Resolves the concept name into
%   \cs{l__spintent_geo_side_cmd_tl} — |ángulo| for |angle|,
%   |medida-del-ángulo| for |mangle| — and calls
%   \myhlc|\__spintent_luafun_geo_angulo_parse_and_set:nnnnn| with the
%   current mode, silent flag, and |recto| flag. When |#1| is |mangle|
%   and |print-m| is |true|, |\operatorname{m}| is typeset, wrapped in
%   \cs{MathMLintent}\marg{:silent}, before
%   \myhlc|\__spintent_spangle_sym:n|.
%    \begin{macrocode}
\cs_new_protected:Nn \__spintent_spangle_process:nn
  {
    \str_if_eq:eeTF {#1} { mangle }
      { \str_set:Nn \l__spintent_geo_side_cmd_tl { medida-del-ángulo } }
      { \str_set:Nn \l__spintent_geo_side_cmd_tl { ángulo } }
    \__spintent_luafun_geo_angulo_parse_and_set:nnnnn
      { \l__spintent_geo_side_cmd_tl }
      { \exp_not:n {#2} }
      { \tl_use:c { l__spintent_sp #1 _mode_str } }
      { \bool_if:cTF { l__spintent_sp #1 _silent_cmd_bool } { true } { false } }
      { \bool_if:cTF { l__spintent_sp #1 _recto_bool }      { true } { false } }
    \str_if_eq:VnT \l__spintent_geo_error_str { true }
      {
        \msg_error:nnee { spintent } { sp #1 -invalid-arg }
          { \c_backslash_str sp #1 } { \exp_not:n {#2} }
      }
    \str_if_eq:eeTF { \tl_use:c { l__spintent_sp #1 _mode_str } } { school }
      {
        \MathMLintent
          {
            \tl_if_empty:cF { l__spintent_sp #1 _prefix_tl }
              { \tl_use:c { l__spintent_sp #1 _prefix_tl } - }
            \l__spintent_geo_intent_str
          }
          {
            \str_if_eq:eeT {#1} { mangle }
              {
                \bool_if:cT { l__spintent_sp #1 _print_m_bool }
                  { \MathMLintent { :silent } { \operatorname { m } } }
              }
            \__spintent_spangle_sym:n {#1}
            \l__spintent_geo_visual_tl
          }
        \__spintent_invisible_sep:
      }
      {
        \MathMLintent
          {
            \tl_if_empty:cF { l__spintent_sp #1 _prefix_tl }
              { \tl_use:c { l__spintent_sp #1 _prefix_tl } - }
            \l__spintent_geo_intent_str
          }
          {
            \str_if_eq:eeT {#1} { mangle }
              {
                \bool_if:cT { l__spintent_sp #1 _print_m_bool }
                  { \MathMLintent { :silent } { \operatorname { m } } }
              }
            \__spintent_spangle_sym:n {#1}
            \MathMLarg { pts } { \l__spintent_geo_visual_tl }
          }
        \__spintent_invisible_sep:
      }
  }
%    \end{macrocode}
% \end{macro}
%
% \begin{macro}[int]{\__spintent_spangle_exec:nnn}
%    \begin{macrocode}
\cs_new_protected:Nn \__spintent_spangle_exec:nnn
  {
    \keys_set:nn { spintent / sp #1 } { #2 }
    \__spintent_spangle_process:nn { #1 } { #3 }
  }
%    \end{macrocode}
% \end{macro}
%
% \begin{macro}{\spangle, \spmangle}
%    \begin{macrocode}
\NewDocumentCommand \spangle { O{} m }
  {
    \mode_if_math:TF
      {
        \group_begin:
          \__spintent_spangle_exec:nnn { angle } { #1 } { #2 }
        \group_end:
      }
      { \msg_error:nnn { spintent } { need-math-mode } { \spangle } }
  }
\NewDocumentCommand \spmangle { O{} m }
  {
    \mode_if_math:TF
      {
        \group_begin:
          \__spintent_spangle_exec:nnn { mangle } { #1 } { #2 }
        \group_end:
      }
      { \msg_error:nnn { spintent } { need-math-mode } { \spmangle } }
  }
%    \end{macrocode}
% \end{macro}
%
% \subsection{The commands \cs{sparc} and \cs{spmarc}}\label{cmd:sparc}
%
% The user-level commands |\sparc| and |\spmarc| — renamed from
% |\sparco|/|\spmarco| — provide a semantic interface for arcs,
% following the same generic |:c| pattern. Unlike |\spmside| and
% |\spmangle|, |\spmarc| has no \mykey{print-m} key: |\operatorname{m}|
% is always present, since $m\widearc{AB}$ (or its Unicode fallback)
% is the only standard notation for the measure of an arc.
%
% \subsubsection{Definition of variables for \cs[no-index]{sparc} and
%   \cs[no-index]{spmarc}}
%
% \begin{variable}[int]{
%   \l__spintent_sparc_mode_str,        \l__spintent_spmarc_mode_str,
%   \l__spintent_sparc_prefix_tl,       \l__spintent_spmarc_prefix_tl,
%   \l__spintent_sparc_silent_cmd_bool, \l__spintent_spmarc_silent_cmd_bool}
%   Declared once for both |arc| and |marc| via
%   \cs{clist_map_inline:nn}, following the same pattern as the rest of
%   the generic pairs.
%    \begin{macrocode}
\clist_map_inline:nn { arc , marc }
  {
    \str_new:c  { l__spintent_sp #1 _mode_str        }
    \tl_new:c   { l__spintent_sp #1 _prefix_tl       }
    \bool_new:c { l__spintent_sp #1 _silent_cmd_bool }
  }
%    \end{macrocode}
% \end{variable}
%
% \subsubsection{Definition of error messages for \cs[no-index]{sparc}
%   and \cs[no-index]{spmarc}}
%
% \begin{macro}[int]{sparc-invalid-points, spmarc-invalid-points}
%    \begin{macrocode}
\clist_map_inline:nn { arc , marc }
  {
    \msg_new:nnnn { spintent } { sp #1 -invalid-points }
      { Invalid~point~format~in~##1:~'##2'. }
      {
        The~argument~must~be~a~comma-separated~list~of~uppercase~
        letters~with~optional~subscripts~(\_{...})~or~primes~(').\\
        Valid~examples:~'A,~B',~'A_\{1\},~B_\{1\}',~'A'',~B'''.
      }
  }
%    \end{macrocode}
% \end{macro}
%
% \subsubsection{Definition of keys for \cs[no-index]{sparc} and
%   \cs[no-index]{spmarc}}
%
% \begin{macro}[int]{read-cmd, prefix, silent-cmd, unknown}
%    \begin{macrocode}
\clist_map_inline:nn { arc , marc }
  {
    \keys_define:nn { spintent / sp #1 }
      {
        read-cmd            .choices:nn = { school , mathcat }
          {
            \int_case:nn { \l_keys_choice_int }
              {
                {1} { \str_set:cn { l__spintent_sp #1 _mode_str } { school  } }
                {2} { \str_set:cn { l__spintent_sp #1 _mode_str } { mathcat } }
              }
          },
        read-cmd            .initial:n  = school,
        read-cmd            .value_required:n = true,
        read-cmd / unknown  .code:n =
          \msg_error:nneee { spintent } { unknown-choice }
            { read-cmd } { school,~mathcat } { \exp_not:n {##1} },
        prefix              .code:n =
          { \__spintent_sanitize_affix:cn { l__spintent_sp #1 _prefix_tl } {##1} },
        prefix              .initial:n  = ,
        prefix              .value_required:n = true,
        silent-cmd          .bool_set:c = { l__spintent_sp #1 _silent_cmd_bool },
        silent-cmd          .initial:n  = false,
        silent-cmd          .value_required:n = true,
        unknown             .code:n =
          {
            \str_set:Nn \l__spintent_command_name_str { sp #1 }
            \__spintent_unknown_keys_cmd:n {##1}
          },
      }
  }
%    \end{macrocode}
% \end{macro}
%
% \subsubsection{Internal functions for \cs[no-index]{sparc} and
%   \cs[no-index]{spmarc}}
%
% \begin{macro}[int]{\__spintent_sparc_deco:n}
%   Applies the arc decorator to its argument: \cs{widearc} when it
%   exists in the current document (e.g.\ provided by \mypkg{yhmath} or
%   \mypkg{MnSymbol}); otherwise \cs{Umathaccent}|~0~0~"23DC|, which
%   renders the extensible U+23DC glyph from the currently loaded math
%   font. The existence check is performed at render time via
%   \cs{cs_if_exist:NTF}, so |\widearc| never needs to be defined
%   globally by \mypkg*[no-index]{spintent}.
%    \begin{macrocode}
\cs_new_protected:Nn \__spintent_sparc_deco:n
  {
    \cs_if_exist:NTF \widearc
      { \widearc {#1} }
      { \Umathaccent 0 0 "23DC {#1} }
  }
%    \end{macrocode}
% \end{macro}
%
% \begin{macro}[int]{\__spintent_sparc_process:nn}
%   Resolves the concept name into
%   \cs{l__spintent_geo_side_cmd_tl} — |arco| for |arc|,
%   |medida-del-arco| for |marc| — and calls
%   \myhlc|\__spintent_luafun_geo_parse_and_set:nnnn|. When |#1| is
%   |marc|, |\operatorname{m}| is always typeset, wrapped in
%   \cs{MathMLintent}\marg{:silent}, before
%   \myhlc|\__spintent_sparc_deco:n|.
%    \begin{macrocode}
\cs_new_protected:Nn \__spintent_sparc_process:nn
  {
    \str_if_eq:eeTF {#1} { marc }
      { \str_set:Nn \l__spintent_geo_side_cmd_tl { medida-del-arco } }
      { \str_set:Nn \l__spintent_geo_side_cmd_tl { arco } }
    \__spintent_luafun_geo_parse_and_set:nnnn
      { \l__spintent_geo_side_cmd_tl }
      { \exp_not:n {#2} }
      { \tl_use:c { l__spintent_sp #1 _mode_str } }
      { \bool_if:cTF { l__spintent_sp #1 _silent_cmd_bool } { true } { false } }
    \str_if_eq:VnT \l__spintent_geo_error_str { true }
      {
        \msg_error:nnee { spintent } { sp #1 -invalid-points }
          { \c_backslash_str sp #1 } { \exp_not:n {#2} }
      }
    \str_if_eq:eeTF { \tl_use:c { l__spintent_sp #1 _mode_str } } { school }
      {
        \MathMLintent
          {
            \tl_if_empty:cF { l__spintent_sp #1 _prefix_tl }
              { \tl_use:c { l__spintent_sp #1 _prefix_tl } - }
            \l__spintent_geo_intent_str
          }
          {
            \c__spintent_invisible_sep_tl
            \str_if_eq:eeT {#1} { marc }
              { \MathMLintent { :silent } { \operatorname { m } } }
            \__spintent_sparc_deco:n { \l__spintent_geo_visual_tl }
          }
        \__spintent_invisible_sep:
      }
      {
        \MathMLintent
          {
            \tl_if_empty:cF { l__spintent_sp #1 _prefix_tl }
              { \tl_use:c { l__spintent_sp #1 _prefix_tl } - }
            \l__spintent_geo_intent_str
          }
          {
            \c__spintent_invisible_sep_tl
            \str_if_eq:eeT {#1} { marc }
              { \MathMLintent { :silent } { \operatorname { m } } }
            \__spintent_sparc_deco:n
              { \MathMLarg { pts } { \l__spintent_geo_visual_tl } }
          }
        \__spintent_invisible_sep:
      }
  }
%    \end{macrocode}
% \end{macro}
%
% \begin{macro}[int]{\__spintent_sparc_exec:nnn}
%    \begin{macrocode}
\cs_new_protected:Nn \__spintent_sparc_exec:nnn
  {
    \keys_set:nn { spintent / sp #1 } { #2 }
    \__spintent_sparc_process:nn { #1 } { #3 }
  }
%    \end{macrocode}
% \end{macro}
%
% \begin{macro}{\sparc, \spmarc}
%    \begin{macrocode}
\NewDocumentCommand \sparc { O{} m }
  {
    \mode_if_math:TF
      {
        \group_begin:
          \__spintent_sparc_exec:nnn { arc } { #1 } { #2 }
        \group_end:
      }
      { \msg_error:nnn { spintent } { need-math-mode } { \sparc } }
  }
\NewDocumentCommand \spmarc { O{} m }
  {
    \mode_if_math:TF
      {
        \group_begin:
          \__spintent_sparc_exec:nnn { marc } { #1 } { #2 }
        \group_end:
      }
      { \msg_error:nnn { spintent } { need-math-mode } { \spmarc } }
  }
%    \end{macrocode}
% \end{macro}
%
% \subsection{The command \cs{spPoly}}\label{cmd:spPoly}
%
% The user-level command |\spPoly| provides a semantic interface for
% typesetting geometric figures defined by their vertices. Its name and
% internal structure are unchanged from the earlier design; it does not
% take part in the generic |:c| pattern, since it has no measure
% counterpart of its own — that role is filled separately by
% |\spAfig|/|\spPfig| (\S\ref{cmd:spAfig}).
%
% \subsubsection{Definition of variables for \cs[no-index]{spPoly}}
%
% \begin{variable}[int]{
%   \l__spintent_spPoly_mode_str,
%   \l__spintent_spPoly_prefix_tl,
%   \l__spintent_spPoly_print_sym_bool}
%   The variables \myhlc|\l__spintent_spPoly_mode_str| and
%   \myhlc|\l__spintent_spPoly_prefix_tl| follow the same semantics as
%   the corresponding variables of |\spside| (\S\ref{cmd:spside}). The
%   boolean \myhlc|\l__spintent_spPoly_print_sym_bool| controls whether
%   the visual symbol (\cs{triangle} or \cs{mdlgwhtsquare}) and the
%   figure name are included in the output; it only has effect when
%   the number of vertices is three or four.
%    \begin{macrocode}
\str_new:N  \l__spintent_spPoly_mode_str
\tl_new:N   \l__spintent_spPoly_prefix_tl
\bool_new:N \l__spintent_spPoly_print_sym_bool
%    \end{macrocode}
% \end{variable}
%
% \subsubsection{Definition of error messages for \cs[no-index]{spPoly}}
%
% \begin{macro}[int]{spPoly-invalid-points}
%    \begin{macrocode}
\msg_new:nnnn { spintent } { spPoly-invalid-points }
  { Invalid~point~format~in~\token_to_str:N \spPoly:~'#1'. }
  {
    The~argument~must~be~a~comma-separated~list~of~at~least~three~
    uppercase~letters~with~optional~subscripts~(\_{...})~or~primes~(').\\
    Valid~examples:~'A,~B,~C',~'A_\{1\},~B_\{1\},~C_\{1\}'.
  }
%    \end{macrocode}
% \end{macro}
%
% \subsubsection{Definition of keys for \cs[no-index]{spPoly}}
%
% \begin{macro}[int]{read-cmd, prefix, print-sym, silent-cmd, unknown}
%   The key \mykey[spPoly]{silent-cmd} is a \cs{keys} alias for
%   \mykey[spPoly]{print-sym}|=false|, provided for consistency with the
%   rest of the \texttt{geo} family.
%    \begin{macrocode}
\keys_define:nn { spintent / spPoly }
  {
    read-cmd            .choices:nn = { school , mathcat }
      {
        \int_case:nn { \l_keys_choice_int }
          {
            {1} { \str_set:Nn \l__spintent_spPoly_mode_str { school  } }
            {2} { \str_set:Nn \l__spintent_spPoly_mode_str { mathcat } }
          }
      },
    read-cmd            .initial:n  = school,
    read-cmd            .value_required:n = true,
    read-cmd / unknown  .code:n =
      \msg_error:nneee { spintent } { unknown-choice }
        { read-cmd } { school,~mathcat } { \exp_not:n {#1} },
    prefix              .code:n =
      { \__spintent_sanitize_affix:Nn \l__spintent_spPoly_prefix_tl {#1} },
    prefix              .initial:n  = ,
    prefix              .value_required:n = true,
    print-sym           .bool_set:N = \l__spintent_spPoly_print_sym_bool,
    print-sym           .initial:n  = true,
    print-sym           .value_required:n = true,
    silent-cmd          .meta:n     = { print-sym = false },
    unknown             .code:n =
      {
        \str_set:Nn \l__spintent_command_name_str { spPoly }
        \__spintent_unknown_keys_cmd:n {#1}
      },
  }
%    \end{macrocode}
% \end{macro}
%
% \subsubsection{Internal functions for \cs[no-index]{spPoly}}
%
% \begin{macro}[int]{\__spintent_spPoly_sym:}
%    \begin{macrocode}
\cs_new_protected:Nn \__spintent_spPoly_sym:
  {
    \str_case:Vn \l__spintent_geo_poly_sym_str
      {
        { triangle } { \triangle      }
        { square   } { \mdlgwhtsquare }
      }
  }
%    \end{macrocode}
% \end{macro}
%
% \begin{macro}[int]{\__spintent_spPoly_print:n}
%   Calls \myhlc|\__spintent_luafun_geo_poly_parse_and_set:nn| with the
%   point list in `|#1|' and the current mode. The Lua parser sets
%   \cs{l__spintent_geo_poly_sym_str}, \cs{l__spintent_geo_intent_str}
%   and \cs{l__spintent_geo_intent_pts_str} according to the number of
%   points. When \cs{l__spintent_spPoly_print_sym_bool} is |true| and a
%   symbol is available, it is wrapped in
%   \cs{MathMLintent}\marg{:silent} and the full |intent| string is
%   used; when |false|, only the vertex tokens are typeset, using
%   \cs{l__spintent_geo_intent_pts_str} as the |intent|.
%    \begin{macrocode}
\cs_new_protected:Nn \__spintent_spPoly_print:n
  {
    \__spintent_luafun_geo_poly_parse_and_set:nn
      { \exp_not:n {#1} }
      { \l__spintent_spPoly_mode_str }
    \str_if_eq:VnT \l__spintent_geo_error_str { true }
      {
        \msg_error:nne { spintent } { spPoly-invalid-points }
          { \exp_not:n {#1} }
      }
    \bool_if:NTF \l__spintent_spPoly_print_sym_bool
      {
        \str_if_eq:VnTF \l__spintent_spPoly_mode_str { school }
          {
            \MathMLintent
              {
                \tl_if_empty:NF \l__spintent_spPoly_prefix_tl
                  { \tl_use:N \l__spintent_spPoly_prefix_tl - }
                \l__spintent_geo_intent_str
              }
              {
                \c__spintent_invisible_sep_tl
                \str_if_empty:NF \l__spintent_geo_poly_sym_str
                  { \MathMLintent { :silent } { \__spintent_spPoly_sym: } }
                \l__spintent_geo_visual_tl
              }
            \__spintent_invisible_sep:
          }
          {
            \MathMLintent
              {
                \tl_if_empty:NF \l__spintent_spPoly_prefix_tl
                  { \tl_use:N \l__spintent_spPoly_prefix_tl - }
                \l__spintent_geo_intent_str
              }
              {
                \c__spintent_invisible_sep_tl
                \str_if_empty:NF \l__spintent_geo_poly_sym_str
                  { \MathMLintent { :silent } { \__spintent_spPoly_sym: } }
                \MathMLarg { pts } { \l__spintent_geo_visual_tl }
              }
            \__spintent_invisible_sep:
          }
      }
      {
        \str_if_eq:VnTF \l__spintent_spPoly_mode_str { school }
          {
            \MathMLintent
              {
                \tl_if_empty:NF \l__spintent_spPoly_prefix_tl
                  { \tl_use:N \l__spintent_spPoly_prefix_tl - }
                \l__spintent_geo_intent_pts_str
              }
              { \c__spintent_invisible_sep_tl \l__spintent_geo_visual_tl }
            \__spintent_invisible_sep:
          }
          {
            \MathMLintent
              {
                \tl_if_empty:NF \l__spintent_spPoly_prefix_tl
                  { \tl_use:N \l__spintent_spPoly_prefix_tl - }
                \l__spintent_geo_intent_pts_str
              }
              {
                \c__spintent_invisible_sep_tl
                \MathMLarg { pts } { \l__spintent_geo_visual_tl }
              }
            \__spintent_invisible_sep:
          }
      }
  }
%    \end{macrocode}
% \end{macro}
%
% \begin{macro}[int]{\__spintent_spPoly_exec:nn}
%    \begin{macrocode}
\cs_new_protected:Nn \__spintent_spPoly_exec:nn
  {
    \keys_set:nn { spintent / spPoly } { #1 }
    \__spintent_spPoly_print:n { #2 }
  }
%    \end{macrocode}
% \end{macro}
%
% \begin{macro}{\spPoly}
%    \begin{macrocode}
\NewDocumentCommand \spPoly { O{} m }
  {
    \mode_if_math:TF
      {
        \group_begin:
          \__spintent_spPoly_exec:nn { #1 } { #2 }
        \group_end:
      }
      { \msg_error:nnn { spintent } { need-math-mode } { \spPoly } }
  }
%    \end{macrocode}
% \end{macro}
%
% \subsection{The commands \cs{spAfig} and \cs{spPfig}}\label{cmd:spAfig}
%
% The user-level commands |\spAfig| and |\spPfig| provide a semantic
% interface for the area and perimeter of a polygonal figure. Their
% names and delimited-optional-argument syntax |d()| — matching
% |\spmcd|/|\spmcm| — are unchanged; what changed is the backend, now
% generated by the same generic |:c| pattern as the rest of the family,
% over the short identifiers |Afig| and |Pfig|.
%
% \subsubsection{Definition of variables for \cs[no-index]{spAfig} and
%   \cs[no-index]{spPfig}}
%
% \begin{variable}[int]{
%   \l__spintent_spAfig_mode_str,      \l__spintent_spPfig_mode_str,
%   \l__spintent_spAfig_prefix_tl,     \l__spintent_spPfig_prefix_tl,
%   \l__spintent_spAfig_print_sym_bool,\l__spintent_spPfig_print_sym_bool,
%   \l__spintent_spAfig_read_arg_str,  \l__spintent_spPfig_read_arg_str}
%   Declared once for both |Afig| and |Pfig| via
%   \cs{clist_map_inline:nn}. The |_print_sym_bool| variable follows the
%   same semantics as \cs{l__spintent_spPoly_print_sym_bool}
%   (\S\ref{cmd:spPoly}); the |_read_arg_str| variable stores the
%   figure-name override set by \mykey[spAfig]{read-arg}, independent
%   of \mykey[spAfig]{print-sym}.
%    \begin{macrocode}
\clist_map_inline:nn { Afig , Pfig }
  {
    \str_new:c  { l__spintent_sp #1 _mode_str        }
    \tl_new:c   { l__spintent_sp #1 _prefix_tl       }
    \bool_new:c { l__spintent_sp #1 _print_sym_bool  }
    \str_new:c  { l__spintent_sp #1 _read_arg_str    }
  }
%    \end{macrocode}
% \end{variable}
%
% \subsubsection{Definition of error messages for \cs[no-index]{spAfig}
%   and \cs[no-index]{spPfig}}
%
% \begin{macro}[int]{spAfig-invalid-points, spPfig-invalid-points}
%    \begin{macrocode}
\clist_map_inline:nn { Afig , Pfig }
  {
    \msg_new:nnnn { spintent } { sp #1 -invalid-points }
      { Invalid~point~format~in~##1:~'##2'. }
      {
        The~argument~must~be~empty,~or~a~comma-separated~list~of~at~
        least~three~uppercase~letters~with~optional~subscripts~
        (\_{...})~or~primes~(').\\
        Valid~examples:~'()',~'(A,~B,~C)',~'(A_\{1\},~B_\{1\},~C_\{1\})'.
      }
  }
%    \end{macrocode}
% \end{macro}
%
% \subsubsection{Definition of keys for \cs[no-index]{spAfig} and
%   \cs[no-index]{spPfig}}
%
% \begin{macro}[int]{read-cmd, prefix, print-sym, read-arg, unknown}
%    \begin{macrocode}
\clist_map_inline:nn { Afig , Pfig }
  {
    \keys_define:nn { spintent / sp #1 }
      {
        read-cmd            .choices:nn = { school , mathcat }
          {
            \int_case:nn { \l_keys_choice_int }
              {
                {1} { \str_set:cn { l__spintent_sp #1 _mode_str } { school  } }
                {2} { \str_set:cn { l__spintent_sp #1 _mode_str } { mathcat } }
              }
          },
        read-cmd            .initial:n  = school,
        read-cmd            .value_required:n = true,
        read-cmd / unknown  .code:n =
          \msg_error:nneee { spintent } { unknown-choice }
            { read-cmd } { school,~mathcat } { \exp_not:n {##1} },
        prefix              .code:n =
          { \__spintent_sanitize_affix:cn { l__spintent_sp #1 _prefix_tl } {##1} },
        prefix              .initial:n  = ,
        prefix              .value_required:n = true,
        print-sym           .bool_set:c = { l__spintent_sp #1 _print_sym_bool },
        print-sym           .initial:n  = true,
        print-sym           .value_required:n = true,
        read-arg            .code:n =
          { \str_set:cn { l__spintent_sp #1 _read_arg_str } {##1} },
        read-arg            .initial:n  = ,
        read-arg            .value_required:n = true,
        unknown             .code:n =
          {
            \str_set:Nn \l__spintent_command_name_str { sp #1 }
            \__spintent_unknown_keys_cmd:n {##1}
          },
      }
  }
%    \end{macrocode}
% \end{macro}
%
% \subsubsection{Internal functions for \cs[no-index]{spAfig} and
%   \cs[no-index]{spPfig}}
%
% \begin{macro}[int]{\__spintent_spfig_sym:}
%    \begin{macrocode}
\cs_new_protected:Nn \__spintent_spfig_sym:
  {
    \str_case:Vn \l__spintent_geo_poly_sym_str
      {
        { triangle } { \triangle      }
        { square   } { \mdlgwhtsquare }
      }
  }
%    \end{macrocode}
% \end{macro}
%
% \begin{macro}[int]{\__spintent_spfig_process:nn}
%   Takes the short identifier |Afig|/|Pfig| in `|#1|' and the point
%   list in `|#2|' (possibly empty). Sets
%   \cs{l__spintent_geo_fig_op_tl} to |área| or |perímetro| — the word
%   Lua needs to build the |intent| — distinct from the visual symbol
%   $\mathcal{A}$/$\mathcal{P}$, which is selected separately when
%   typesetting. Calls
%   \myhlc|\__spintent_luafun_geo_fig_parse_and_set:nnnnn| with the
%   operator word, the point list, the mode, the print-symbol flag, and
%   the read-arg override.
%
%   \smallskip
%
%   When the point list is non-empty, the subscript is built with
%   \cs{c_math_subscript_token} rather than a literal |_|, since the
%   latter is not guaranteed to be read as a subscript operator inside
%   the token stream constructed by \cs{MathMLintent}.
%    \begin{macrocode}
\tl_new:N \l__spintent_geo_fig_op_tl
\cs_new_protected:Nn \__spintent_spfig_process:nn
  {
    \str_if_eq:eeTF {#1} { Afig }
      { \tl_set:Nn \l__spintent_geo_fig_op_tl { área } }
      { \tl_set:Nn \l__spintent_geo_fig_op_tl { perímetro } }
    \__spintent_luafun_geo_fig_parse_and_set:nnnnn
      { \l__spintent_geo_fig_op_tl }
      { \exp_not:n {#2} }
      { \tl_use:c { l__spintent_sp #1 _mode_str } }
      { \bool_if:cTF { l__spintent_sp #1 _print_sym_bool } { true } { false } }
      { \tl_use:c { l__spintent_sp #1 _read_arg_str } }
    \str_if_eq:VnT \l__spintent_geo_error_str { true }
      {
        \msg_error:nnee { spintent } { sp #1 -invalid-points }
          { \c_backslash_str sp #1 } { \exp_not:n {#2} }
      }
    \str_if_eq:eeTF { \tl_use:c { l__spintent_sp #1 _mode_str } } { school }
      {
        \MathMLintent
          {
            \tl_if_empty:cF { l__spintent_sp #1 _prefix_tl }
              { \tl_use:c { l__spintent_sp #1 _prefix_tl } - }
            \l__spintent_geo_intent_str
          }
          {
            \c__spintent_invisible_sep_tl
            \str_if_eq:eeTF {#1} { Afig }
              { \mathcal { A } } { \mathcal { P } }
            \tl_if_empty:NF \l__spintent_geo_visual_tl
              {
                \c_math_subscript_token
                  { \__spintent_spfig_sym: \l__spintent_geo_visual_tl }
              }
          }
        \__spintent_invisible_sep:
      }
      {
        \MathMLintent
          {
            \tl_if_empty:cF { l__spintent_sp #1 _prefix_tl }
              { \tl_use:c { l__spintent_sp #1 _prefix_tl } - }
            \l__spintent_geo_intent_str
          }
          {
            \c__spintent_invisible_sep_tl
            \str_if_eq:eeTF {#1} { Afig }
              { \mathcal { A } } { \mathcal { P } }
            \tl_if_empty:NF \l__spintent_geo_visual_tl
              {
                \c_math_subscript_token
                  {
                    \__spintent_spfig_sym:
                    \MathMLarg { pts } { \l__spintent_geo_visual_tl }
                  }
              }
          }
        \__spintent_invisible_sep:
      }
  }
%    \end{macrocode}
% \end{macro}
%
% \begin{macro}[int]{\__spintent_spfig_exec:nnn}
%    \begin{macrocode}
\cs_new_protected:Nn \__spintent_spfig_exec:nnn
  {
    \keys_set:nn { spintent / sp #1 } { #2 }
    \__spintent_spfig_process:nn { #1 } { #3 }
  }
%    \end{macrocode}
% \end{macro}
%
% \begin{macro}{\spAfig, \spPfig}
%   Following the same \cs{d()}-argument pattern as |\spmcd|/|\spmcm|,
%   both public commands convert the |-NoValue-| marker produced by an
%   omitted delimited argument into an explicit empty group before
%   forwarding to \myhlc|\__spintent_spfig_exec:nnn|, since the Lua
%   parser interprets an empty point list as \enquote{no figure — just
%   the bare symbol}.
%    \begin{macrocode}
\NewDocumentCommand \spAfig { O{} d() }
  {
    \mode_if_math:TF
      {
        \group_begin:
          \tl_if_novalue:nTF {#2}
            { \__spintent_spfig_exec:nnn { Afig } { #1 } { } }
            { \__spintent_spfig_exec:nnn { Afig } { #1 } { #2 } }
        \group_end:
      }
      { \msg_error:nnn { spintent } { need-math-mode } { \spAfig } }
  }
\NewDocumentCommand \spPfig { O{} d() }
  {
    \mode_if_math:TF
      {
        \group_begin:
          \tl_if_novalue:nTF {#2}
            { \__spintent_spfig_exec:nnn { Pfig } { #1 } { } }
            { \__spintent_spfig_exec:nnn { Pfig } { #1 } { #2 } }
        \group_end:
      }
      { \msg_error:nnn { spintent } { need-math-mode } { \spPfig } }
  }
%    \end{macrocode}
% \end{macro}
%
% \subsection{Finish package}
%
% Finish package implementation.
%
%    \begin{macrocode}
\file_input_stop:
%</package>
%    \end{macrocode}
% \Finale
% \endinput
% \iffalse
%<*spintentlua>
% \fi
--[[
     Lua module spintent.lua for spintent package - v0.98 [2026-08-20]
--]]

-- CACHÉ, LPEG Y HERRAMIENTAS GLOBALES

local tonumber   = tonumber
local tostring   = tostring
local math_floor = math.floor
local math_sqrt  = math.sqrt
local math_min   = math.min
local t_sort     = table.sort
local s_format   = string.format
local t_concat   = table.concat
local t_unpack   = table.unpack or unpack
local s_match    = string.match
local s_gmatch   = string.gmatch
local s_sub      = string.sub
local s_gsub     = string.gsub
local s_lower    = string.lower
local s_upper    = string.upper
local s_byte     = string.byte

-- Nuevas referencias locales para rendimiento
local pairs      = pairs

local token_set_macro = token.set_macro

-- APIs node.direct cacheadas a nivel de módulo (una sola vez, no en cada
-- invocación de \spmQ) para máximo rendimiento, mismo criterio que mylhmc.lua.
local d_todirect = node.direct.todirect
local d_tonode   = node.direct.tonode
local d_traverse = node.direct.traverse_glyph
local d_getfield = node.direct.getfield
local d_new      = node.direct.new
local d_setfield = node.direct.setfield

local lpeg_base = lpeg or require("lpeg")
local P, R, S, Cs, Ct, Cg, C, Cc =
    lpeg_base.P,
    lpeg_base.R,
    lpeg_base.S,
    lpeg_base.Cs,
    lpeg_base.Ct,
    lpeg_base.Cg,
    lpeg_base.C,
    lpeg_base.Cc

local spintent_spshort_V  = lpeg_base.V
local spintent_digit = R"09"

local spintent_math_space =
    P"\\," + P"\\-" + P"\\;" + P"\\:" + P"\\!" +
    P"\\>" + P"\\quad" + P"\\qquad" + P"~" +
    P"\\ " + P"\\thinspace" + P"\\nobreakspace"

local spintent_discard_space = (S" \t\r\n" + spintent_math_space) / ""
local spintent_opt_spaces = spintent_discard_space^0

-- Trim de espacios centralizado en LPeg (sin backtracking), reemplaza el
-- patrón s_match(x, "^%s*(.-)%s*$") repetido en cada función de registro.
local spintent_ws = S" \t\r\n"
local spintent_trim_pattern =
    spintent_ws^0 * C((P(1) - spintent_ws^0 * P(-1))^0) * spintent_ws^0
local function spintent_trim(str)
    return spintent_trim_pattern:match(str)
end

-- Combina trim + minúsculas para los casos que no necesitan conservar
-- el case original después (a diferencia de luafun_spmoney_lookup_data,
-- que sí reutiliza la versión sin s_lower para la detección ISO).
local function spintent_normalize_key(str)
    return str and s_lower(spintent_trim(str)) or ""
end

-- Eliminación de espacios matemáticos para cálculos
local spintent_num_chunk = Cs(
    spintent_digit * (spintent_opt_spaces * spintent_digit)^0
)

local spintent_num_chunk_keep = Cs(
    spintent_digit * (
        ((S" \t\r\n" + spintent_math_space)^0 * #spintent_digit) * spintent_digit
    )^0
)

-- Registro de las funciones para expl3
local function register_tex_cmd(name, func, args)
    name = "__spintent_" .. name .. ":" .. string.rep("n", #args)
    local scanners = {}
    for i = 1, #args do
        local scan_type = (args[i] == "string" and "scan_argument") or "scan_" .. args[i]
        scanners[i] = token[scan_type]
    end

    local scanning_func
    if #scanners == 0 then
        scanning_func = func
    elseif #scanners == 1 then
        local s1 = scanners[1]
        scanning_func = function() func(s1()) end
    elseif #scanners == 2 then
        local s1, s2 = scanners[1], scanners[2]
        scanning_func = function() func(s1(), s2()) end
    else
        scanning_func = function()
            local values = {}
            for i = 1, #scanners do values[i] = scanners[i]() end
            func(t_unpack(values))
        end
    end

    local index = luatexbase.new_luafunction(name)
    lua.get_functions_table()[index] = scanning_func
    token.set_lua(name, index, "global", "protected")
end

-- 1. MOTOR CENTRAL Y NÚMEROS (\spnum y \spunit)

local spintent_num_sign = S "+-"
local spintent_num_decimal = S ".," + P "{.}" + P "{,}"
local spintent_num_semi = P ";"
local spintent_forbidden_in_extra = spintent_num_decimal + spintent_num_semi
local spintent_custom_spunit_aliases = {}

local spintent_extra_char = P(1) - spintent_forbidden_in_extra

local spintent_number_pattern = Ct(
    Cg(spintent_num_sign^-1, "sign")
    * spintent_opt_spaces
    * Cg(spintent_num_chunk_keep^-1, "integer")
    * (
        spintent_opt_spaces
        * Cg(C(spintent_num_decimal), "decimal")
        * spintent_opt_spaces
        * Cg(spintent_num_chunk_keep, "fraction")
    )^-1
    * (
        spintent_opt_spaces
        * spintent_num_semi
        * spintent_opt_spaces
        * Cg(spintent_num_chunk_keep, "period")
    )^-1
    * (
        spintent_opt_spaces
        * S"eE"
        * spintent_opt_spaces
        * Cg(Cs(spintent_num_sign^-1 * spintent_num_chunk), "exponent")
    )^-1
    * spintent_opt_spaces
    * Cg(Cs(spintent_extra_char^0), "extra")
    * P(-1)
)

-- Optimización: uso exclusivo de longitud de bytes (#) y s_sub para dígitos (O(1) vs O(N))
local function spintent_rae_format_digits(str_num, reverse)
    if not str_num or str_num == "" then return "" end

    str_num = tostring(str_num)
    local len = #str_num -- Los números en Lua son 1 byte por carácter (ASCII)
    if len <= 4 then return str_num end

    local chunks = {}
    if reverse then
        for i = 1, len, 3 do
            chunks[#chunks + 1] = s_sub(str_num, i, i + 2)
        end
    else
        local first = len % 3
        if first == 0 then first = 3 end
        chunks[#chunks + 1] = s_sub(str_num, 1, first)

        for i = first + 1, len, 3 do
            chunks[#chunks + 1] = s_sub(str_num, i, i + 2)
        end
    end

    return t_concat(chunks, "\\,")
end

-- Patrón de una sola pasada: reemplaza el loop de N gsub por un único
-- recorrido LPeg que elimina cualquier separador matemático o espacio.
local spintent_strip_math_spaces = Cs((spintent_discard_space + P(1)) ^ 0)

register_tex_cmd("luafun_clean_split_and_set", function(raw_string)
    raw_string = spintent_trim(raw_string)

    local result = spintent_number_pattern:match(raw_string)

    -- Prevención de creación de tabla si falla el match (nil or {})
    local r_sign = (result and result.sign) or ""
    local r_int  = (result and result.integer) or ""
    local r_frac = (result and result.fraction) or ""

    local r_int_clean  = spintent_strip_math_spaces:match(r_int)
    local r_frac_clean = spintent_strip_math_spaces:match(r_frac)

    local r_dec    = (result and result.decimal) or ""
    local r_period = (result and result.period) or ""
    local r_extra  = (result and result.extra)
    local r_exp    = (result and result.exponent) or ""

    local has_exponent = (r_exp ~= "") and "true" or "false"
    local above = ""
    local below = ""
    local unit_status = "valid"
    local has_units = "false"
    local denom_has_numeric = "false"

    if r_extra then
        local extra = spintent_trim(r_extra)

        if spintent_custom_spunit_aliases[extra] then
            extra = spintent_custom_spunit_aliases[extra]
        end

        if extra ~= "" then
            has_units = "true"
            local parts = {}
            for part in s_gmatch(extra, "[^/]+") do
                parts[#parts + 1] = part
            end

            if #parts > 2 or s_match(extra, "/%s*/") then
                unit_status = "multislash"
            else
                above = parts[1] or ""
                below = parts[2] or ""

                if below ~= "" and s_match(below, "^%s*%d") then
                    denom_has_numeric = "true"
                end
            end
        end
    end

    local is_million_clean = "false"
    if r_int_clean ~= "" then
        local int_value = tonumber(r_int_clean)
        if int_value and int_value > 999999 and (int_value % 1000000 == 0) then
            is_million_clean = "true"
        end
    end

    local has_dec = r_dec ~= ""
    local has_per = r_period ~= ""

    local dec_and_per      = "false"
    local not_dec_and_per  = "false"
    local not_dec_not_per  = "false"

    if has_dec and has_per then
        dec_and_per = "true"
    elseif not has_dec and has_per then
        not_dec_and_per = "true"
    elseif not has_dec and not has_per then
        not_dec_not_per = "true"
    end

    local has_integer = "false"
    local has_natural = "false"

    if r_int_clean ~= "" and r_dec == "" and r_frac_clean == "" and r_period == "" and has_exponent == "false" and has_units == "false" then
        has_integer = "true"
        if r_sign == "" then
            local int_val = tonumber(r_int_clean)
            if int_val and int_val > 0 then
                has_natural = "true"
            end
        end
    end

    token_set_macro("l__spintent_luaset_part_int_tl", r_int)
    token_set_macro("l__spintent_luaset_part_dec_tl", r_frac)
    token_set_macro("l__spintent_luaset_sign_tl", r_sign)
    token_set_macro("l__spintent_luaset_dec_sep_tl", r_dec)
    token_set_macro("l__spintent_luaset_part_period_tl", r_period)
    token_set_macro("l__spintent_luaset_only_part_int_str", r_int_clean)
    token_set_macro("l__spintent_luaset_only_part_dec_str", r_frac_clean)
    token_set_macro("l__spintent_luaset_only_part_period_str", (result and result.period) or "")
    token_set_macro("l__spintent_luaset_format_part_int_str", spintent_rae_format_digits(r_int_clean, false))
    token_set_macro("l__spintent_luaset_format_part_dec_str", spintent_rae_format_digits(r_frac_clean, true))
    token_set_macro("l__spintent_luaset_millons_str", is_million_clean)
    token_set_macro("l__spintent_luaset_decimal_period_str", dec_and_per)
    token_set_macro("l__spintent_luaset_not_decimal_period_str", not_dec_and_per)
    token_set_macro("l__spintent_luaset_not_decimal_not_period_str", not_dec_not_per)
    token_set_macro("l__spintent_luaset_has_exponent_str", has_exponent)
    token_set_macro("l__spintent_luaset_exponent_tl", r_exp)
    token_set_macro("l__spintent_luaset_has_units_str", has_units)
    token_set_macro("l__spintent_luaset_status_str", unit_status)
    token_set_macro("l__spintent_luaset_arg_above_tl", above)
    token_set_macro("l__spintent_luaset_arg_below_tl", below)
    token_set_macro("l__spintent_luaset_denom_has_number_str", denom_has_numeric)
    token_set_macro("l__spintent_luaset_has_integer_str", has_integer)
    token_set_macro("l__spintent_luaset_has_natural_str", has_natural)

    if r_frac_clean ~= "" and s_match(r_frac_clean, "^0+$") then
        token_set_macro("l__spintent_luaset_dec_is_all_zeros_str", "true")
    else
        token_set_macro("l__spintent_luaset_dec_is_all_zeros_str", "false")
    end
end, { "string" })

-- 2. UNIDADES FÍSÍCAS Y CIENTÍFICAS (\spunit)

local spintent_units = {
    ["°K"]   = "°K",
    ["A"]    = "A",    ["B"]    = "B",     ["Bq"]   = "Bq",   ["C"]    = "C",
    ["Da"]   = "Da",   ["F"]    = "F",     ["Gy"]   = "Gy",   ["H"]    = "H",
    ["Hz"]   = "Hz",   ["J"]    = "J",     ["K"]    = "K",    ["l"]    = "l",
    ["L"]    = "L",    ["ℓ"]    = "ℓ",     ["N"]    = "N",    ["Np"]   = "Np",
    ["Pa"]   = "Pa",   ["S"]    = "S",     ["Sv"]   = "Sv",   ["T"]    = "T",
    ["V"]    = "V",    ["W"]    = "W",     ["Wb"]   = "Wb",   ["a"]    = "a",
    ["atm"]  = "atm",  ["au"]   = "au",    ["b"]    = "b",    ["bar"]  = "bar",
    ["cal"]  = "cal",  ["cd"]   = "cd",    ["d"]    = "d",    ["dB"]   = "dB",
    ["dyn"]  = "dyn",  ["eV"]   = "eV",    ["erg"]  = "erg",  ["g"]    = "g",
    ["h"]    = "h",    ["ha"]   = "ha",    ["kat"]  = "kat",  ["kg"]   = "kg",
    ["km"]   = "km",   ["lm"]   = "lm",    ["lx"]   = "lx",   ["m"]    = "m",
    ["cm"]   = "cm",   ["min"]  = "min",   ["mm"]   = "mm",   ["mol"]  = "mol",
    ["pc"]   = "pc",   ["rad"]  = "rad",   ["rpm"]  = "rpm",  ["s"]    = "s",
    ["sr"]   = "sr",   ["t"]    = "t",     ["u"]    = "u",    ["w"]    = "w",
    ["y"]    = "y",    ["°"]    = "°",     ["°C"]   = "°C",   ["°F"]   = "°F",
    ["Ω"]    = "Ω",    ["℧"]    = "℧",     ["′"]    = "′",    ["″"]    = "″",
    ["Å"]    = "Å",    ["µ"]    = "µ",
    ["in"]  = "in",    ["pulgada"]  = "in",   ["pulgadas"]  = "in",
    ["ft"]  = "ft",    ["pie"]      = "ft",   ["pies"]    = "ft",
    ["yd"]  = "yd",    ["yarda"]    = "yd",   ["yardas"]  = "yd",
    ["mi"]  = "mi",    ["milla"]    = "mi",   ["millas"]  = "mi",
    ["lb"]  = "lb",    ["libra"]    = "lb",   ["libras"]  = "lb",
    ["oz"]  = "oz",    ["onza"]     = "oz",   ["onzas"]   = "oz",
    ["gal"] = "gal",   ["galon"]    = "gal",  ["galones"] = "gal",
    ["hora"]     = "h",   ["horas"]      = "h",
    ["dia"]      = "d",   ["dias"]       = "d",
    ["caloria"]  = "cal", ["calorias"]   = "cal",
    ["faradio"]  = "F",   ["culombio"]   = "C",
    ["bit"]      = "b",   ["byte"]       = "B",
    ["bytes"]    = "B",   ["fahrenheit"] = "°F",
    ["kelvin"]   = "K",   ["Kelvin"]   = "K",
    ["ohmio"]    = "Ω",   ["ohm"]      = "Ω",   ["Omega"]    = "Ω",
    ["amperio"]  = "A",   ["ampere"]   = "A",
    ["voltio"]   = "V",   ["volt"]     = "V",
    ["vatio"]    = "W",   ["watt"]     = "W",
    ["julio"]    = "J",   ["joule"]    = "J",
    ["newton"]   = "N",
    ["pascal"]   = "Pa",
    ["hercio"]   = "Hz",  ["hertz"]    = "Hz",
    ["metro"]    = "m",   ["metros"]   = "m",
    ["segundo"]  = "s",   ["segundos"] = "s",
    ["litro"]    = "l",   ["l-litro"]  = "l",
    ["gramo"]    = "g",   ["gramos"]   = "g",
    ["angstrom"] = "Å",
    ["celsius"]  = "°C",  ["degC"]     = "°C",  ["℃"]        = "°C",
    ["degF"]     = "°F",  ["℉"]        = "°F",  ["micro"]    = "µ",
    ["minmin"]   = "′",   ["segseg"]   = "″",   ["u-masa"]   = "u",
    ["mho"]      = "℧",
    ["㎡"]   = "m^2",      ["㎢"]   = "km^2",     ["㎠"]   = "cm^2",     ["㎜²"]  = "mm^2",
    ["㎥"]   = "m^3",      ["㎦"]   = "km^3",     ["㎤"]   = "cm^3",     ["㎜³"]  = "mm^3",
    ["㎐"]   = "Hz",       ["㎑"]   = "kHz",      ["㎒"]   = "MHz",      ["㎓"]   = "GHz",      ["㎔"]   = "THz",
    ["㎽"]   = "mW",       ["㎾"]   = "kW",       ["㎿"]   = "MW",       ["㎴"]   = "V",        ["㎵"]   = "kV",
    ["㎶"]   = "µV",       ["㎷"]   = "mV",       ["㎸"]   = "kV",       ["㎀"]   = "pA",       ["㎁"]   = "nA",
    ["㎂"]   = "µA",       ["㎃"]   = "mA",       ["㎄"]   = "kA",       ["㎫"]   = "MPa",      ["㎬"]   = "GPa",
    ["㎺"]   = "pW",       ["㎻"]   = "nW",       ["㎼"]   = "µW",
    ["㎍"]   = "µg",       ["㎎"]   = "mg",       ["㎏"]   = "kg",
    ["㎜"]   = "mm",       ["㎝"]   = "cm",       ["㎞"]   = "km",       ["㎛"]   = "µm",       ["㎕"]   = "µL",
    ["㎖"]   = "mL",       ["㎗"]   = "dL",       ["㎘"]   = "kL",       ["¼"]   = "cm^3",     ["㎧"]   = "m/s",
    ["㎨"]   = "m/s^2",    ["㎭"]   = "rad",      ["㎮"]   = "rad/s",    ["㎯"]   = "rad/s^2"
}

local spintent_unit_compact_spoken_names = {
    ["㎡"]     = "metro-cuadrado",
    ["㎥"]     = "metro-cúbico",
    ["㎢"]     = "kilómetro-cuadrado",
    ["㎟"]     = "milímetro-cuadrado",
    ["㎠"]     = "centímetro-cuadrado",
    ["㎣"]     = "milímetro-cúbico",
    ["㎤"]     = "centímetro-cúbico",
    ["㎦"]     = "kilómetro-cúbico",
    ["¼"]      = "centímetro-cúbico",
    ["㎧"]     = "metros-por-segundo",
    ["㎨"]     = "metros-por-segundo-al-cuadrado",
    ["㎐"]     = "hercio",
    ["℃"]     = "grados-celsius",
    ["℉"]     = "grados-fahrenheit",
    ["'"]      = "minutos",
    ["''"]     = "segundos",
    ['"']      = "segundos",
    ["minmin"] = "minutos",
    ["segseg"] = "segundos",
    ["′"]      = "minutos",
    ["″"]      = "segundos",
}

local spintent_custom_spunit_spoken_names = {}

local spintent_normalizations = {
    ["K"] = "K", ["Ω"] = "Ω", ["ℓ"] = "\\ell", ["μ"] = "µ", ["°K"] = "°K",
    ["'"]  = "′", ["''"] = "″", ['"']  = "″", ["minmin"] = "′", ["segseg"] = "″"
}

local spintent_tex_accents = {
    ["\\'a"] = "á", ["\\'e"] = "é", ["\\'i"] = "í", ["\\'o"] = "ó", ["\\'u"] = "ú",
    ["\\'A"] = "Á", ["\\'E"] = "É", ["\\'I"] = "Í", ["\\'O"] = "Ó", ["\\'U"] = "Ú",
    ["\\\"u"] = "ü", ["\\\"U"] = "Ü", ["\\~n"] = "ñ", ["\\~N"] = "Ñ"
}

local function spintent_sanitize_for_screen_reader(raw_string)
    if not raw_string then return "" end
    local clean = s_gsub(raw_string, "[{}]", "")
    clean = s_gsub(clean, "\\[\'~\x22].", spintent_tex_accents)
    clean = s_gsub(clean, "%s+", "-")
    return clean
end

register_tex_cmd("luafun_define_custom_unit", function(unit_symbol, spoken_name)
    if spintent_units[unit_symbol] or spintent_custom_spunit_spoken_names[unit_symbol] then
        token_set_macro("l__spintent_spunit_luaset_status_str", "duplicate")
        return
    end

    -- Sanitizamos y aglutinamos para MathCAT
    local clean_spoken_name = spintent_sanitize_for_screen_reader(spoken_name)

    spintent_custom_spunit_spoken_names[unit_symbol] = clean_spoken_name
    spintent_custom_spunit_aliases[unit_symbol] = true

    token_set_macro("l__spintent_spunit_luaset_status_str", "success")
end, { "string", "string" })

register_tex_cmd("luafun_spunit_lookup_alias", function(raw_unit_name)
    raw_unit_name = s_gsub(raw_unit_name, "%s+", "")
    local clean_exp_format = s_gsub(raw_unit_name, "[{}]", "")
    token_set_macro("l__spintent_spunit_luaset_is_sexagesimal_str", "false")

    local search_name = spintent_normalizations[clean_exp_format] or spintent_normalizations[raw_unit_name] or raw_unit_name
    local canonical = spintent_units[search_name] or spintent_units[clean_exp_format]

    if canonical then
        local base_clean = s_gsub(canonical, "%^%-?%d+", "")
        local spoken = spintent_custom_spunit_spoken_names[base_clean]
            or spintent_unit_compact_spoken_names[raw_unit_name]
            or spintent_unit_compact_spoken_names[clean_exp_format]
            or spintent_unit_compact_spoken_names[search_name]
            or ":unit"

        token_set_macro("l__spintent_spunit_luaset_read_str", spoken)

        if s_match(canonical, "%^") then
            local base, exp = s_match(canonical, "([^%^]+)%^(%-?%d+)")
            if base and exp then
                token_set_macro("l__spintent_spunit_luaset_canonical_str", base)
                token_set_macro("l__spintent_spunit_luaset_compact_exp_str", exp)
                token_set_macro("l__spintent_spunit_luaset_is_compact_str", "true")
            else
                token_set_macro("l__spintent_spunit_luaset_canonical_str", canonical)
                token_set_macro("l__spintent_spunit_luaset_is_compact_str", "false")
            end
        else
            token_set_macro("l__spintent_spunit_luaset_canonical_str", canonical)
            token_set_macro("l__spintent_spunit_luaset_is_compact_str", "false")
        end

        if canonical == "°" or canonical == "′" or canonical == "″" then
            token_set_macro("l__spintent_spunit_luaset_is_sexagesimal_str", "true")
        else
            token_set_macro("l__spintent_spunit_luaset_is_sexagesimal_str", "false")
        end
        token_set_macro("l__spintent_spunit_luaset_lookup_status_str", "found")
    else
        local literal_base, literal_exp = s_match(clean_exp_format, "([^%^]+)%^(%-?%d+)")
        if literal_base and literal_exp then
            local canonical_base = spintent_units[spintent_normalizations[literal_base] or literal_base]
            if canonical_base then
                token_set_macro("l__spintent_spunit_luaset_canonical_str", canonical_base)
                token_set_macro("l__spintent_spunit_luaset_compact_exp_str", literal_exp)
                token_set_macro("l__spintent_spunit_luaset_is_compact_str", "true")
                local spoken = spintent_custom_spunit_spoken_names[canonical_base] or ":unit"
                token_set_macro("l__spintent_spunit_luaset_read_str", spoken)
                if canonical_base == "°" or canonical_base == "′" or canonical_base == "″" then
                    token_set_macro("l__spintent_spunit_luaset_is_sexagesimal_str", "true")
                end
                token_set_macro("l__spintent_spunit_luaset_lookup_status_str", "found")
            else
                token_set_macro("l__spintent_spunit_luaset_is_compact_str", "false")
                token_set_macro("l__spintent_spunit_luaset_lookup_status_str", "notfound")
            end
        else
            token_set_macro("l__spintent_spunit_luaset_is_compact_str", "false")
            token_set_macro("l__spintent_spunit_luaset_lookup_status_str", "notfound")
        end
    end
end, { "string" })

register_tex_cmd("luafun_define_spunit_alias", function(alias_name, unit_expression)
    alias_name = s_gsub(alias_name, "%s+", "")
    unit_expression = spintent_trim(unit_expression)
    if spintent_units[alias_name] or spintent_normalizations[alias_name] or spintent_custom_spunit_aliases[alias_name] then
        token_set_macro("l__spintent_spunit_luaset_status_str", "duplicate")
        return
    end
    local is_valid = true
    local _, slash_count = s_gsub(unit_expression, "/", "")
    if slash_count > 1 then is_valid = false end

    if s_match(unit_expression, "[^%a°ΩµÅℓ′″℧%*%/%^%-%d%s%'%\x22{}]") then is_valid = false end
    if is_valid then
        for unit_base in s_gmatch(unit_expression, "[%a°ΩµÅℓ′″℧%'%\x22]+") do
            if not (spintent_units[unit_base] or spintent_normalizations[unit_base]) then is_valid = false; break end
        end
    end
    if not is_valid then
        token_set_macro("l__spintent_spunit_luaset_status_str", "invalid-expr")
    else
        spintent_custom_spunit_aliases[alias_name] = unit_expression
        token_set_macro("l__spintent_spunit_luaset_status_str", "success")
    end
end, { "string", "string" })

-- 3. DIVISAS Y MONEDAS (\spmoney)

local spintent_internal_currencies = {
    ["$"]        = "$",         ["€"]        = "€",         ["£"]        = "£",         ["¥"]        = "¥",
    ["¢"]        = "¢",         ["₩"]        = "₩",         ["₪"]        = "₪",         ["₹"]        = "₹",
    ["₽"]        = "₽",         ["₺"]        = "₺",         ["₴"]        = "₴",
    ["dolar"]    = "$",         ["dolares"]  = "$",         ["usd"]      = "$",
    ["peso"]     = "$",         ["pesos"]    = "$",         ["clp"]      = "$",         ["mxn"]      = "$",
    ["euro"]     = "€",         ["euros"]    = "€",         ["eur"]      = "€",
    ["libra"]    = "£",         ["libras"]   = "£",         ["gbp"]      = "£",
    ["yen"]      = "¥",         ["yenes"]    = "¥",         ["yuan"]     = "¥",         ["yuanes"]   = "¥",
    ["jpy"]      = "¥",         ["cny"]      = "¥",         ["centavo"]  = "¢",         ["centavos"] = "¢",
    ["won"]      = "₩",         ["wons"]     = "₩",         ["krw"]      = "₩",
    ["shekel"]   = "₪",         ["shekels"]  = "₪",         ["sequel"]   = "₪",         ["sequeles"] = "₪",
    ["ils"]      = "₪",         ["rupia"]    = "₹",         ["rupias"]   = "₹",         ["inr"]      = "₹",
    ["rublo"]    = "₽",         ["rublos"]   = "₽",         ["rub"]      = "₽",
    ["lira"]     = "₺",         ["liras"]    = "₺",         ["try"]      = "₺",
    ["grivna"]   = "₴",         ["grivnas"]  = "₴",         ["uah"]      = "₴",
}

local spintent_currency_grammatical_dict = {
    ["$"]        = { sing = "peso",                  plur = "pesos",                  conde = "de-pesos" },
    ["peso"]     = { sing = "peso",                  plur = "pesos",                  conde = "de-pesos" },
    ["pesos"]    = { sing = "peso",                  plur = "pesos",                  conde = "de-pesos" },
    ["clp"]      = { sing = "peso-chileno",          plur = "pesos-chilenos",          conde = "de-pesos-chilenos" },
    ["mxn"]      = { sing = "peso-mexicano",         plur = "pesos-mexicanos",         conde = "de-pesos-mexicanos" },
    ["dolar"]    = { sing = "dólar",                 plur = "dólares",                conde = "de-dólares" },
    ["dolares"]  = { sing = "dólar",                 plur = "dólares",                conde = "de-dólares" },
    ["usd"]      = { sing = "dólar-estadounidense",  plur = "dólares-estadounidenses", conde = "de-dólares-estadounidenses" },
    ["euro"]     = { sing = "euro",                  plur = "euros",                  conde = "de-euros" },
    ["euros"]    = { sing = "euro",                  plur = "euros",                  conde = "de-euros" },
    ["eur"]      = { sing = "euro",                  plur = "euros",                  conde = "de-euros" },
    ["libra"]    = { sing = "libra",                 plur = "libras",                 conde = "de-libras" },
    ["libras"]   = { sing = "libra",                 plur = "libras",                 conde = "de-libras" },
    ["gbp"]      = { sing = "libra",                 plur = "libras",                 conde = "de-libras" },
    ["yen"]      = { sing = "yen",                   plur = "yenes",                  conde = "de-yenes" },
    ["yenes"]    = { sing = "yen",                   plur = "yenes",                  conde = "de-yenes" },
    ["jpy"]      = { sing = "yen",                   plur = "yenes",                  conde = "de-yenes" },
    ["yuan"]     = { sing = "yuan",                  plur = "yuanes",                 conde = "de-yuanes" },
    ["yuanes"]   = { sing = "yuan",                  plur = "yuanes",                 conde = "de-yuanes" },
    ["cny"]      = { sing = "yuan",                  plur = "yuanes",                 conde = "de-yuanes" },
    ["won"]      = { sing = "won",                   plur = "wons",                   conde = "de-wons" },
    ["krw"]      = { sing = "won",                   plur = "wons",                   conde = "de-wons" },
    ["shekel"]   = { sing = "séquel",                plur = "séqueles",               conde = "de-séqueles" },
    ["ils"]      = { sing = "séquel",                plur = "séqueles",               conde = "de-séqueles" },
    ["rupia"]    = { sing = "rupia",                 plur = "rupias",                 conde = "de-rupias" },
    ["inr"]      = { sing = "rupia",                 plur = "rupias",                 conde = "de-rupias" },
    ["rublo"]    = { sing = "rublo",                 plur = "rublos",                 conde = "de-rublos" },
    ["rub"]      = { sing = "rublo",                 plur = "rublos",                 conde = "de-rublos" },
    ["lira"]     = { sing = "lira",                  plur = "liras",                  conde = "de-liras" },
    ["try"]      = { sing = "lira",                  plur = "liras",                  conde = "de-liras" },
    ["grivna"]   = { sing = "grivna",                plur = "grivnas",                conde = "de-grivnas" },
    ["uah"]      = { sing = "grivna",                plur = "grivnas",                conde = "de-grivnas" },
    -- nuevos
    ["€"]        = { sing = "euro",                  plur = "euros",                  conde = "de-euros" },
    ["£"]        = { sing = "libra",                 plur = "libras",                 conde = "de-libras" },
    ["¥"]        = { sing = "yen",                   plur = "yenes",                  conde = "de-yenes" },
    ["₩"]        = { sing = "won",                   plur = "wons",                   conde = "de-wons" },
    ["₪"]        = { sing = "séquel",                plur = "séqueles",               conde = "de-séqueles" },
    ["₹"]        = { sing = "rupia",                 plur = "rupias",                 conde = "de-rupias" },
    ["₽"]        = { sing = "rublo",                 plur = "rublos",                 conde = "de-rublos" },
    ["₺"]        = { sing = "lira",                  plur = "liras",                  conde = "de-liras" },
    ["₴"]        = { sing = "grivna",                plur = "grivnas",                conde = "de-grivnas" },
}

local spintent_currency_subunits_matrix = {
    ["pesos"]                   = { sing = "centavo", plur = "centavos" },
    ["pesos-chilenos"]          = nil,
    ["pesos-colombianos"]       = nil,
    ["pesos-mexicanos"]         = { sing = "centavo", plur = "centavos" },
    ["dólares"]                 = { sing = "centavo", plur = "centavos" },
    ["dólares-estadounidenses"] = { sing = "centavo", plur = "centavos" },
    ["euros"]                   = { sing = "céntimo", plur = "céntimos" },
    ["libras"]                  = { sing = "penique", plur = "peniques" },
    ["yenes"]                   = nil,
    ["yuanes"]                  = { sing = "fen",     plur = "fen" },
    ["wons"]                    = nil,
    ["séqueles"]                = { sing = "ágora",   plur = "agorot" },
    ["rupias"]                  = { sing = "paisa",   plur = "paisas" },
    ["rublos"]                  = { sing = "kopek",   plur = "kopeks" },
    ["liras"]                   = { sing = "kuruş",   plur = "kuruş" },
    ["grivnas"]                 = { sing = "kopek",   plur = "kopeks" }
}

local spintent_currency_spoken_names = {
    ["$"]        = "pesos",                    ["peso"]    = "pesos",                    ["pesos"]   = "pesos",
    ["clp"]      = "pesos-chilenos",           ["mxn"]     = "pesos-mexicanos",
    ["dolar"]    = "dólares-estadounidenses",  ["dolares"] = "dólares-estadounidenses",  ["usd"]     = "dólares-estadounidenses",
    ["euro"]     = "euros",                    ["euros"]   = "euros",                    ["eur"]     = "euros",
    ["libra"]    = "libras",                   ["libras"]  = "libras",                   ["gbp"]     = "libras",
    ["yen"]      = "yenes",                    ["yenes"]   = "yenes",                    ["jpy"]     = "yenes",
    ["yuan"]     = "yuanes",                   ["yuanes"]  = "yuanes",                   ["cny"]     = "yuanes",
    ["centavo"]  = "centavos",                 ["centavos"]= "centavos",
    ["won"]      = "wons",                     ["krw"]     = "wons",
    ["shekel"]   = "séqueles",                 ["ils"]     = "séqueles",
    ["rupia"]    = "rupias",                   ["inr"]     = "rupias",
    ["rublo"]    = "rublos",                   ["rub"]     = "rublos",
    ["lira"]     = "liras",                    ["try"]     = "liras",
    ["grivna"]   = "grivnas",                  ["uah"]     = "grivnas",
    -- nuevos
    ["€"]        = "euros",
    ["£"]        = "libras",
    ["¥"]        = "yenes",
    ["¢"]        = "centavos",
    ["₩"]        = "wons",
    ["₪"]        = "séqueles",
    ["₹"]        = "rupias",
    ["₽"]        = "rublos",
    ["₺"]        = "liras",
    ["₴"]        = "grivnas",
}

register_tex_cmd("luafun_spmoney_lookup_data", function(currency_name)
    local trimmed = spintent_trim(currency_name)
    local clean = s_lower(trimmed)
    local resolved_symbol = spintent_internal_currencies[clean] or "$"

    local gram_entry = spintent_currency_grammatical_dict[clean] or spintent_currency_grammatical_dict[resolved_symbol]
      or { sing = "peso", plur = "pesos", conde = "de-pesos" }

    local position = "pre"
    if resolved_symbol == "€" then position = "post" end

    if trimmed == s_upper(trimmed) and s_match(trimmed, "^%a+$") then
        token_set_macro("l__spintent_spmoney_luaset_print_iso_str", "true")
    else
        token_set_macro("l__spintent_spmoney_luaset_print_iso_str", "false")
    end

    local sub_db = spintent_currency_subunits_matrix[gram_entry.plur]

    token_set_macro("l__spintent_spmoney_luaset_symbol_str", resolved_symbol)
    token_set_macro("l__spintent_spmoney_luaset_position_str", position)
    token_set_macro("l__spintent_spmoney_luaset_sing_str", gram_entry.sing)
    token_set_macro("l__spintent_spmoney_luaset_plur_str", gram_entry.plur)
    token_set_macro("l__spintent_spmoney_luaset_conde_str", gram_entry.conde)

    if sub_db then
        token_set_macro("l__spintent_spmoney_luaset_subsing_str", sub_db.sing)
        token_set_macro("l__spintent_spmoney_luaset_subplur_str", sub_db.plur)
    else
        token_set_macro("l__spintent_spmoney_luaset_subsing_str", "")
        token_set_macro("l__spintent_spmoney_luaset_subplur_str", "")
    end
end, { "string" })

register_tex_cmd("luafun_spmoney_normalize_key", function(raw_input)
    local clean = spintent_normalize_key(raw_input)
    local resolved_symbol = spintent_internal_currencies[clean]
    local resolved_spoken = nil

    if resolved_symbol then
        resolved_spoken = spintent_currency_spoken_names[clean] or spintent_currency_spoken_names[resolved_symbol]
    else
        resolved_spoken = spintent_currency_spoken_names[clean]
    end

    if resolved_spoken then
        token_set_macro("l__spintent_spmoney_luaset_currency_arg_str", resolved_spoken)
        token_set_macro("l__spintent_spmoney_luaset_status_str", "found")
    else
        token_set_macro("l__spintent_spmoney_luaset_currency_arg_str", raw_input)
        token_set_macro("l__spintent_spmoney_luaset_status_str", "notfound")
    end
end, { "string" })

register_tex_cmd("luafun_spmoney_prepare_input", function(raw_input)
    -- Extrae solo la moneda (borra dígitos, signos, espacios y separadores numéricos)
    local curr_part = s_gsub(raw_input, "[%d%s%+%-%.,{}%:%;]", "")

    -- Extrae solo el número matemático puro
    local num_part  = s_gsub(raw_input, "[^%d%s%+%-%.,{}%:%;]", "")

    token_set_macro("l__spintent_spmoney_extracted_curr_str", curr_part)
    token_set_macro("l__spintent_spmoney_extracted_num_tl", num_part)
end, { "string" })

-- 4. ABREVIATURAS Y ORDINALES (\spshort)

local spintent_spshort_dict = {
  ["dr.a"]   = { actualtext = "doctora",   layout_type = "superscript", base = "Dr.", suffix = "a" },
  ["dr.as"]  = { actualtext = "doctoras",  layout_type = "superscript", base = "Dr.", suffix = "as" },
  ["prof.a"] = { actualtext = "profesora", layout_type = "superscript", base = "Prof.", suffix = "a" },
  ["d.a"]    = { actualtext = "doña",      layout_type = "superscript", base = "D.", suffix = "a" },
  ["m.a"]    = { actualtext = "maría",     layout_type = "superscript", base = "M.", suffix = "a" },
  ["n.o"]    = { actualtext = "número",    layout_type = "superscript", base = "N.", suffix = "o" },
  ["n.os"]   = { actualtext = "números",   layout_type = "superscript", base = "N.", suffix = "os" },
  ["c.ia"]   = { actualtext = "compañía",  layout_type = "superscript", base = "C.", suffix = "ia" },

  ["dr.ª"] = { actualtext = "doctora", layout_type = "superscript", base = "Dr.", suffix = "a" },
  ["d.ª"]  = { actualtext = "doña",    layout_type = "superscript", base = "D.", suffix = "a" },
  ["n.º"]  = { actualtext = "número",  layout_type = "superscript", base = "N.", suffix = "o" },

  ["pág."] = { actualtext = "página",   layout_type = "linear_regular", output = "pág." },
  ["pag."] = { actualtext = "página",   layout_type = "linear_regular", output = "pág." },
  ["vol."] = { actualtext = "volumen",  layout_type = "linear_regular", output = "vol." },
  ["etc."] = { actualtext = "etcétera", layout_type = "linear_regular", output = "etc." },

  ["et al."]    = { actualtext = "y otros",      layout_type = "linear_regular", output = "et\u{00A0}al." },
  ["et. al."]   = { actualtext = "y otros",      layout_type = "linear_regular", output = "et\u{00A0}al." },
  ["ibíd."]     = { actualtext = "ibídem",       layout_type = "linear_regular", output = "ibíd." },
  ["ibid."]     = { actualtext = "ibídem",       layout_type = "linear_regular", output = "ibíd." },
  ["op. cit."]  = { actualtext = "obra citada",  layout_type = "linear_regular", output = "op.\u{00A0}cit." },
  ["op.cit."]   = { actualtext = "obra citada",  layout_type = "linear_regular", output = "op.\u{00A0}cit." },
  ["loc. cit."] = { actualtext = "lugar citado", layout_type = "linear_regular", output = "loc.\u{00A0}cit." },
  ["loc.cit."]  = { actualtext = "lugar citado", layout_type = "linear_regular", output = "loc.\u{00A0}cit." },
  ["v. gr."]    = { actualtext = "verbigracia",  layout_type = "linear_regular", output = "v.\u{00A0}gr." },
  ["v.gr."]     = { actualtext = "verbigracia",  layout_type = "linear_regular", output = "v.\u{00A0}gr." },
  ["i. e."]     = { actualtext = "esto es",      layout_type = "linear_regular", output = "i.\u{00A0}e." },
  ["i.e."]      = { actualtext = "esto es",      layout_type = "linear_regular", output = "i.\u{00A0}e." },
  ["e. g."]     = { actualtext = "por ejemplo",  layout_type = "linear_regular", output = "e.\u{00A0}g." },
  ["e.g."]      = { actualtext = "por ejemplo",  layout_type = "linear_regular", output = "e.\u{00A0}g." },
  ["p. ej."]    = { actualtext = "por ejemplo",  layout_type = "linear_regular", output = "p.\u{00A0}ej." },
  ["p.ej."]     = { actualtext = "por ejemplo",  layout_type = "linear_regular", output = "p.\u{00A0}ej." },

  ["sr."]    = { actualtext = "señor",       layout_type = "linear_regular", output = "Sr." },
  ["sra."]   = { actualtext = "señora",      layout_type = "linear_regular", output = "Sra." },
  ["srta."]  = { actualtext = "señorita",    layout_type = "linear_regular", output = "Srta." },
  ["dr."]    = { actualtext = "doctor",      layout_type = "linear_regular", output = "Dr." },
  ["dra."]   = { actualtext = "doctora",     layout_type = "linear_regular", output = "Dra." },
  ["dres."]  = { actualtext = "doctores",    layout_type = "linear_regular", output = "Dres." },
  ["dras."]  = { actualtext = "doctoras",    layout_type = "linear_regular", output = "Dras." },
  ["prof."]  = { actualtext = "profesor",    layout_type = "linear_regular", output = "Prof." },
  ["profa."] = { actualtext = "profesora",   layout_type = "linear_regular", output = "Profa." },
  ["profs."] = { actualtext = "profesores",  layout_type = "linear_regular", output = "Profs." },
  ["ing."]   = { actualtext = "ingeniero",   layout_type = "linear_regular", output = "Ing." },
  ["ings."]  = { actualtext = "ingenieros",  layout_type = "linear_regular", output = "Ings." },
  ["lic."]   = { actualtext = "licenciado",  layout_type = "linear_regular", output = "Lic." },
  ["v. b."]  = { actualtext = "visto bueno", layout_type = "linear_regular", output = "V.\u{00A0}B." },
  ["v.b."]   = { actualtext = "visto bueno", layout_type = "linear_regular", output = "V.\u{00A0}B." },

  ["s. a."]    = { actualtext = "sociedad anónima",         layout_type = "linear_caps",    output = "S.\u{00A0}A." },
  ["s.a."]     = { actualtext = "sociedad anónima",         layout_type = "linear_caps",    output = "S.\u{00A0}A." },
  ["ee. uu."]  = { actualtext = "estados unidos",          layout_type = "linear_caps",    output = "EE.\u{00A0}UU." },
  ["ee.uu."]   = { actualtext = "estados unidos",           layout_type = "linear_caps",    output = "EE.\u{00A0}UU." },
  ["dd. hh."]  = { actualtext = "derechos humanos",        layout_type = "linear_caps",    output = "DD.\u{00A0}HH." },
  ["dd.hh."]   = { actualtext = "derechos humanos",        layout_type = "linear_caps",    output = "DD.\u{00A0}HH." },
  ["jj. oo."]  = { actualtext = "juegos olímpicos",        layout_type = "linear_caps",    output = "JJ.\u{00A0}OO." },
  ["jj.oo."]   = { actualtext = "juegos olímpicos",        layout_type = "linear_caps",    output = "JJ.\u{00A0}OO." },
  ["ff. aa."]  = { actualtext = "fuerzas armadas",         layout_type = "linear_caps",    output = "FF.\u{00A0}AA." },
  ["ff.aa."]   = { actualtext = "fuerzas armadas",         layout_type = "linear_caps",    output = "FF.\u{00A0}AA." },
  ["rr. ee."]  = { actualtext = "relaciones exteriores",   layout_type = "linear_caps",    output = "RR.\u{00A0}EE." },
  ["rr.ee."]   = { actualtext = "relaciones exteriores",   layout_type = "linear_caps",    output = "RR.\u{00A0}EE." },
  ["cc. aa."]  = { actualtext = "comunidades autónomas",   layout_type = "linear_caps",    output = "CC.\u{00A0}AA." },
  ["cc.aa."]   = { actualtext = "comunidades autónomas",   layout_type = "linear_caps",    output = "CC.\u{00A0}AA." },
  ["p. d."]    = { actualtext = "posdata",                 layout_type = "linear_caps",    output = "P.\u{00A0}D." },
  ["p.d."]     = { actualtext = "posdata",                 layout_type = "linear_caps",    output = "P.\u{00A0}D." },

-- Solo por compatibilidad (no deben llevar punto acorde a la RAE)
  ["d. n. i."] = { actualtext = "documento nacional de identidad", layout_type = "linear_caps", output = "D.\u{00A0}N.\u{00A0}I." },
  ["d.n.i."]   = { actualtext = "documento nacional de identidad", layout_type = "linear_caps", output = "D.\u{00A0}N.\u{00A0}I." },
  ["r. u. t."] = { actualtext = "rol único tributario", layout_type = "linear_caps", output = "R.\u{00A0}U.\u{00A0}T." },
  ["r.u.t."]   = { actualtext = "rol único tributario", layout_type = "linear_caps", output = "R.\u{00A0}U.\u{00A0}T." },
  ["r. u. n."] = { actualtext = "rol único nacional", layout_type = "linear_caps", output = "R.\u{00A0}U.\u{00A0}N." },
  ["r.u.n."]   = { actualtext = "rol único nacional", layout_type = "linear_caps", output = "R.\u{00A0}U.\u{00A0}N." },

  ["a. c."]    = { actualtext = "antes de Cristo",          layout_type = "linear_mixed",   output = "a.\u{00A0}C." },
  ["a.c."]     = { actualtext = "antes de Cristo",          layout_type = "linear_mixed",   output = "a.\u{00A0}C." },
  ["d. c."]    = { actualtext = "después de Cristo",        layout_type = "linear_mixed",   output = "d.\u{00A0}C." },
  ["d.c."]     = { actualtext = "después de Cristo",        layout_type = "linear_mixed",   output = "d.\u{00A0}C." },

  ["dni"]      = { actualtext = "documento nacional de identidad",     layout_type = "small_caps_pure", output = "DNI" },
  ["run"]      = { actualtext = "rol único nacional",                  layout_type = "small_caps_pure", output = "RUN" },
  ["rut"]      = { actualtext = "rol único tributario",                layout_type = "small_caps_pure", output = "RUT" },
  ["onu"]      = { actualtext = "organización de las naciones unidas", layout_type = "small_caps_pure", output = "ONU" },
  ["rae"]      = { actualtext = "real academia española",             layout_type = "small_caps_pure", output = "RAE" },
  ["ong"]      = { actualtext = "organización no gubernamental",      layout_type = "small_caps_pure", output = "ONG" },
  ["ongs"]     = { actualtext = "organizaciones no gubernamentales",  layout_type = "small_caps_pure", output = "ONGs" },
  ["urss"]     = { actualtext = "unión de repúblicas socialistas soviéticas", layout_type = "small_caps_pure", output = "URSS" },
  ["oea"]      = { actualtext = "organización de los estados americanos", layout_type = "small_caps_pure", output = "OEA" },
  ["oms"]      = { actualtext = "organización mundial de la salud",    layout_type = "small_caps_pure", output = "OMS" },
  ["fmi"]      = { actualtext = "fondo monetario internacional",      layout_type = "small_caps_pure", output = "FMI" },
  ["bid"]      = { actualtext = "banco interamericano de desarrollo", layout_type = "small_caps_pure", output = "BID" },
  ["otan"]     = { actualtext = "organización del tratado del atlántico norte", layout_type = "small_caps_pure", output = "OTAN" },
  ["unam"]     = { actualtext = "universidad nacional autónoma de méxico", layout_type = "small_caps_pure", output = "UNAM" },
  ["pib"]      = { actualtext = "producto interno bruto",             layout_type = "small_caps_pure", output = "PIB" },
  ["ue"]       = { actualtext = "unión europea",                       layout_type = "small_caps_pure", output = "UE" },
  ["eau"]      = { actualtext = "emiratos árabes unidos",              layout_type = "small_caps_pure", output = "EAU" },
  ["ru"]       = { actualtext = "reino unido",                         layout_type = "small_caps_pure", output = "RU" },
  ["rca"]      = { actualtext = "república centroafricana",            layout_type = "small_caps_pure", output = "RCA" },
  ["rdc"]      = { actualtext = "república democrática del congo",     layout_type = "small_caps_pure", output = "RDC" },
  ["rfa"]      = { actualtext = "república federal alemana",           layout_type = "small_caps_pure", output = "RFA" },
  ["rda"]      = { actualtext = "república democrática alemana",       layout_type = "small_caps_pure", output = "RDA" },

  ["unicef"]   = { actualtext = "fondo de las naciones unidas para la infancia", layout_type = "acronym_long",   output = "Unicef" },
  ["unesco"]   = { actualtext = "organización de las naciones unidas para la educación, la ciencia y la cultura", layout_type = "acronym_long",   output = "Unesco" },
  ["mercosur"] = { actualtext = "mercado común del sur",               layout_type = "acronym_long",   output = "Mercosur" },
  ["cepal"]    = { actualtext = "comisión económica para américa latina", layout_type = "acronym_long",  output = "Cepal" },
  ["celac"]    = { actualtext = "comunidad de estados latinoamericanos", layout_type = "acronym_long",  output = "Celac" },
  ["unasur"]   = { actualtext = "unión de naciones suramericanas",     layout_type = "acronym_long",   output = "Unasur" },
  ["acnur"]    = { actualtext = "alto comisionado de las naciones unidas", layout_type = "acronym_long",  output = "Acnur" },
  ["mineduc"]  = { actualtext = "ministerio de educación",            layout_type = "acronym_long",   output = "Mineduc" },
  ["senadis"]  = { actualtext = "servicio nacional de la discapacidad", layout_type = "acronym_long",  output = "Senadis" },
  ["conicet"]  = { actualtext = "consejo nacional de investigaciones", layout_type = "acronym_long",   output = "Conicet" }
}

local spintent_spshort_ord_suffixes = {
  ["a"]  = "ª",
  ["o"]  = "º",
  ["er"] = "er",
  ["os"] = "os",
  ["as"] = "as"
}

local spintent_ord_units = {
  [1] = "primer", [2] = "segund", [3] = "tercer", [4] = "cuart", [5] = "quint",
  [6] = "sext",   [7] = "séptim", [8] = "octav",  [9] = "noven"
}

local spintent_ord_tens = {
  [1] = "décim",       [2] = "vigésim",     [3] = "trigésim",    [4] = "cuadragésim",
  [5] = "quincuagésim",[6] = "sexagésim",   [7] = "septuagésim", [8] = "octogésim",
  [9] = "nonagésim"
}

local function spintent_get_semantic_ordinal(val, suffix)
  if val < 1 or val > 99 then
    return val .. (spintent_spshort_ord_suffixes[suffix] or suffix)
  end

  local end_char = (suffix == "a" or suffix == "as") and "a" or "o"
  local plural = (suffix == "as" or suffix == "os") and "s" or ""

  if suffix == "er" then
    end_char = "o"
    plural = ""
  end

  local ten_val = math_floor(val / 10)
  local unit_val = val % 10

  if val == 11 and suffix ~= "er" then return "undécim" .. end_char .. plural end
  if val == 12 then return "duodécim" .. end_char .. plural end

  local result = ""
  if ten_val > 0 then
    result = spintent_ord_tens[ten_val] .. end_char .. plural
  end

  if unit_val > 0 then
    local unit_str = spintent_ord_units[unit_val]

    if unit_val == 1 and suffix ~= "er" then unit_str = "primer" end
    if unit_val == 3 and suffix ~= "er" then unit_str = "tercer" end

    local unit_end = (suffix == "er") and "" or (end_char .. plural)
    local full_unit = unit_str .. unit_end

    if ten_val > 0 then
      result = result .. " " .. full_unit
    else
      result = full_unit
    end
  end

  return result
end

local spintent_spshort_digits       = spintent_digit^1
local spintent_spshort_dot          = P(".")
local spintent_spshort_raw_suffix   = C(P("er") + P("os") + P("as") + P("a") + P("o"))
local spintent_spshort_illegal_suff = C((R("az") + R("AZ"))^1)

local spintent_spshort_ord_grammar = P({
  "Entry",
  Entry  = #spintent_digit * spintent_spshort_V("Main"),
  Main   = C(spintent_spshort_digits) * spintent_spshort_dot * (spintent_spshort_raw_suffix + spintent_spshort_illegal_suff * Cg(Cc(true)))
})

local function spintent_spshort_execute_analysis(raw_input)
  local clean_input = spintent_normalize_key(raw_input)

  local dict_match = spintent_spshort_dict[clean_input]
  if not dict_match then
    if s_sub(clean_input, -1) == "." then
      dict_match = spintent_spshort_dict[s_sub(clean_input, 1, -2)]
    else
      dict_match = spintent_spshort_dict[clean_input .. "."]
    end
  end

  if dict_match then
    token_set_macro("l__spintent_spshort_luaset_status_str", "success")
    token_set_macro("l__spintent_spshort_luaset_layout_str", dict_match.layout_type)
    token_set_macro("l__spintent_spshort_luaset_expanded_str", dict_match.actualtext)

    if dict_match.layout_type == "superscript" then
      token_set_macro("l__spintent_spshort_luaset_base_str", dict_match.base)
      token_set_macro("l__spintent_spshort_luaset_suffix_str", dict_match.suffix)
    else
      token_set_macro("l__spintent_spshort_luaset_output_str", dict_match.output)
    end
    return
  end

  local ordinal_input = clean_input
  if s_sub(ordinal_input, -1) == "." and not s_match(ordinal_input, "%d+%.%a+$") and not s_match(ordinal_input, "^%d+%.") then
    ordinal_input = s_sub(ordinal_input, 1, -2)
  end

  local num_part, suffix_part, is_illegal = spintent_spshort_ord_grammar:match(ordinal_input)

  if num_part and not is_illegal and spintent_spshort_ord_suffixes[suffix_part] then

    if suffix_part == "er" then
      local val = tonumber(num_part)
      local last_digit = val % 10
      if last_digit ~= 1 and last_digit ~= 3 then
        token_set_macro("l__spintent_spshort_luaset_status_str", "error")
        return
      end
    end

    local semantic_read = spintent_get_semantic_ordinal(tonumber(num_part), suffix_part)

    token_set_macro("l__spintent_spshort_luaset_status_str", "success")
    token_set_macro("l__spintent_spshort_luaset_layout_str", "superscript")
    token_set_macro("l__spintent_spshort_luaset_expanded_str", semantic_read)
    token_set_macro("l__spintent_spshort_luaset_base_str", num_part .. ".")
    token_set_macro("l__spintent_spshort_luaset_suffix_str", spintent_spshort_ord_suffixes[suffix_part])
    return
  end

  if s_match(raw_input, "^%a+$") then
    token_set_macro("l__spintent_spshort_luaset_status_str", "fallback")
    token_set_macro("l__spintent_spshort_luaset_layout_str", "none")
    token_set_macro("l__spintent_spshort_luaset_expanded_str", raw_input)
    token_set_macro("l__spintent_spshort_luaset_output_str", raw_input)
  else
    token_set_macro("l__spintent_spshort_luaset_status_str", "error")
  end
end

register_tex_cmd("luafun_spshort_process", function(raw_input)
  spintent_spshort_execute_analysis(raw_input)
end, { "string" })

-- 5. SÍGLOS Y NÚMEROS ROMANOS (\spsiglo)

local spintent_arab_to_roman_map = {
  {1000, "m"}, {900, "cm"}, {500, "d"}, {400, "cd"},
  {100, "c"}, {90, "xc"}, {50, "l"}, {40, "xl"},
  {10, "x"}, {9, "ix"}, {5, "v"}, {4, "iv"}, {1, "i"}
}

local spintent_roman_to_arab_byte_map = {
  [105] = 1,   -- 'i'
  [118] = 5,   -- 'v'
  [120] = 10,  -- 'x'
  [108] = 50,  -- 'l'
  [99]  = 100, -- 'c'
  [100] = 500, -- 'd'
  [109] = 1000 -- 'm'
}

local function spintent_arabic_to_roman(num)
  local result = {}
  for i = 1, #spintent_arab_to_roman_map do
    local pair = spintent_arab_to_roman_map[i]
    while num >= pair[1] do
      result[#result + 1] = pair[2]
      num = num - pair[1]
    end
  end
  return t_concat(result)
end

local function spintent_roman_to_arabic(str_roman)
  local total = 0
  local i = 1
  local len = #str_roman

  while i <= len do
    local v1 = spintent_roman_to_arab_byte_map[s_byte(str_roman, i)] or 0

    if i + 1 <= len then
      local v2 = spintent_roman_to_arab_byte_map[s_byte(str_roman, i + 1)] or 0

      if v1 < v2 then
        total = total + (v2 - v1)
        i = i + 2
      else
        total = total + v1
        i = i + 1
      end
    else
      total = total + v1
      i = i + 1
    end
  end
  return total
end

register_tex_cmd("luafun_spsiglo_parse", function(raw_siglo_input)
  local clean = s_lower(s_gsub(raw_siglo_input, "%s+", ""))
  local arabic_val = nil
  local roman_val = nil
  local is_error = false

  if s_match(clean, "^%d+$") then
    arabic_val = tonumber(clean)
    if arabic_val > 0 and arabic_val <= 4000 then
      roman_val = spintent_arabic_to_roman(arabic_val)
    else
      is_error = true
    end
  elseif s_match(clean, "^[ivxlcdm]+$") and clean ~= "" then
    roman_val = clean
    arabic_val = spintent_roman_to_arabic(clean)

    if spintent_arabic_to_roman(arabic_val) ~= clean then
      is_error = true
    end
  else
    is_error = true
  end

  if is_error then
    token_set_macro("l__spintent_spsiglo_luaset_error_str", "true")
    token_set_macro("l__spintent_spsiglo_luaset_arabic_str", "")
    token_set_macro("l__spintent_spsiglo_luaset_roman_min_str", "")
  else
    token_set_macro("l__spintent_spsiglo_luaset_error_str", "false")
    token_set_macro("l__spintent_spsiglo_luaset_arabic_str", tostring(arabic_val))
    token_set_macro("l__spintent_spsiglo_luaset_roman_min_str", roman_val)
  end
end, { "string" })

-- 6. FECHA Y HORA (\spdate, \sptime)

local spintent_date_sep = C(S"/-")

-- Patrón neutro: 3 partes numéricas divididas por 2 separadores
local spintent_date_pattern = Ct(
    Cg(spintent_num_chunk, "p1") * Cg(spintent_date_sep, "sep1") * Cg(spintent_num_chunk, "p2") * Cg(spintent_date_sep, "sep2") * Cg(spintent_num_chunk, "p3") * P(-1)
)

-- Función auxiliar: Validar que el día y mes existan (¡incluyendo bisiestos!)
local days_in_month = {31, 28, 31, 30, 31, 30, 31, 31, 30, 31, 30, 31}

local function spintent_is_valid_date(d, m, y)
    if m < 1 or m > 12 then return false end
    if d < 1 then return false end

    -- Ajuste para febrero en años bisiestos
    if m == 2 then
        local is_leap = (y % 4 == 0 and y % 100 ~= 0) or (y % 400 == 0)
        if is_leap then
            return d <= 29
        end
    end
    return d <= days_in_month[m]
end

register_tex_cmd("luafun_spdate_parse", function(raw_date_input)
    if not raw_date_input then return end
    local clean_str = s_gsub(raw_date_input, "%s+", "")
    local result = spintent_date_pattern:match(clean_str)

    -- Si no encaja el patrón, o si el usuario mezcla separadores (ej. 2024-01/02) -> error
    if not result or result.sep1 ~= result.sep2 then
        token_set_macro("l__spintent_spdate_luaset_error_str", "true")
        return
    end

    local p1, p2, p3 = result.p1, result.p2, result.p3
    local sep = result.sep1
    local year, month, day

    -- Como spintent_num_chunk elimina espacios, las longitudes son exactas
    local p1_len = #p1
    local p3_len = #p3

    -- Lógica 1: Detección del Año
    if p1_len == 4 and p3_len <= 2 then
        year, month, day = p1, p2, p3
    elseif p3_len == 4 and p1_len <= 2 then
        day, month, year = p1, p2, p3
    else
        -- Si ninguno tiene 4 dígitos (ej. 12-10-81) o ambos lo tienen, hay ambigüedad
        token_set_macro("l__spintent_spdate_luaset_error_str", "true")
        return
    end

    local num_y, num_m, num_d = tonumber(year), tonumber(month), tonumber(day)

    -- Lógica 2: Validación de calendario real
    if not (num_y and num_m and num_d) or not spintent_is_valid_date(num_d, num_m, num_y) then
        token_set_macro("l__spintent_spdate_luaset_error_str", "true")
        return
    end

    -- Lógica 3: Reglas de renderizado visual
    local output_str
    if p1_len == 4 and sep == "/" then
        -- Único caso a corregir: YYYY/MM/DD se invierte a DD/MM/YYYY
        output_str = day .. "/" .. month .. "/" .. year
    else
        -- Para el resto, se devuelve EXACTAMENTE el mismo formato que ingresó
        output_str = p1 .. sep .. p2 .. sep .. p3
    end

    -- Si sobrevivió a todo, pasamos las variables limpias a TeX
    token_set_macro("l__spintent_spdate_luaset_error_str", "false")
    token_set_macro("l__spintent_spdate_luaset_day_str", day)
    token_set_macro("l__spintent_spdate_luaset_month_str", month)
    token_set_macro("l__spintent_spdate_luaset_year_str", year)
    token_set_macro("l__spintent_spdate_luaset_output_str", output_str)
end, { "string" })

local spintent_time_pattern = Ct(
    Cg(spintent_num_chunk, "h") * P":" * Cg(spintent_num_chunk, "m") * (P":" * Cg(spintent_num_chunk, "s"))^-1 * P(-1)
)

register_tex_cmd("luafun_sptime_parse", function(raw_time_input)
    if not raw_time_input then return end
    local clean_str = s_gsub(raw_time_input, "%s+", "")
    local result = spintent_time_pattern:match(clean_str)

    if not result then
        token_set_macro("l__spintent_sptime_luaset_error_str", "true")
        return
    end

    local h = tonumber(result.h)
    local m = tonumber(result.m)
    local s = result.s and tonumber(result.s) or nil

    if not h or not m or h >= 24 or m >= 60 or (s and s >= 60) then
        token_set_macro("l__spintent_sptime_luaset_error_str", "true")
        return
    end

    token_set_macro("l__spintent_sptime_luaset_error_str", "false")
    token_set_macro("l__spintent_sptime_luaset_base_str", result.h .. ":" .. result.m)

    -- Solo es "fracción especial" si la hora es < 13 y los minutos son 15 o 30
    if h < 13 and (m == 15 or m == 30) then
        token_set_macro("l__spintent_sptime_luaset_is_fraction_str", "true")
    else
        token_set_macro("l__spintent_sptime_luaset_is_fraction_str", "false")
    end

    if result.s then
        token_set_macro("l__spintent_sptime_luaset_has_seconds_str", "true")
        token_set_macro("l__spintent_sptime_luaset_seconds_str", result.s)
    else
        token_set_macro("l__spintent_sptime_luaset_has_seconds_str", "false")
        token_set_macro("l__spintent_sptime_luaset_seconds_str", "")
    end
end, { "string" })

-- 7. MATEMÁTICAS: MCM Y MCD (\spmcm, \spmcd)

local function spintent_gcd_algorithm(val_a, val_b)
    while val_b ~= 0 do val_a, val_b = val_b, val_a % val_b end
    return val_a
end

local function spintent_lcm_algorithm(val_a, val_b)
    if val_a == 0 or val_b == 0 then return 0 end
    return math_floor((val_a * val_b) / spintent_gcd_algorithm(val_a, val_b))
end

local function spintent_execute_mcm_mcd_result(raw_csv_list, tl_out, operation_fn)
    local numbers = {}
    token_set_macro("l__spintent_spmcm_spmcd_luaset_error_str", "false")

    for item in s_gmatch(raw_csv_list, "([^,]+)") do
        local clean_item = spintent_trim(item)
        local result = spintent_number_pattern:match(clean_item) or {}

        local r_int  = result.integer
        local r_sign = result.sign
        local r_dec  = result.decimal
        local r_per  = result.period
        local r_ext  = result.extra

        local es_natural = r_int and (not r_sign or r_sign == "")
          and (not r_dec or r_dec == "") and (not r_per or r_per == "")
          and (not r_ext or s_gsub(r_ext, "%s+", "") == "")

        if es_natural then
            numbers[#numbers + 1] = tonumber(r_int)
        else
            token_set_macro("l__spintent_spmcm_spmcd_luaset_error_str", "true")
            return
        end
    end
    if #numbers == 0 then
        token_set_macro("l__spintent_spmcm_spmcd_luaset_error_str", "true")
        return
    end
    local final_result = numbers[1]
    for i = 2, #numbers do final_result = operation_fn(final_result, numbers[i]) end
    token_set_macro(tl_out, s_format("%d", final_result))
end

register_tex_cmd("luafun_calculate_mcd", function(raw_csv_list)
    spintent_execute_mcm_mcd_result(raw_csv_list, "l__spintent_spmcd_luaset_mcd_value_tl", spintent_gcd_algorithm)
end, { "string" })

register_tex_cmd("luafun_calculate_mcm", function(raw_csv_list)
    spintent_execute_mcm_mcd_result(raw_csv_list, "l__spintent_spmcm_luaset_mcm_value_tl", spintent_lcm_algorithm)
end, { "string" })

-- 8. SÍSTEMA SEXAGESIMAL (\spang)

local spintent_angle_sexag_pattern = Ct(
    Cg(spintent_num_chunk^-1, "a") * P ":" * Cg(spintent_num_chunk^-1, "b") * (P ":" * Cg(spintent_num_chunk^-1, "c"))^-1 * P(-1)
)

register_tex_cmd("luafun_spang_parse", function(raw_sexag_str)
    raw_sexag_str = s_gsub(raw_sexag_str, "%s+", "")
    local result = spintent_angle_sexag_pattern:match(raw_sexag_str)

    if not result then
        token_set_macro("l__spintent_spang_luaset_error_str", "true")
        return
    end

    token_set_macro("l__spintent_spang_luaset_error_str", "false")
    token_set_macro("l__spintent_spang_luaset_grado_str", result.a or "")
    token_set_macro("l__spintent_spang_luaset_minuto_str", result.b or "")
    token_set_macro("l__spintent_spang_luaset_segundo_str", result.c or "")
end, { "string" })

-- 9. MATEMÁTICAS: MÚLTIPLOS Y DIVISORES (\spnM, \spnD)

local function spintent_calculate_multiples(n_val, count)
    local results = {}
    for i = 1, count do
        results[#results + 1] = tostring(n_val * i)
    end
    return t_concat(results, ", ")
end

local function spintent_calculate_divisors(n_val, limit)
    local divisors = {}
    -- Algoritmo optimizado de raíz cuadrada para divisores rápidos
    for i = 1, math_floor(math_sqrt(n_val)) do
        if n_val % i == 0 then
            divisors[#divisors + 1] = i
            -- Si el cociente es diferente, agregarlo también
            local quotient = n_val // i
            if i ~= quotient then
                divisors[#divisors + 1] = quotient
            end
        end
    end
    -- Ordenar los divisores de menor a mayor
    t_sort(divisors)

    local results = {}
    local actual_count = #divisors
    local count_to_show = (limit > 0) and math_min(limit, actual_count) or actual_count

    for i = 1, count_to_show do
        results[#results + 1] = tostring(divisors[i])
    end

    local is_truncated = (count_to_show < actual_count) and "true" or "false"
    return t_concat(results, ", "), is_truncated
end

local function spintent_execute_spnM_spnD_result(raw_n, raw_limit, is_multiple)
    -- Asumimos que expl3 ya usó luafun_clean_split_and_set y garantizó que raw_n
    -- es un número natural puro. ¡Solo convertimos y calculamos!
    local num_n = tonumber(raw_n)

    local clean_limit = spintent_normalize_key(raw_limit)
    local num_limit = tonumber(clean_limit) or 0
    if clean_limit == "all" or clean_limit == "" then
        num_limit = 0
    end

    if is_multiple then
        if num_limit <= 0 then num_limit = 2 end -- 2 múltiplos por defecto
        local result_str = spintent_calculate_multiples(num_n, num_limit)

        token_set_macro("l__spintent_spnM_luaset_result_tl", result_str)
        token_set_macro("l__spintent_spnM_luaset_is_truncated_str", "true") -- Siempre infinitos
    else
        local result_str, is_truncated = spintent_calculate_divisors(num_n, num_limit)

        token_set_macro("l__spintent_spnD_luaset_result_tl", result_str)
        token_set_macro("l__spintent_spnD_luaset_is_truncated_str", is_truncated)
    end
end

register_tex_cmd("luafun_calcular_nM", function(raw_n, raw_limit)
    spintent_execute_spnM_spnD_result(raw_n, raw_limit, true)
end, { "string", "string" })

register_tex_cmd("luafun_calcular_nD", function(raw_n, raw_limit)
    spintent_execute_spnM_spnD_result(raw_n, raw_limit, false)
end, { "string", "string" })

-- 10. NÚMEROS MIXTOS, code for \spmQ (see https://tex.stackexchange.com/q/764828)

local spintent_font_cache         = {}
local spintent_measuring = false

-- Tabla constante: nunca cambia entre llamadas, se construye una sola vez
-- en vez de en cada invocación de \spmQ.
local spintent_mathstyle_names = {
    [0] = "display",      [1] = "display",
    [2] = "text",         [3] = "text",
    [4] = "script",       [5] = "script",
    [6] = "scriptscript", [7] = "scriptscript",
}

local function spintent_is_digit(n)
    return n.char >= 48 and n.char <= 57
end

local function spintent_get_scaled_font(id, factor)
    local key = tostring(id) .. ":" .. tostring(factor)
    if spintent_font_cache[key] then return spintent_font_cache[key] end
    local old_font = font.getfont(id)
    if not old_font then return id end
    local spec = {}
    for k, v in pairs(old_font.specification) do spec[k] = v end
    spec.size = math_floor(old_font.size * factor)
    local ok, data = pcall(fonts.definers.loadfont, spec)
    if not ok or not data then return id end
    local new_id = font.define(data)
    fonts.hashes.identifiers[new_id] = data
    spintent_font_cache[key] = new_id
    return new_id
end

local function spintent_scale_glyphs(head, factor)
    for n in node.traverse_glyph(head) do
        if spintent_is_digit(n) then
            local new_id = spintent_get_scaled_font(n.font, factor)
            n.font = new_id
            local f = font.getfont(new_id)
            if f and f.characters then
                local c = f.characters[n.char]
                if c then
                    n.width  = c.width
                    n.height = c.height
                    n.depth  = c.depth
                end
            end
        end
    end
    return head
end

local function spintent_build_char_noads(str)
    local head, tail
    for i = 1, #str do
        local kernel = node.new("math_char")
        kernel.char  = str:byte(i)
        kernel.fam   = 0
        local n      = node.new("noad")
        n.subtype    = 0
        n.nucleus    = kernel
        if not head then head = n tail = n
        else tail.next = n tail = n end
    end
    return head
end

local function spintent_build_submlist_noad(str)
    local mlist  = node.new("sub_mlist")
    mlist.head   = spintent_build_char_noads(str)
    local noad   = node.new("noad")
    noad.subtype = 0
    noad.nucleus = mlist
    return noad
end

-- Movida a nivel de módulo: no captura nada específico de una llamada a
-- \spmQ (dt/cmd son parámetros explícitos, no upvalues), así que crear
-- un closure nuevo por invocación era trabajo evitable.
local function spintent_medir_frac(dt, cmd)
    spintent_measuring = true
    tex.runtoks(function()
        tex.sprint(
            "\\setbox0\\hbox{$\\" .. dt .. "style" ..
            "\\" .. cmd .. "{0}{0}$}"
        )
    end)
    spintent_measuring = false
    local h  = tex.box[0].height
    local dp = tex.box[0].depth
    tex.box[0] = nil
    return h, dp
end

register_tex_cmd("luafun_spmQ_scale_int", function(entero, frac_cmd)
    local display_type = (frac_cmd == "dfrac") and "display" or
        (spintent_mathstyle_names[tonumber(tex.mathstyle)] or "text")

    -- 1. Medir fracción
    local frac_height, frac_depth = spintent_medir_frac(display_type, frac_cmd)
    local ht_frac_real = frac_height + frac_depth

    -- 2. Construir y medir la hlist del entero UNA sola vez, antes de
    --    escalarla. Antes se llamaba node.mlist_to_hlist dos veces sobre
    --    la misma estructura (una para medir y descartar, otra para
    --    escalar) — la operación más costosa de esta función. Como
    --    spintent_build_submlist_noad/spintent_build_char_noads son puras (sin efectos
    --    secundarios) y spintent_measuring no se lee en ningún lado,
    --    medir antes de escalar produce los mismos valores sin pagar
    --    el costo de una segunda conversión mlist->hlist.
    local h_esc = node.mlist_to_hlist(spintent_build_submlist_noad(entero), display_type, false)
    local ht_entero = h_esc.height + h_esc.depth

    -- 3. Factor
    local factor = (ht_entero == 0) and 1.0 or (ht_frac_real / ht_entero)

    -- 4. Escalar la misma hlist ya medida
    spintent_scale_glyphs(h_esc.list, factor)

    -- 5. Ancho y Glyph nodes físicos (Bucle fusionado y acelerado por C)
    local ancho = 0
    local glyph_nodes = {}

    for n in d_traverse(d_todirect(h_esc.list)) do
        ancho = ancho + d_getfield(n, "width")
        -- Reconvertimos a userdata temporalmente porque spintent_is_digit lo espera así
        local n_user = d_tonode(n)
        if spintent_is_digit(n_user) then
            glyph_nodes[#glyph_nodes + 1] = n_user
        end
    end

    h_esc.width  = ancho
    h_esc.height = frac_height
    h_esc.depth  = frac_depth
    h_esc.shift  = math_floor(frac_depth + 0.5)

    -- 6. sub_box con mathml_filter usando node.direct
    local sub = d_new("sub_box")
    d_setfield(sub, "head", d_todirect(h_esc))

    local noad = d_new("noad")
    d_setfield(noad, "subtype", 0)
    d_setfield(noad, "nucleus", sub)

    -- Devolvemos a userdata para interactuar con las properties de luamml y node.write
    local sub_user = d_tonode(sub)
    local noad_user = d_tonode(noad)

    local mml_mn = {
        [0]         = 'mn',
        entero,
        [':nodes']  = glyph_nodes,
        [':actual'] = entero,
    }

    local properties = node.get_properties_table()
    local p = properties[sub_user] or {}
    p.mathml_filter = function(_mml, _core)
        return mml_mn, mml_mn
    end
    properties[sub_user] = p

    -- 7. Inyectar el noad con luamml deshabilitado temporalmente.
    --    node.write aquí solo deposita el noad en la lista matemática
    --    actual; el callback de luamml se dispara al CERRAR el grupo
    --    matemático (al final de $ o \]). Por tanto deshabilitar el
    --    flag aquí NO suprime el procesamiento del noad — luamml lo
    --    verá igual cuando procese la lista completa.
    --    Lo que sí suprime es cualquier callback intermedio espurio
    --    que pudiera dispararse durante el node.write mismo.
    node.write(noad_user)

end, {"string", "string"})

-- ============================================================
-- §12  GEOMETRÍA PLANA
--      Orden: Punto · Recta · Rayo · Segmento · Ángulo · Arco ·
--      Polígono (incluye \spAfig / \spPfig)
--
--      Comandos cubiertos: \spPoint, \sprecta, \sprayo, \spside,
--      \spmside, \spangle, \spmangle, \sparc, \spmarc, \spPoly,
--      \spAfig, \spPfig
-- ============================================================

-- ------------------------------------------------------------
-- Gramática y utilidades compartidas por toda la familia
-- ------------------------------------------------------------

local spintent_geo_cont_sub = C((1 - P"}")^1)
local spintent_geo_primas   = C(P"'"^1)

local spintent_geo_punto = Ct(
    C(R"AZ") *
    (
        P"_{" * spintent_geo_cont_sub * P"}" * Cc"sub"
        + spintent_geo_primas             * Cc"prima"
        + Cc""                            * Cc""
    )
)

-- Separador de coma entre puntos geométricos.
-- IMPORTANTE: no usa spintent_opt_spaces aquí, porque esa función
-- descarta espacios vía "/ ''" (captura de sustitución). Esa captura,
-- usada DENTRO de un Ct externo, no se descarta — se inserta como
-- elemento string "" suelto en la tabla resultante, duplicando el
-- conteo real de puntos. Se usa un patrón puro sin captura.
local spintent_geo_ws     = S" \t\r\n"
local spintent_geo_ws_opt = spintent_geo_ws^0
local spintent_geo_sep    = spintent_geo_ws_opt * P"," * spintent_geo_ws_opt

local spintent_geo_grammar =
    Ct( spintent_geo_punto * (spintent_geo_sep * spintent_geo_punto)^0 * P(-1) )

-- Gramática sin anclaje final, usada donde se necesita re-chequear
-- cardinalidad exacta (\spPoly, \spAfig, \spPfig)
local spintent_geo_poly_grammar =
    Ct( spintent_geo_punto * (spintent_geo_sep * spintent_geo_punto)^0 )

-- Gramática para nombre propio de recta (mayúsculas Y minúsculas)
local spintent_geo_nombre_recta = Ct(
    C(R"AZ" + R"az") *
    (
        P"_{" * spintent_geo_cont_sub * P"}" * Cc"sub"
        + spintent_geo_primas             * Cc"prima"
        + Cc""                            * Cc""
    )
)

-- Tabla de conversión de subíndices numéricos a texto
local spintent_geo_num_a_texto = {
    ["0"] = "cero",  ["1"] = "uno",   ["2"] = "dos",
    ["3"] = "tres",  ["4"] = "cuatro",["5"] = "cinco",
    ["6"] = "seis",  ["7"] = "siete", ["8"] = "ocho",
    ["9"] = "nueve",
}

-- Prefijos para primas
local spintent_geo_prima_prefijo = {
    "prima", "doble-prima", "triple-prima", "cuádruple-prima"
}

-- Constructor de intent literal
-- Si cmd == "" (silent-cmd=true en school) omite el concepto
local function spintent_geo_build_intent(cmd, puntos)
    local partes = {}
    if cmd ~= "" then
        partes[#partes + 1] = cmd
    end
    for _, pt in ipairs(puntos) do
        local letra, mod, tipo = pt[1], pt[2], pt[3]
        partes[#partes + 1] = letra
        if tipo == "sub" then
            partes[#partes + 1] = "sub"
            partes[#partes + 1] = spintent_geo_num_a_texto[mod] or mod
        elseif tipo == "prima" then
            local n = #mod
            partes[#partes + 1] = spintent_geo_prima_prefijo[n] or (n .. "-prima")
        end
    end
    return t_concat(partes, "-")
end

-- Constructor de intent para nombre propio de recta
local function spintent_geo_build_intent_nombre(cmd, nombre)
    return spintent_geo_build_intent(cmd, { nombre })
end

-- Constructor del visual para \overline/\overrightarrow/\overleftrightarrow/\widearc
local function spintent_geo_build_visual(puntos)
    local partes = {}
    for _, pt in ipairs(puntos) do
        local letra, mod, tipo = pt[1], pt[2], pt[3]
        if tipo == "sub" then
            partes[#partes + 1] = letra .. "_{" .. mod .. "}"
        elseif tipo == "prima" then
            local buf = { letra }
            for _ = 1, #mod do buf[#buf + 1] = "'" end
            partes[#partes + 1] = t_concat(buf)
        else
            partes[#partes + 1] = letra
        end
    end
    return t_concat(partes)
end

-- Constructor del visual para nombre propio de recta
local function spintent_geo_build_visual_nombre(nombre)
    local letra, mod, tipo = nombre[1], nombre[2], nombre[3]
    if tipo == "sub" then
        return letra .. "_{" .. mod .. "}"
    elseif tipo == "prima" then
        local buf = { letra }
        for _ = 1, #mod do buf[#buf + 1] = "'" end
        return t_concat(buf)
    else
        return letra
    end
end

-- ------------------------------------------------------------
-- §12.1  PUNTO — \spPoint
-- ------------------------------------------------------------
local spintent_geo_single_point_grammar = spintent_geo_punto * P(-1)

register_tex_cmd("luafun_geo_point_parse_and_set",
    function(csv, mode, silent)
    csv = spintent_trim(csv)

    local cmd_base   = "punto"
    local cmd_intent = (silent == "true") and "" or cmd_base

    local punto = spintent_geo_single_point_grammar:match(csv)
    if not punto then
        token_set_macro("l__spintent_geo_error_str",  "true")
        token_set_macro("l__spintent_geo_intent_str", "")
        token_set_macro("l__spintent_geo_visual_tl",  "")
        return
    end

    token_set_macro("l__spintent_geo_error_str", "false")

    if mode == "school" then
        token_set_macro("l__spintent_geo_intent_str",
                        spintent_geo_build_intent_nombre(cmd_intent, punto))
        token_set_macro("l__spintent_geo_visual_tl",
                        spintent_geo_build_visual_nombre(punto))
    else
        local mathcat_intent
        if silent == "true" then
            mathcat_intent = "_" .. cmd_base .. ":prefix($pt):silent"
        else
            mathcat_intent = "_" .. cmd_base .. ":prefix($pt)"
        end
        token_set_macro("l__spintent_geo_intent_str", mathcat_intent)
        token_set_macro("l__spintent_geo_visual_tl",
                        spintent_geo_build_visual_nombre(punto))
    end
end, { "string", "string", "string" })

-- ------------------------------------------------------------
-- §12.2  RECTA, RAYO, SEGMENTO Y ARCO
--        \sprecta, \sprayo, \spside, \spmside, \sparc, \spmarc
--
-- Los seis comparten la misma función Lua genérica de dos puntos,
-- registrada una sola vez a continuación. El parámetro "cmd" recibido
-- desde expl3 determina el concepto verbalizado ("recta", "rayo",
-- "lado", "segmento", "trazo", "arco", ...); el visual (\overline,
-- \overrightarrow, \overleftrightarrow, \widearc) se decide del lado
-- expl3 según qué comando público invocó la función.
-- ------------------------------------------------------------
register_tex_cmd("luafun_geo_parse_and_set", function(cmd, csv, mode, silent)
    csv = spintent_trim(csv)

    local cmd_intent = cmd
    if mode == "school" and silent == "true" then
        cmd_intent = ""
    end

    local mathcat_intent
    if silent == "true" then
        mathcat_intent = "_" .. cmd .. ":prefix($pts):silent"
    else
        mathcat_intent = "_" .. cmd .. ":prefix($pts)"
    end

    local tiene_coma = s_match(csv, ",")

    if tiene_coma then
        local puntos = spintent_geo_grammar:match(csv)
        if not puntos or #puntos == 0 then
            token_set_macro("l__spintent_geo_error_str",     "true")
            token_set_macro("l__spintent_geo_intent_str",    "")
            token_set_macro("l__spintent_geo_visual_tl",     "")
            token_set_macro("l__spintent_geo_is_nombre_str", "false")
            return
        end
        token_set_macro("l__spintent_geo_error_str",     "false")
        token_set_macro("l__spintent_geo_is_nombre_str", "false")
        if mode == "school" then
            token_set_macro("l__spintent_geo_intent_str",
                            spintent_geo_build_intent(cmd_intent, puntos))
            token_set_macro("l__spintent_geo_visual_tl",
                            spintent_geo_build_visual(puntos))
        else
            token_set_macro("l__spintent_geo_intent_str", mathcat_intent)
            token_set_macro("l__spintent_geo_visual_tl",
                            spintent_geo_build_visual(puntos))
        end
    else
        local nombre = spintent_geo_nombre_recta:match(csv)
        if not nombre then
            token_set_macro("l__spintent_geo_error_str",     "true")
            token_set_macro("l__spintent_geo_intent_str",    "")
            token_set_macro("l__spintent_geo_visual_tl",     "")
            token_set_macro("l__spintent_geo_is_nombre_str", "false")
            return
        end
        token_set_macro("l__spintent_geo_error_str",     "false")
        token_set_macro("l__spintent_geo_is_nombre_str", "true")
        if mode == "school" then
            token_set_macro("l__spintent_geo_intent_str",
                            spintent_geo_build_intent_nombre(cmd_intent, nombre))
            token_set_macro("l__spintent_geo_visual_tl",
                            spintent_geo_build_visual_nombre(nombre))
        else
            token_set_macro("l__spintent_geo_intent_str", mathcat_intent)
            token_set_macro("l__spintent_geo_visual_tl",
                            spintent_geo_build_visual_nombre(nombre))
        end
    end
end, { "string", "string", "string", "string" })

-- ------------------------------------------------------------
-- §12.3  ÁNGULO — \spangle, \spmangle
-- ------------------------------------------------------------
local spintent_geo_angulo_un_punto = spintent_geo_punto * P(-1)

local function spintent_geo_build_intent_angulo(cmd, csv, silent)
    local cmd_intent = (silent == "true") and "" or cmd

    local tiene_coma = s_match(csv, ",")

    if tiene_coma then
        local puntos = spintent_geo_grammar:match(csv)
        if not puntos or #puntos == 0 then
            return nil, nil, nil
        end
        local intent = spintent_geo_build_intent(cmd_intent, puntos)
        local visual = spintent_geo_build_visual(puntos)
        return intent, visual, "tres-puntos"
    else
        local un_punto = spintent_geo_angulo_un_punto:match(csv)
        if un_punto then
            local letra  = un_punto[1]
            local mod    = un_punto[2]
            local tipo   = un_punto[3]
            local partes_en = { "en" }
            partes_en[#partes_en + 1] = letra
            if tipo == "sub" then
                partes_en[#partes_en + 1] = "sub"
                partes_en[#partes_en + 1] = spintent_geo_num_a_texto[mod] or mod
            elseif tipo == "prima" then
                local n = #mod
                partes_en[#partes_en + 1] =
                    spintent_geo_prima_prefijo[n] or (n .. "-prima")
            end
            local sufijo = t_concat(partes_en, "-")
            local intent
            if cmd_intent == "" then
                intent = sufijo
            else
                intent = cmd_intent .. "-" .. sufijo
            end
            local visual = spintent_geo_build_visual_nombre(un_punto)
            return intent, visual, "un-punto"
        else
            local nombre = spintent_trim(csv)
            if nombre == "" then return nil, nil, nil end
            local nombre_intent = s_gsub(nombre, "[\\{}%s]", "-")
            nombre_intent = s_gsub(nombre_intent, "%-%-+", "-")
            local intent
            if cmd_intent == "" then
                intent = nombre_intent
            else
                intent = cmd_intent .. "-" .. nombre_intent
            end
            return intent, nombre, "nombre"
        end
    end
end

register_tex_cmd("luafun_geo_angulo_parse_and_set",
    function(cmd, csv, mode, silent, recto)
    csv = spintent_trim(csv)

    local cmd_base
    if recto == "true" then
        cmd_base = cmd .. "-recto"
    else
        cmd_base = cmd
    end

    local intent, visual, caso =
        spintent_geo_build_intent_angulo(cmd_base, csv, silent)

    if not intent then
        token_set_macro("l__spintent_geo_error_str",       "true")
        token_set_macro("l__spintent_geo_intent_str",      "")
        token_set_macro("l__spintent_geo_visual_tl",       "")
        token_set_macro("l__spintent_geo_angulo_caso_str", "")
        return
    end

    token_set_macro("l__spintent_geo_error_str", "false")
    token_set_macro("l__spintent_geo_angulo_caso_str", caso)

    if mode == "school" then
        token_set_macro("l__spintent_geo_intent_str", intent)
        token_set_macro("l__spintent_geo_visual_tl",  visual)
    else
        local mathcat_intent
        if silent == "true" then
            mathcat_intent = "_" .. cmd_base .. ":prefix($pts):silent"
        else
            mathcat_intent = "_" .. cmd_base .. ":prefix($pts)"
        end
        token_set_macro("l__spintent_geo_intent_str", mathcat_intent)
        token_set_macro("l__spintent_geo_visual_tl",  visual)
    end
end, { "string", "string", "string", "string", "string" })

-- ------------------------------------------------------------
-- §12.4  POLÍGONO — \spPoly, \spAfig, \spPfig
-- ------------------------------------------------------------
local spintent_geo_poly_npts = {
    [3] = { cmd = "triángulo",  sym = "triangle" },
    [4] = { cmd = "cuadrado",   sym = "square"   },
    [5] = { cmd = "pentágono",  sym = ""         },
    [6] = { cmd = "hexágono",   sym = ""         },
    [7] = { cmd = "heptágono",  sym = ""         },
    [8] = { cmd = "octágono",   sym = ""         },
}

register_tex_cmd("luafun_geo_poly_parse_and_set",
    function(csv, mode)
    csv = spintent_trim(csv)

    local tiene_coma = s_match(csv, ",")
    if not tiene_coma then
        token_set_macro("l__spintent_geo_error_str",      "true")
        token_set_macro("l__spintent_geo_intent_str",     "")
        token_set_macro("l__spintent_geo_intent_pts_str", "")
        token_set_macro("l__spintent_geo_visual_tl",      "")
        token_set_macro("l__spintent_geo_poly_sym_str",   "")
        return
    end

    local puntos = spintent_geo_poly_grammar:match(csv)
    if not puntos or #puntos < 3 then
        token_set_macro("l__spintent_geo_error_str",      "true")
        token_set_macro("l__spintent_geo_intent_str",     "")
        token_set_macro("l__spintent_geo_intent_pts_str", "")
        token_set_macro("l__spintent_geo_visual_tl",      "")
        token_set_macro("l__spintent_geo_poly_sym_str",   "")
        return
    end

    local info     = spintent_geo_poly_npts[#puntos]
    local cmd_base = info and info.cmd or "polígono"
    local sym      = info and info.sym or ""

    token_set_macro("l__spintent_geo_error_str",      "false")
    token_set_macro("l__spintent_geo_poly_sym_str",   sym)
    token_set_macro("l__spintent_geo_intent_pts_str",
                    spintent_geo_build_intent("", puntos))

    if mode == "school" then
        token_set_macro("l__spintent_geo_intent_str",
                        spintent_geo_build_intent(cmd_base, puntos))
        token_set_macro("l__spintent_geo_visual_tl",
                        spintent_geo_build_visual(puntos))
    else
        token_set_macro("l__spintent_geo_intent_str",
                        "_" .. cmd_base .. ":prefix($pts)")
        token_set_macro("l__spintent_geo_visual_tl",
                        spintent_geo_build_visual(puntos))
    end
end, { "string", "string" })

register_tex_cmd("luafun_geo_fig_parse_and_set",
    function(op, csv, mode, print_sym, read_arg)
    csv = spintent_trim(csv)
    read_arg = spintent_trim(read_arg)

    if csv == "" then
        token_set_macro("l__spintent_geo_error_str",    "false")
        token_set_macro("l__spintent_geo_poly_sym_str", "")
        token_set_macro("l__spintent_geo_intent_str",   op)
        token_set_macro("l__spintent_geo_visual_tl",    "")
        return
    end

    local tiene_coma = s_match(csv, ",")
    if not tiene_coma then
        token_set_macro("l__spintent_geo_error_str",    "true")
        token_set_macro("l__spintent_geo_intent_str",   "")
        token_set_macro("l__spintent_geo_visual_tl",    "")
        token_set_macro("l__spintent_geo_poly_sym_str", "")
        return
    end

    local puntos = spintent_geo_poly_grammar:match(csv)
    if not puntos or #puntos < 3 then
        token_set_macro("l__spintent_geo_error_str",    "true")
        token_set_macro("l__spintent_geo_intent_str",   "")
        token_set_macro("l__spintent_geo_visual_tl",    "")
        token_set_macro("l__spintent_geo_poly_sym_str", "")
        return
    end

    local info    = spintent_geo_poly_npts[#puntos]
    local cmd_fig = (read_arg ~= "" and read_arg)
                    or (info and info.cmd)
                    or "polígono"
    local sym     = (info and info.sym) or ""

    token_set_macro("l__spintent_geo_error_str", "false")
    token_set_macro("l__spintent_geo_poly_sym_str",
                    (print_sym == "true") and sym or "")
    token_set_macro("l__spintent_geo_visual_tl",
                    spintent_geo_build_visual(puntos))

    local conector = "del"
    local partes = { op, conector, cmd_fig }
    for _, pt in ipairs(puntos) do
        local letra, mod, tipo = pt[1], pt[2], pt[3]
        partes[#partes + 1] = letra
        if tipo == "sub" then
            partes[#partes + 1] = "sub"
            partes[#partes + 1] = spintent_geo_num_a_texto[mod] or mod
        elseif tipo == "prima" then
            local n = #mod
            partes[#partes + 1] =
                spintent_geo_prima_prefijo[n] or (n .. "-prima")
        end
    end

    if mode == "school" then
        token_set_macro("l__spintent_geo_intent_str", t_concat(partes, "-"))
    else
        token_set_macro("l__spintent_geo_intent_str",
                        "_" .. op .. "-" .. conector .. "-" .. cmd_fig ..
                        ":prefix($pts)")
    end
end, { "string", "string", "string", "string", "string" })
% \iffalse
%</spintentlua>
% \fi
%
% \iffalse
%<@@=>
%<*mylhmcsty>
% \fi
\NeedsTeXFormat{LaTeX2e}[2026-11-01]
\ProvidesExplPackage{mylhmc}{2025-08-19}{2.0}
  {Line-by-line macrocode reader for l3doc in Lua}
\sys_if_engine_luatex:F
  {
    \msg_new:nnn { mylhmc } { luatex-required } { mylhmc~requires~LuaTeX. }
    \msg_fatal:nn { mylhmc } { luatex-required }
  }
\tl_if_exist:NF \g__codedoc_module_name_tl
  {
    \msg_new:nnn { mylhmc } { l3doc-required } { mylhmc~must~be~loaded~via~l3doc.cls. }
    \msg_fatal:nn { mylhmc } { l3doc-required }
  }
\lua_load_module:n { mylhmc.lua }
\cs_generate_variant:Nn \__mylhmc_set_color:nn { VV }

\seq_new:N \g__mylhmc_colors_export_seq
\seq_new:N \l__mylhmc_path_seq
\tl_new:N  \l__mylhmc_color_export_tl
\tl_new:N  \l__mylhmc_export_name_tl

\cs_new_protected:Npn \__mylhmc_define_xcolor:nnn #1#2#3
  {
    \definecolor { #1 } { #2 } { #3 }
  }
\cs_generate_variant:Nn \__mylhmc_define_xcolor:nnn { Vnn }

\cs_new_protected:Npn \__mylhmc_color_export_name:n #1
  {
    % Clave Lua = path COMPLETO unido con '_' (no solo el último
    % segmento): así el path sirve para organizar la UI de colores
    % y a la vez garantiza claves únicas en color_data sin colisiones
    % entre paths distintos que compartan el último segmento
    % (p.ej. guard/name vs module/name).
    \seq_set_split:Nnn \l__mylhmc_path_seq { / } { #1 }
    \seq_set_map:NNn \l__mylhmc_path_seq \l__mylhmc_path_seq { \tl_trim_spaces:n {##1} }
    \tl_set:Ne \l__mylhmc_export_name_tl { \seq_use:Nn \l__mylhmc_path_seq { _ } }
  }

\cs_new_protected:Npn \__mylhmc_export_color_aux:nnn #1#2#3
  {
    \__mylhmc_color_export_name:n { #1 }
    \__mylhmc_define_xcolor:Vnn \l__mylhmc_export_name_tl { #2 } { #3 }
  }
\cs_new_protected:Npn \__mylhmc_process_color_export:n #1
  {
    \__mylhmc_export_color_aux:nnn #1
  }

\cs_new_protected:Npn \__mylhmc_define_color:nnn #1#2#3
  {
    % Native expl3 color
    \color_set:nnn { mylhmc / #1 } { #2 } { #3 }
    % Export to Lua using the last path component as the name
    \__mylhmc_color_export_name:n { #1 }
    \color_export:nnN { mylhmc / #1 } { space-sep-rgb } \l__mylhmc_color_export_tl
    \__mylhmc_set_color:VV \l__mylhmc_export_name_tl \l__mylhmc_color_export_tl
    % Defer xcolor definition until xcolor is available
    \seq_gput_right:Nn \g__mylhmc_colors_export_seq { { #1 } { #2 } { #3 } }
  }

\__mylhmc_define_color:nnn { argument   } { HTML } { 48467C } % Violáceo Oscuro   -> # (sin número)
\__mylhmc_define_color:nnn { argnum     } { HTML } { 6A1B9A } % Púrpura           -> número en #1
\__mylhmc_define_color:nnn { privatevar } { HTML } { 3D7A7A } % Verde azulado     -> \g__var \l__var
\__mylhmc_define_color:nnn { bracket    } { HTML } { 4A5568 } % Gris Carbón       -> []
\__mylhmc_define_color:nnn { publicvar  } { HTML } { 0066CC } % Azul Medio        -> \l_var \g_var
\__mylhmc_define_color:nnn { privatefun } { HTML } { 00695C } % Menta Oscuro      -> \__foo:Nn
\__mylhmc_define_color:nnn { publicfun  } { HTML } { 004085 } % Azul Denim        -> \cs_new:Nn
\__mylhmc_define_color:nnn { brace      } { HTML } { 4A5568 } % Gris Carbón       -> { }
\__mylhmc_define_color:nnn { comment    } { HTML } { 7F8C8D } % Gris Medio        -> %
\__mylhmc_define_color:nnn { signature  } { HTML } { 5C6BC0 } % Azul Acero        -> :
\__mylhmc_define_color:nnn { sigargs    } { HTML } { 7F8C8D } % Gris Medio        -> Nn nTF
\__mylhmc_define_color:nnn { number     } { HTML } { 6A1B9A } % Púrpura           -> 10 100
\__mylhmc_define_color:nnn { constant   } { HTML } { 5E2A84 } % Púrpura Oscuro    -> \c_foo \c__foo
\__mylhmc_define_color:nnn { ltcmd      } { HTML } { 002EA6 } % Azul Eléctrico    -> \NewDocumentCommand
\__mylhmc_define_color:nnn { ltcmd_arg  } { HTML } { 5C658A } % Azul Pizarra      -> { O{} m +m t }
\__mylhmc_define_color:nnn { legacy_at  } { HTML } { 607D8B } % Gris Azulado      -> \mi@cmd
\__mylhmc_define_color:nnn { scan       } { HTML } { A64426 } % Rojo Ladrillo     -> \s_stop
\__mylhmc_define_color:nnn { quark      } { HTML } { 800040 } % Vino              -> \q_nil
\__mylhmc_define_color:nnn { keyfun     } { HTML } { B27A00 } % Oro Viejo         -> .meta:n
\__mylhmc_define_color:nnn { backtick   } { HTML } { 558B2F } % Verde Manzana     -> `A
\__mylhmc_define_color:nnn { escape     } { HTML } { D81B60 } % Rosa Eléctrico    -> ^^M ~
\__mylhmc_define_color:nnn { danger     } { HTML } { D32F2F } % Rojo Sangre       -> :D
\__mylhmc_define_color:nnn { quote      } { HTML } { E65100 } % Naranja Profundo  -> ' "
\__mylhmc_define_color:nnn { math       } { HTML } { FF1818 } % Rojo Neón         -> $ $$ \( \) \[ \]
\__mylhmc_define_color:nnn { compare    } { HTML } { 808080 } % Gris Medio        -> < > = <= >= !=
\__mylhmc_define_color:nnn { latex_cmd  } { HTML } { 3F3FFF } % Azul Brillante    -> \mathrm \textbf \,
\__mylhmc_define_color:nnn { structure  } { HTML } { 004085 } % Azul Denim (bold) -> \NeedsTeXFormat \ProvidesExplPackage
\__mylhmc_define_color:nnn { unit       } { HTML } { E67E22 } % Naranja Calabaza  -> pt em cm mu
\__mylhmc_define_color:nnn { pkgstruct  } { HTML } { 064FA5 } % Azul cuerpo       -> \RequirePackage \IfPackageLoadedTF
\__mylhmc_define_color:nnn { support    } { HTML } { 3587B1 } % Azul Boston       -> hyperref, fontspec, unicode-math
\__mylhmc_define_color:nnn { mathtag    } { HTML } { 2036D2 } % Azul MathML       -> \MathMLintent \MathMLarg \luamml_annotate:en \tagpdfsetup etc
\__mylhmc_define_color:nnn { mathml_var } { HTML } { 008000 } % Verde Puro        -> $foo identifiers in MathMLintent/MathMLarg
\__mylhmc_define_color:nnn { intent     } { HTML } { 5C658A } % Azul Pizarra      -> :prefix :postfix :infix :silent etc (global)
\__mylhmc_define_color:nnn { filename   } { HTML } { 4A5568 } % Gris Carbón       -> \file_... \input \include \IfFileExists (nombre de archivo)
\__mylhmc_define_color:nnn { percent    } { HTML } { 4A5568 } % Gris Carbón       -> %
\__mylhmc_define_color:nnn { arroba     } { HTML } { D81B60 } % Rosa Eléctrico    -> @
\__mylhmc_define_color:nnn { star       } { HTML } { D81B60 } % Rosa Eléctrico    -> *
\__mylhmc_define_color:nnn { keyname    } { HTML } { 483D8B } % Gris Índigo       -> cifras-sep read-sign terse clear
\__mylhmc_define_color:nnn { other_np   } { HTML } { D32F2F } % Rojo Sangre       -> \__other_mod: privados ajenos
\__mylhmc_define_color:nnn { pkgname    } { HTML } { 6A5ACD } % SlateBlue         -> nombre de paquete (negrita)
\__mylhmc_define_color:nnn { clsname    } { HTML } { 8470FF } % LightSlateBlue    -> nombre de clase (negrita)

\__mylhmc_define_color:nnn { guard / angle } { HTML } { 4A5568 } % Gris Carbón  -> 〈 〉
\__mylhmc_define_color:nnn { guard / name  } { HTML } { 5A6370 } % Gris Pizarra -> nombre módulo
\__mylhmc_define_color:nnn { guard / star  } { HTML } { 95A5A6 } % Gris Plata   -> * (guardas)
\__mylhmc_define_color:nnn { guard / slash } { HTML } { 95A5A6 } % Gris Plata   -> / (guardas)
\__mylhmc_define_color:nnn { guard / paren } { HTML } { 95A5A6 } % Gris Plata   -> ( ) (guardas)
\__mylhmc_define_color:nnn { guard / amp   } { HTML } { 95A5A6 } % Gris Plata   -> & (guardas)
\__mylhmc_define_color:nnn { guard / pipe  } { HTML } { 95A5A6 } % Gris Plata   -> | (guardas)
\__mylhmc_define_color:nnn { guard / bang  } { HTML } { 95A5A6 } % Gris Plata   -> ! (guardas)
\__mylhmc_define_color:nnn { module / at   } { HTML } { 008080 } % Teal         -> __spintent
\__mylhmc_define_color:nnn { module / name } { HTML } { 008080 } % Teal         -> nombre módulo en __spintent=
\__mylhmc_define_color:nnn { module / eq   } { HTML } { 808080 } % Gris Medio   -> =

\hook_gput_code:nnn { begindocument } { mylhmc }
  {
    \IfPackageLoadedT { xcolor }
      {
        \seq_map_function:NN \g__mylhmc_colors_export_seq \__mylhmc_process_color_export:n
      }
  }

\NewDocumentCommand \myhlc { v }
  {
    \mode_leave_vertical:
    \group_begin:
      \ttfamily
      \__mylhmc_inline:n {#1}
    \group_end:
  }
\cs_new_protected:Npn \__mylhmc_inline:n #1
  {
    \tl_set:Nn \l__mylhmc_verbatim_line_tl {#1}
    \__mylhmc_inline:
  }

\hook_gput_code:nnn { begindocument } { mylhmc }
  {
    \renewcommand \theCodelineNo
      {
        \__mylhmc_actual_text_begin:
        \group_begin:
          \color_select:n { mylhmc / guard / angle }
          \ttfamily \tiny \arabic { CodelineNo }
        \group_end:
        \__mylhmc_actual_text_end:
      }
  }

\group_begin:
  \int_set:Nn \tex_endlinechar:D { -1 }
  \use:n
    {
      \char_set_catcode_active:n { 32 }
      \tl_const:Nn \c__mylhmc_active_space_tl
    }{ }
\group_end:

\tl_new:N \l__mylhmc_verbatim_line_tl
\tl_new:N \l__mylhmc_verbatim_finish_tl

\group_begin:
  \char_set_catcode_active:n { 13 }
  \cs_new_protected:Npe \__mylhmc_set_end_pattern:n #1
    {
      \tl_set:Nn \exp_not:N \l__mylhmc_verbatim_finish_tl
        {
          \c_percent_str
          \prg_replicate:nn { 4 } { \exp_not:o { \c__mylhmc_active_space_tl } }
          \exp_not:o { \active@escape@char } end
          \c_left_brace_str #1 \c_right_brace_str
          \exp_not:N ^^M
        }
    }
  \cs_new_protected:Npn \__mylhmc_verbatim_start:w #1
    {
      \str_if_eq:nnTF { #1 } { ^^M }
        { \__mylhmc_verbatim_read_line:w }
        { \__mylhmc_verbatim_read_line:w #1 }
    }
  \cs_new_protected:Npn \__mylhmc_verbatim_read_line:w #1 ^^M
    {
      \tl_set:Nn \l__mylhmc_verbatim_line_tl { #1 ^^M }
      \tl_if_eq:NNTF \l__mylhmc_verbatim_line_tl \l__mylhmc_verbatim_finish_tl
        { \exp_args:Ne \end { \@currenvir } }
        {
          \__mylhmc_output_line:
          \__mylhmc_verbatim_read_line:w
        }
    }
\group_end:

\cs_new_protected_nopar:Nn \__mylhmc_output_line:
  {
    \tex_noindent:D
    \__mylhmc_process_line:
    \tex_par:D
  }

\cs_generate_variant:Nn \__mylhmc_set_end_pattern:n { V }
\cs_set_protected_nopar:Npn \xmacro@code
  {
    \__mylhmc_set_end_pattern:V \@currenvir
    \__mylhmc_verbatim_start:w
  }
\cs_set_eq:cN { xmacro@code* } \xmacro@code

% ============================================================
% Interfaz de usuario: \setmylhmc { key = val }
% ============================================================

\msg_new:nnnn { mylhmc } { invalid-hex-color }
  { Invalid~color~value~'#2'~for~key~'#1':~expected~6~hex~digits~(e.g.~FF0000). }
  { The~color~was~not~changed~and~keeps~its~previous~value. }
\msg_new:nnnn { mylhmc } { unknown-key }
  { Unknown~mylhmc~key~'#1'. }
  { Check~the~color~name~(use~its~full~path,~e.g.~guard/star)~or~use~
    provide~to~register~module~commands. }

\keys_define:nn { mylhmc }
  {
    unknown .code:n =
      { \msg_error:nne { mylhmc } { unknown-key } { \l_keys_key_str } } ,
    provide .code:n =
      { \clist_map_function:nN {#1} \__mylhmc_register_module_cmd:n } ,
    provide .value_required:n = true ,
  }

% Genera automáticamente una clave por cada color ya registrado (mismo
% path que su nombre, p.ej. \setmylhmc{guard/star=FF0000}). Debe ir
% DESPUÉS de todas las llamadas a \__mylhmc_define_color:nnn de arriba,
% ya que recorre \g__mylhmc_colors_export_seq tal como quedó poblada.
\cs_new_protected:Npn \__mylhmc_define_color_key:nnn #1#2#3
  {
    \keys_define:nn { mylhmc }
      {
        #1 .code:n =
          {
            \regex_match:nnTF { ^ [0-9A-Fa-f] {6} $ } {##1}
              { \__mylhmc_define_color:nnn { #1 } { HTML } {##1} }
              { \msg_error:nnnn { mylhmc } { invalid-hex-color } { #1 } {##1} }
          } ,
        #1 .value_required:n = true ,
      }
  }
\seq_map_inline:Nn \g__mylhmc_colors_export_seq
  { \__mylhmc_define_color_key:nnn #1 }

\NewDocumentCommand \setmylhmc { m }
  { \keys_set:nn { mylhmc } { #1 } }

\file_input_stop:
% \iffalse
%</mylhmcsty>
% \fi
%
% \iffalse
%<*mylhmclua>
% \fi
-- mylhmc.lua v2.0 2025-08-19 (node.direct + Loop/Memory Micro-Optimized)
-- LuaLaTeX syntax highlighter for l3doc macrocode environments

local require      = require
local pcall        = pcall
local ipairs       = ipairs
local pairs        = pairs
local tonumber     = tonumber
local t_unpack     = table.unpack or unpack

local s_sub        = string.sub
local s_find       = string.find
local s_match      = string.match
local s_format     = string.format
local s_gsub       = string.gsub
local s_gmatch     = string.gmatch
local s_byte       = string.byte
local s_lower      = string.lower
local s_rep        = string.rep

local u_codes      = utf8.codes

local d_new        = node.direct.new
local d_copy       = node.direct.copy
local d_copylist   = node.direct.copy_list
local d_tonode     = node.direct.tonode
local d_setlink    = node.direct.setlink
local d_getnext    = node.direct.getnext
local d_setfield   = node.direct.setfield
local n_write      = node.write

local f_current    = font.current
local f_getfont    = font.getfont
local f_max        = font.max

local t_get_macro    = token.get_macro
local t_set_macro    = token.set_macro
local t_scan_arg     = token.scan_argument
local t_scan_int     = token.scan_int
local t_scan_dimen   = token.scan_dimen
local t_scan_glue    = token.scan_glue
local t_scan_toks    = token.scan_toks
local t_scan_word    = token.scan_word
local token_set_lua  = token.set_lua

local kpse_find      = kpse.find_file
local f_read         = fonts.constructors.readanddefine
local f_hashes       = fonts.hashes.identifiers

local luatexbase_new = luatexbase.new_luafunction
local lua_get_funcs  = lua.get_functions_table

local current_module_name   = ''
local current_module_prefix = ''

local lpeg_base  = lpeg or require('lpeg')
local P, S, V    = lpeg_base.P, lpeg_base.S, lpeg_base.V
local Cp, Cc     = lpeg_base.Cp, lpeg_base.Cc
local lpeg_match = lpeg_base.match
lpeg_base.locale(lpeg_base)
local alpha      = lpeg_base.alpha
local digit      = lpeg_base.digit

local scan_map = {
  string = t_scan_arg,
  int    = t_scan_int,
  dimen  = t_scan_dimen,
  glue   = t_scan_glue,
  toks   = t_scan_toks,
  word   = t_scan_word
}

local function register_tex_cmd(name, func, args)
  name = '__mylhmc_' .. name .. ':' .. s_rep('n', #args)
  local scanners = {}
  for i = 1, #args do
    scanners[i] = scan_map[args[i]] or t_scan_arg
  end
  local scanning_func
  if #scanners == 0 then
    scanning_func = func
  elseif #scanners == 1 then
    local s1 = scanners[1]
    scanning_func = function() func(s1()) end
  elseif #scanners == 2 then
    local s1, s2 = scanners[1], scanners[2]
    scanning_func = function() func(s1(), s2()) end
  else
    scanning_func = function()
      local values = {}
      for i = 1, #scanners do values[i] = scanners[i]() end
      func(t_unpack(values))
    end
  end
  local index = luatexbase_new(name)
  lua_get_funcs()[index] = scanning_func
  token_set_lua(name, index, 'global', 'protected')
end

local sig_chars  = S'NnVvoxefpTFwDcq'
local hex_chars  = S'0123456789abcdefABCDEF'
local name_chars = alpha + S'_@'
local var_prefix = P'\\' * S'lg'
local name_chars_no_at = alpha + S'_'

local structure_names =
  P'NeedsTeXFormat' + P'ProvidesExplPackage' +
  P'ExplSyntaxOn' + P'ExplSyntaxOff' +
  P'makeatletter' + P'makeatother'

local pkgstruct_names =
  P'ProvidesPackage' + P'ProvidesFile' + P'ProvidesClass' +
  P'RequirePackageWithOptions' + P'RequirePackage' +
  P'LoadClassWithOptions' + P'LoadClass' +
  P'PassOptionsToPackage' + P'PassOptionsToClass' +
  P'DocumentMetadata' +
  P'IfFormatAtLeast'    * (P'TF' + P'T' + P'F') ^ -1 +
  P'IfPackageLoaded'    * (P'WithOptions') ^ -1 * (P'TF' + P'T' + P'F') ^ -1 +
  P'IfFileExists'       * (P'TF' + P'T' + P'F') ^ -1 +
  P'IfDocumentMetadata' * (P'TF' + P'T' + P'F') ^ -1 +
  P'NewHook' + P'NewReversedHook' + P'NewMirroredHookPair' +
  P'AddToHook' * (P'WithArguments') ^ -1 +
  P'RemoveFromHook' + P'UseHook' + P'UseOneTimeHook' +
  P'IfHookEmpty' * (P'TF' + P'T' + P'F') +
  P'DeclareHookRule' + P'ClearHookRule' + P'ClearHook' +
  P'SetDefaultHookLabel' + P'PushDefaultHookLabel' + P'PopDefaultHookLabel' +
  P'ShowHook' + P'LogHook' +
  P'NewProperty' + P'SetProperty' +
  P'RecordProperties' + P'RecordProperty' + P'RefProperty' +
  P'IfPropertyExists'   * (P'TF' + P'T' + P'F') +
  P'IfLabelExists'      * (P'TF' + P'T' + P'F') +
  P'IfPropertyRecorded' * (P'TF' + P'T' + P'F')

-- Familia hyperref: comandos públicos, ninguno con prefijo real de otro
-- (todos divergen justo tras "hyper" o son nombres completamente
-- distintos), así que el orden dentro de la alternancia no es crítico.
local support_names =
  P'hyperlink' + P'hypertarget' + P'hyperbaseurl' + P'hyperimage' +
  P'hyperdef' + P'hyperget' + P'hypersetup' + P'hyperref' +
  P'href' + P'nolinkurl' + P'url' +
  P'autopageref' + P'autoref' +
  P'pdfstringdef' +
  P'currentpdfbookmark' + P'subpdfbookmark' + P'belowpdfbookmark' + P'pdfbookmark' +
  P'texorpdfstring' + P'phantomsection' +
  -- fontspec: sin colisiones de prefijo entre sí (verificado con texlua).
  P'setmainfont' + P'setsansfont' + P'setmonofont' + P'fontspec' +
  P'setboldmathrm' + P'setmathrm' + P'setmathsf' + P'setmathtt' +
  P'IfFontExistsTF' + P'addfontfeature' + P'defaultfontfeatures' +
  -- unicode-math:
  P'setmathfont' + P'unimathsetup'

local mathtag_names =
  P'MathMLintent' + P'MathMLarg' +
  P'luamml_annotate:nen' + P'luamml_annotate:en' +
  P'luamml_attribute:een' + P'luamml_attribute_core:een' +
  P'luamml_begin_single_file:' + P'luamml_end_single_file:' +
  P'luamml_flag_ignore:' + P'luamml_flag_process:' +
  P'luamml_flag_save:nNn' + P'luamml_flag_save:nN' +
  P'luamml_flag_save:nn' + P'luamml_flag_save:n' +
  P'luamml_get_last_mathml_stream:e' +
  P'luamml_ignore:' + P'luamml_pdf_write:' + P'luamml_process:' +
  P'luamml_register_output_hook:N' +
  P'luamml_save:nNn' + P'luamml_save:nN' +
  P'luamml_save:nn' + P'luamml_save:n' +
  P'luamml_set_filename:n' + P'luamml_structelem:' +
  P'tag_start:n' + P'tag_start:' +
  P'tag_stop:n' + P'tag_stop:' +
  P'tag_suspend:n' + P'tag_resume:n' +
  P'tag_get:n' +
  P'tag_if_active_p:' + P'tag_if_active:' +
  P'tag_if_box_tagged:N' +
  P'tag_if_in:n' +
  P'tag_mc_add_missing_to_stream:Nn' +
  P'tag_mc_artifact_group_begin:n' + P'tag_mc_artifact_group_end:' +
  P'tag_mc_begin_pop:n' + P'tag_mc_begin:n' +
  P'tag_mc_end_push:' + P'tag_mc_end:' +
  P'tag_mc_if_in:' +
  P'tag_mc_new_stream:n' +
  P'tag_mc_reset_box:N' +
  P'tag_mc_use:n' +
  P'tag_socket_use:nnn' + P'tag_socket_use:nn' + P'tag_socket_use:n' +
  P'tag_socket_use_expandable:n' +
  P'tag_spacechar_off:' + P'tag_spacechar_on:' +
  P'tag_check_child:nn' +
  P'tag_struct_begin:n' +
  P'tag_struct_end:n' + P'tag_struct_end:' +
  P'tag_struct_use_num:n' + P'tag_struct_use:n' +
  P'tag_struct_object_ref:n' +
  P'tag_struct_insert_annot:nn' +
  P'tag_struct_parent_int:' +
  P'tag_struct_gput_ref:nnn' + P'tag_struct_gput:nnn' +
  -- latex-lab-math / latex-lab-mathintent: comandos públicos del sistema
  -- de tagging de matemáticas -- mismo color 'mathtag', sin firma expl3
  -- ni estado especial (ninguno activa in_mathml_intent/in_mathml_arg).
  P'MathCollectTrue' + P'MathCollectFalse' + P'm@th' +
  P'SuspendTagging' + P'UseMathForPositioningText' + P'UseStructureName' +
  P'invisibletimes' + P'functionapplication' +
  -- tagpdf: comandos públicos estilo LaTeX2e (sin guión bajo, no
  -- colisionan con los \tag_*:n de arriba, que sí llevan guión bajo).
  P'tagpdfsetup' + P'tagtool' +
  P'tagmcifinTF' + P'tagmcuse' + P'tagmcbegin' + P'tagmcend' +
  P'tagstructbegin' + P'tagstructend' + P'tagstructuse' +
  P'ShowTagging'

local keys_names =
  P'keys_define:' * S'Ncn' * P'n' +
  P'keys_set:'    * S'Ncn' * P'n'

-- Resto de la familia \keys_...:n... (l3keys) cuyo PRIMER argumento es
-- siempre {módulo} o {módulo/path}, pero cuyo SEGUNDO argumento cambia
-- de significado según la función (lista de claves, lista de grupos,
-- nombre de clave único, etc.) -- por eso solo coloreamos el primero,
-- sin extender el tratamiento de "cuerpo de claves" como en .meta/keys_define.
-- Orden: nombres más largos antes que sus prefijos (set_known/set_groups/
-- set_exclude_groups antes de cualquier 'set' bare, que ya cubre keys_names).
local keys_module_arg_names =
  P'keys_' *
  (P'set_exclude_groups' + P'set_groups' + P'set_known' +
   P'precompile' +
   P'if_choice_exist_p' + P'if_choice_exist' +
   P'if_exist_p' + P'if_exist' +
   P'show' + P'log') *
  P':' * P'n' * sig_chars ^ 0

-- \file_..., \lua_load_module:n, y la familia clásica \input/\include/
-- \InputIfFileExists: primer argumento = nombre de archivo. Cada literal
-- incluye su propio ':' (los nombres expl3 no comparten prefijo real
-- una vez incluido el ':', porque divergen en la letra siguiente).
local file_arg_names =
  P'file_input_raw:'          * sig_chars ^ 1 +
  P'file_input:'              * sig_chars ^ 1 +
  P'file_if_exist_input:'     * sig_chars ^ 1 +
  P'file_if_exist_p:'         * sig_chars ^ 1 +
  P'file_if_exist:'           * sig_chars ^ 1 +
  P'file_forget:'             * sig_chars ^ 1 +
  P'file_hex_dump:'           * sig_chars ^ 1 +
  P'file_get_full_name:'      * sig_chars ^ 1 +
  P'file_get:'                * sig_chars ^ 1 +
  P'file_full_name:'          * sig_chars ^ 1 +
  P'file_parse_full_name_apply:' * sig_chars ^ 1 +
  P'file_parse_full_name:'    * sig_chars ^ 1 +
  P'lua_load_module:'         * sig_chars ^ 1 +
  P'includeonly' + P'input' + P'include' + P'InputIfFileExists'

-- \hook_gput_code:nnn, \hook_gput_code_with_args:nnn, \hook_gremove_code:nn
-- (label en arg2) y \hook_gset_rule:nnnn (labels en arg2 Y arg4). El
-- argumento 'hook' (arg1) y 'code'/'relation' (arg3) NO se tocan aquí:
-- siguen su coloreado normal (código real en el caso de 'code').
local hook_label_names =
  P'hook_gput_code_with_args:' * sig_chars ^ 1 +
  P'hook_gput_code:'           * sig_chars ^ 1 +
  P'hook_gremove_code:'        * sig_chars ^ 1 +
  P'hook_gset_rule:'           * sig_chars ^ 1

-- Convención de nombres de hook predefinidos: package/NOMBRE/before,
-- class/NOMBRE/after, file/NOMBRE.ext/before, etc. Global: aparece como
-- texto plano dentro del argumento {hook} de \UseHook, \AddToHook,
-- \hook_use:n, \hook_new:n, el propio arg1 de la familia de arriba, etc.
-- Solo coloreamos el NOMBRE de en medio; el resto queda sin tocar.
local hook_name_convention =
  (P'package' + P'class' + P'file') * P'/' *
  (P(1) - P'/') ^ 1 * P'/' * (P'before' + P'after')

local msgerror_names =
  P'msg_error:nneee'

local msg_generic_names =
  P'msg_' * (P'new' + P'error' + P'warning' + P'info' + P'note' +
             P'critical' + P'fatal' + P'expandable_error' +
             P'redirect_name' + P'redirect_class' + P'log' + P'term') *
  P':' * sig_chars ^ 1

-- Paréntesis balanceados (con soporte de anidación, p.ej. e:infix($x,$y)
-- dentro de otro grupo): mismo idioma recursivo P{...V(1)...} que usa
-- combined_pattern más abajo en este archivo.
local balanced_parens
balanced_parens = P{ P'(' * ((P(1) - S'()') + V(1)) ^ 0 * P')' }

local intent_names =
  P':' * S' \t' ^ 0 *
  (P'prefix' + P'postfix' + P'function' + P'nofix' +
   P'silent' + P'time' + P'date' + P'infix') *
  balanced_parens ^ -1

local ltcmd_def_names =
  P'NewDocumentCommand' + P'NewDocumentEnvironment' +
  P'DeclareDocumentCommand' + P'DeclareDocumentEnvironment' +
  P'NewExpandableDocumentCommand' + P'DeclareExpandableDocumentCommand' +
  P'RenewDocumentCommand' + P'RenewDocumentEnvironment' +
  P'ProvideDocumentCommand' + P'ProvideDocumentEnvironment'

local ltcmd_use_names =
  P'BooleanTrue' + P'BooleanFalse' +
  P'IfNoValue' * (P'TF' + P'T' + P'F') +
  P'IfValue'   * (P'TF' + P'T' + P'F') +
  P'IfBoolean' * (P'TF' + P'T' + P'F') +
  P'ProcessedArgument' + P'ReverseBoolean' +
  P'SplitArgument' + P'SplitList' +
  P'ProcessList' + P'TrimSpaces'

local gen_variant_suffix =
  P'cs_generate_variant:' * S'Nc' * P'n' +
  P'prg_generate_conditional_variant:' * S'Nc' * P'nn' +
  P'prg_new_conditional:' * S'N' * P'pnn' +
  P'prg_new_conditional:' * S'N' * P'nn' +
  P'prg_new_protected_conditional:' * S'N' * P'pnn' +
  P'prg_new_protected_conditional:' * S'N' * P'nn' +
  P'prg_new_eq_conditional:' * S'N' * P'Nn'

local legacy_names =
  P'legacy_if_set_true:n'  + P'legacy_if_set_false:n'  +
  P'legacy_if_gset_true:n' + P'legacy_if_gset_false:n' +
  P'legacy_if_set:nn'      + P'legacy_if_gset:nn'      +
  P'legacy_if_p:n'         + P'legacy_if:n' * (P'TF' + P'T' + P'F')

local cs_w     = P'cs:w'
local cs_end   = P'cs_end:'
local dim_unit = P'pt' + P'pc' + P'in' + P'cm' + P'mm' + P'bp' +
                 P'dd' + P'cc' + P'sp' + P'em' + P'ex' + P'px' + P'mu'
local fill_kw  = P'filll' + P'fill' + P'fil'
local skip_kw  = P'minus' + P'plus'
local num_part = digit ^ 1 * (P'.' * digit ^ 1) ^ -1

local code_rules = {
  {24, P'#' ^ 1 * digit},
  {1,  P'#' ^ 1},
  {42, intent_names},
  {21, P'$$'}, {21, P'$'},
  {21, P'\\' * S'(['}, {21, P'\\' * S')]'},
  {18, P'\x5e\x5e\x5e\x5e' * hex_chars * hex_chars * hex_chars * hex_chars},
  {18, P'\\' ^ -1 * P'\x5e\x5e' * P(1) * P(1) ^ -1},
  {18, P'~'},
  {22, P' ' * (P'<=' + P'>=' + P'!=' + P'<' + P'>' + P'=' + P'/') * P' '},
  {22, P' ' * P'=' * (P'\\' + -P(1))},
  {29, P','},
  {13, P'\\' * (alpha + S'_@') ^ 1 * P':D'},
  {10, P'\\' * P's__' * (alpha + S'_') ^ 1},
  {10, P'\\' * P's_'  * (alpha + S'_') ^ 1},
  {11, P'\\' * P'q__' * (alpha + S'_') ^ 1},
  {11, P'\\' * P'q_'  * (alpha + S'_') ^ 1},
  {38, P'\\' * mathtag_names},
  {39, P'\\' * keys_names},
  {44, P'\\' * keys_module_arg_names},
  {40, P'\\' * msgerror_names},
  {41, P'\\' * msg_generic_names},
  {46, P'\\' * file_arg_names},
  {47, P'\\' * hook_label_names},
  {48, hook_name_convention},
  {37, P'\\' * pkgstruct_names},
  {49, P'\\' * support_names},
  {25, P'\\' * structure_names},
  {14, P'\\' * ltcmd_def_names},
  {33, P'\\' * ltcmd_use_names},
  {19, P'\\' * P'c__' * name_chars ^ 1},
  {20, P'\\' * P'c_'  * (alpha + S'@') * name_chars ^ 0},
  {2,  var_prefix * P'__' * name_chars ^ 1},
  {4,  var_prefix * P'_' * (alpha + S'@') * name_chars ^ 0},
  {15, P'\\' * (alpha ^ 0 * (P'@' + P'!' + P'&' + P'*') * alpha ^ 0) ^ 1},
  {5,  P'\\' * P'__' * name_chars ^ 1 * (P':' * sig_chars ^ 0) ^ -1},
  {34, P'\\' * legacy_names},
  {26, P'\\' * gen_variant_suffix},
  {30, P'\\' * cs_w},
  {31, P'\\' * cs_end},
  {16, P'.' * (alpha + S'_') ^ 1 * P':' * (alpha + S'_') ^ 0},
  {28, P'``'},
  {17, P'`' * P'\\' * (P(1) - alpha)},
  {17, P'`' * P'\\' * alpha * -alpha},
  {17, P'`' * (P(1) - alpha - P'\\')},
  {17, P'`' * alpha * -alpha},
  {17, P'`' * -alpha},
  {28, P'`'},
  {28, P'\''},
  {28, P'\x22'},
  {23, P'\\\\'},
  {23, P'\\' * S',:;!>~|/*+-={}#$_&%\x5e'},
  {35, num_part * dim_unit},
  {36, (fill_kw + skip_kw) * -alpha},
  {12, digit ^ 1},
  {6,  P'\\' * alpha * name_chars_no_at ^ 0 * P':' * sig_chars ^ 0},
  {23, P'\\' * alpha * name_chars_no_at ^ 0},
  {8,  S'{}'},
  {3,  S'[]'},
  {9,  P'%' * (P(1) - S'\r\n') ^ 0 * (S'\r\n' + -1)},
}

local function build_combined()
  local alts = nil
  for _, rule in ipairs(code_rules) do
    local p = Cp() * rule[2] * Cp() * Cc(rule[1])
    alts = alts and (alts + p) or p
  end
  return P{ alts + P(1) * V(1) }
end

local combined_pattern = build_combined()
local LM_LANGLE = 0x2329
local LM_RANGLE = 0x232A
local bold_cache, lm_cache, italic_cache = {}, {}, {}

local LM_ITALIC_PATH, LM_REGULAR_PATH, LM_BOLD_PATH = nil, nil, nil

local function ensure_lm_paths()
  if not LM_REGULAR_PATH then
    LM_REGULAR_PATH = kpse_find('lmroman10-regular.otf', 'opentype fonts')
    LM_ITALIC_PATH  = kpse_find('lmmono10-italic.otf',   'opentype fonts')
    LM_BOLD_PATH    = kpse_find('lmmonolt10-bold.otf',   'opentype fonts')
  end
end

local function load_font_at_size(filepath, target_size)
  if not filepath then return nil end
  local font_spec = '[' .. filepath .. ']'
  local ok, data, new_id = pcall(f_read, font_spec, target_size)
  if not ok or not data or not new_id or new_id <= 0 then return nil end
  f_hashes[new_id] = data
  return new_id
end

local function get_bold_fid(fid)
  if bold_cache[fid] then return bold_cache[fid] end
  local f = f_getfont(fid)
  if not f or not f.name or not f.size then
    bold_cache[fid] = fid; return fid
  end
  local base_family = s_match(f.name, '^%[?([^%-]+)') or f.name
  local target_size = f.size
  local best_id, best_score, best_path = nil, 0, nil

  -- Optimización: cacheamos f_max() para el for
  local max_f = f_max()
  for i = 1, max_f do
    local bf = f_getfont(i)
    if bf and bf.name and bf.name ~= f.name and
       s_find(bf.name, base_family, 1, true) then
      local lower_name = s_lower(bf.name)
      local score = 0
      if s_find(lower_name, 'medium') then score = 3
      elseif s_find(lower_name, 'semibold') or s_find(lower_name, 'semi') or
             s_find(lower_name, 'demi') then score = 2
      elseif s_find(lower_name, 'bold') then score = 1
      end
      if score > best_score then
        best_score = score
        if bf.size == target_size then
          bold_cache[fid] = i; return i
        else
          best_id = i
          best_path = bf.filename or bf.name
        end
      elseif score > 0 and score == best_score and bf.size == target_size then
        bold_cache[fid] = i; return i
      end
    end
  end
  if best_path and best_id then
    local clean_path = s_match(best_path, '^%[?(.-)%]?$') or best_path
    local new_id = load_font_at_size(clean_path, target_size)
    if new_id then bold_cache[fid] = new_id; return new_id end
    bold_cache[fid] = best_id; return best_id
  end
  ensure_lm_paths()
  local fallback_id = load_font_at_size(LM_BOLD_PATH, target_size)
  if fallback_id then bold_cache[fid] = fallback_id; return fallback_id end
  bold_cache[fid] = fid; return fid
end

local function get_italic_fid(fid)
  if italic_cache[fid] then return italic_cache[fid] end
  local f = f_getfont(fid)
  if not f or not f.size then italic_cache[fid] = fid; return fid end
  ensure_lm_paths()
  local fallback_id = load_font_at_size(LM_ITALIC_PATH, f.size)
  if fallback_id then italic_cache[fid] = fallback_id; return fallback_id end
  italic_cache[fid] = fid; return fid
end

local function get_lm_fid(fid)
  fid = fid or f_current()
  if lm_cache[fid] then return lm_cache[fid] end
  local cf = f_getfont(fid)
  if not cf or not cf.size then lm_cache[fid] = fid; return fid end

  -- Optimización: cacheamos f_max() para el for
  local max_f = f_max()
  for i = 1, max_f do
    local f = f_getfont(i)
    local fname = f and (f.filename or f.name) or ''
    if f and f.size == cf.size and
       s_find(s_lower(fname), 'lmroman10%-regular') then
      lm_cache[fid] = i; return i
    end
  end
  ensure_lm_paths()
  local new_id = load_font_at_size(LM_REGULAR_PATH, cf.size)
  if new_id then lm_cache[fid] = new_id; return new_id end
  lm_cache[fid] = fid; return fid
end

local color_data, push_templates, module_user_cmds = {}, {}, {}

register_tex_cmd('register_module_cmd', function(name)
  name = s_match(name, '^%s*(.-)%s*$') or name
  local clean = s_sub(name, 1, 1) == '\\' and s_sub(name, 2) or name
  -- quitar un '*' final (variante con asterisco, estilo \section*): el
  -- colorizer siempre separa nombre + '*' antes de consultar
  -- module_user_cmds, así que registrar "nuevo*" tal cual no serviría.
  clean = s_match(clean, '^(.-)%*$') or clean
  local bare  = s_match(clean, '^(.-):[NnVvoxefpTFwDcq]+$') or clean
  if bare and bare ~= '' then
    module_user_cmds[bare], module_user_cmds[clean] = true, true
  end
end, {'string'})

register_tex_cmd('set_color', function(name, spec)
  local r, g, b = s_match(spec, '([%d%.]+)%s+([%d%.]+)%s+([%d%.]+)')
  if r then
    local nr, ng, nb = tonumber(r), tonumber(g), tonumber(b)
    color_data[name] = s_format(
      '%.4f %.4f %.4f rg %.4f %.4f %.4f RG', nr, ng, nb, nr, ng, nb)
  end
end, {'string', 'string'})

local function colorstack_push(color_name)
  local tpl = push_templates[color_name]
  if not tpl then
    tpl = d_new('whatsit', 'pdf_colorstack')
    d_setfield(tpl, 'stack', 0)
    d_setfield(tpl, 'command', 1)
    d_setfield(tpl, 'data', color_data[color_name] or '')
    push_templates[color_name] = tpl
  end
  return d_copy(tpl)
end

local pop_template = d_new('whatsit', 'pdf_colorstack')
d_setfield(pop_template, 'stack', 0)
d_setfield(pop_template, 'command', 2)
d_setfield(pop_template, 'data', '')

local function colorstack_pop() return d_copy(pop_template) end

local function append(rh, rt, h, t)
  if not h then return rh, rt end
  if not t then
    t = h; while d_getnext(t) do t = d_getnext(t) end
  end
  if rt then d_setlink(rt, h) else rh = h end
  return rh, t
end

-- Bounded glyph cache
local raw_glyph_cache = {}
local CACHE_MAX_LEN = 16

local function str_to_nodes(str, fid, fallback_fn)
  if not str or str == '' then return nil, nil end
  fid = fid or f_current()

  local cacheable = #str <= CACHE_MAX_LEN
  local key
  if cacheable then
    -- Optimización: concatenación implícita de número en vez de tostring()
    key = str .. '_' .. fid
    local cached = raw_glyph_cache[key]
    if cached then
      local h = d_copylist(cached)
      local t = h
      while d_getnext(t) do t = d_getnext(t) end
      return h, t
    end
  end

  fallback_fn = fallback_fn or get_lm_fid
  local head, tail, lm_fid, fallback_fid = nil, nil, nil, nil
  local chars = f_getfont(fid) and f_getfont(fid).characters

  for _, cp in u_codes(str) do
    local current_fid = fid
    if cp == LM_LANGLE or cp == LM_RANGLE then
      if not lm_fid then lm_fid = get_lm_fid(fid) end
      current_fid = lm_fid
    elseif chars and not chars[cp] then
      if not fallback_fid then fallback_fid = fallback_fn(fid) end
      current_fid = fallback_fid
    end
    local g = d_new('glyph')
    d_setfield(g, 'font', current_fid)
    d_setfield(g, 'char', cp)
    if head then d_setlink(tail, g) else head = g end
    tail = g
  end

  if head and cacheable then
    raw_glyph_cache[key] = d_copylist(head)
  end
  return head, tail
end

-- Envuelve texto en un Span con ActualText = mismo texto: los glyphs se
-- muestran igual visualmente, pero el PDF garantiza que al copiar/pegar
-- se extraigan exactamente esos caracteres, sin depender de cómo el
-- backend termine serializando corridas de glyphs (que a veces colapsan
-- espacios repetidos en puro kerning sin carácter extraíble, o incluso
-- provocan que el lector de PDF inserte un salto de línea espurio ante
-- un hueco horizontal ancho sin texto extraíble de por medio). Uso
-- previsto: SOLO indentación/espacios de alineación (siempre ASCII de
-- un byte) -- por eso convertir byte a byte a UTF-16BE es seguro aquí.
-- Si esta función se reutiliza alguna vez para texto arbitrario con
-- caracteres UTF-8 multibyte, este bucle byte-a-byte NO sería
-- suficiente (habría que decodificar codepoints Unicode reales, no
-- bytes crudos) -- habría que revisarlo antes de reutilizar.
--
-- EXPERIMENTO 2 (el 1, sustituir U+0020 por U+00A0, no funcionó en
-- Acrobat -- la heurística de recorte no distingue el TIPO de espacio):
-- se antepone un centinela U+200B (ZERO WIDTH SPACE) ANTES de los
-- espacios reales dentro del valor de ActualText. A diferencia de
-- U+00A0, Unicode NO clasifica U+200B como espacio en blanco (categoría
-- Zs) sino como carácter de formato (Cf) -- si la heurística de Acrobat
-- recorta basándose en "¿el primer carácter extraído de la fila es
-- espacio?", este centinela invisible rompe esa condición sin alterar
-- el glifo visible (que sigue siendo el espacio normal).
local function actual_text_wrap(text, fid)
  local begin_n = d_new('whatsit', 'pdf_literal')
  d_setfield(begin_n, 'mode', 2)
  -- Notación hex nativa del PDF (<...>): BOM UTF-16BE (FEFF) + centinela
  -- ZWSP (200B) + cada carácter real como 4 dígitos hex.
  local hex_content = 'FEFF200B'
  for i = 1, #text do
    hex_content = hex_content .. s_format('%04X', s_byte(text, i))
  end
  d_setfield(begin_n, 'data', '/Span << /ActualText <' .. hex_content .. '> >> BDC')
  local end_n = d_new('whatsit', 'pdf_literal')
  d_setfield(end_n, 'mode', 2)
  d_setfield(end_n, 'data', 'EMC')
  local h, t = str_to_nodes(text, fid)
  local rh, rt = append(nil, nil, begin_n, begin_n)
  rh, rt = append(rh, rt, h, t)
  rh, rt = append(rh, rt, end_n, end_n)
  return rh, rt
end

-- Como str_to_nodes, pero envuelve en ActualText cualquier corrida de
-- 2+ espacios consecutivos (típicas de alineación de columnas en
-- código expl3/TeX bien formateado). Texto normal (sin corridas
-- largas) no paga overhead extra -- pasa directo por str_to_nodes.
local function str_to_nodes_ws_safe(str, fid)
  local rh, rt, i, len = nil, nil, 1, #str
  while i <= len do
    local ws_start, ws_end = s_find(str, '  +', i)
    if not ws_start then
      local h, t = str_to_nodes(s_sub(str, i), fid)
      rh, rt = append(rh, rt, h, t)
      break
    end
    if ws_start > i then
      local h, t = str_to_nodes(s_sub(str, i, ws_start - 1), fid)
      rh, rt = append(rh, rt, h, t)
    end
    local h, t = actual_text_wrap(s_sub(str, ws_start, ws_end), fid)
    rh, rt = append(rh, rt, h, t)
    i = ws_end + 1
  end
  return rh, rt
end

local function colored_str(str, color_name, fid)
  local h, t = str_to_nodes(str, fid)
  if not h then return nil, nil end
  if color_data[color_name] then
    local push, pop = colorstack_push(color_name), colorstack_pop()
    d_setlink(push, h)
    d_setlink(t, pop)
    return push, pop
  end
  return h, t
end

local function colored_str_bold(str, color_name, fid)
  return colored_str(str, color_name, get_bold_fid(fid))
end

local function colored_str_italic(str, color_name, fid)
  local italic_fid = get_italic_fid(fid)
  local h, t = str_to_nodes(str, italic_fid, get_italic_fid)
  if not h then return nil, nil end
  if color_data[color_name] then
    local push, pop = colorstack_push(color_name), colorstack_pop()
    d_setlink(push, h)
    d_setlink(t, pop)
    return push, pop
  end
  return h, t
end

local type_colors = {
  [1] ='argument',   [2] ='privatevar', [3] ='bracket',
  [4] ='publicvar',  [5] ='privatefun', [6] ='publicfun',
  [8] ='brace',      [9] ='comment',    [10]='scan',
  [11]='quark',      [12]='number',     [13]='danger',
  [14]='ltcmd',      [16]='keyfun',
  [17]='escape',     [18]='escape',     [19]='constant',
  [20]='constant',   [21]='math',       [22]='compare',
  [23]='latex_cmd',  [24]='argument',   [25]='structure',
  [26]='publicfun',  [28]='quote',      [29]='brace',
  [30]='publicfun',  [31]='publicfun',  [33]='ltcmd_arg',
  [34]='publicfun',  [35]='unit',       [36]='compare',
  [37]='pkgstruct',  [38]='mathtag',    [39]='publicfun',  [40]='publicfun',
  [41]='publicfun',  [42]='intent',     [44]='publicfun',
  [46]='publicfun',
  [47]='publicfun',  [48]='module_name', [49]='support',
}

local csname_positions = {}
local arg_counter = 0
local in_csname, csname_base_color, csname_depth = false, nil, 0
local variant_braces_remaining, variant_depth = 0, 0
local in_variant_group = false
local ltcmd_brace_count, ltcmd_name_pending, ltcmd_in_args = 0, false, false
local in_mathml_intent, in_mathml_arg = false, false
local mathml_brace_count, mathml_after_dollar = 0, false
local in_keys_context, keys_body_opened, keys_depth = false, false, 0
local keys_in_choices = false

-- .meta:nn (2 args: ruta módulo/clave + lista de claves) y .meta:n
-- (1 arg: lista de claves). Reutiliza colorize_keys_path para la ruta
-- y colorize_trimmed_keyname para la lista, igual que el resto de keys.
local keys_meta_pending      = false
local keys_meta_arity        = 0  -- 1 = .meta:n, 2 = .meta:nn
local keys_meta_argnum       = 0  -- argumento actual (1 o 2)
local keys_meta_brace_count  = 0

local in_msgerror       = false
local msgerror_argnum   = 0
local msgerror_depth    = 0
local msgerror_is_unknown = false

local in_msg_simple    = false
local msg_simple_depth = 0
local msg_simple_done  = false

-- Resto de la familia \keys_...:n... (set_known, set_groups,
-- set_exclude_groups, precompile, if_exist, if_choice_exist, show, log):
-- arg1 (primer 'n' de la firma) = ruta módulo/path -> colorize_keys_path;
-- args 2..N (resto de las 'n' minúsculas de la firma) = nombre de
-- clave/grupo/opción -> colorize_trimmed_keyname; lo que venga después
-- (ramas T/F de un sufijo TF, argumentos N sin llaves, etc.) no se toca.
local in_keys_module_arg     = false
local keys_module_arg_depth  = 0
local keys_module_arg_done   = false
local keys_module_arg_argnum = 0
local keys_module_arg_total_n = 0

-- \ProvidesExplPackage {nombre} {...} {...} {...}: el PRIMER argumento
-- (el nombre del paquete/módulo) va en negrita + color module_name.
local provides_pkg_pending = false
local provides_pkg_brace_count = 0

-- Nombres de paquete/clase (\ProvidesPackage, \RequirePackage,
-- \ProvidesClass, \LoadClass, \PassOptionsTo(Package|Class),
-- \IfPackageLoaded...) -> negrita + color package_name/class_name.
-- pkgcls_target_arg indica en qué argumento {} está el nombre (1 o 2);
-- [opciones] opcional no afecta el conteo porque [] no toca estas
-- variables (solo {} las incrementa).
local pkgcls_pending     = false
local pkgcls_target_arg  = 0
local pkgcls_argnum      = 0
local pkgcls_brace_depth = 0
local pkgcls_color       = nil

-- \RequirePackage [opciones] {pkg} y \LoadClass [opciones] {cls}: los
-- corchetes opcionales, si aparecen, contienen una lista de opciones
-- (key=value, ...) -> mismo tratamiento keyname. Si en vez de '[' llega
-- '{' directamente (sin opciones), pkgcls_bracket_pending se limpia sin
-- activar nada.
local pkgcls_bracket_pending = false
local pkgcls_bracket_active  = false
local pkgcls_bracket_depth   = 0

-- Nombre de archivo: \file_..., \lua_load_module:n, \IfFileExists
-- (activado desde el handler de tipo 37), \input/\include/
-- \InputIfFileExists. Solo el primer argumento -> colorize_filename.
-- \includeonly recibe una LISTA separada por comas -> file_arg_is_list.
local in_file_arg     = false
local file_arg_depth  = 0
local file_arg_done   = false
local file_arg_is_list = false

-- \hook_gput_code:nnn/_with_args:nnn/\hook_gremove_code:nn (label en
-- arg2) y \hook_gset_rule:nnnn (labels en arg2 Y arg4, con 'relation'
-- en arg3 sin colorear). hook_label_positions guarda qué números de
-- argumento deben colorearse como module_name.
local in_hook_label        = false
local hook_label_depth     = 0
local hook_label_argnum    = 0
local hook_label_total_n   = 0
local hook_label_positions = nil

-- \hypersetup{key=val,...}, \tagpdfsetup{...}, \tagtool{...},
-- \tagmcbegin{...}, \tagstructbegin{...}, \ShowTagging{...}, y también
-- \unimathsetup{...}/\addfontfeature{...}/\defaultfontfeatures{...}:
-- único argumento tratado como lista de claves, mismo tratamiento
-- visual que el cuerpo de \keys_define:nn/\keys_set:nn
-- (colorize_trimmed_keyname).
local in_single_keyval    = false
local single_keyval_depth = 0

-- \setmainfont{fuente}[features] y familia (fontspec/unicode-math): el
-- ORDEN es inverso a \RequirePackage[opciones]{pkg} (llave primero,
-- corchete opcional después) -- mecanismo independiente de
-- pkgcls_bracket_*. font_brace_pending/depth trackea la primera llave
-- (nombre de fuente, SIN color especial); al cerrarse, se activa
-- font_bracket_pending esperando un '[' opcional. Salvaguarda: se
-- limpia al final de cada línea (en process_line/inline) para no
-- "contaminar" un [...] no relacionado en una línea posterior si el
-- comando se usó sin features.
local font_brace_pending   = false
local font_brace_depth     = 0
local font_bracket_pending = false
local font_bracket_active  = false
local font_bracket_depth   = 0

local colorize_code

local function reset_csname_state()
  for k in pairs(csname_positions) do csname_positions[k] = nil end
  arg_counter, in_csname, csname_base_color, csname_depth = 0, false, nil, 0
end

local function module_aware_color(chunk, base_color)
  local name = s_sub(chunk, 1, 1) == '\\' and s_sub(chunk, 2) or chunk
  local colon = s_find(name, ':', 1, true)
  local bare = colon and s_sub(name, 1, colon - 1) or name
  if module_user_cmds[bare] or module_user_cmds[name] then
    return 'module_name'
  end
  if current_module_name ~= '' and s_find(name, current_module_name, 1, true)
     and not s_find(name, '__', 1, true) then
    return 'module_name'
  end
  return base_color
end

-- Detecta \l__foo, \g__foo, \c__foo, \__foo cuyo segmento de módulo NO
-- coincide con current_module_name: acceso explícito a privados de OTRO
-- módulo (hackeo de internos ajenos) -> color de advertencia 'other_np'.
local function private_owner_color(chunk, fallback_color)
  local name = s_sub(chunk, 1, 1) == '\\' and s_sub(chunk, 2) or chunk
  local c3 = s_sub(name, 1, 3)
  local rest
  if c3 == 'l__' or c3 == 'g__' or c3 == 'c__' then
    rest = s_sub(name, 4)
  elseif s_sub(name, 1, 2) == '__' then
    rest = s_sub(name, 3)
  else
    return fallback_color
  end
  if current_module_name == '' then return fallback_color end
  if module_user_cmds[rest] then return fallback_color end
  if s_find(rest, current_module_name, 1, true) then return fallback_color end
  return 'other_np'
end

local function emit_with_signature(chunk, base_color, fid)
  local colon = s_find(chunk, ':', 2, true)
  if not colon then
    reset_csname_state()
    return colored_str(chunk, base_color, fid)
  end
  local name, sig = s_sub(chunk, 1, colon - 1), s_sub(chunk, colon + 1)
  reset_csname_state()
  local pos = 0
  for ch in s_gmatch(sig, '.') do
    if s_match(ch, '[NnVvoxefpTFwDcq]') then
      pos = pos + 1
      if ch == 'c' or ch == 'v' then csname_positions[pos] = true end
    end
  end
  local rh, rt = nil, nil
  local h, t = colored_str(name, base_color, fid); rh, rt = append(rh, rt, h, t)
  h, t = colored_str(':', 'signature', fid); rh, rt = append(rh, rt, h, t)
  if sig ~= '' then
    h, t = colored_str(sig, 'sigargs', fid); rh, rt = append(rh, rt, h, t)
  end
  return rh, rt
end

local function colorize_ltcmd_args(text, fid, rh, rt)
  local i, depth, text_len = 1, 0, #text
  while i <= text_len do
    local b = s_byte(text, i)
    if b == 123 then
      depth = depth + 1
      local h, t = colored_str('{', 'brace', fid); rh, rt = append(rh, rt, h, t)
      i = i + 1
    elseif b == 125 then
      depth = depth - 1
      local h, t = colored_str('}', 'brace', fid); rh, rt = append(rh, rt, h, t)
      i = i + 1
    elseif depth > 0 then
      local j, d = i, depth
      while j <= text_len do
        local bb = s_byte(text, j)
        if bb == 123 then d = d + 1
        elseif bb == 125 then d = d - 1; if d < depth then break end end
        j = j + 1
      end
      local content = s_sub(text, i, j - 1)
      if content ~= '' then
        local ch, ct = colorize_code(content, fid); rh, rt = append(rh, rt, ch, ct)
      end
      i = j
    elseif b == 32 then
      local h, t = str_to_nodes(' ', fid); rh, rt = append(rh, rt, h, t)
      i = i + 1
    else
      local h, t = colored_str(s_sub(text, i, i), 'ltcmd_arg', fid)
      rh, rt = append(rh, rt, h, t); i = i + 1
    end
  end
  return rh, rt
end

local function csname_color_from_prefix(text)
  text = s_match(text, '^%s*(.-)%s*$') or text
  local p2, p3 = s_sub(text, 1, 2), s_sub(text, 1, 3)
  local base, rest
  if p3 == 'l__' or p3 == 'g__' then base, rest = 'privatevar', s_sub(text, 4)
  elseif p3 == 'c__' then base, rest = 'constant', s_sub(text, 4)
  elseif p2 == '__' then base, rest = 'privatefun', s_sub(text, 3)
  elseif p2 == 'l_' or p2 == 'g_' then return 'publicvar'
  elseif p2 == 'c_' then return 'constant'
  else return 'latex_cmd'
  end
  if current_module_name == '' then return base end
  if module_user_cmds[rest] then return base end
  if s_find(rest, current_module_name, 1, true) then return base end
  return 'other_np'
end

local function emit_csname_text(text, base_color, fid, rh, rt)
  local has_special, len = false, #text
  for i = 1, len do
    local b = s_byte(text, i)
    if b == 64 or b == 33 or b == 38 or b == 94 or b == 42 or b == 58 then
      has_special = true; break
    end
  end
  if not has_special then
    local h, t = colored_str(text, base_color, fid); return append(rh, rt, h, t)
  end
  local cur = 1
  for i = 1, len do
    local b = s_byte(text, i)
    if b == 64 or b == 33 or b == 38 or b == 94 or b == 42 then
      if i > cur then
        local h, t = colored_str(s_sub(text, cur, i-1), base_color, fid)
        rh, rt = append(rh, rt, h, t)
      end
      local color_name = (b == 64 and 'arroba') or (b == 42 and 'star') or 'escape'
      local h, t = colored_str(s_sub(text, i, i), color_name, fid)
      rh, rt = append(rh, rt, h, t); cur = i + 1
    elseif b == 58 then
      if i > cur then
        local h, t = colored_str(s_sub(text, cur, i-1), base_color, fid)
        rh, rt = append(rh, rt, h, t)
      end
      local h, t = colored_str(':', 'signature', fid)
      rh, rt = append(rh, rt, h, t); local rest = s_sub(text, i+1)
      if rest ~= '' then
        h, t = colored_str(rest, 'sigargs', fid); rh, rt = append(rh, rt, h, t)
      end
      cur = len + 1; break
    end
  end
  if cur <= len then
    local h, t = colored_str(s_sub(text, cur), base_color, fid)
    rh, rt = append(rh, rt, h, t)
  end
  return rh, rt
end

local function colorize_mathml_intent(text, fid, rh, rt)
  local i, len = 1, #text
  while i <= len do
    local c, b = s_sub(text, i, i), s_byte(text, i)
    if c == ':' then
      local j = i + 1
      -- saltar espacios opcionales entre ':' y la palabra de intent
      while j <= len and s_byte(text, j) == 32 do j = j + 1 end
      local word_start = j
      while j <= len do
        local cb = s_byte(text, j)
        if (cb >= 65 and cb <= 90) or (cb >= 97 and cb <= 122) or
           (cb >= 48 and cb <= 57) or cb == 95 or cb == 45 then j = j + 1
        else break end
      end
      local h, t = colored_str(':', 'signature', fid)
      rh, rt = append(rh, rt, h, t)
      if word_start > i + 1 then
        h, t = str_to_nodes(s_sub(text, i+1, word_start-1), fid)
        rh, rt = append(rh, rt, h, t)
      end
      if j > word_start then
        h, t = colored_str(s_sub(text, word_start, j-1), 'ltcmd_arg', fid)
        rh, rt = append(rh, rt, h, t)
      end
      i = j
    elseif c == '$' then
      local h, t = colored_str('$', 'math', fid); rh, rt = append(rh, rt, h, t)
      local j = i + 1
      while j <= len do
        local cb = s_byte(text, j)
        if (cb >= 65 and cb <= 90) or (cb >= 97 and cb <= 122) or
           (cb >= 48 and cb <= 57) or cb == 95 then j = j + 1
        else break end
      end
      if j > i + 1 then
        h, t = colored_str(s_sub(text, i+1, j-1), 'mathml_var', fid)
        rh, rt = append(rh, rt, h, t)
      end
      i = j
    elseif c == '(' or c == ')' or c == ',' then
      local h, t = colored_str(c, 'brace', fid); rh, rt = append(rh, rt, h, t)
      i = i + 1
    elseif c == ' ' then
      local h, t = str_to_nodes(c, fid); rh, rt = append(rh, rt, h, t); i = i + 1
    else
      local j = i + 1
      while j <= len do
        local nc = s_sub(text, j, j)
        if nc == ':' or nc == '$' or nc == '(' or nc == ')' or
           nc == ',' or nc == ' ' then break end
        j = j + 1
      end
      local chunk = s_sub(text, i, j-1)
      local saved_intent, saved_arg = in_mathml_intent, in_mathml_arg
      in_mathml_intent, in_mathml_arg = false, false
      local ch, ct = colorize_code(chunk, fid)
      in_mathml_intent, in_mathml_arg = saved_intent, saved_arg
      rh, rt = append(rh, rt, ch, ct); i = j
    end
  end
  return rh, rt
end

-- Nombre de archivo: divide en '.' (color 'signature', ya existente)
-- y el resto del texto (color 'filename', independiente y editable).
local function colorize_filename(text, fid, rh, rt)
  local i, len, start = 1, #text, 1
  while i <= len do
    if s_sub(text, i, i) == '.' then
      if i > start then
        local h, t = colored_str(s_sub(text, start, i - 1), 'filename', fid)
        rh, rt = append(rh, rt, h, t)
      end
      local h, t = colored_str('.', 'signature', fid)
      rh, rt = append(rh, rt, h, t)
      start = i + 1
    end
    i = i + 1
  end
  if start <= len then
    local h, t = colored_str(s_sub(text, start), 'filename', fid)
    rh, rt = append(rh, rt, h, t)
  end
  return rh, rt
end

-- Lista de archivos separados por comas (\includeonly{cap1,cap2,cap3}):
-- mismo tratamiento de '.' que colorize_filename, y además separa por
-- ',' (color 'brace', igual que el resto de comas separadoras del archivo).
local function colorize_filename_list(text, fid, rh, rt)
  local i, len, start = 1, #text, 1
  while i <= len do
    local c = s_sub(text, i, i)
    if c == '.' or c == ',' then
      if i > start then
        local h, t = colored_str(s_sub(text, start, i - 1), 'filename', fid)
        rh, rt = append(rh, rt, h, t)
      end
      local h, t = colored_str(c, c == '.' and 'signature' or 'brace', fid)
      rh, rt = append(rh, rt, h, t)
      start = i + 1
    end
    i = i + 1
  end
  if start <= len then
    local h, t = colored_str(s_sub(text, start), 'filename', fid)
    rh, rt = append(rh, rt, h, t)
  end
  return rh, rt
end

local function colorize_keys_path(text, fid, rh, rt)
  local i, len = 1, #text
  while i <= len do
    local b = s_byte(text, i)
    if (b >= 65 and b <= 90) or (b >= 97 and b <= 122) or b == 95 or b == 45 then
      local j = i + 1
      while j <= len do
        local bb = s_byte(text, j)
        if (bb >= 65 and bb <= 90) or (bb >= 97 and bb <= 122) or
           bb == 95 or bb == 45 then j = j + 1 else break end
      end
      local word = s_sub(text, i, j - 1)
      if module_user_cmds[word] or
         (current_module_name ~= '' and word == current_module_name) then
        local h, t = colored_str(word, 'module_name', fid)
        rh, rt = append(rh, rt, h, t)
      else
        local h, t = str_to_nodes(word, fid)
        rh, rt = append(rh, rt, h, t)
      end
      i = j
    else
      local h, t = str_to_nodes(s_sub(text, i, i), fid)
      rh, rt = append(rh, rt, h, t)
      i = i + 1
    end
  end
  return rh, rt
end

local function colorize_trimmed_keyname(text, fid, rh, rt)
  local lead = s_match(text, '^(%s*)')
  local core = s_sub(text, #lead + 1)
  local trimmed = s_match(core, '^(.-)%s*$')
  local trail = s_sub(core, #trimmed + 1)
  if lead ~= '' then
    local h2, t2 = str_to_nodes_ws_safe(lead, fid)
    rh, rt = append(rh, rt, h2, t2)
  end
  if trimmed ~= '' then
    local h, t = colored_str(trimmed, 'keyname', fid)
    rh, rt = append(rh, rt, h, t)
  end
  if trail ~= '' then
    local h, t = str_to_nodes_ws_safe(trail, fid)
    rh, rt = append(rh, rt, h, t)
  end
  return rh, rt
end

local function ltcmd_name_color(text)
  text = s_match(text, '^%s*(.-)%s*$') or text
  local name = s_sub(text, 1, 1) == '\\' and s_sub(text, 2) or text
  local colon = s_find(name, ':', 1, true)
  local bare = colon and s_sub(name, 1, colon - 1) or name
  if module_user_cmds[bare] or module_user_cmds[text] then return 'module_name' end
  if current_module_name ~= '' and s_find(bare, current_module_name, 1, true) then
    return 'module_name'
  end
  return 'latex_cmd'
end

local function colorize_plain_text(text, fid, rh, rt)
  local h, tl
  if in_variant_group then h, tl = colored_str(text, 'sigargs', fid)
  elseif in_csname then
    csname_base_color = csname_base_color or csname_color_from_prefix(text)
    rh, rt = emit_csname_text(text, csname_base_color, fid, rh, rt)
    h, tl = nil, nil
  elseif in_mathml_intent then
    if mathml_after_dollar then
      mathml_after_dollar = false
      local ident = s_match(text, '^%s*([%a%d_]+)')
      if ident then
        h, tl = colored_str(ident, 'mathml_var', fid)
        rh, rt = append(rh, rt, h, tl)
        local fnd = s_find(text, ident, 1, true) or 1
        local rest = s_sub(text, fnd + #ident)
        if rest ~= '' then rh, rt = colorize_mathml_intent(rest, fid, rh, rt) end
        h, tl = nil, nil
      else
        rh, rt = colorize_mathml_intent(text, fid, rh, rt); h, tl = nil, nil
      end
    else
      rh, rt = colorize_mathml_intent(text, fid, rh, rt); h, tl = nil, nil
    end
  elseif in_mathml_arg then h, tl = colored_str(text, 'mathml_var', fid)
  elseif keys_meta_pending and keys_meta_brace_count == 1
         and keys_meta_argnum == 1 and keys_meta_arity == 2 then
    -- .meta:nn primer argumento: ruta módulo/clave (spintent/spnum)
    rh, rt = colorize_keys_path(text, fid, rh, rt); h, tl = nil, nil
  elseif keys_meta_pending and keys_meta_brace_count == 1 then
    -- .meta:n único argumento, o .meta:nn segundo argumento: lista de claves
    rh, rt = colorize_trimmed_keyname(text, fid, rh, rt); h, tl = nil, nil
  elseif in_keys_context and keys_body_opened and keys_depth == 1 then
    rh, rt = colorize_trimmed_keyname(text, fid, rh, rt); h, tl = nil, nil
  elseif in_keys_context and keys_body_opened and keys_in_choices and keys_depth == 2 then
    rh, rt = colorize_trimmed_keyname(text, fid, rh, rt); h, tl = nil, nil
  elseif in_msgerror and msgerror_depth == 1 and msgerror_argnum == 1 then
    h, tl = colored_str(text, 'module_name', fid)
  elseif in_msgerror and msgerror_depth == 1 and msgerror_argnum == 2 then
    local trimmed = s_match(text, '^%s*(.-)%s*$')
    msgerror_is_unknown = (trimmed == 'unknown-choice')
    h, tl = str_to_nodes(text, fid)
  elseif in_msgerror and msgerror_depth == 1 and msgerror_is_unknown
         and (msgerror_argnum == 3 or msgerror_argnum == 4) then
    rh, rt = colorize_trimmed_keyname(text, fid, rh, rt); h, tl = nil, nil
  elseif ltcmd_in_args then
    if ltcmd_brace_count > 1 then h, tl = str_to_nodes(text, fid)
    else rh, rt = colorize_ltcmd_args(text, fid, rh, rt); h, tl = nil, nil end
  elseif ltcmd_name_pending and ltcmd_brace_count == 1 then
    rh, rt = emit_csname_text(text, ltcmd_name_color(text), fid, rh, rt)
    h, tl = nil, nil
  elseif in_keys_context and not keys_body_opened then
    rh, rt = colorize_keys_path(text, fid, rh, rt); h, tl = nil, nil
  elseif in_msg_simple and not msg_simple_done and msg_simple_depth == 1 then
    rh, rt = colorize_keys_path(text, fid, rh, rt); h, tl = nil, nil
  elseif in_keys_module_arg and not keys_module_arg_done
         and keys_module_arg_depth == 1 then
    if keys_module_arg_argnum == 1 then
      rh, rt = colorize_keys_path(text, fid, rh, rt); h, tl = nil, nil
    else
      rh, rt = colorize_trimmed_keyname(text, fid, rh, rt); h, tl = nil, nil
    end
  elseif provides_pkg_pending and provides_pkg_brace_count == 1 then
    h, tl = colored_str_bold(text, 'module_name', fid)
  elseif pkgcls_pending and pkgcls_brace_depth == 1
         and pkgcls_argnum == pkgcls_target_arg then
    h, tl = colored_str_bold(text, pkgcls_color, fid)
  elseif pkgcls_pending and pkgcls_brace_depth == 1
         and pkgcls_target_arg == 2 and pkgcls_argnum == 1 then
    -- \PassOptionsToPackage/\PassOptionsToClass: primer argumento =
    -- lista de opciones (key=value, ...), mismo tratamiento que keyname.
    rh, rt = colorize_trimmed_keyname(text, fid, rh, rt); h, tl = nil, nil
  elseif pkgcls_bracket_active and pkgcls_bracket_depth == 1 then
    -- \RequirePackage/\LoadClass [opciones]: mismo tratamiento keyname.
    rh, rt = colorize_trimmed_keyname(text, fid, rh, rt); h, tl = nil, nil
  elseif in_file_arg and not file_arg_done and file_arg_depth == 1 then
    if file_arg_is_list then
      rh, rt = colorize_filename_list(text, fid, rh, rt)
    else
      rh, rt = colorize_filename(text, fid, rh, rt)
    end
    h, tl = nil, nil
  elseif in_hook_label and hook_label_depth == 1
         and hook_label_positions[hook_label_argnum] then
    h, tl = colored_str(text, 'module_name', fid)
  elseif in_single_keyval and single_keyval_depth == 1 then
    rh, rt = colorize_trimmed_keyname(text, fid, rh, rt); h, tl = nil, nil
  elseif font_bracket_active and font_bracket_depth == 1 then
    rh, rt = colorize_trimmed_keyname(text, fid, rh, rt); h, tl = nil, nil
  else h, tl = str_to_nodes_ws_safe(text, fid) end
  return append(rh, rt, h, tl)
end

colorize_code = function(line, fid)
  local rh, rt, pos, line_len = nil, nil, 1, #line
  while pos <= line_len do
    local b, e, t = lpeg_match(combined_pattern, line, pos)
    if not b then
      local remainder = s_sub(line, pos)
      rh, rt = colorize_plain_text(remainder, fid, rh, rt)
      break
    end
    if b > pos then
      local prefix = s_sub(line, pos, b - 1)
      rh, rt = colorize_plain_text(prefix, fid, rh, rt)
    end
    local chunk = s_sub(line, b, e - 1)
    local h, tl
    if t == 8 then
      if chunk == '{' then
        arg_counter = arg_counter + 1
        local is_csname = csname_positions[arg_counter]
        if not is_csname and arg_counter == 1 then
          for k in pairs(csname_positions) do
            if k >= 1 then is_csname = true; break end
          end
        end
        if is_csname then
          in_csname, csname_base_color, csname_depth = true, nil, 0
          in_variant_group, variant_braces_remaining, variant_depth = false, 0, 0
        elseif in_csname then csname_depth = csname_depth + 1 end
        if in_variant_group then variant_depth = variant_depth + 1
        elseif variant_braces_remaining > 0 then in_variant_group = true end
        if ltcmd_name_pending or ltcmd_in_args then
          ltcmd_brace_count = ltcmd_brace_count + 1
        end
        if in_mathml_intent or in_mathml_arg then
          mathml_brace_count = mathml_brace_count + 1
        end
        if in_keys_context then
          keys_depth = keys_depth + 1
        end
        if keys_meta_pending then
          keys_meta_brace_count = keys_meta_brace_count + 1
          if keys_meta_brace_count == 1 then
            keys_meta_argnum = keys_meta_argnum + 1
          end
        end
        if in_msgerror then
          if msgerror_depth == 0 then
            msgerror_argnum = msgerror_argnum + 1
          end
          msgerror_depth = msgerror_depth + 1
        end
        if in_msg_simple and not msg_simple_done then
          msg_simple_depth = msg_simple_depth + 1
        end
        if in_keys_module_arg and not keys_module_arg_done then
          keys_module_arg_depth = keys_module_arg_depth + 1
          if keys_module_arg_depth == 1 then
            keys_module_arg_argnum = keys_module_arg_argnum + 1
          end
        end
        if provides_pkg_pending then
          provides_pkg_brace_count = provides_pkg_brace_count + 1
        end
        if pkgcls_pending then
          pkgcls_brace_depth = pkgcls_brace_depth + 1
          if pkgcls_brace_depth == 1 then
            pkgcls_argnum = pkgcls_argnum + 1
          end
        end
        pkgcls_bracket_pending = false
        if in_file_arg and not file_arg_done then
          file_arg_depth = file_arg_depth + 1
        end
        if in_hook_label then
          hook_label_depth = hook_label_depth + 1
          if hook_label_depth == 1 then
            hook_label_argnum = hook_label_argnum + 1
          end
        end
        if in_single_keyval then
          single_keyval_depth = single_keyval_depth + 1
        end
        if font_brace_pending then
          font_brace_depth = font_brace_depth + 1
        end
        -- Si esperábamos un '[' opcional y en cambio llega otra '{'
        -- (sin corchete de por medio), descartamos la espera.
        font_bracket_pending = false
      elseif chunk == '}' then
        if in_csname then
          if csname_depth > 0 then csname_depth = csname_depth - 1
          else in_csname, csname_depth = false, 0 end
        end
        if in_variant_group then
          if variant_depth > 0 then variant_depth = variant_depth - 1
          else
            in_variant_group = false
            variant_braces_remaining = variant_braces_remaining - 1
          end
        end
        if ltcmd_in_args and ltcmd_brace_count > 0 then
          ltcmd_brace_count = ltcmd_brace_count - 1
          if ltcmd_brace_count == 0 then ltcmd_in_args = false end
        elseif ltcmd_name_pending and ltcmd_brace_count > 0 then
          ltcmd_brace_count = ltcmd_brace_count - 1
          if ltcmd_brace_count == 0 then
            ltcmd_name_pending, ltcmd_in_args = false, true
          end
        end
        if mathml_brace_count > 0 then
          mathml_brace_count = mathml_brace_count - 1
          if mathml_brace_count == 0 then
            in_mathml_intent, in_mathml_arg = false, false
          end
        end
        if in_keys_context and keys_depth > 0 then
          keys_depth = keys_depth - 1
          if keys_in_choices and keys_depth == 1 then
            keys_in_choices = false
          end
          if keys_depth == 0 then
            if not keys_body_opened then
              keys_body_opened = true
            else
              in_keys_context, keys_body_opened = false, false
            end
          end
        end
        if keys_meta_pending and keys_meta_brace_count > 0 then
          keys_meta_brace_count = keys_meta_brace_count - 1
          if keys_meta_brace_count == 0 then
            if keys_meta_arity == 1 or keys_meta_argnum >= 2 then
              keys_meta_pending = false
            end
          end
        end
        if in_msgerror and msgerror_depth > 0 then
          msgerror_depth = msgerror_depth - 1
          if msgerror_depth == 0 and msgerror_argnum >= 5 then
            in_msgerror, msgerror_argnum, msgerror_is_unknown = false, 0, false
          end
        end
        if in_msg_simple and not msg_simple_done and msg_simple_depth > 0 then
          msg_simple_depth = msg_simple_depth - 1
          if msg_simple_depth == 0 then
            msg_simple_done, in_msg_simple = true, false
          end
        end
        if in_keys_module_arg and not keys_module_arg_done
           and keys_module_arg_depth > 0 then
          keys_module_arg_depth = keys_module_arg_depth - 1
          if keys_module_arg_depth == 0
             and keys_module_arg_argnum >= keys_module_arg_total_n then
            keys_module_arg_done, in_keys_module_arg = true, false
          end
        end
        if provides_pkg_pending and provides_pkg_brace_count > 0 then
          provides_pkg_brace_count = provides_pkg_brace_count - 1
          if provides_pkg_brace_count == 0 then
            provides_pkg_pending = false
          end
        end
        if pkgcls_pending and pkgcls_brace_depth > 0 then
          pkgcls_brace_depth = pkgcls_brace_depth - 1
          if pkgcls_brace_depth == 0 and pkgcls_argnum >= pkgcls_target_arg then
            pkgcls_pending = false
          end
        end
        if in_file_arg and not file_arg_done and file_arg_depth > 0 then
          file_arg_depth = file_arg_depth - 1
          if file_arg_depth == 0 then
            file_arg_done, in_file_arg = true, false
          end
        end
        if in_hook_label and hook_label_depth > 0 then
          hook_label_depth = hook_label_depth - 1
          if hook_label_depth == 0 and hook_label_argnum >= hook_label_total_n then
            in_hook_label = false
          end
        end
        if in_single_keyval and single_keyval_depth > 0 then
          single_keyval_depth = single_keyval_depth - 1
          if single_keyval_depth == 0 then
            in_single_keyval = false
          end
        end
        if font_brace_pending and font_brace_depth > 0 then
          font_brace_depth = font_brace_depth - 1
          if font_brace_depth == 0 then
            font_brace_pending = false
            font_bracket_pending = true
          end
        end
      end
      h, tl = colored_str(chunk, 'brace', fid)
    elseif t == 3 then
      if chunk == '[' then
        if pkgcls_bracket_pending then
          pkgcls_bracket_active, pkgcls_bracket_depth, pkgcls_bracket_pending =
            true, 1, false
        elseif pkgcls_bracket_active then
          pkgcls_bracket_depth = pkgcls_bracket_depth + 1
        end
        if font_bracket_pending then
          font_bracket_active, font_bracket_depth, font_bracket_pending =
            true, 1, false
        elseif font_bracket_active then
          font_bracket_depth = font_bracket_depth + 1
        end
      elseif chunk == ']' then
        if pkgcls_bracket_active and pkgcls_bracket_depth > 0 then
          pkgcls_bracket_depth = pkgcls_bracket_depth - 1
          if pkgcls_bracket_depth == 0 then
            pkgcls_bracket_active = false
          end
        end
        if font_bracket_active and font_bracket_depth > 0 then
          font_bracket_depth = font_bracket_depth - 1
          if font_bracket_depth == 0 then
            font_bracket_active = false
          end
        end
      end
      h, tl = colored_str(chunk, 'bracket', fid)
    elseif t == 24 then
      local hashes = s_match(chunk, '^#+')
      local num    = s_sub(chunk, #hashes + 1)
      h, tl = colored_str(hashes, 'argument', fid); rh, rt = append(rh, rt, h, tl)
      h, tl = colored_str(num, 'argnum', fid)
    elseif t == 26 then
      arg_counter, variant_depth, in_variant_group = 0, 0, false
      if s_find(chunk, 'prg_generate_conditional_variant', 1, true) then
        variant_braces_remaining = 2
      else variant_braces_remaining = 1 end
      h, tl = emit_with_signature(chunk, 'publicfun', fid)
    elseif t == 30 then
      in_csname, csname_base_color, csname_depth = true, nil, 0
      h, tl = colored_str(chunk, 'publicfun', fid)
    elseif t == 31 then
      in_csname = false; h, tl = colored_str(chunk, 'publicfun', fid)
    elseif t == 16 then
      local dot  = s_sub(chunk, 1, 1)
      local rest = s_sub(chunk, 2)
      local colon = s_find(rest, ':', 1, true)
      h, tl = colored_str(dot, 'signature', fid); rh, rt = append(rh, rt, h, tl)
      if colon then
        local kname, kargs = s_sub(rest, 1, colon - 1), s_sub(rest, colon + 1)
        if in_keys_context and keys_body_opened and keys_depth == 1 then
          keys_in_choices = (kname == 'choices')
        end
        -- .meta:nn/.meta:n es autocontenido (sus dos argumentos son ruta
        -- de módulo + lista de claves, sin depender de estar dentro de
        -- un \keys_define:nn real) -- se reconoce sin importar el
        -- contexto, para que \myhlc también lo coloree correctamente.
        if kname == 'meta' then
          keys_meta_pending, keys_meta_argnum, keys_meta_brace_count =
            true, 0, 0
          keys_meta_arity = (kargs == 'n') and 1 or 2
        else
          keys_meta_pending = false
        end
        reset_csname_state(); local cpos = 0
        for i = 1, #kargs do
          local ch = s_sub(kargs, i, i)
          if s_match(ch, '[NnVvoxefpTFwDcq]') then
            cpos = cpos + 1
            if ch == 'c' or ch == 'v' then csname_positions[cpos] = true end
          end
        end
        h, tl = colored_str(kname, 'keyfun', fid); rh, rt = append(rh, rt, h, tl)
        h, tl = colored_str(':', 'signature', fid); rh, rt = append(rh, rt, h, tl)
        h, tl = colored_str(kargs, 'sigargs', fid)
      else h, tl = colored_str(rest, 'keyfun', fid) end
    elseif t == 9 then
      h, tl = colored_str(s_sub(chunk, 1, 1), 'guard_angle', fid)
      rh, rt = append(rh, rt, h, tl)
      h, tl = colored_str_italic(s_sub(chunk, 2), 'comment', fid)
    elseif t == 14 then
      arg_counter, ltcmd_brace_count, ltcmd_name_pending, ltcmd_in_args = 0, 0, true, false
      h, tl = emit_with_signature(chunk, 'ltcmd', fid)
    elseif t == 39 then
      in_keys_context, keys_body_opened, keys_depth = true, false, 0
      keys_in_choices = false
      h, tl = emit_with_signature(chunk, 'publicfun', fid)
    elseif t == 40 then
      in_msgerror, msgerror_argnum, msgerror_depth, msgerror_is_unknown =
        true, 0, 0, false
      h, tl = emit_with_signature(chunk, 'publicfun', fid)
    elseif t == 41 then
      in_msg_simple, msg_simple_depth, msg_simple_done = true, 0, false
      h, tl = emit_with_signature(chunk, 'publicfun', fid)
    elseif t == 44 then
      in_keys_module_arg, keys_module_arg_depth, keys_module_arg_done =
        true, 0, false
      keys_module_arg_argnum = 0
      local colon = s_find(chunk, ':', 1, true)
      local sig = colon and s_sub(chunk, colon + 1) or ''
      local n_count = 0
      for ch in s_gmatch(sig, '.') do
        if ch == 'n' then n_count = n_count + 1 end
      end
      keys_module_arg_total_n = n_count
      h, tl = emit_with_signature(chunk, 'publicfun', fid)
    elseif t == 46 then
      in_file_arg, file_arg_depth, file_arg_done = true, 0, false
      -- Coincidencia de PREFIJO exacto justo tras el '\' (no substring en
      -- cualquier parte): evita que \file_input:n "contenga" 'input' y
      -- se confunda con el \input clásico.
      local name_start = s_sub(chunk, 2)
      local function starts_with(prefix)
        return s_sub(name_start, 1, #prefix) == prefix
      end
      file_arg_is_list = starts_with('includeonly')
      local base
      if file_arg_is_list
         or starts_with('InputIfFileExists')
         or starts_with('include')
         or starts_with('input') then
        base = 'pkgstruct'
      else
        base = module_aware_color(chunk, 'publicfun')
      end
      h, tl = emit_with_signature(chunk, base, fid)
    elseif t == 47 then
      in_hook_label, hook_label_depth, hook_label_argnum = true, 0, 0
      local colon = s_find(chunk, ':', 1, true)
      local sig = colon and s_sub(chunk, colon + 1) or ''
      local n_count = 0
      for ch in s_gmatch(sig, '.') do
        if ch == 'n' then n_count = n_count + 1 end
      end
      hook_label_total_n = n_count
      if s_find(chunk, 'hook_gset_rule', 1, true) then
        hook_label_positions = { [2] = true, [4] = true }
      else
        hook_label_positions = { [2] = true }
      end
      h, tl = emit_with_signature(chunk, module_aware_color(chunk, 'publicfun'), fid)
    elseif t == 48 then
      -- package/NOMBRE/before|after (también class/... y file/...):
      -- colorea solo el NOMBRE de en medio.
      local slash1 = s_find(chunk, '/', 1, true)
      local slash2 = s_find(chunk, '/', slash1 + 1, true)
      h, tl = str_to_nodes(s_sub(chunk, 1, slash1), fid); rh, rt = append(rh, rt, h, tl)
      h, tl = colored_str(s_sub(chunk, slash1 + 1, slash2 - 1), 'module_name', fid)
      rh, rt = append(rh, rt, h, tl)
      h, tl = str_to_nodes(s_sub(chunk, slash2), fid)
    elseif t == 34 then
      h, tl = emit_with_signature(chunk, 'publicfun', fid)
      csname_positions[1] = true
    elseif t == 33 then h, tl = emit_with_signature(chunk, 'ltcmd_arg', fid)
    elseif ltcmd_name_pending and ltcmd_brace_count == 1 then
      rh, rt = emit_csname_text(chunk, ltcmd_name_color(chunk), fid, rh, rt)
      h, tl = nil, nil
    elseif ltcmd_name_pending and ltcmd_brace_count == 0 and s_sub(chunk, 1, 1) == '\\' then
      h, tl = colored_str(chunk, ltcmd_name_color(chunk), fid)
      ltcmd_name_pending, ltcmd_in_args = false, true
    elseif ltcmd_in_args then
      if ltcmd_brace_count > 1 then h, tl = colored_str(chunk, type_colors[t], fid)
      else rh, rt = colorize_ltcmd_args(chunk, fid, rh, rt); h, tl = nil, nil end
    elseif in_csname then
      csname_base_color = csname_base_color or csname_color_from_prefix(chunk)
      if s_sub(chunk, 1, 1) == '\\' then
        h, tl = colored_str(chunk, type_colors[t], fid)
      else rh, rt = emit_csname_text(chunk, csname_base_color, fid, rh, rt); h, tl = nil, nil end
    elseif t == 5 or t == 6 then
      arg_counter = 0
      if ltcmd_in_args or ltcmd_name_pending then
        ltcmd_in_args, ltcmd_name_pending, ltcmd_brace_count = false, false, 0
      end
      local eff
      if t == 6 then
        eff = module_aware_color(chunk, 'publicfun')
      else
        eff = private_owner_color(chunk, 'privatefun')
      end
      h, tl = emit_with_signature(chunk, eff, fid)
    elseif t == 25 then
      arg_counter = 0
      if s_find(chunk, 'ProvidesExplPackage', 1, true) then
        provides_pkg_pending, provides_pkg_brace_count = true, 0
      end
      h, tl = colored_str_bold(chunk, 'structure', fid)
    elseif t == 37 then
      arg_counter = 0
      pkgcls_bracket_pending = false
      if s_find(chunk, 'ProvidesPackage', 1, true) then
        pkgcls_pending, pkgcls_target_arg, pkgcls_color = true, 1, 'pkgname'
      elseif s_find(chunk, 'ProvidesClass', 1, true) then
        pkgcls_pending, pkgcls_target_arg, pkgcls_color = true, 1, 'clsname'
      elseif s_find(chunk, 'RequirePackageWithOptions', 1, true) then
        pkgcls_pending, pkgcls_target_arg, pkgcls_color = true, 1, 'pkgname'
      elseif s_find(chunk, 'RequirePackage', 1, true) then
        pkgcls_pending, pkgcls_target_arg, pkgcls_color = true, 1, 'pkgname'
        pkgcls_bracket_pending = true
      elseif s_find(chunk, 'LoadClassWithOptions', 1, true) then
        pkgcls_pending, pkgcls_target_arg, pkgcls_color = true, 1, 'clsname'
      elseif s_find(chunk, 'LoadClass', 1, true) then
        pkgcls_pending, pkgcls_target_arg, pkgcls_color = true, 1, 'clsname'
        pkgcls_bracket_pending = true
      elseif s_find(chunk, 'PassOptionsToPackage', 1, true) then
        pkgcls_pending, pkgcls_target_arg, pkgcls_color = true, 2, 'pkgname'
      elseif s_find(chunk, 'PassOptionsToClass', 1, true) then
        pkgcls_pending, pkgcls_target_arg, pkgcls_color = true, 2, 'clsname'
      elseif s_find(chunk, 'IfPackageLoaded', 1, true) then
        pkgcls_pending, pkgcls_target_arg, pkgcls_color = true, 1, 'pkgname'
      else
        pkgcls_pending = false
      end
      if pkgcls_pending then pkgcls_argnum, pkgcls_brace_depth = 0, 0 end
      if s_find(chunk, 'IfFileExists', 1, true) then
        in_file_arg, file_arg_depth, file_arg_done = true, 0, false
      end
      h, tl = emit_with_signature(chunk, 'pkgstruct', fid)
    elseif t == 49 then
      if s_find(chunk, 'hypersetup', 1, true)
         or s_find(chunk, 'unimathsetup', 1, true)
         or s_find(chunk, 'addfontfeature', 1, true)
         or s_find(chunk, 'defaultfontfeatures', 1, true) then
        in_single_keyval, single_keyval_depth = true, 0
      elseif s_find(chunk, 'setmainfont', 1, true)
         or s_find(chunk, 'setsansfont', 1, true)
         or s_find(chunk, 'setmonofont', 1, true)
         or s_find(chunk, 'setboldmathrm', 1, true)
         or s_find(chunk, 'setmathrm', 1, true)
         or s_find(chunk, 'setmathsf', 1, true)
         or s_find(chunk, 'setmathtt', 1, true)
         or s_find(chunk, 'setmathfont', 1, true)
         or s_find(chunk, 'fontspec', 1, true) then
        font_brace_pending, font_brace_depth = true, 0
      end
      h, tl = emit_with_signature(chunk, 'support', fid)
    elseif t == 38 then
      mathml_brace_count, mathml_after_dollar = 0, false
      if s_find(chunk, 'MathMLintent', 1, true) then
        in_mathml_intent, in_mathml_arg = true, false
      elseif s_find(chunk, 'MathMLarg', 1, true) then
        in_mathml_arg, in_mathml_intent = true, false
      else in_mathml_intent, in_mathml_arg = false, false end
      if s_find(chunk, 'tagpdfsetup', 1, true)
         or s_find(chunk, 'tagtool', 1, true)
         or s_find(chunk, 'tagmcbegin', 1, true)
         or s_find(chunk, 'tagstructbegin', 1, true)
         or s_find(chunk, 'ShowTagging', 1, true)
         or s_find(chunk, 'tag_struct_begin:n', 1, true)
         or s_find(chunk, 'tag_mc_begin:n', 1, true) then
        in_single_keyval, single_keyval_depth = true, 0
      end
      h, tl = emit_with_signature(chunk, 'mathtag', fid)
    elseif in_mathml_intent and t ~= 8 then
      if t == 21 then mathml_after_dollar = true; h, tl = colored_str(s_sub(chunk, 1, 1), 'math', fid)
      elseif mathml_after_dollar then
        mathml_after_dollar = false
        rh, rt = colorize_mathml_intent(chunk, fid, rh, rt); h, tl = nil, nil
      else rh, rt = colorize_mathml_intent(chunk, fid, rh, rt); h, tl = nil, nil end
    elseif in_mathml_arg and t ~= 8 then h, tl = colored_str(chunk, 'mathml_var', fid)
    elseif t == 42 then
      h, tl = colored_str(':', 'signature', fid); rh, rt = append(rh, rt, h, tl)
      local rest = s_sub(chunk, 2)
      local ws = s_match(rest, '^([ \t]*)')
      local after_ws = s_sub(rest, #ws + 1)
      if ws ~= '' then
        h, tl = str_to_nodes(ws, fid); rh, rt = append(rh, rt, h, tl)
      end
      local paren_start = s_find(after_ws, '(', 1, true)
      local word = paren_start and s_sub(after_ws, 1, paren_start - 1) or after_ws
      h, tl = colored_str(word, 'intent', fid); rh, rt = append(rh, rt, h, tl)
      if paren_start then
        local paren_content = s_sub(after_ws, paren_start)
        rh, rt = colorize_mathml_intent(paren_content, fid, rh, rt)
      end
      h, tl = nil, nil
    elseif t == 35 then
      local split = s_find(chunk, '%a')
      if split then
        h, tl = colored_str(s_sub(chunk, 1, split-1), 'number', fid)
        rh, rt = append(rh, rt, h, tl)
        h, tl = colored_str(s_sub(chunk, split), 'unit', fid)
      else h, tl = colored_str(chunk, 'number', fid) end
    elseif t == 18 then
      local four_caret = '\x5e\x5e\x5e\x5e'
      if s_sub(chunk, 1, 4) == four_caret then
        h, tl = colored_str(four_caret, 'escape', fid); rh, rt = append(rh, rt, h, tl)
        h, tl = colored_str(s_sub(chunk, 5), 'number', fid)
      else h, tl = colored_str(chunk, 'escape', fid) end
    elseif t == 13 then h, tl = emit_with_signature(chunk, 'danger', fid)
    elseif t == 15 then
      local cur, len = 2, #chunk
      for i = 2, len do
        local b = s_byte(chunk, i)
        if b == 64 or b == 33 or b == 38 or b == 94 or b == 42 then
          local part = (cur == 2 and '\\' or '') .. s_sub(chunk, cur, i - 1)
          if #part > 0 then
            local pcolor = module_aware_color(part, 'latex_cmd')
            if s_find(part, 'nuevo', 1, true) then
              texio.write_nl('[[CALL-DEBUG]] part=[' .. part .. '] pcolor=[' .. pcolor .. '] color_data_module_name=' .. tostring(color_data['module_name'] ~= nil))
            end
            h, tl = colored_str(part, pcolor, fid)
            rh, rt = append(rh, rt, h, tl)
          end
          local color_name = (b == 64 and 'arroba') or (b == 42 and 'star') or 'escape'
          h, tl = colored_str(s_sub(chunk, i, i), color_name, fid)
          rh, rt = append(rh, rt, h, tl); cur = i + 1
        end
      end
      if cur <= len then
        local trailing = (cur == 2 and '\\' or '') .. s_sub(chunk, cur)
        h, tl = colored_str(trailing, module_aware_color(trailing, 'latex_cmd'), fid)
      else h, tl = nil, nil end
    elseif t == 4 then
      h, tl = colored_str(chunk, module_aware_color(chunk, 'publicvar'), fid)
    elseif t == 2 then
      h, tl = colored_str(chunk, private_owner_color(chunk, 'privatevar'), fid)
    elseif t == 19 then
      h, tl = colored_str(chunk, private_owner_color(chunk, 'constant'), fid)
    else
      local effective = type_colors[t]
      if effective then
        local bare = s_sub(chunk, 1, 1) == '\\' and s_sub(chunk, 2) or chunk
        if module_user_cmds[bare] or module_user_cmds[chunk] then effective = 'module_name' end
      end
      h, tl = colored_str(chunk, effective, fid)
    end
    rh, rt = append(rh, rt, h, tl); pos = e
  end
  return rh, rt
end

local function make_angle_nodes(content_head, content_tail, angle_color, fid)
  local lm_fid = get_lm_fid(fid)
  local cname = (color_data[angle_color] and angle_color) or 'guard_angle'
  local lg, rg = d_new('glyph'), d_new('glyph')
  d_setfield(lg, 'font', lm_fid)
  d_setfield(lg, 'char', LM_LANGLE)
  d_setfield(rg, 'font', lm_fid)
  d_setfield(rg, 'char', LM_RANGLE)
  local rh, rt = nil, nil
  if color_data[cname] then rh, rt = append(rh, rt, colorstack_push(cname)) end
  rh, rt = append(rh, rt, lg)
  if color_data[cname] then rh, rt = append(rh, rt, colorstack_pop()) end
  if content_head then rh, rt = append(rh, rt, content_head, content_tail) end
  if color_data[cname] then rh, rt = append(rh, rt, colorstack_push(cname)) end
  rh, rt = append(rh, rt, rg)
  if color_data[cname] then rh, rt = append(rh, rt, colorstack_pop()) end
  return rh, rt
end

-- Mapea cada símbolo de guarda a su propio color independiente
-- (antes todos compartían 'operator'; ahora cada uno es editable aparte).
local function guard_symbol_color(b)
  if b == 42 then return 'guard_star'    -- *
  elseif b == 47 then return 'guard_slash'  -- /
  elseif b == 40 or b == 41 then return 'guard_paren' -- ( )
  elseif b == 38 then return 'guard_amp'    -- &
  elseif b == 124 then return 'guard_pipe'  -- |
  elseif b == 33 then return 'guard_bang'   -- !
  end
  return 'guard_star' -- fallback (inalcanzable: los 6 casos cubren todos
                       -- los bytes que colorize_guard_expr le pasa a esta función)
end

local function colorize_guard_expr(expr, fid)
  local rh, rt, name_start, len = nil, nil, 1, #expr
  for i = 1, len do
    local b = s_byte(expr, i)
    if b == 33 or b == 32 or b == 38 or b == 40 or b == 41 or b == 42 or
       b == 47 or b == 124 then
      if name_start <= i - 1 then
        local h, t = colored_str(s_sub(expr, name_start, i - 1), 'guard_name', fid)
        rh, rt = append(rh, rt, h, t)
      end
      name_start = i + 1
      if b ~= 32 then
        local h, t = colored_str(s_sub(expr, i, i), guard_symbol_color(b), fid)
        rh, rt = append(rh, rt, h, t)
      end
    end
  end
  if name_start <= len then
    local h, t = colored_str(s_sub(expr, name_start, len), 'guard_name', fid)
    rh, rt = append(rh, rt, h, t)
  end
  return rh, rt
end

local function render_guard(line, fid)
  local rest = s_gsub(line, '^%s*', '')
  local close = s_find(rest, '>', 2)
  if not close then
    local h, t = colored_str('%', 'guard_angle', fid)
    if h then n_write(d_tonode(h)) end
    h, t = colored_str_italic(line, 'comment', fid)
    if h then n_write(d_tonode(h)) end
    return
  end
  local inner = s_sub(rest, 2, close - 1)
  local after = s_sub(rest, close + 1)
  local ch, ct, h, t = nil, nil, nil, nil

  if s_sub(inner, 1, 3) == '@@=' then
    local modname = s_sub(inner, 4)
    t_set_macro('g__codedoc_module_name_tl', modname, 'global')
    current_module_name = modname
    current_module_prefix = '__' .. modname
    h, t = colored_str('@@', 'module_at', fid); ch, ct = append(ch, ct, h, t)
    h, t = colored_str('=', 'module_eq', fid); ch, ct = append(ch, ct, h, t)
    h, t = colored_str_bold(modname, 'module_name', fid); ch, ct = append(ch, ct, h, t)
  else
    local sigil, expr, first = '', inner, s_sub(inner, 1, 1)
    if first == '*' or first == '/' then sigil = first; expr = s_sub(inner, 2) end
    if sigil ~= '' then
      h, t = colored_str(sigil, guard_symbol_color(s_byte(sigil)), fid)
      ch, ct = append(ch, ct, h, t)
    end
    if expr ~= '' then
      h, t = colorize_guard_expr(expr, fid); ch, ct = append(ch, ct, h, t)
    end
  end

  local rh, rt = make_angle_nodes(ch, ct, 'guard_angle', fid)
  if rh then n_write(d_tonode(rh)) end
  if after and after ~= '' then
    local ah, at = colorize_code(after, fid); if ah then n_write(d_tonode(ah)) end
  end
end

local function apply_at_at(line)
  if not s_find(line, '@@', 1, true) then return line end
  if current_module_name ~= '' then
    line = s_gsub(line, '@@@@', '\0')
    line = s_gsub(line, '__@@', current_module_prefix)
    line = s_gsub(line, '_@@',  current_module_prefix)
    line = s_gsub(line, '@@',   current_module_prefix)
    line = s_gsub(line, '\0',   '@@')
  end
  return line
end

-- El nombre de módulo activo (current_module_name/current_module_prefix) se
-- sincroniza desde \g__codedoc_module_name_tl al inicio de cada línea, pero
-- SOLO cuando el valor leído no está vacío: así se mantiene sincronía con
-- el mecanismo nativo de l3doc.cls (y con cualquier otro código externo
-- que también fije esa variable), sin que una lectura vacía intermedia
-- borre el estado que Lua ya tiene cacheado desde un %<@@=modulo> anterior.
local function init_line_state()
  local m = t_get_macro('g__codedoc_module_name_tl')
  if m and m ~= '' then
    current_module_name = m
    current_module_prefix = '__' .. m
  end
end

-- Reemplazo mínimo de \BeginAccSupp{ActualText={}}/\EndAccSupp{} (accsupp)
-- vía whatsit pdf_literal directo: envuelve el número de línea en un
-- Span PDF con ActualText vacío, para que no se copie al extraer texto.
-- mode=2 ('direct'): el literal se inserta exactamente en el punto del
-- content stream, sin reordenamiento -- necesario para que BDC/EMC
-- queden bien balanceados alrededor del contenido que envuelven.
--
-- HISTORIAL de intentos descartados (ver conversación para detalle):
-- - <FEFF200B> (BOM + centinela de ancho cero U+200B): pensado para
--   evitar un caso límite teórico de ActualText vacío -- nunca
--   confirmado como problema real, y en cambio introdujo un espacio
--   extra visible al inicio de cada línea copiada (el U+200B se
--   extrae como carácter real y algunos visores lo renderizan con
--   ancho). Descartado.
-- - /Artifact BMC/EMC: probado durante la investigación del recorte de
--   indentación en Adobe Acrobat -- no resolvió ese problema (quedó
--   documentado como limitación conocida) y además dejó de ocultar el
--   número de línea correctamente. Descartado.
-- Versión actual: ActualText vacío simple, la única confirmada sin
-- efectos secundarios.
register_tex_cmd('actual_text_begin', function()
  local n = d_new('whatsit', 'pdf_literal')
  d_setfield(n, 'mode', 2)
  d_setfield(n, 'data', '/Span << /ActualText () >> BDC')
  n_write(d_tonode(n))
end, {})

register_tex_cmd('actual_text_end', function()
  local n = d_new('whatsit', 'pdf_literal')
  d_setfield(n, 'mode', 2)
  d_setfield(n, 'data', 'EMC')
  n_write(d_tonode(n))
end, {})


register_tex_cmd('inline', function()
  local text = t_get_macro('l__mylhmc_verbatim_line_tl')
  if not text or text == '' then return end
  text = s_gsub(text, '[\r\n]+$', '')
  init_line_state()
  text = apply_at_at(text)
  local fid = f_current()
  local rh, rt = colorize_code(text, fid)
  if rh then n_write(d_tonode(rh)) end
  font_bracket_pending = false
end, {})

register_tex_cmd('process_line', function()
  local line = t_get_macro('l__mylhmc_verbatim_line_tl')
  if not line or line == '' then return end
  line = s_gsub(line, '[\r\n]+$', '')
  if line == '' then return end
  init_line_state()
  line = apply_at_at(line)
  local fid = f_current()
  if s_sub(line, 1, 1) == '%' then
    local rest = s_sub(line, 2)
    local trimmed = s_match(rest, '^%s*(.*)')
    if s_sub(trimmed, 1, 1) == '<' then render_guard(rest, fid)
    else
      local h, t = colored_str(s_sub(line, 1, 1), 'guard_angle', fid)
      if h then n_write(d_tonode(h)) end
      h, t = colored_str_italic(s_sub(line, 2), 'comment', fid)
      if h then n_write(d_tonode(h)) end
    end
  else
    local indent = s_match(line, '^(%s+)')
    local rest = line
    local rh, rt
    if indent and indent ~= '' then
      rest = s_sub(line, #indent + 1)
      rh, rt = actual_text_wrap(indent, fid)
    end
    local ch, ct = colorize_code(rest, fid)
    rh, rt = append(rh, rt, ch, ct)
    if rh then n_write(d_tonode(rh)) end
  end
  font_bracket_pending = false
end, {})
% \iffalse
%</mylhmclua>
% \fi
